% Paper V-D v1.58 / Framework v3.058: fixed rooted Hopf return technical checks, no action-source occurrence.
% Paper source Paper5D-TechnicalMaterial-v1_56.tex: v3.040 S3 literature mathematical/conditional comparison; physical EM packet unresolved.
% Framework source integration v3.024 (audit v1.7--1.8): explicit action-jet source-port diagnostic, guards retained.
% Paper-series v2.382: audit-identified standalone repair integration; mathematical scopes and open physical owners preserved.
% Paper-series v2.381: complete algebraic exceptional torsion atlas and all-cocycle twist-stability; original claims retained.
% Paper-series v2.380: unified bibliography, exact exceptional double/length/gauge mathematics in I-B; previous theorem labels retained.
% Paper-series v2.379: full equivariant Dirac/Hodge/eta/torsion atlas in Paper I-B; paper-root references synchronized.
% Paper-series v2.378: 2026 exceptional literature consolidation and reference synchronization. The mathematical additions are located in Paper I-B.
% Paper-series v2.377: reference-only rebind to current Paper I-B v1.1; original mathematics retained except new Paper I-B comparison chapter.
% Framework v2.370: derived exceptional-root/Hurwitz comparison in I-A/I-C; other papers reference-only synchronized.
% Framework v2.368: preserve prior II-B convention/unital fixes and verified II-D corrections; current references synchronized.
% Framework v2.365: scoped Hurwitz/600-cell hadron-source integration; physical APBI/decay frontier remains open.
% Framework v2.358: current 30-root source reference synchronization; selectively updated III-E, III-F, VI-B, VI-C.
% Framework v2.356: bibliography/source-registry synchronization; mathematical body unchanged.
% Framework v2.348: Paper V audit, terminal observable scope, canonical Gram eigenvalues and exact weak-memory variance.
% Framework v2.321: Paper V Petersen rooted 3+12 induction, exact sector compression, conjugation orbit refinement and source-incidence firewall. Original source theorem bodies preserved.
% Framework v2.275: exact positive dense-twin word cubature, finite repeatability and eighth-moment classification, nested sampling frames, joint action and POVM-jet exports; preserves v2.274.
% Framework v2.274: two-scalar positive convolution certificates, first-leaked exceptional coset detection, branch-symmetric twin mixing, and finite endpoint moments; all v2.273 branches preserved.
% Framework v2.273: restored missing Problem 3 theorems from v2.268; exact exceptional frames, twin geometry, convolution repeatability and projector/occurrence firewalls; merged on v2.272 without overwriting independent branches.
% Framework v2.272: merged v2.270 covariant native-readout mathematics with v2.271 Problem 4 finite-action Hopf hypercharge/holonomy corrections; no universal bulk-gauge uniqueness claim.
% Independent branches: v2.270 readout/guard exports and v2.270 K16 action-consumed charge closure; all prior scoped assumptions retained.
% Framework v2.267: exact Paper II-A converse/observable certificates, III-F action tests, V-B right orbit and V-C native gauge scope; prior v2.266 preserved.
% Framework v2.266: exact Paper II-C outcome-count/readout converse results and scoped downstream imports; 30-root references synchronized.
% Framework v2.265: II-C coordinate-stage constitution and exact first-symbol host inheritance; prior mathematical domains retained.
% Framework v2.264: II-B reverse and action-bridge integration; preserves prior mathematical domains; active 30-root references synchronized.
% Framework v2.263: Paper I-D adds bridge-controlled root-carry parity and source-relative old-cap cancellation; other root mathematical bodies retained; active 30-root references synchronized.
% Framework v2.262: merges Paper II-F three-gap closure from v2.260 into v2.261, preserving newer II-B/II-D repairs; 30-root references synchronized.
% Framework v2.261: merged Paper II-B GP2-level/geometric repairs; preserves v2.260 II-D repairs; active references synchronized.
% Framework v2.260: Paper II-D signed RS, finite-symbol typing, and conditional curvature/Jordan repairs; references synchronized.
% Framework v2.259: conditional native SU(2) atlas and explicit later discharge; v2.258 mathematical bodies retained elsewhere.
% Framework v2.258: preserves parallel II-A/dynamics update and adds complete readout/frequency closures.
% Framework v2.257: exact minimal complete position readout and arbitrary-preparation frequency law.
% Framework v2.256: exact native inertia and sharp bounded ordered-source closure; preserves concurrent v2.255 work.
% Framework v2.255: preserves parallel detector integration and adds exact finite readout closures.
% Framework v2.254: shared bibliography/reference sync; substantive Pasterski detector-Ward additions only IV-B/IV-C/VII-C.
% Framework v2.253: typed raw coordinate defect in I-A and synchronized VI-A import; active references synchronized.
% Framework v2.252: optional shell-frame source-return calculation in III-G; active references synchronized.
% Framework v2.251: forward common-action execution, joint source jets and scoped coupled control.
% Framework v2.251: exact full-heat variations, fermion source jets and explicit physical completion frontier.
% Framework v2.251: certified interacting source jets, pole stability and scoped completion proofs.
% Framework v2.251: interacting finite-source evaluation and coupled local vacuum classification.
% Framework v2.251: exact finite source execution and native causal common-parent bridge.
% Framework v2.251: constructive parent and scoped source-incidence audit; see release report.
% Framework v2.251: common-action two-jet, corrected source/response closure and synchronized references.
% Framework v2.244: scoped rooted source-order Schur diagnostics and synchronized paper references.
% Framework v2.242: signed rooted phase lift, complete transverse return and exact family cubic.
% Framework v2.241: proved rooted conormal geometry, conditional cap realization and scoped action tests.
% Framework v2.240: observable-minimal action and response certificates integrated; inherited v2.239 geometry retained.
% Framework v2.234: technical proof companion to V-A, V-B and V-C.
\documentclass[11pt,a4paper]{article}
\usepackage[T1]{fontenc}
\usepackage[utf8]{inputenc}
\usepackage{lmodern}
\setlength{\emergencystretch}{3em}

\usepackage{mathtools}
\usepackage{graphicx}
\usepackage{float}
\usepackage{caption}
\usepackage{microtype}
\PassOptionsToPackage{hidelinks}{hyperref}

% Canonical series preamble.
% Paper-series v2.381: new I-B torsion/modular results, bibliography and current source roots.
% Paper-series v2.380: new exceptional bibliography and synchronized paper references.
% Paper-series v2.379 current 31-paper macro/aux registry.
% Paper-series v2.378: current 31-paper AUX/bibliography registry updated for Paper I-B v1.2 literature additions.
% Framework v2.371: current-source reference registry refreshed; old theorem labels retained.
% Framework v2.370: current 30-paper source/AUX root rebinding; native mathematical macros unchanged.
% Framework v2.368: synchronized merged 30-root registry; mathematical notation unchanged.
% Framework v2.358: active 30-paper source registry synchronized; existing mathematical macros unchanged.
% Framework v2.357: selective new source-owned certificates; active reference roots updated for five revised papers.
% Framework v2.356: active 30-root registry synchronized to the 30 papers in this release.
% Framework v2.275: exact positive dense-twin word cubature, finite repeatability and eighth-moment classification, nested sampling frames, joint action and POVM-jet exports; preserves v2.274.
% Framework v2.274: filename and root-source synchronization.
% Framework v2.273: restored missing Problem 3 theorems from v2.268; exact exceptional frames, twin geometry, convolution repeatability and projector/occurrence firewalls; merged on v2.272 without overwriting independent branches.
% Framework v2.272: merged v2.270 covariant native-readout mathematics with v2.271 Problem 4 finite-action Hopf hypercharge/holonomy corrections; no universal bulk-gauge uniqueness claim.
% Independent branches: v2.270 readout/guard exports and v2.270 K16 action-consumed charge closure; all prior scoped assumptions retained.
% Framework v2.269: synchronized 30 active AUX roots; mathematics and bibliography preserved from Macros-v2_161.
% Framework v2.266: Paper II-C converse-derived outcome resources; II-E, III-C, IV-A/B, V-A and VI-B downstream mathematical closures; 30-root registry synchronized.
% Framework v2.265: II-C coordinate-stage constitution and exact first-symbol host inheritance; prior mathematical domains retained.
% Framework v2.264: paper II-B reverse/audit and IV-A observable-minimal exchange; synchronized external references to 30 active roots.
% Framework v2.263: Paper I-C adds bridge-controlled root-carry parity and source-relative old-cap cancellation; other root mathematical bodies retained; active 30-root references synchronized.
% Framework v2.262: integrates Paper II-F radius, Casimir and readout results; active 30-root references synchronized.
% Framework v2.261: merged Paper II-B GP2-level/geometric repairs; preserves v2.260 II-D repairs; active references synchronized.
% Framework v2.260: Paper II-D signed RS, finite-symbol typing, and conditional curvature/Jordan repairs; references synchronized.
% Framework v2.259: conditional native SU(2) atlas and explicit later discharge; v2.258 mathematical bodies retained elsewhere.
% Framework v2.258: preserves parallel II-A/dynamics update and adds complete readout/frequency closures.
% Framework v2.257: exact minimal complete position readout and arbitrary-preparation frequency law.
% Framework v2.256: exact native inertia and sharp bounded ordered-source closure; preserves concurrent v2.255 work.
% Framework v2.255: preserves parallel detector integration and adds exact finite readout closures.
% Framework v2.254: Pasterski bibliography and synchronized 30-paper external-reference registry.
% Framework v2.253: typed raw coordinate defect in I-A and synchronized VI-A import; active references synchronized.
% Framework v2.252: synchronized source registry; optional III-G shell-frame calculation.
% Framework v2.251: forward common-action execution, joint source jets and scoped coupled control.
% Framework v2.251: exact full-heat variations, fermion source jets and explicit physical completion frontier.
% Framework v2.251: certified interacting source jets, pole stability and scoped completion proofs.
% Framework v2.251: synchronized active source registry; mathematical notation unchanged.
% Framework v2.251: synchronized active source registry; mathematical notation unchanged.
% Framework v2.251: synchronized active source registry; mathematical notation unchanged.
% Framework v2.251: synchronized source registry; no mathematical notation changes.
% Framework v2.244: synchronized 30-paper AUX registry and rooted phase-source action-order tests.
% Framework v2.242: signed rooted phase lift, complete transverse return and exact family cubic.
% Framework v2.241: proved rooted conormal geometry, conditional cap realization and scoped action tests.
% Framework v2.239: all 30 active source roots synchronized; no mathematical macro changes.
% Framework v2.238: Paper IV A/B/C split and 30-root reference registry.
% Framework v2.237: Paper VII-D technical split and 28-root reference registry.
% Framework v2.236: Paper VI-C technical-volume reference added; all current root bindings advanced.
% Framework v2.235: 26-root AUX and bibliography bindings synchronized; mathematical macros unchanged.
% BEGIN INLINE SOURCE: Paper00-TeXBuildCompatibility-v1_0.tex
% Rendering-only repairs for inherited long appendices and bitmap text fonts.
% Integer counters through 26 retain the original labels; later values use AA,
% AB, ... . This file changes no action coefficient, state or theorem premise.
\RequirePackage{alphalph}
\makeatletter
\let\BHPrimitiveAlph\@Alph
\newalphalph{\BHExtendedAlph}{\BHPrimitiveAlph}{26}
\let\@Alph\BHExtendedAlph
\AtBeginDocument{%
  \@ifpackageloaded{microtype}{\microtypesetup{expansion=false}}{}%
}
\makeatother

% END INLINE SOURCE: Paper00-TeXBuildCompatibility-v1_0.tex

% Framework v2.164: current reference/bibliography profile; mathematical definitions retained.
% Framework v2.162: finite metric, functional Hilbert and scoped Ward integration.
% Framework v2.161: inherits Macros-v2_57.tex exactly except the current external-reference binding.
% Macros v2.57 / framework v2.156: current reference registry only; mathematical definitions unchanged.
% Framework v2.143 / Macros v2.55: current external-reference bindings; mathematical notation unchanged.
% Framework v2.110 / Macros v2.54: current external-reference registry and bibliography; mathematical notation inherited.
% Framework v2.67 / Macros-v2_53: adds Born-action-cell and joint-EH readout symbols; absolute unit realization is metrological; MATCH_{J/scal}^{EH} remains the sole G24 normalization bridge.
% Framework v2.35 / Macros-v2_52: grouped paper numbering, active AUX registry, and Born-v6_20 bibliography binding synchronized; notation unchanged.
% Framework v2.26: active-source AUX and bibliography bindings synchronized; notation unchanged.
% Framework v2.09: reviewed early-reciprocity integration; scoped physical imports.
% Framework v2.08: Born--Euler mathematics and conditional complete-action potential criterion.
% Framework v2.07: active AUX registry includes Paper I-B v1.2; shared bibliography Born-v6_18.
% Macros-v2_46 / framework v2.04: mathematical atlas and downstream realization correction.
% Macros-v2_46 / framework v2.03: scoped structure integration and current dependency bindings.
% Macros-v2_46 / framework v2.01: shared source bindings; prior revision comments are historical.
% v2.42 / framework v1.99: Synchronized RG/AS owner and bibliography bindings.
% v2.41 / framework v1.99: active manuscript AUX bindings and geometric-closure bibliography synchronized.
% Release-v1.80 centralized source synchronization for scalar-circle/comparison-band integration.
% Release v1.79: Hurwitz hierarchy and kinetic-normalized action response integration; physical mass matching remains conditional.
% Release v1.78: approved M8 closure-batch integration; physical and analytic domains retained.
% Current v1.78 registry: all active TeX roots are synchronized below.
% Earlier release comments are historical scope records, not current status.
% Release-v1.77 synchronization: native same-owner Noether closure, generator--flow/history typing, analytic-certification firewalls, scalaron/K16, LF/QCD, quantum-Ward, and RAR boundary-owner transport are registered; shared mathematical notation is unchanged.
% Release-v1.76 synchronization: registry points to the v1.76 active paper roots.
% Release-v1.76 synchronization: registers the local A+Phi+g Noether/constraint closure and dynamics-frontier propagation; shared mathematical notation is unchanged.
% Release-v1.75 synchronization: registry points to the v1.75 active paper roots.
% Release-v1.75 synchronization: integrates the joint-Dirac Einstein bulk real form and the local odd-scalar heat-coefficient action on explicit analytic first-jet/interior domains; v1.74 Yang--Mills closure is retained. Fermionic reality, pure-metric curvature-squared/exact-nonlocal terms, noncanonical boundaries, physical vacuum selection, and actual solution-history tokens remain open.
% Release-v1.74 synchronization: registry points to the v1.74 active paper roots and retains the v1.73 theorem labels.
% Release-v1.74 synchronization: canonical Macros-v2_34.tex; generated local nonlinear Yang--Mills action, constrained connection-field phase space, source-free Gauss momentum map, and Legendre closure are integrated on their proved domains. The remaining non-gauge ACT_cl sectors and actual solution-history selection remain open.
% Release-v1.71 synchronization: closes the former JOINT_x/JOINT_5
% connection-occurrence gap by canonical structure-group extension of the
% nonflat Born connection, faithful associated representation, and the local
% cotangent--Clifford operator constructor.  Actual fluctuation members,
% actions, vacua, solutions, calibrations, and analytic QFT domains remain
% separately typed.  The active-AUX registry is advanced once.
% Release-v1.71 canonical macro and active-AUX registry.
% ======================================================================
% Macros-v2_39.tex -- canonical publication macro file
% Release v1.71 retains the v1.70 null-ray/QCD boundaries and additionally
% removes the duplicate structural connection actuality premise.  It does not
% claim global bundle topology, select a gauge/Higgs member, or promote the
% finite construction to an interacting continuum theory.
% Release v1.68 constructs the finite light-front total-fraction operators and
% parton number/momentum operator measures, proves their simplex and Hausdorff-limit
% theorems, and separates positive cut measures from scheme-renormalized
% light-ray distributions.  It also separates finite real Schur poles, stable
% norm-resolvent/Riesz pole-residue limits, and continued resonance
% pole/residue limits.  The Paper-VII-C scattering package is claim-indexed:
% local QFT and spectral claims import no S-matrix or infrared premise, while
% asymptotic amplitudes, LSZ, S-matrix poles, completeness, and cross sections
% import only their named massive or massless/long-range clauses.
% Release v1.67 factorizes BND48_H into a same-action Hessian/Schur
% descendant and still-open typed boundary-trace, critical-lift, and
% refinement-limit conditions.  It proves the finite-collar Dirichlet-to-
% Neumann law sqrt(y)tanh(L sqrt(y)), its nonuniform critical limit, and the
% finite-return ultra-infrared turnover.  The Paper-I-A Hardy constant mode is
% an exact geometric near-hit, not yet the physical metric boundary mode.
% The current local f(R) Hessian remains massive and therefore still refutes
% a RAR prediction on the minimal action branch.
% Release v1.66 derives the transverse depth-six normalizer bridge from the
% already generated C6 Reynolds image and Paper-I-A binary-dihedral normalizer.
% The bridge is the unique least 2D3-stable scalar-containing realization of
% the faithful quotient 2D3/C6=C2, has trace defect six and hence K=48, but is
% transverse to the nested active carrier.  Exact K=44 and K=52 contraction,
% semigroup, and energy-midpoint identities are registered as structural
% readouts rather than projector events.  The active E48 carrier gap remains
% unchanged.  A same-action boundary-symbol/Calderon intertwiner is isolated
% as the necessary and still unproved owner bridge to the static DtN sort;
% representation equivalence alone neither derives RAR nor creates a K=48
% action junction, pole, occupation, or cosmological transition.
% Release v1.65 derives the native E_48-centered carrier mass gap from the
% GP1--GP3-generated active spectrum, distinguishes its carrier-free interval
% from a particle located at E_48, and withdraws the unsupported promotion to
% a positive continuum Yang--Mills mass gap.  The latter remains a valid
% conditional transfer theorem on an enlarged continuum branch, but it is not
% an active framework prediction or an unfinished native-gap proof.
% Release v1.64 synchronizes Derivation Batches 14--22.  It registers the
% GP1--GP3 heat-profile orbit, same-Hessian electroweak response and pole
% ownership, K16/K48 action-response corrections, perturbation state and
% non-Gaussianity typing, finite-vacuum/gap/QCD limit repairs, factorized
% position-area ownership, the actual static metric response, and the ordered
% Wilson/string-tension limits.  Strict continuum convergence, positive
% physical gap, surviving higher cumulants, and positive tension remain open.
% Release v1.61 closes TIMEWIDTH16 as an independent direct-production datum.
% JUNC_{2|7} continues to own the finite intrinsic microscopic profile/width,
% while the actual same-action background/production history determines its
% physical-time pullback, support, and duration.  The current Paper-I-A--VII-C AUX
% registry is advanced once; shared mathematical notation is unchanged.
% Release v1.60 derives the final-octet primitive-remnant depletion, canonical
% Hopf-current representative, exact Hopf descent to dOmega/(4pi), and
% closed-FLRW a^{-3} neutral-current shape.  It retypes K=48 as the onset
% of the final one-unit geometric transfer and records cDM/RAR as distinct
% physical projections of that common neutral parent while preserving
% occurrence, Hilbert-owner, stress-matching, and static-kernel premises.
% Release v1.59 transports the v1.58 exceptional-shell/K16 closure through
% Papers III-F and VI-B--VII-C.  It replaces Paper VI-B's obsolete mixed-Hessian
% output typing by the third-variation pair vertex and Bogoliubov occupation
% support, derives the post-seed Higgs-pair feedback vertex and its exact
% alpha--delta kernel factors, propagates the unique K16/K48 action-junction
% split, identifies the QCD delta count with the upstream 2O shell atom only
% as a common-scalar numerical compatibility, and narrows THERM16 to dynamical
% approach/beta selection because the finite Gibbs/KMS target is already
% unique on a supplied simple observable algebra once its Hamiltonian and
% inverse temperature are supplied; direct-sum central weights remain typed.
% Release v1.58 registers the exceptional nearest-shell closure, the derived
% transverse K=16 action event, the unique interior action junction, the
% type-correct ordered pair vertex, the thin-layer squeeze law, and the
% radiation-like/matter-like kinematic split.  REG_0, transfer magnitude,
% physical layer width, and thermalization remain explicitly conditional.
% Release v1.57 registers the active-depth-six obstruction, the native
% K=48 carrier desert, the typed 7|9 owner reduction, and the dynamical
% reheating handoff across Papers III-B--VI-B.
% Release v1.56 registers the least-support n_2 inflation dressing of the
% persistent n_7 Laurent response, the analytic global-convexity and positive
% endpoint-measure proofs, the derived inflation core and K=16 coefficient
% freeze, and the no-particle dark-sector split into static RAR plus the
% conditional homogeneous effective-curvature response at reflected K=48.
% There is no n_6 action, scalaron-DM mass, or second Einstein--Hilbert term.
% Release v1.55 registers the protected-octet resolved/residual reflection
% q_sp+Z_W=8, its j<->8-j and K<->64-K scale covariance, and the derived
% rank-one-exit quadratic grade C_1^2/U=C_2 in Papers III-C, III-E, and VIII.
% The reflection is a typed trace/scale theorem, not a time reversal or a
% physical event-occurrence theorem; reheating transfer remains conditional.
% Release v1.54 registers the exact positive-part Laurent indices, the
% active-domain/terminal-endpoint action split, the no-second-EH firewall, and
% the q_DM=6/K_DM=48 scalaron readout typing across Papers III-D/E/F, VI,
% VIII, and IX.
% Release v1.53 registers the common nonnegative Laurent obstacle projection
% for n_1^+ and n_7^+, discharges the three former exit/freeze assumptions,
% and advances the Paper-III-D/E/F, VIII, and IX consumers.  Shared notation
% and bibliography ownership are unchanged.
% Release v1.51 registers the PQTI derivation batch and its updated
% Paper-III-A/C/D/E/F, VIII, and IX consumers.  Shared notation and
% bibliography ownership are unchanged.
% Release v1.50 advances the Paper-III-B AUX target to v1.40 after the
% rigorous audit repair; shared notation and bibliography ownership are unchanged.
% Release v1.49 advances the Paper-II-E/F AUX target to v3.14 after the
% minimal audit repair; shared notation and bibliography ownership are unchanged.
% Release v1.46 synchronizes Paper II-A's minimal frozen constructor-origin
% theorem and its five dependent AUX targets.  Shared notation is unchanged.
% Release v1.45 updates the Paper-II-B AUX target to v2.23 after the
% GP3-rigidity and appendix-placement repair; shared notation is unchanged.
% Release v1.44 advances the centralized Paper-I-A--VII-C AUX registry after the
% Paper-I-A v5.16 protected-completion scope repair.  Every current
% paper imports this one file; no paper-level opt-in or \externaldocument list
% is permitted.  The bibliography remains centralized here through Born-v6_13.
% Shared bibliography, theorem infrastructure, colors, notation alignment,
% and current Paper-I-A--VII-C external-reference registration.
% This file centralizes common imports, bibliography sorting, colors, theorem
% infrastructure, and the notation shared by the complete series.
% v2.9 (release v1.43) made the Paper-I-A--VII-C AUX registry automatic and
% centralized the bibliography and all external-reference ownership.
% v2.8 release-registry refresh (v1.33) registers the rooted generated-parent
% cascade and the POL/REG/free-law/11+4/Born-group closure revisions.  No
% shared mathematical notation or Paper-II-A fundamental principle is changed.
% v2.8 release-registry refresh (v1.32) registers the universal generated-
% continuation revisions in Papers II-A, III-A, III-C, IV, and VIII.  No
% shared mathematical notation or Paper-II-A fundamental principle is changed.
% v2.8 release-registry refresh (v1.30) registers the CRP4/projective-state
% source revisions and their direct registry-only consumers.  No shared
% mathematical notation or Paper-II-A axiom is changed.
% v2.7 release-registry refresh (v1.27) updates only the current external
% Paper-I-A--VII-C basenames after the placement/assumption synchronization; no
% shared mathematical notation or Paper-II-A axiom is changed.
% v2.7 synchronizes the geometric defect-repair revision: Paper-II-A v3.42,
% Paper-III-C v1.26, Paper-III-D v1.27, Paper-III-E v1.32, and Paper-III-F
% v1.34.  It registers the common trace-defect/spectral-filtration theorem,
% its owner-typed response pullback, and the reduced GRED and TERM interfaces.
% No shared mathematical notation is changed.
% v2.6 synchronizes Paper-III-E v1.31's cone-type repair, minimal-response
% classification, and explicit classical action-role boundary, together with
% Paper-III-F v1.32's downstream response-atlas synchronization.  No shared
% mathematical notation is changed.
% v2.5 synchronizes the Paper-III-D v1.26 faithful-trace-duality proof,
% exact ordered-owner joint locus, and typed alternative-presentation
% bookkeeping repair, together with its Paper-III-E and closure-frontier
% references.  No shared mathematical notation is changed.
% v2.4 synchronizes the Paper-III-C v1.25 orbit/process/owner repair, its
% downstream Papers III-B--F, V, VI, VIII, and IX revisions, and Born-v6_13.
% v2.3 synchronizes the Paper-II-A v3.41 witness/support/principal inheritance
% layer, the Paper-III-A v1.24 joint-realizability/readout firewall, all dependent
% host-rate revisions, and the current Paper-I-A--VII-B source registry.
% v2.2 synchronizes the generated active-realization, determinant-depth,
% dimension-weighted hierarchy, and typed cosmological-event interfaces.
% v2.1 synchronizes the exceptional-support interface revisions: Paper I-A
% v5.9, Paper II-A v3.30, Paper II-E/F v3.7, and Paper IV v4.16.
% v2.0 introduces the stage-safe G0 occurrence notation, synchronizes the
% active Paper-I-A--IV source registry (current rigorous Paper-III closure:
% III-A v1.17, III-B v1.22, III-C--E v1.17, III-F v1.14, IV v4.14), and makes
% external-reference self-skip depend on logical paper role rather than a
% versioned job name.
% ======================================================================

% ======================================================================
% macros.tex — Shared macros for the Born Geometry series
% ======================================================================

% ======================================================================
% 1. PACKAGES, BIBLIOGRAPHY, COLORS, AND UNIT DEFINITIONS
% ======================================================================
% Core packages shared across the complete Paper-I-A--VII-C series.  Loading them here makes
% the macro file a genuine series preamble while remaining compatible with existing
% paper preambles that load the same packages before \input.
\usepackage{amsmath}
\usepackage{amssymb}
\usepackage{amsthm}
\usepackage{mathrsfs}
\usepackage{array}
\usepackage{booktabs}
\usepackage{xcolor}
\makeatletter
\@ifpackageloaded{geometry}{}{\usepackage{geometry}\geometry{margin=1in}}
\makeatother
\usepackage{mdframed}

% Shared support packages required by the macros below and by most papers.
\usepackage{ragged2e}
\usepackage{longtable}
\usepackage{tabularx}
\usepackage{enumitem}
\usepackage[most]{tcolorbox}

% Series-wide numeric bibliography convention using the portable BibTeX path.
% All Papers I-A--VII-C use the ordinary \cite interface.
\usepackage[numbers,sort&compress]{natbib}
\providecommand{\printbibliography}{%
  \bibliographystyle{plainnat}%
  \bibliography{Born-v6_59}%
}

% Shared colors.  \providecolor avoids clobbering a paper-specific override.
\providecolor{postulatecolor}{RGB}{230,242,255}
\providecolor{compare}{RGB}{230,245,230}
\providecolor{definitioncolor}{RGB}{230,242,255}

\providecolor{openriskcolor}{RGB}{255,235,235}
\providecolor{bridgecolor}{RGB}{240,240,255}
\providecolor{gold}{RGB}{255,215,0}

% ======================================================================
% 2. CUSTOM ENVIRONMENTS
% ======================================================================
% Common table column types used in all papers.  Guarded so that a paper may
% define its own versions before loading this file.
\makeatletter
\@ifundefined{NC@find@L}{%
  \newcolumntype{L}[1]{>{\RaggedRight\arraybackslash}p{#1}}%
}{}
\@ifundefined{NC@find@C}{%
  \newcolumntype{C}[1]{>{\Centering\arraybackslash}p{#1}}%
}{}
\@ifundefined{NC@find@Y}{%
  \newcolumntype{Y}{>{\RaggedRight\arraybackslash}X}%
}{}
\makeatother

% ======================================================================
% 3. THEOREM STYLES AND ENVIRONMENTS
% ======================================================================
% Custom underlined theorem style (useful for remarks or notes)
\newtheoremstyle{underlined}
  {3pt}   % Space above
  {3pt}   % Space below
  {}      % Body font (empty means normal/roman text, typical for remarks)
  {}      % Indent amount
  {}      % Theorem head font (can be \itshape, \bfseries, etc.)
  {.}     % Punctuation after theorem head
  { }     % Space after theorem head
  {\underline{\thmname{#1}\thmnumber{ #2}}\thmnote{ (#3)}} % Theorem head formatting

% Standard plain theorem style environments.  Each environment is guarded so
% the shared macro file can be loaded after a paper-specific theorem setup.
\theoremstyle{plain}
\makeatletter
\@ifundefined{theorem}{\newtheorem{theorem}{Theorem}[section]}{}
\@ifundefined{lemma}{\newtheorem{lemma}{Lemma}[section]}{}
\@ifundefined{proposition}{\newtheorem{proposition}{Proposition}[section]}{}
\@ifundefined{corollary}{\newtheorem{corollary}{Corollary}[section]}{}
\makeatother

% Definition style environments (typically non-italicized body)
\theoremstyle{definition}
\makeatletter
\@ifundefined{definition}{\newtheorem{definition}{Definition}[section]}{}
\@ifundefined{postulate}{\newtheorem{postulate}{Postulate}[section]}{}
\@ifundefined{closure}{\newtheorem{closure}{Closure}[section]}{}
\makeatother

% Applying the custom underlined style
\theoremstyle{underlined}
\makeatletter
\@ifundefined{remark}{\newtheorem{remark}{Remark}[section]}{}
\@ifundefined{convention}{\newtheorem{convention}{Convention}[section]}{}
\@ifundefined{principle}{\newtheorem{principle}{Principle}[section]}{}
\makeatother

% Framed proof environment.
% Every \begin{proof} ... \end{proof} is automatically framed/shaded.
% This avoids mdframed/trivlist nesting, which can drop the top border after
% theorem-like environments.
\makeatletter
\renewenvironment{proof}[1][\proofname]{%
  \par\addvspace{6pt}%
  \noindent
  \begin{tcolorbox}[%
    enhanced,
    breakable,
    colback=black!4,
    colframe=black!35,
    boxrule=0.35pt,
    arc=1pt,
    left=6pt,
    right=6pt,
    top=4pt,
    bottom=4pt,
    boxsep=0pt,
    before skip=0pt,
    after skip=6pt,
    parbox=false
  ]%
  \pushQED{\qed}%
  \normalfont\itshape #1\@addpunct{.}\enspace
  \normalfont\ignorespaces
}{%
  \unskip\nobreak\hfill\popQED
  \end{tcolorbox}%
  \par\addvspace{2pt}%
}
\makeatother

\mdfdefinestyle{postulatebox}{
  backgroundcolor=yellow!10,
  linecolor=black,
  linewidth=0.4pt,
  skipabove=\baselineskip,
  skipbelow=\baselineskip,
  nobreak=true
}

% ======================================================================
% PROOF-READABILITY HELPERS
% ======================================================================

\newcommand{\ProofPlan}[1]{%
  \par\smallskip
  \noindent\textbf{Proof plan. }#1%
  \par\smallskip
}

\newlist{proofsteps}{enumerate}{2}
\setlist[proofsteps,1]{
  label=\textbf{(\alph*)},
  ref=\alph*,
  leftmargin=2.2em,
  itemsep=0.45\baselineskip,
  topsep=0.35\baselineskip
}
\setlist[proofsteps,2]{
  label=\textbf{(\roman*)},
  ref=\roman*,
  leftmargin=2.0em,
  itemsep=0.25\baselineskip,
  topsep=0.2\baselineskip
}


% ======================================================================
% 4. NUMBER SYSTEMS AND ALPHABETS
% ======================================================================
\providecommand{\RR}{\mathbb{R}} % Real numbers
\providecommand{\PP}{\mathbb{P}} % Projective space / Probability
\providecommand{\CC}{\mathbb{C}} % Complex numbers
\providecommand{\ZZ}{\mathbb{Z}} % Integers
\providecommand{\NN}{\mathbb{N}} % Natural numbers
\providecommand{\HH}{\mathbb{H}} % Quaternions / Hyperbolic space
\providecommand{\LFRayGeom}{\mathsf{LF\!\mbox{-}RAY\!\mbox{-}GEOM}}
\providecommand{\LFRep}{\mathsf{LF\!\mbox{-}REP}}
\providecommand{\LFState}{\mathsf{LF\!\mbox{-}STATE}}
\providecommand{\LFLim}{\mathsf{LF\!\mbox{-}LIM}}
\providecommand{\LFScat}{\mathsf{LF\!\mbox{-}SCAT}}
% Number systems / projective spaces
\providecommand{\CP}{\mathbb{CP}}
\providecommand{\RP}{\mathbb{RP}}

% ======================================================================
% 5. GROUPS AND LIE THEORY
% ======================================================================
\providecommand{\SO}{\ensuremath{\mathrm{SO}}}     % Special Orthogonal group
\providecommand{\SU}{\ensuremath{\mathrm{SU}}}     % Special Unitary group
\providecommand{\Spin}{\ensuremath{\mathrm{Spin}}} % Spin group
\providecommand{\CS}{\ensuremath{\mathrm{CS}}}     % Chern-Simons (capitalized)
\providecommand{\cs}{\ensuremath{\mathrm{cs}}}     % Chern-Simons (lowercase)

% ======================================================================
% 6. COMMON MANIFOLDS AND TOPOLOGY
% ======================================================================
\providecommand{\Sone}{S^1}                 % Standard circle

% ======================================================================
% 7. MATH OPERATORS (Declared via \DeclareMathOperator)
% ======================================================================
\DeclareMathOperator{\End}{End}   % Endomorphisms
\DeclareMathOperator{\Aut}{Aut}   % Automorphisms
\DeclareMathOperator{\spec}{spec} % Lowercase spectrum
\DeclareMathOperator{\tr}{tr}     % Lowercase trace (if you already have \Tr)
\DeclareMathOperator{\supp}{supp} % Support
\DeclareMathOperator{\TrA}{Tr_{A}} % Partial trace over subsystem A
\DeclareMathOperator{\TrB}{Tr_{B}} % Partial trace over subsystem B
% Lie/algebra operators
\DeclareMathOperator{\Ad}{Ad}
\DeclareMathOperator{\Sym}{Sym}
\DeclareMathOperator{\su}{su}
% ======================================================================
% 8. STANDARD OPERATOR-LIKE SYMBOLS (Safe definitions via \providecommand)
% ======================================================================
\providecommand{\vol}{\operatorname{vol}}   % Volume form
\providecommand{\Vol}{\operatorname{Vol}}   % Total volume
\providecommand{\Ric}{\operatorname{Ric}}   % Ricci tensor
\providecommand{\Rm}{\operatorname{Rm}}     % Riemann tensor
\providecommand{\Tr}{\operatorname{Tr}}     % Uppercase trace
\providecommand{\Spec}{\operatorname{Spec}} % Uppercase spectrum
\providecommand{\diag}{\operatorname{diag}} % Diagonal matrix
\providecommand{\rank}{\operatorname{rank}} % Rank
\providecommand{\scal}{\mathcal R} % Scalar curvature
\providecommand{\Cl}{\operatorname{Cl}}     % Clifford algebra
\providecommand{\diam}{\operatorname{diam}} % Diameter

% ======================================================================
% 9. FUNCTIONS, TRIGONOMETRY, AND CALCULUS
% ======================================================================

% Codifferential naming conventions
% (Avoid \ddag: commonly defined by physics.sty and symbol packages)
\providecommand{\codiff}{\mathrm{d}^{\dagger}} % Codifferential operator

% Quantum mechanics bra-ket notation (needed for basis tables)
\providecommand{\ket}[1]{\lvert #1 \rangle}
\providecommand{\bra}[1]{\langle #1 \rvert}
\providecommand{\braket}[2]{\langle #1 \mid #2 \rangle}

% ======================================================================
% 10. PHYSICAL SCALES AND CONSTANTS (Paper 1c / Born Geometry)
% ======================================================================
\providecommand{\Rp}{\ensuremath{R_{\mathrm P}}}     % matched common EH curvature scale (not defined from observed G)
\providecommand{\hBorn}{\ensuremath{h_{\mathrm B}}} % physical representative 2\pi\hbar of one primitive Born scalar action cell
\providecommand{\UEHjoint}{\ensuremath{U_{J}}} % joint-Dirac EH curvature-coordinate scale 8 kappa_J
\providecommand{\Gread}{\ensuremath{G_{\rm scal}}} % scalar-action coordinate c_op^3/(h_B U); physical matching explicit
\providecommand{\GreadJoint}{\ensuremath{G_{J,\mathrm{raw}}}} % raw ordinary-trace coordinate = G_scal/96
\providecommand{\Rds}{R_{\mathrm{dS}}}    % de Sitter Curvature
\providecommand{\Mpl}{M_{\mathrm{Pl}}}    % Planck mass
\providecommand{\lp}{\ell_{\mathrm{Pl}}}  % Planck length
\providecommand{\Pmax}{P_{\mathrm{max}}}  % Maximum momentum
\providecommand{\Mplbar}{\bar{M}_{\mathrm{Pl}}} % Reduced Planck mass
\providecommand{\eff}{\mathrm{eff}}       % Subscript for effective quantities
\providecommand{\Cosmo}{\ensuremath{\Lambda}}                % Cosmological constant (use \Cosmo in source)
\providecommand{\LCDM}{\ensuremath{\Cosmo\mathrm{CDM}}}  % \LambdaCDM (typesets with \Cosmo)
\providecommand{\Lcut}{\ensuremath{\Lambda_{\mathrm{cut}}}}  % Cutoff scale (brace when used as script)
\providecommand{\Rcs}{\ensuremath{R_{\mathrm{cs}}}}      % Configuration-space radius
\providecommand{\CosmoEff}{\ensuremath{\Cosmo_{\eff}}}   % Effective cosmological constant

% Disambiguation: \Rmom is the momentum-space (S^3) radius. 
% Reserve plain R for the 4D curvature scalar in f(R).
\providecommand{\Rmom}{\ensuremath{R_{\mathrm{mom}}}} 
\providecommand{\Punit}{P_{0}} % Momentum unit set by \Rmom (= \hbar/\Rmom)
\providecommand{\Rvec}{\ensuremath{R_{\mathrm{vec}}}} % invariant vector-space radius in the reduced round sector
% Specific numeric values and expressions
\providecommand{\Vlogcap}{L_{\mathrm{Hopf}}^{(1)}}
\providecommand{\Hopf}{\operatorname{Hopf}}
% Scalar phase (S^1) sector + universal cover notation (Paper I-A)
\providecommand{\tper}{T_\star}      % Time period/scale
\providecommand{\Eper}{E_\star}      % Energy period/scale
\providecommand{\ttilde}{\tilde t}   % Universal cover time
\providecommand{\Etilde}{\tilde E}   % Universal cover energy

% ======================================================================
% 10A. SERIES-WIDE LOGICAL / NOTATIONAL DISAMBIGUATIONS
% ======================================================================
% Use these symbols to avoid collisions among: (i) the Paper--III
% no-insertion/geometric-closure class, (ii) Paper--I-A trimer diagonals,
% (iii) Paper--V-A/V-C family-space diagonals, (iv) scalar-index running,
% and (v) protected shell-clock labels.
\providecommand{\Rprim}{\ensuremath{\mathsf R_{\mathrm{prim}}^{\Sigma}}} % primitive finite structural readouts
\providecommand{\Rloop}{\ensuremath{\mathsf R_{\circ}^{\Sigma}}} % closed-loop readout role
\providecommand{\Rbdry}{\ensuremath{\mathsf R_{\partial}^{\Sigma}}} % boundary-conditioned readout role
\providecommand{\Rgen}{\ensuremath{\mathsf R_{\mathrm{fin}}^{\Sigma}}} % finite generated closure of primitive roles
\providecommand{\PhysOcc}{\operatorname{PhysOcc}}
\providecommand{\Actual}{\operatorname{Actual}}
\providecommand{\GeomLic}{\operatorname{GeomLic}}
\providecommand{\InvDist}{\operatorname{InvDist}}
\providecommand{\Dcl}{\ensuremath{\mathsf{D}_1}}      % unmarked positive Dirac-even closed spectral-readout form class
\providecommand{\Dheat}{\ensuremath{\mathsf D_{\mathrm{heat}}}} % one-parameter heat-semigroup subclass
\providecommand{\DtriI}{\ensuremath{D^{\mathrm{tri}}_1}} % Paper--I-A cyclic-trimer diagonal
\providecommand{\DtriII}{\ensuremath{D^{\mathrm{tri}}_2}}
\providecommand{\DfamI}{\ensuremath{D^{\mathrm{fam}}_1}} % Family-space diagonal selector
\providecommand{\DfamII}{\ensuremath{D^{\mathrm{fam}}_2}}
\providecommand{\alpharun}{\ensuremath{\alpha_{\mathrm{run}}}} % scalar spectral-index running
\providecommand{\Ksh}{\ensuremath{K_{\mathrm{sh}}}} % protected shell-clock coordinate
\providecommand{\ddef}{\ensuremath{\delta_{\mathrm{def}}}} % universal defect scalar after Paper--III realization

% ======================================================================
% 10B. SERIES-WIDE LOGICAL STATUS TAGS
% ======================================================================
% Upright small-caps wrapper for audit/status labels.  It prevents LaTeX from
% trying to use unavailable small-caps italic font shapes inside theorem bodies
% such as TU/lmr/m/scit.
\providecommand{\AuditLabel}[1]{}
\renewcommand{\AuditLabel}[1]{{\normalfont\scshape #1}}

% These tags are bookkeeping labels for consistency audits. They do not
% create new mathematical hypotheses; they record how a datum is being used
% in the proof chain.
\providecommand{\StatusPrimitive}{}
\renewcommand{\StatusPrimitive}{\AuditLabel{Primitive}}
\providecommand{\StatusDerived}{}
\renewcommand{\StatusDerived}{\AuditLabel{Derived}}
\providecommand{\StatusImported}{}
\renewcommand{\StatusImported}{\AuditLabel{Imported}}
\providecommand{\StatusConditional}{}
\renewcommand{\StatusConditional}{\AuditLabel{Conditional}}
\providecommand{\StatusBridge}{}
\renewcommand{\StatusBridge}{\AuditLabel{Bridge-derived}}
\providecommand{\StatusComparison}{}
\renewcommand{\StatusComparison}{\AuditLabel{Comparison-only}}

% ======================================================================
% 10C. SERIES-WIDE CLAIM-WORDING AUDIT LABELS
% ======================================================================
% These are textual audit labels used to keep modal language consistent
% across the series.  They introduce no new mathematical assumptions.
\providecommand{\ClaimForced}{}
\renewcommand{\ClaimForced}{\AuditLabel{forced}}
\providecommand{\ClaimUnique}{}
\renewcommand{\ClaimUnique}{\AuditLabel{unique}}
\providecommand{\ClaimCanonical}{}
\renewcommand{\ClaimCanonical}{\AuditLabel{canonical}}
\providecommand{\ClaimMinimal}{}
\renewcommand{\ClaimMinimal}{\AuditLabel{minimal}}
\providecommand{\ClaimClosed}{}
\renewcommand{\ClaimClosed}{\AuditLabel{closed}}
\providecommand{\ClaimPrediction}{}
\renewcommand{\ClaimPrediction}{\AuditLabel{prediction}}

% ======================================================================
% 10D. SERIES-WIDE THEOREM-CHAIN AUDIT LABELS
% ======================================================================
% These labels mark the direction of load-bearing theorem chains. They are
% audit annotations only; they introduce no new mathematical assumptions.
\providecommand{\ChainInput}{}
\renewcommand{\ChainInput}{\AuditLabel{Chain input}}
\providecommand{\ChainOutput}{}
\renewcommand{\ChainOutput}{\AuditLabel{Chain output}}
\providecommand{\ChainGuard}{}
\renewcommand{\ChainGuard}{\AuditLabel{Chain guard}}
\providecommand{\ChainNoBackflow}{}
\renewcommand{\ChainNoBackflow}{\AuditLabel{No backflow}}

% ======================================================================
% 10E. SERIES-WIDE LOCAL THEOREM-AUDIT LABELS
% ======================================================================
% These labels mark a statement-level dependency audit around individual
% theorem/definition nodes. They are bookkeeping labels only.
\providecommand{\LocalAudit}{}
\renewcommand{\LocalAudit}{\AuditLabel{Local audit}}
\providecommand{\UsesOnly}{}
\renewcommand{\UsesOnly}{\AuditLabel{Uses only}}
\providecommand{\Exports}{}
\renewcommand{\Exports}{\AuditLabel{Exports}}
\providecommand{\DoesNotUse}{}
\renewcommand{\DoesNotUse}{\AuditLabel{Does not use}}

% ======================================================================
% 10F. SERIES-WIDE NUMERICAL-ROW AUDIT LABELS
% ======================================================================
% These labels classify individual rows in numerical, benchmark, status,
% and comparison tables.  They are audit annotations only and do not change
% any theorem, value, formula, or comparison number.
\providecommand{\RowAudit}{}
\renewcommand{\RowAudit}{\AuditLabel{Numerical-row audit}}
\providecommand{\RowTheorem}{}
\renewcommand{\RowTheorem}{\AuditLabel{Theorem-fixed row}}
\providecommand{\RowBridge}{}
\renewcommand{\RowBridge}{\AuditLabel{Bridge-fixed row}}
\providecommand{\RowRepresentative}{}
\renewcommand{\RowRepresentative}{\AuditLabel{Representative row}}
\providecommand{\RowComparison}{}
\renewcommand{\RowComparison}{\AuditLabel{Comparison row}}
\providecommand{\RowCalibrated}{}
\renewcommand{\RowCalibrated}{\AuditLabel{Calibrated row}}
\providecommand{\RowEstimate}{}
\renewcommand{\RowEstimate}{\AuditLabel{Estimate row}}
\providecommand{\RowOpen}{}
\renewcommand{\RowOpen}{\AuditLabel{Open row}}
\providecommand{\RowNoSelector}{}
\renewcommand{\RowNoSelector}{\AuditLabel{No selector}}

% ======================================================================
% 10G. TABLE HELPERS
% ======================================================================
% Table helpers used by Paper V-C.  The unused readout-label preset has been
% removed; only active table helpers remain.
\providecommand{\gpsection}[1]{%
  \addlinespace[0.55em]
  \multicolumn{3}{@{}l}{\textsc{#1}}\\[-0.25em]
  \midrule
}
\providecommand{\mathcell}[1]{%
  \(\begin{gathered}#1\end{gathered}\)%
}

\newcommand{\DerBy}[1]{%
  \begin{tabular}[t]{@{}l@{}}#1\end{tabular}%
}

\newcommand{\TableNotes}[1]{%
  \par\medskip
  \begin{minipage}{0.94\linewidth}
  \footnotesize\emph{Notes.} #1
  \end{minipage}\par
}
% ======================================================================
% 10H. PAPER VI-A / QCD NOTATION
% ======================================================================
% Shared symbols used by Paper VI-A and downstream QCD/Yang--Mills handoffs.
\providecommand{\Pqcd}{\mathfrak P_{\rm VIII}^{\rm QCD}}
\providecommand{\Iqcd}{\mathfrak I_{\rm QCD}^{(7)}}
\providecommand{\muQCD}{\mu_{\rm QCD}^{\rm int}}
\providecommand{\LambdaQCD}{\Lambda_{\rm QCD}}
\providecommand{\LambdaBH}{\Lambda_{\rm BH}^{\rm QCD}}
\providecommand{\HQCD}{\mathcal H_{\rm QCD}^{\rm BH}}
\providecommand{\HQCDfin}{H_{\rm QCD}^{\rm fin}}
\providecommand{\Hraw}{\mathcal H_{\rm raw}}
\providecommand{\HphysQCD}{\mathcal H_{\rm phys}^{\rm QCD}}
\providecommand{\Vcol}{V_C}
\providecommand{\Ps}{P_{\rm sing}}
\providecommand{\PiOct}{\Pi_8}
\providecommand{\MSbar}{\overline{\rm MS}}
\providecommand{\YM}{\mathrm{YM}}
\providecommand{\Ktri}{K_{\triangle}}
\providecommand{\KYM}{K_{\YM}^{\rm BH}}
\providecommand{\aBH}{a_{\rm BH}}
\providecommand{\epsBH}{\varepsilon_{\rm BH}}
\newcommand{\RowConstructive}{\textnormal{\textsc{Constructive}}}
\newcommand{\RowObligation}{\textnormal{\textsc{Bridge obligation}}}
\newcommand{\RowTarget}{\textnormal{\textsc{Target row}}}
\newcommand{\RowValidation}{\textnormal{\textsc{Validation row}}}
\providecommand{\slashed}[1]{\not{#1}}

% ======================================================================
% 10I. PAPER VI-B / CONTINUOUS-QFT NOTATION
% ======================================================================
% Shared symbols used by Paper VI-B and any downstream continuous-QFT or
% Yang--Mills handoff paper.  All definitions are guarded so individual
% papers may override them locally before loading this file if needed.
\providecommand{\A}{\mathfrak A}
\providecommand{\Hphys}{\mathcal H_{\rm phys}}
\providecommand{\QCD}{\mathrm{QCD}}
\providecommand{\BH}{\mathrm{BH}}
\providecommand{\loc}{\mathrm{loc}}

\providecommand{\RBH}{\mathcal R_{\rm BH}}
\providecommand{\DG}{\mathfrak D_G}
\providecommand{\PW}{\mathcal{P\!W}} % Peter--Weyl spectral-symbol subspace
\providecommand{\Toe}{\operatorname{Toe}} % Toeplitz compression/representative
\providecommand{\Hol}{\operatorname{Hol}} % Holonomy

% ======================================================================
% 11. LATE-LOADED PACKAGES (Must be at the end of the preamble)
% ======================================================================
% Load late packages only if a paper has not already loaded them, avoiding
% option clashes in paper-specific preambles.
\makeatletter
\@ifpackageloaded{hyperref}{}{\usepackage{hyperref}}
\makeatother
% --- Series-wide disambiguations (Pass 3) ---
\providecommand{\diracdelta}{\delta_{\rm D}} % Dirac delta distribution (delta reserved for hierarchy)
\providecommand{\Hg}{H_{\mathrm g}} % Hodge harmonic projector (keep H for Hubble)
\providecommand{\Gr}{G_{\mathrm r}} % Green operator (keep G for Newton)
\providecommand{\FD}{\mathrm{D}} % functional/variational derivative symbol (delta reserved)
\providecommand{\Id}{\mathbf{1}} % identity/Kronecker tensor (avoid \delta_{ij})


% BEGIN INLINE SOURCE: Paper00-ActiveExternalReferences-v2_130.tex
% Framework v2.191: current reference binding; existing definitions retained.
% Framework v2.190: current22-root registry with native clock-evolution proofs.
% Framework v2.189: exact current22-root external registry, including retained v2.188 hardening owners.
% Framework v2.187: current registry synchronization only; source release v2.186 preserved.
% Framework v2.186: authoritative current 22-paper registry.
% Framework v2.182: all current 22 logical paper AUX bindings.
% Framework v2.170: current I-B v1.43 external AUX binding.
% Framework v2.167: exact current 22-root registry after CP--BL2 integration.
% Framework v2.166: exact current 22-root compact-release AUX registry.
% Framework v2.165: authoritative current root bindings.
% Framework v2.164: exact current 22 manuscript AUX bindings.
% Framework v2.162: finite metric, functional Hilbert and scoped Ward integration.
% Framework v2.161: current 22 manuscript AUX bindings; source only.
% Framework v2.156 / registry v2.113: synchronized current 22 manuscript roots; source-only registry.
% Framework v2.143 / registry v2.111: active manuscript AUX basenames; no generated AUX or PDF included.
% ======================================================================
% Central external-reference imports, numbers only, no PDF jumping
% ======================================================================

% xr-hyper is loaded centrally so that papers may retain their existing
% \usepackage{xr-hyper} lines without redefining \externaldocument after this
% registry is installed.  If plain xr was loaded earlier, retain it.
\makeatletter
\@ifpackageloaded{xr-hyper}{}{%
  \@ifpackageloaded{xr}{}{\usepackage{xr-hyper}}%
}
\makeatother

% Make external-document registration idempotent.  The registry identifies a
% logical paper by its Paper<number><optional letter> prefix.  Prefix-qualified
% imports remain independent.  Current paper roots contain no local
% \externaldocument calls; this guard also makes accidental duplicates harmless.
\ExplSyntaxOn
\prop_new:N \g_bh_external_refs_prop
\cs_new_eq:NN \BHOriginalExternalDocument \externaldocument

\cs_new_protected:Npn \bh_external_document:nnnn #1#2#3#4
 {
  \tl_set:Nn \l_tmpa_tl {#3}
  \regex_replace_once:nnN
   { \A (?: xref/ )? (Paper[0-9]+[A-F]?) .* \Z }
   { \1 }
   \l_tmpa_tl
  \tl_set:Nx \l_tmpb_tl { #1 | \tl_use:N \l_tmpa_tl }
  \prop_if_in:NVTF \g_bh_external_refs_prop \l_tmpb_tl
   {
    \PackageInfo{BornHopf}
     {Skipping~duplicate~external~refs~for~#3}
   }
   {
    \prop_gput:NVn \g_bh_external_refs_prop \l_tmpb_tl {#3}
    \IfNoValueTF{#4}
     { \BHOriginalExternalDocument[#1][#2]{#3} }
     { \BHOriginalExternalDocument[#1][#2]{#3}[#4] }
   }
 }

\RenewDocumentCommand{\externaldocument}{O{} O{} m o}
 { \bh_external_document:nnnn {#1}{#2}{#3}{#4} }

\NewDocumentCommand{\BHExternalRefs}{m}
 {
  % Do not import the aux file of the logical paper currently being compiled,
  % even when the job and registry use different version suffixes.
  \tl_set:Nx \l_tmpa_tl { \c_sys_jobname_str }
  \regex_replace_once:nnN
   { \A (Paper[0-9]+[A-F]?) .* \Z }
   { \1 }
   \l_tmpa_tl
  \tl_set:Nn \l_tmpb_tl { #1 }
  \regex_replace_once:nnN
   { \A (Paper[0-9]+[A-F]?) .* \Z }
   { \1 }
   \l_tmpb_tl
  \str_if_eq:eeTF { \tl_use:N \l_tmpa_tl } { \tl_use:N \l_tmpb_tl }
   {
    \PackageInfo{BornHopf}
     {Skipping~external~refs~for~#1~because~this~is~the~current~document}
   }
   {
    \IfFileExists{xref/#1.aux}
     { \externaldocument[][nocite]{xref/#1}[#1.pdf] }
     {
      \IfFileExists{#1.aux}
       { \externaldocument[][nocite]{#1} }
       {
        \PackageInfo{BornHopf}
         {External~aux~file~#1.aux~not~found;~no~labels~imported}
       }
     }
   }
 }
\ExplSyntaxOff

% External references are automatic.  This file is the sole owner of both the
% bibliography and the current Paper-I-A--VII-D AUX registry (30 roots including I-C).  Missing auxiliaries
% are informational on the first pass and are populated by the source-only
% two-pass series build.  Basenames deliberately omit temporary download/copy
% suffixes such as ``(1)'' or ``-1''.
  % Paper I-A: intrinsic S3 and exceptional group introduction.
  \BHExternalRefs{Paper1A-Geometry-v5_175}

  % Paper I-B: complete exceptional geometry and finite spectral structures.
  \BHExternalRefs{Paper1B-ExceptionalGeometry-v1_5}
  % Paper I-C: reciprocal product M8 geometry.
  \BHExternalRefs{Paper1D-TechnicalMaterial-v1_65}
  % Paper I-D: terminal technical proofs and calculations.
  \BHExternalRefs{Paper1C-M8Geometry-v1_137}

\ifdefined\BHMathOnly
  % Papers I-A and I-B import only the mathematical atlas labels.
\else
  % Paper II: five physical-mathematical main papers plus one technical proof volume.
  \BHExternalRefs{Paper2A-GeometricConstitution-v3_221}
  \BHExternalRefs{Paper2B-MomentumCarrier-v2_144}
  \BHExternalRefs{Paper2C-ReciprocalLocalization-v2_169}
  \BHExternalRefs{Paper2D-SymmetrySpectra-v2_141}
  % The former II-F export labels remain in II-E; current II-F is the technical volume.
  \BHExternalRefs{Paper2E-KinematicExports-v3_138}
  % Paper II-F: complete technical proof and calculation owner for II-A--II-E.
  \BHExternalRefs{Paper2F-TechnicalMaterial-v1_67}

  % Paper III: Dynamics and Intrinsic Geometric Constitution (six components).
  \BHExternalRefs{Paper3A-GeneratedDynamics-v1_176}
  \BHExternalRefs{Paper3B-SpectralAction-v1_187}
  \BHExternalRefs{Paper3C-ProtectedHierarchy-v1_187}
  \BHExternalRefs{Paper3D-RunningResponse-v1_165}
  \BHExternalRefs{Paper3E-GravityAction-v1_197}
  \BHExternalRefs{Paper3F-CompleteDynamics-v1_263}
  % Paper III-G: complete preserved technical proof owner.
  \BHExternalRefs{Paper3G-TechnicalMaterial-v1_53}

  % Paper IV-A: generated finite quantum mechanics.
  \BHExternalRefs{Paper4A-FiniteQuantumMechanics-v4_214}
  % Paper IV-B: complete finite quantum processes and physical response.
  \BHExternalRefs{Paper4B-QuantumResponse-v1_49}
  % Paper IV-C: extended proofs and finite technical calculations.
  \BHExternalRefs{Paper4C-TechnicalMaterial-v1_49}

  % Paper V-A: The Standard Model.
  \BHExternalRefs{Paper5A-FiniteMatterGeometry-v3_209}

  % Paper V-B: Finite Spectrum and Flavor Closure.
  \BHExternalRefs{Paper5B-TerminalSpectrum-v5_199}

  % Paper V-C: Finite Interaction Boundaries of the Standard Model.
  \BHExternalRefs{Paper5C-InteractionResponse-v2_177}
  % Paper V-D: preserved technical proof and calculation owner.
  \BHExternalRefs{Paper5D-TechnicalMaterial-v1_53}

  % Paper VI-A: Gravity and Closed Cosmology.
  \BHExternalRefs{Paper6A-GravityCosmology-v2_260}

  % Paper VI-B: Geometric Dark Sector and Cosmological Readouts.
  \BHExternalRefs{Paper6B-CosmologyReadouts-v1_206}

  % Paper VI-C: complete technical proof and comparison owner.
  \BHExternalRefs{Paper6C-TechnicalMaterial-v1_50}

  % Paper VII-A: Finite QCD Kernel and Nonperturbative Closure.
\BHExternalRefs{Paper7A-QCDKernel-v1_172}

  % Paper VII-B: Hadronic, Thermal, and Nuclear QCD Readouts.
  \BHExternalRefs{Paper7B-QCDReadouts-v1_157}

  % Paper VII-C: Continuous Quantum Field Theory.
  \BHExternalRefs{Paper7C-ContinuousQFT-v2_202}

  % Paper VII-D: complete technical material for VII-A through VII-C.
  \BHExternalRefs{Paper7D-TechnicalMaterial-v1_53}

\fi

% END INLINE SOURCE: Paper00-ActiveExternalReferences-v2_130.tex


\newcolumntype{R}[1]{>{\raggedleft\arraybackslash}p{#1}}


% Combined local notation, with V-C-first priority on shared names.
\newcommand{\PaperSevenRevision}{v2.106 gauge variation and source normalization}
\newcommand{\Hkin}{\mathsf{H}_{\rm kin}}
\newcommand{\Hqft}{\mathsf{H}_{\rm QFT}}
\newcommand{\Hem}{\mathsf{H}_{\rm EM}}
\newcommand{\Hvac}{\mathsf{H}_{\rm vac}}
\newcommand{\Hcp}{\mathsf{H}_{\rm CP}}
\newcommand{\Hnu}{\mathsf{H}_{\nu}}
\newcommand{\Hgauge}{\mathsf{H}_{\rm gauge}}
\newcommand{\Aem}{\mathcal A_{\rm EM}}
\newcommand{\YN}{Y_\nu}
\newcommand{\Gfive}{\mathsf G_5^{\rm fin}}
\newcommand{\Rrep}{\mathsf R_5^{\rm rep}}
\newcommand{\Rconn}{\mathsf J_5^{\rm conn,der}}
\newcommand{\Rssqft}{\mathsf R_5^{\rm SSQFT}}
\newcommand{\JOINTfiveocc}{\mathsf{JOINT}_{5}^{\rm der}}
\newcommand{\CALgauge}{\mathsf{CAL}_{\rm gauge}}
\newcommand{\REGsig}{\mathsf{REG}_{\sigma}}
\newcommand{\REGzero}{\mathsf{REG}_{0}}
\newcommand{\Ustr}{\mathcal U_6^{\rm str}}
\newcommand{\Urole}{\mathcal U_6^{\rm role}}
\newcommand{\Uratio}{\mathcal U_6^{\rm ratio}}
\newcommand{\Uabs}{\mathcal U_6^{\rm abs}}
\newcommand{\Uroleratio}{\mathcal U_6^{\rm role,ratio}}
\newcommand{\Uroleabs}{\mathcal U_6^{\rm role,abs}}
\newcommand{\Istr}{\mathfrak I_{5C}^{\rm str}}
\newcommand{\Irole}{\mathfrak I_{5C}^{\rm role}}
\newcommand{\Brawsix}{\mathsf B_{6}^{\rm hist}}
\newcommand{\Arawsix}{\mathfrak A_{6,{\rm fin}}^{\rm hist}}
\newcommand{\Acompsix}{\mathfrak A_{6,{\rm fin}}^{\rm comp}}
\newcommand{\CMPexp}{\mathsf{CMP}_{\rm exp}}
\newcommand{\Hresopt}{\mathsf H_{\rm res}^{\rm opt}}
\newcommand{\Hhopfopt}{\mathsf H_{\rm Hopf}^{\rm opt}}
\newcommand{\PaperSixRevision}{v5.115 neutral geometry and source-scope integration}
\newcommand{\CoreSix}{\mathsf C_6}
\newcommand{\BridgeSix}{\mathsf B_6}
\newcommand{\BridgeLive}{\mathsf B_6^{\rm live}}
\newcommand{\BridgeRaw}{\mathsf B_6^{\rm hist}}
\newcommand{\BridgeObs}{\mathsf B_6^{\rm obs}}
\newcommand{\CLKrem}{\mathsf{CLK}_{\rm rem}}
\newcommand{\Ubase}{\mathcal U_6^{\rm str}}
\newcommand{\UU}{\mathrm U}
\newcommand{\Ran}{\operatorname{Ran}}
\newcommand{\HomG}{\operatorname{Hom}}
\newcommand{\sgnline}{\mathbf 1_{\rm sgn}}
\newcommand{\etaH}{\eta_{\rm H}}
\newcommand{\PVIStatus}[1]{\textnormal{\textsc{#1}}}
\newcommand{\pvcite}[1]{\cite{#1}}
\title{\textbf{Born--Hopf Geometry in a Closed Universe V-D: Technical Material}\\[0.45em]
\large Exact representation, spectrum, source and interaction calculations}
\author{Peter Kugelmann}
\date{Version 1.57 --- 19 November 2026 \\ Framework v2.379}
\begin{document}
% Publication line-breaking flexibility for technical prose (no math scaling).
\emergencystretch=2em
\maketitle
\begin{abstract}
Paper V-D preserves the full proofs and calculations moved from Papers V-A, V-B, and V-C. The original theorem, proof, equation, and label ownership is retained and the source-specific hypotheses are not relaxed. In particular, positive finite geometry, optional spectral branches, and negative comparison tests do not independently imply a physical interacting particle pole, measured flavor matrix, continuum anomaly or baryogenesis history.
\end{abstract}
\tableofcontents
\clearpage

\section{Finite representations, quotient characters and anomaly computations}
\label{p5d:sec:A}
\noindent\textit{Technical owner.} These exact inherited calculations
support the corresponding statements in Papers~V-A--V-C;
the appropriate action/observable occurrence conditions remain in those papers.
\subsection{Gauge data from the structural carrier}
\label{p5:sec:part4-gauge-data}

\subsubsection{Gauge-group and hypercharge inputs}
On \(\mathsf G_5^{\rm fin}\), the represented gauge
group is
\[
SU(3)_c\times SU(2)_L\times U(1)_Y
\]
together with the normalized hypercharge ratio
\begin{equation}
\label{p5:eq:part4-hypercharge-normalization}
\boxed{k_3:k_2:k_Y=1:1:\frac53.}
\end{equation}
In the coefficient ordering used below, the unbroken couplings are
\[
(\hat\alpha_1,\hat\alpha_2,\hat\alpha_3),
\]
where $\hat\alpha_1$ denotes the $5/3$-normalized hypercharge coupling. Throughout the one-loop bridge, the matter index is written as the Dynkin index $T_i(R)$ for the representation $R$ with respect to the factor $G_i$; no second index notation is introduced. The adjoint self-response is written separately as $C_2(G_i)$. The required indices are
\begin{equation}
\label{p5:eq:part4-casimirs}
\boxed{C_2(SU(3))=3,\qquad C_2(SU(2))=2,\qquad C_2(U(1))=0,}
\end{equation}
\begin{equation}
\label{p5:eq:part4-fundamental-indices}
\boxed{T(\mathbf 3)=\frac12,\qquad T(\mathbf 2)=\frac12,\qquad T_1(R)=\frac35\,Y^2\times(\text{spectator multiplicity}).}
\end{equation}
\subsection{Per-representation one-loop blocks}
\label{p5:sec:part4-representation-blocks}

\subsubsection{One-generation representation census}
For one generation of the Weyl carrier on the generated structural hypercharge branch, the
matter content is
\[
Q_L:(\mathbf 3,\mathbf 2,\tfrac16),\qquad u_R:(\mathbf 3,\mathbf 1,\tfrac23),\qquad d_R:(\mathbf 3,\mathbf 1,-\tfrac13),
\]
\[
L_L:(\mathbf 1,\mathbf 2,-\tfrac12),\qquad e_R:(\mathbf 1,\mathbf 1,-1),\qquad \nu_R:(\mathbf 1,\mathbf 1,0),
\]
together with one Higgs doublet
\[
H:(\mathbf 1,\mathbf 2,\tfrac12).
\]
The right-handed neutrino is neutral and therefore does not contribute to gauge running.

\subsubsection{Per-representation coefficient blocks}
\begin{proposition}[Evaluation of the per-representation one-loop blocks]
\label{p5:prop:part4-representation-blocks}
Assume
\(\mathsf G_5^{\rm fin}
+\mathsf R_5^{\rm SSQFT}\).  The graded heat-boundary theorem evaluated on the
reconstructed representation gives, in \((b_1,b_2,b_3)\) order,
\begin{equation}
\label{p5:eq:part4-bQL}
b(Q_L)=\left(\frac{1}{15},\,1,\,\frac23\right),
\end{equation}
\begin{equation}
\label{p5:eq:part4-buR}
b(u_R)=\left(\frac{8}{15},\,0,\,\frac13\right),
\end{equation}
\begin{equation}
\label{p5:eq:part4-bdR}
b(d_R)=\left(\frac{2}{15},\,0,\,\frac13\right),
\end{equation}
\begin{equation}
\label{p5:eq:part4-bLL}
b(L_L)=\left(\frac15,\,\frac13,\,0\right),
\end{equation}
\begin{equation}
\label{p5:eq:part4-beR}
b(e_R)=\left(\frac25,\,0,\,0\right),
\end{equation}
\begin{equation}
\label{p5:eq:part4-bH}
b(H)=\left(\frac{1}{10},\,\frac16,\,0\right).
\end{equation}
\end{proposition}

\begin{proof}
\ProofPlan{We evaluate the one-loop matter blocks representation by representation: left quark doublet, right quark singlets, lepton blocks, and the Higgs scalar block.}

\begin{proofsteps}
\item \textbf{Left quark doublet block.}
\leavevmode\par\noindent
Apply \eqref{p5:eq:part4-standard-bi} representation by representation.  For
the left quark doublet,
\[
\begin{aligned}
T_1(Q_L)
&=
\frac35\left(\frac16\right)^2(3\cdot 2)
=
\frac{1}{10},
\\
T_2(Q_L)
&=
3\,T(\mathbf 2)
=
\frac32,
\\
T_3(Q_L)
&=
2\,T(\mathbf 3)
=
1.
\end{aligned}
\]
This gives
\[
b(Q_L)=\left(\frac{1}{15},1,\frac23\right).
\]

\item \textbf{Right quark singlet blocks.}
Similarly,
\[
T_1(u_R)=\frac45,\quad T_3(u_R)=\frac12\quad\Longrightarrow\quad b(u_R)=\left(\frac{8}{15},0,\frac13\right),
\]
\[
T_1(d_R)=\frac15,\quad T_3(d_R)=\frac12\quad\Longrightarrow\quad b(d_R)=\left(\frac{2}{15},0,\frac13\right),
\]

\item \textbf{Lepton blocks.}
\[
T_1(L_L)=\frac{3}{10},\quad T_2(L_L)=\frac12\quad\Longrightarrow\quad b(L_L)=\left(\frac15,\frac13,0\right),
\]
\[
T_1(e_R)=\frac35\quad\Longrightarrow\quad b(e_R)=\left(\frac25,0,0\right).
\]

\item \textbf{Higgs scalar block.}
For the Higgs doublet,
\[
T_1(H)=\frac{3}{10},\qquad T_2(H)=\frac12,
\]
and since $H$ is a complex scalar one obtains
\[
b(H)=\left(\frac{1}{10},\frac16,0\right).
\]
\end{proofsteps}
\end{proof}

\subsection{Balanced generation theorem}
\label{p5:sec:part4-balanced-generation}

\subsubsection{Balanced one-generation block}
\begin{theorem}[Gauge-balanced generation]
\label{p5:prop:part4-balanced-generation}
\label{p5:thm:part4-balanced-generation}% compatibility alias with environment prefix
Assume
\(\mathsf G_5^{\rm fin}
+\mathsf R_5^{\rm SSQFT}\).  One complete generation contributes equally to all
three gauge sectors:
\begin{equation}
\label{p5:eq:part4-balanced-generation}
\boxed{b_{\rm 1\,gen}=\left(\frac43,\frac43,\frac43\right).}
\end{equation}
Hence one complete generation is gauge-balanced in the \(5/3\)-normalized
Abelian convention.
\end{theorem}

\begin{proof}
Summing Eqs.~\eqref{p5:eq:part4-bQL}--\eqref{p5:eq:part4-beR} gives
\[
\frac{1}{15}+\frac{8}{15}+\frac{2}{15}+\frac15+\frac25=\frac43,
\]
for the hypercharge sector,
\[
1+\frac13=\frac43,
\]
for the weak sector, and
\[
\frac23+\frac13+\frac13=\frac43,
\]
for the color sector. Therefore
\[
b_{\rm 1\,gen}=\left(\frac43,\frac43,\frac43\right).
\]
\end{proof}

\subsubsection{Three-generation matter block}
\begin{corollary}[Three-generation matter block]
\label{p5:cor:part4-three-generation-matter}
On the same
\(\mathsf G_5^{\rm fin}
+\mathsf R_5^{\rm SSQFT}\) domain, three complete generations give
\begin{equation}
\label{p5:eq:part4-three-generation-matter}
\boxed{b_{\rm 3\,gen\,matter}=(4,4,4).}
\end{equation}
\end{corollary}

\begin{proof}
Theorem~\ref{p5:prop:part4-balanced-generation} gives one complete generation
as \((4/3,4/3,4/3)\).  The generated finite package supplies the rank-three
replication carrier, so multiplying the generation block by three gives
\((4,4,4)\).
\end{proof}

\subsection{Normalized pair coefficients on the existing physical fibre}
\label{p5:app:v2174-field-coefficient-dictionary}
This appendix supplies the explicit basis details of
Corollary~\ref{p5:cor:v2174-pair-field-dictionary}. Its only ancestry is
the existing internal spin transport and generated Higgs/gauge fibre; no
new physical carrier or quantum-field identification is assumed.

For $t_j=\sigma_j/2$ and $t_\pm=t_1\pm it_2$, the symmetric map
$K\mapsto K\epsilon$ has normalized basis images
\[
 e_1e_1\mapsto t_+,\qquad
 (e_1e_2+e_2e_1)/\sqrt2\mapsto-\sqrt2t_3,\qquad
 e_2e_2\mapsto-t_-.
\]
They are orthonormal for the complex Hilbert--Schmidt pairing, and
$[t_3,t_\pm]=\pm t_\pm$. In particular the left $L_2L_2$ coefficient
has the weak lowering type; its minus sign is a transported presentation
phase. On mixed pairs, $K\mapsto-K\epsilon$ intertwines left/right
transport since $g_R^T\epsilon=\epsilon g_R^{-1}$.
The normalized geometric coefficients of
Proposition~\ref{p2b:prop:v2174-geometric-pair-presentation} are
$K_a=A_a\epsilon/\sqrt2$, hence map to $A_a/\sqrt2$.
This factor normalizes the pair basis; it is not a fitted field scale.

\begin{lemma}[The dictionary preserves the inherited fibre charges]
\label{p5:lem:v2174-field-dictionary-charges}
On the existing $48$-dimensional fibre
$F=(W_L\oplus W_R)\otimes Q\otimes G$, with
$Q=\ell\oplus C$ and $\dim G=3$, retain
\[
 Y=T_R^3+\tfrac12(\tfrac13P_C-P_\ell),\qquad Q_{\rm el}=T_L^3+Y.
\]
The upper scalar slots, extended by $1_Q\otimes1_G$, obey
\[
 [Y,H]=\tfrac12H,\qquad [Y,\widetilde H]=-\tfrac12\widetilde H,
 \qquad [T_c,H]=[T_c,\widetilde H]=0,
\]
where $H:W_{R-}\to W_L$, $\widetilde H:W_{R+}\to W_L$, and $T_c$
ranges over color generators. Also
\[
 [Q_{\rm el},t_-\otimes1_Q\otimes1_G]
       =-t_-\otimes1_Q\otimes1_G.
\]
\end{lemma}
\begin{proof}
On $W_L$, $T_R^3=0$; on the two right input lines its eigenvalues are
$\mp1/2$ respectively. Taking output minus input gives $\pm1/2$.
The component character commutes with the identity on $Q$, as do all color
generators. The last equation is the left $[t_3,t_-]=-t_-$ relation;
the other fibre factors commute. These calculations retain both upper
scalar columns; the self-adjoint scalar also retains their lower adjoints.
\end{proof}
The active gauge algebra remains
$\mathfrak{su}(2)_L\oplus\mathfrak u(1)_Y\oplus\mathfrak{su}(3)_c$,
of real dimension twelve. The right off-Cartan adjoint is outside that
algebra; its Cartan enters the already generated $Y$. Color comes from
$C\subset Q$, not from the four auxiliary spin modes. These role statements
delete no geometric sector or complete action term. Coefficient matrices
are not identified with quantum creation maps, and the ancestry does not
identify internal spin labels with Lorentz chirality.

% END INLINE SOURCE: Paper5A-FieldCoefficientDictionary-Appendix-v1_0.tex

% END C01 V2174 FIELD DICTIONARY APPENDIX



% BEGIN INLINE SOURCE: Paper5A-GaugeExportScope-Appendix-v1_0.tex
% Framework v2.171: integrated conversation proof on its stated domain.
\subsection{Cartan transgression and physical export scope}
\label{p5a:v2171:app:gaugehistoryscope}
This appendix compares already owned mathematical gauge paths with the
additional physical current and history map consumed by baryogenesis.
I-B, \texttt{p1b:prop:g11-transgression}, supplies the precise mathematical
test on its stated Cartan family:
\begin{equation}
 A_f=f\Theta,\qquad N_{\rm CS}(f)=3f^2-2f^3,\qquad
 \nu=N(f_f)-N(f_i).
 \label{p5a:v2171:eq:csflux}
\end{equation}
The sole doublet crossing is $\lambda_0=(3/2-3f)/R$ on
\texttt{p1b:thm:dirac-spectrum}; its upward-positive spectral flow is $-1$
from $f=0$ to $1$. The physical Weyl/time convention must be kept when
relating that sign to \eqref{p5a:v2171:eq:unitexport}. A constant midpoint has
$N_{\rm CS}=1/2$ but $\nu=0$ and no time crossing. The Euclidean cylinder
kink is explicitly an additional mathematical branch in I-B, rather than
the native Lorentzian baryogenesis history. In particular
$f(t)=\eta_{\rm reh}(t)$ or a selected kink is not derived here.
This missing history and its regulated physical charge map are distinct
from the already closed native K8--K16 reheating owner.

\subsection{External anomaly literature checked without importing its history}
\label{p5a:v2171:app:externalanomaly}
The supplied bibliography already cites Gibbons and Steif's
\emph{Yang--Mills Cosmologies and Collapsing Gravitational Sphalerons},
arXiv:hep-th/9311098, especially its round-sphere level crossing.
Its Sections 1, 2 and 4 consider actual evolving gauge backgrounds and
their fermion equation; its vacuum/spectral-flow interpretation uses its
own quantum convention. That background and convention are not the native
cyclic state/handoff merely because their three-sphere spectra are related.
The primary paper was checked at
\texttt{https://arxiv.org/pdf/hep-th/9311098}.

Annala and Rummukainen, \emph{Electroweak sphaleron in a magnetic field},
Physical Review D 107, 073006 (2023), DOI
\texttt{10.1103/PhysRevD.107.073006}, gives the weak and hypercharge
Chern--Simons contributions to the anomaly relation. This supplies an
external normalization cross-check; its lattice rates, Higgs background,
plasma and thermodynamic history are not adopted as native input.
The coefficients in \eqref{p5a:v2171:eq:exportcoefficients} were calculated
independently from the supplied framework's representation table.

The mathematical midpoint, its half-integer Chern--Simons coordinate,
the optional Euclidean kink, and a physical Lorentzian transition are
different objects. Neither external paper supplies the missing native
regulated current/history map by citation alone.

% END INLINE SOURCE: Paper5A-GaugeExportScope-Appendix-v1_0.tex


\label{p5:part:appendices}

\subsection{Conditional transverse electroweak propagation from the common Hessian}
\label{p5:app:v244-transverse-ew-propagation}

The common-Hessian theorem above is already conditional on a supplied nonzero
Higgs orbit.  On that same optional branch it can be composed with the
physical transverse gauge operator generated in Paper~III-F.  This appendix
records that composition.  Paper~V-C owns only the claim-scoped scalar
stationary/effective-action boundary; native electroweak reduction itself is
already fixed by the occurring chiral solder and its stabilizer.

\begin{corollary}[Common-Hessian transverse propagation]
\label{p5:cor:v244-common-hessian-transverse-propagation}
On the hypotheses of
Theorem~\ref{p5:thm:v164-common-ew-hessian}, after the source-free physical
coexact reduction of Paper~III-F,
Corollary~\ref{p3f:cor:v244-transverse-gauge-propagation}, the charged,
neutral massive, and photon eigendirections of the one electroweak Hessian
obey
\begin{align}
 K_{\gamma,T}&=\partial_t^2+c_{\rm op}^2\Delta_1^{\rm co},\\
 K_{W,T}&=\partial_t^2+c_{\rm op}^2\Delta_1^{\rm co}
 +\frac{m_W^2c_{\rm op}^4}{\hbar^2},\\
 K_{Z,T}&=\partial_t^2+c_{\rm op}^2\Delta_1^{\rm co}
 +\frac{m_Z^2c_{\rm op}^4}{\hbar^2},
\end{align}
where the mass-to-frequency conversion uses the same physical-unit convention
as Paper~IV.  The mass matrix is zeroth differential order, so all three
operators retain the same principal Lorentzian cone.

On an oriented round closed position slice for which a physical radius \(R_x\)
has separately been constituted, Paper~I-A's coexact curl spectrum gives
\[
 \omega_{V,L}^2
 =c_{\rm op}^2\frac{(L+1)^2}{R_x^2}
 +\frac{m_V^2c_{\rm op}^4}{\hbar^2},
 \qquad V\in\{\gamma,W,Z\},\quad L\ge1,
\]
with \(m_\gamma=0\).  Without that dimensionful calibration the same result
is retained spectrally with the eigenvalues of \(\Delta_1^{\rm co}\).
This corollary treats transverse physical polarizations only; the massive
longitudinal \(W/Z\) branch requires the coupled Higgs/Goldstone/Gauss
reduction.
\end{corollary}

\begin{corollary}[Photon helicity factorization and Hopf lowest mode]
\label{p5:cor:v244-photon-helicity-hopf-mode}
On the oriented coexact photon block let \(C=*d\).  Then
\[
 \boxed{K_{\gamma,T}
 =(\partial_t-i c_{\rm op}C)(\partial_t+i c_{\rm op}C)}.
\]
Thus the two curl signs are the two first-order helicity branches.  On a
round slice with physical radius \(R_x\), the Hopf contact form obeys
\(C\eta=(2/R_x)\eta\), so
\[
 a_\eta(t)=e^{-2ic_{\rm op}t/R_x}\eta
\]
is an exact lowest positive-curl transverse Maxwell mode.  This is a mode
existence statement; it does not assert that the electromagnetic state
occupies that mode.
\end{corollary}

\clearpage


% BEGIN M8 v1.78 SM03
% Results: B04-STABILIZERS, B05-EXACT-ENDPOINT, B06-WEAK-ROLE, B07-FULL-STABILIZER, B08-INTERMEDIATE-TYPES, B08-MULTIPLICATION-COLOR, B08-FULL-TORUS-AVERAGE, B09-SUPPORT-CLOSURE, B09-FIRST-ORDER-TEST, B10-STRONGER-PHYSICAL-MARK, B11-PHYSICAL-AXIAL-PORT, B11-DISPLAYED-JOINT-WORD
\subsection{Parent-relation hypercharge scope reminder}
\label{p5:sec:v187-pcad-yfix}
Proposition~\ref{p5:prop:v187-pcad-yfix}, placed before its first structural
consumer, derives the hypercharge stabilizer from the generated parent
quadratic.  The remaining late-port issue is therefore not hypercharge
selection but the typing of any separately claimed field, vacuum, source, or
background member.

\begin{remark}[Scope of the late-port firewall]
\label{p5:rem:v187-yfix-late-port-firewall}
The earlier fixed-input/late-Hom-space no-go remains valid when one starts from
already separated descendants and asks for a new physical transport into
$L_0$.  The present proof changes the primitive object: it descends the
pre-existing parent quadratic before the descendants are separated.  No
firewall is weakened.
\end{remark}


% BEGIN CHARGED-PRIMITIVE-OWNER-G15 v3.72
\subsubsection{Pre-Y trace, faithful characters and global stabilizers}

\begin{proposition}[Trace and anomaly calculation before Y selection]
\label{p5:prop:pre-y-trace-anomaly-plane}
\label{p5:prop:pre-y-abelian-trace-and-anomalies}
Write \(T=T_R^3\), \(B=B-L\), retaining the same one-field/conjugate
candidate domain.  Per family, the six blocks have
\[
\begin{array}{c|rrrrrr}
 &LC&L\ell&R+C&R-C&R+\ell&R-\ell\\ \hline
 \dim&6&2&3&3&1&1\\
 T&0&0&1/2&-1/2&1/2&-1/2\\
 B&1/3&-1&1/3&1/3&-1&-1
\end{array}
\]
Every coefficient of their local chiral two-torus anomaly
polynomial vanishes.  Their ordinary, unsigned trace form is
\[
 \eta=N_g\begin{pmatrix}2&0\\0&16/3\end{pmatrix}.
\]
The direction trace-orthogonal to \(Y=T+B/2\) is
\(X=4T-3B\), with
\[
 \eta(Y,Y)=10N_g/3,\quad \eta(X,X)=80N_g,\quad \chi_0(X)=5.
\]
For a retained quadratic of norm \(\sigma\), the infinitesimal
orbit form is \(4\sigma^2\chi_0\otimes\chi_0\).  Relative to
\(\eta\), it has kernel \(\RR Y\), image \(\RR X\), and
nonzero eigenvalue \(5\sigma^2/(4N_g)\).
\end{proposition}
\begin{proof}
Take positive chirality sign for the two left blocks and
negative sign for the four right blocks.  The paired right
weights cancel \(T,T^3,TB^2\).  The remaining mixed cubic
sum is proportional to \(3(1/3)-1=0\); the left and right
\(B,B^3\) sums cancel.  The \(SU(3)^2B\) sum is
\(2/3-1/3-1/3=0\); the \(SU(2)^2B\) sum is
\(3(1/3)-1=0\).  The analogous \(T\) sums vanish, and the
pure color cubic anomaly cancels between the two left color
copies and two right copies.  These are local polynomial
calculations, not global anomaly or regulator statements.
Unsigned dimension-weighted summation gives \(\eta\).
Direct substitution proves the orthogonality and norms.
Since \(\eta^{-1}\chi_0=X/(16N_g)\), the rank-one
endomorphism associated to the orbit form sends \(X\) to
\(5\sigma^2X/(4N_g)\).
\end{proof}

This trace eigenvalue is a geometric quadratic response; it is
not a physical mass without an action, kinetic normalization
and occurrence of that response.  The anomaly polynomial
vanishes on the whole candidate plane and hence does not select Y.

\begin{proposition}[Faithful candidate kernel and primitive character]
\label{p5:prop:pre-y-faithful-candidate-kernel}
Use integer charges \(r=2T\), \(b=3B\) on \(F\), and the
coordinate doublet with \((r,b)=(1,0)\).  The faithful linear
candidate group on these roles is
\[
 G_{\mathrm{cand}}=
 (SU(3)\times SU(2)\times U(1)_r\times U(1)_b)/\Gamma,
\]
where, for \(\omega_3=e^{2\pi i/3}\),
\[
 \Gamma=\{(\omega_3^a,(-1)^c,(-1)^c,
                 (-1)^c\omega_3^{-a}):a\in\ZZ_3,\ c\in\ZZ_2\}
 \simeq\ZZ_6.
\]
The character of \(L_0\),
\(\chi_N(z,w)=zw^{-3}\), descends and is primitive.
Its connected kernel \(K_1\) is the faithful Standard-Model
candidate subgroup with the inherited \(\ZZ_6\) quotient.
For \(q_Y=6Y=3r+b\) and \(X=2r-b\), the torus map
\[
 (y,x)\longmapsto(z,w)=(y^3x^2,yx^{-1})
\]
has diagonal \(\ZZ_5\) kernel.  Thus
\(G_{\mathrm{cand}}=(K_1\times U(1)_X)/\ZZ_5\).
\end{proposition}
\begin{proof}
A kernel element must be central on both nonabelian defining
roles.  The two right weights force \(z^2=1\); the lepton
blocks give \(z=w^3=(-1)^c\), and the color blocks then give
\(w=(-1)^c\omega_3^{-a}\).  The coordinate doublet has
center product \((-1)^cz=1\).  This proves the kernel.
The character \(zw^{-3}\) equals one on it.  On the covering
torus its exponent vector \((1,-3)\) is primitive; its kernel
\(z=w^3\) is connected, and the included finite quotient
does not change connectedness.  Substitution gives the
integer hypercharge weights \(3r+b\), without assuming a
post-Y group as an antecedent.  The exponent matrix of
the final torus map has determinant \(-5\), and its kernel
is \(y=x,\ x^5=1\), proving the last statement.
\end{proof}

\begin{proposition}[Quadratic stabilizer and center parity]
\label{p5:prop:quadratic-global-stabilizer-parity}
\label{p5:thm:faithful-quadratic-incidence-stabilizer}
For a coherently fixed physical \(m\ne0\),
\[
 K_2=\ker\chi_N^2=K_1\times\ZZ_{2,M}.
\]
The splitting element
\(p=(-1)^{3(B-L)}\) is central, acts as \(-I\) on \(F\)
and \(+I\) on the coordinate/Higgs role, and has
\(\chi_N(p)=-1\).  A primitive degree-\(n\) pointing has
\(n\) components and the same connected Y kernel.
The relative determinant of Paper I-A has physical character
\[
 z^{-2}w^6=\chi_N^{-2},\qquad (T,B)=(-1,+2),\quad
 Y=0,\quad X=-10,
\]
so a nonzero retained map in that line has exactly the same
stabilizer.  The geometric spin-center representative
\(\xi=(1,-1,-1,1)\) and \(p=(1,1,1,-1)\) agree on \(F,H\).
\end{proposition}
\begin{proof}
The fibre of \(\chi_N\) over \(+1\) is \(K_1\), and that
over \(-1\) is \(pK_1\).  The central order-two element
\(p\notin K_1\) splits the extension.  More generally the
preimage of the \(n\)-th roots of unity has \(n\) components.
In the integer torus presentation, \(\chi_N=x^5\), so the
quadratic equation is \(x^{10}=1\); dividing by the existing
\(\ZZ_5\) overlap leaves two components.  The natural
relative determinant character is the square dual
character of \(L_0\), proving its weights.  Finally
\(\xi p^{-1}=(1,-1,-1,-1)\in\Gamma\), proving equality on
the stated faithful tensor roles.
\end{proof}

These kernel statements are role-dependent.  The separately
retained spin and axial factors need not have the same kernel
as their matter tensor product.  In particular,
\(\bar L_0\simeq W_{R-}\otimes\bar\ell\), whereas
\(W_{R-}\otimes\ell\) has a different \(B\)-weight.  The
clock half-turn moves \(x\) to \(-x\) and acts on degree-\(k\)
position harmonics by \((-1)^k\); the internal \(p\) fixes
the position base and has a uniform matter sign.  Their
geometric ancestry does not identify these two actions.
% END M8 v1.78 SM02




Let \(T_R^3\) denote the convention-fixed right half-unit on \(W_R\), with eigenvalues
\(\pm\tfrac12\), and with zero action on \(W_L\).  On the full carrier,
\begin{equation}
\label{p5:eq:part1-TR3-full}
T_R^3 = T_R^3\otimes 1_{\ell\oplus C}\otimes 1_G.
\end{equation}
By Lemma~\ref{p5:lem:one-abelian-slot-derived}, the unrestrained quotient
centralizer is five-dimensional.

The one-field/conjugate condition in
Lemma~\ref{p5:lem:intrinsic-abelian-centralizer-plane} reduces it to the
two-dimensional plane, so any candidate residual generator has the form
\begin{equation}
\label{p5:eq:part1-Y-general}
Y=a\,T_R^3+b(B-L).
\end{equation}

\begin{lemma}[Distinguished-line character-kernel ray]
\label{p5:lem:hypercharge-normalization-unique}
On \(\mathsf G_5^{\rm fin}\), inside the
one-field/conjugate plane
\[
\mathrm{span}\{T_R^3,B-L\},
\]
the character kernel of \(L_0=W_{R+}\otimes\ell\) fixes the unique
nonzero line
\[
\RR\!\left(T_R^3+\frac12(B-L)\right).
\]
Choosing the representative whose coefficient on the convention-fixed
\(T_R^3\) is one gives
\[
Y=T_R^3+\frac12(B-L).
\]
The ray is derived from coordinate covariance and the distinguished-line
character kernel; the last equality fixes a generator convention.  Physical structural identification of this ray follows from
Proposition~\ref{p5:prop:v187-pcad-yfix}.
\end{lemma}

\begin{proof}
For \(X=aT_R^3+b(B-L)\), the character on
\(L_0=W_{R+}\otimes\ell\) is \(a/2-b\).
Corollary~\ref{p5:lem:geometric-majorana-line} therefore gives
\(b=a/2\).  Every nonzero solution spans the displayed line.  Setting
\(a=1\) chooses the representative normalized in the inherited
\(T_R^3\)-coordinate.
\end{proof}

\begin{remark}[Why the module alone does not select the physical subgroup]
\label{p5:rem:yfix-no-go}
If the parent relation is omitted after descendant separation, the same finite data admit distinct connected
one-dimensional subgroups of \(\mathbb T_5\) as the surviving Abelian factor.
For example, the \(T_R^3\) and \(B-L\) rays both commute with the retained
nonabelian factors, are generation blind, and lie in the
coordinate-conjugate plane.  Moreover the full parent torus preserves the
one-dimensional subspace \(L_0\) setwise.  Hence naming the line as a module,
without a prior nonzero pointing/background tensor in it, does not imply that
the physical subgroup acts trivially on a vector.  GP3 cannot distinguish
these realizations.  Proposition~\ref{p5:prop:v187-pcad-yfix} supplies the
missing structural relation by descending the generated parent quadratic.
\end{remark}

\begin{theorem}[Geometric hypercharge reconstruction from the parent relation]
\label{p5:thm:conditional-hypercharge-reconstruction}
\label{p5:thm:part1-anchor-free-hypercharge}
On \(\mathsf G_5^{\rm fin}\) and
\[
F=(W_L\oplus W_R)\otimes(\ell\oplus C)\otimes G,
\]
coordinate/conjugate covariance reduces the diagonal centralizer modulo the
identity to \(\operatorname{span}\{T_R^3,B-L\}\).  The generated parent
quadratic has the nonzero projected component of
Proposition~\ref{p5:prop:v187-pcad-yfix}; its connected stabilizer has Lie
algebra
\[
\RR Y
=
\RR\!\left(T_R^3+\frac12(B-L)\right).
\]
This is the finite represented \(U(1)_Y\) action generated by the structural
constructor chain.  Its realization as a local connection component uses
\(\mathsf J_5^{\rm conn,der}\).  In the representative normalized on the
provenance-fixed right Cartan,
\begin{equation}
\label{p5:eq:part1-Y}
Y=T_R^3+\frac12(B-L).
\end{equation}
No observed Standard--Model charge value, independently selected pointing, or
Majorana condition is used.  Projective/traceless representation is generated
by ray faithfulness; without removing the common identity direction, the
character-kernel relation would not select a unique represented nonzero line.
\end{theorem}

\begin{proof}
Lemma~\ref{p5:lem:intrinsic-abelian-centralizer-plane} gives the candidate
plane.  Proposition~\ref{p5:prop:v187-pcad-yfix} supplies the nonzero projected
parent relation of weight \((-1,+2)\), whose connected stabilizer is
\(-a+2b=0\).  Lemma~\ref{p5:lem:hypercharge-normalization-unique} fixes the
inherited generator convention.  No member-selection operation enters the
proof.
\end{proof}

\begin{corollary}[Scalar weight after hypercharge reconstruction]
\label{p5:cor:scalar-hypercharge-after-reconstruction}
On the domain of
Theorem~\ref{p5:thm:part1-anchor-free-hypercharge}, the one-field modules of
Theorem~\ref{p5:thm:minimal-scalar-left-right-intertwiner} evaluate to
\[
H\sim(\mathbf1,\mathbf2)_{1/2},
\qquad
\widetilde H\sim(\mathbf1,\mathbf2)_{-1/2}.
\]
\end{corollary}

\begin{proof}
For either \(q\in\{\ell,C\}\), the \(B-L\) contribution cancels between
source and target.  The left block has \(T_R^3=0\), whereas \(W_{R-}\) and
\(W_{R+}\) have \(T_R^3=-1/2\) and \(+1/2\), respectively.  Target charge
minus source charge is therefore \(+1/2\) for \(H\) and \(-1/2\) for its
conjugate.
\end{proof}

\begin{theorem}[Yukawa--anomaly cross-check of the hypercharge plane]
\label{p5:thm:yukawa-anomaly-hypercharge-crosscheck}
Assume one Lorentz-chiral generation, one scalar field \(H\) and its
conjugate, all four Dirac slots, and perturbative anomaly freedom.  Write
\[
Y(Q_L)=q,\quad Y(L_L)=l,\quad Y(H)=h,
\]
\[
Y(u_R)=u,\quad Y(d_R)=d,\quad
Y(\nu_R)=n,\quad Y(e_R)=e.
\]
Yukawa covariance gives
\[
u=q+h,\qquad d=q-h,\qquad
n=l+h,\qquad e=l-h.
\]
All perturbative anomaly conditions then reduce to
\[
3q+l=0.
\]
The resulting two-parameter charge family is exactly
\(aT_R^3+b(B-L)\), with \(a=2h\), \(b=3q\).  If the derived distinguished right singlet is
neutral, \(n=0\), then \(b=a/2\).  Choosing \(h=1/2\) gives the charge values
in \eqref{p5:eq:projector-hypercharge-discharge}.
\end{theorem}

\begin{proof}
Using right-handed fields with the opposite sign in left-chiral anomaly
sums,
\[
\mathcal A_{33Y}=2q-u-d=0,
\qquad
\mathcal A_{22Y}=3q+l,
\]
\[
\mathcal A_{{\rm grav}^2Y}
=6q-3u-3d+2l-n-e=0,
\]
and direct expansion gives
\[
\mathcal A_{Y^3}
=6q^3-3u^3-3d^3+2l^3-n^3-e^3
=-6h^2(3q+l).
\]
Thus the weak mixed anomaly is the only independent condition.  Substituting
\(l=-3q\) into the six charges and setting \(a=2h,b=3q\) reproduces the
eigenvalues of \(aT_R^3+b(B-L)\).  Finally \(n=l+h=0\) gives \(h=3q\), or
\(b=a/2\).
\end{proof}



\begin{corollary}[Projector hypercharge and electric-charge discharge]
\label{p5:cor:projector-hypercharge-electric-discharge}
On \(\mathsf G_5^{\rm fin}\), applying the geometrically
reconstructed operator
\[
Y=T_R^3+\frac12(B-L)
\]
to the fixed projector sectors gives, generationwise,
\begin{equation}
\label{p5:eq:projector-hypercharge-discharge}
\begin{aligned}
Y(P_LP_C)&=\frac16,
&
Y(P_{R+}P_C)&=\frac23,
&
Y(P_{R-}P_C)&=-\frac13,\\
Y(P_LP_\ell)&=-\frac12,
&
Y(P_{R+}P_\ell)&=0,
&
Y(P_{R-}P_\ell)&=-1.
\end{aligned}
\end{equation}
With the scalar-stabilizer generator
\[
Q=T_L^3+Y,
\]
the weak-doublet components and their right-handed singlet partners have electric charges
\[
\frac23,
\qquad
-\frac13,
\qquad
0,
\qquad
-1.
\]
Thus the conventional matter names are eigenvalue labels read after the
geometric hypercharge operator has been fixed.
\end{corollary}

\begin{proof}
On left-handed blocks the right Hopf Cartan satisfies \(T_R^3=0\).  On color
blocks \(B-L=1/3\), while on lepton blocks \(B-L=-1\).  Hence
\(Y(P_LP_C)=1/6\) and \(Y(P_LP_\ell)=-1/2\).  On right-handed blocks the
Hopf eigenvalues are \(+1/2\) on \(P_{R+}\) and \(-1/2\) on \(P_{R-}\).
Substituting these eigenvalues into \(T_R^3+\frac12(B-L)\) gives
\(2/3,-1/3,0,-1\) on
\(P_{R+}P_C,P_{R-}P_C,P_{R+}P_\ell,P_{R-}P_\ell\), respectively.  The
scalar-stabilizer theorem gives \(Q=T_L^3+Y\); adding
\(T_L^3=\pm1/2\) on weak doublets reproduces the displayed electric-charge
spectrum.  No charge eigenvalue is used in the preceding normalization
proof; all entries are evaluations of the reconstructed operator on the
already fixed projectors.
\end{proof}

\begin{corollary}[Projector-defined matter names after the derivation]
\label{p5:cor:part1-matter-names-after-Y}
After \(Y\) has been fixed, define
\[
Q_L=P_LP_CF_{\rm core},\quad L_L=P_LP_\ell F_{\rm core},\quad
\nu_R=P_{R+}P_\ell F_{\rm core},\quad e_R=P_{R-}P_\ell F_{\rm core},
\]
\[
u_R=P_{R+}P_CF_{\rm core},\qquad d_R=P_{R-}P_CF_{\rm core}.
\]
These are names for the resulting charge eigenspaces; in particular
\(\nu_R\) is the distinguished right-handed lepton module \(L_0\), and its
neutrality is a consequence of the character-kernel construction.
\end{corollary}

\subsubsection{Finite-algebra taxonomy and the represented gauge action}
\label{p5:sec:part1-finite-algebra}

Two real finite \(*\)-algebras provide a useful component presentation,
\[
\mathcal A_{\chi}=\mathbb H_w\oplus\mathbb C_R,
\qquad
\mathcal A_{q\ell}=\mathbb C_\ell\oplus M_3(\mathbb C)_C,
\]
whose independent algebraic composition is
\begin{equation}
\label{p5:eq:part1-finite-algebra}
\mathcal A_F^{\rm comp}
=
\mathcal A_{\chi}\,\widehat\otimes_{\mathbb R}\,\mathcal A_{q\ell}.
\end{equation}
This algebra must not be confused with the four-summand direct sum
\[
\mathcal A_F^\oplus
=
\mathbb H_w\oplus\mathbb C_R\oplus\mathbb C_\ell\oplus M_3(\mathbb C)_C.
\]
Indeed, as real semisimple algebras,
\begin{equation}
\label{p5:eq:part1-finite-algebra-decomposition}
\mathcal A_F^{\rm comp}
\cong
M_2(\mathbb C)\oplus M_6(\mathbb C)
\oplus\mathbb C\oplus\mathbb C
\oplus M_3(\mathbb C)\oplus M_3(\mathbb C),
\end{equation}
so
\[
\dim_{\mathbb R}\mathcal A_F^{\rm comp}=120,
\qquad
\dim_{\mathbb R}\mathcal A_F^\oplus=26.
\]
They are not isomorphic and play no interchangeable role.

Define genuine real-linear \(*\)-representations
\[
\rho_\chi(q,z)=qP_L+zP_{R+}+\bar z P_{R-},
\qquad
\rho_{q\ell}(\lambda,m)=\lambda P_\ell+mP_C.
\]
The full action is specified on simple tensors by
\begin{equation}
\label{p5:eq:part1-finite-rep}
\rho(a_\chi\otimes a_{q\ell})
=
\rho_\chi(a_\chi)\otimes\rho_{q\ell}(a_{q\ell})\otimes1_G,
\end{equation}
and extended real-linearly.  This is a linear multiplicative representation
of the tensor-product algebra.  The former product formula is therefore
retained only as a product-group action, not as a representation of a
four-summand direct-sum algebra.  By the generation-blindness theorem, the
specified core gauge action is trivial on \(G\).
The physical representation
\[
\rho_{\rm SM}:SU(3)_c\times SU(2)_L\times U(1)_Y\longrightarrow\mathcal U(F)
\]
is the finite action obtained from the generated derivations and the pointed
parent-relation stabilizer theorem.  It is neither the full unitary group
nor the full inner-automorphism group of \(\mathcal A_F^{\rm comp}\).  Its
spacetime-local realization uses \(\mathsf J_5^{\rm conn,der}\).

\begin{proof}[Proof of the algebra decomposition]
Use the real-algebra identities
\[
\mathbb H\widehat\otimes_{\mathbb R}\mathbb C\cong M_2(\mathbb C),
\qquad
\mathbb C\widehat\otimes_{\mathbb R}\mathbb C\cong\mathbb C\oplus\mathbb C.
\]
Tensoring the first identity with \(M_3(\mathbb C)\) gives
\(M_6(\mathbb C)\), while tensoring the second complex factor with
\(M_3(\mathbb C)\) gives two copies of \(M_3(\mathbb C)\).
Distributivity over the two direct sums gives
\eqref{p5:eq:part1-finite-algebra-decomposition}.  The stated dimensions
follow directly.
\end{proof}

\begin{proposition}[Faithful global group of the reconstructed representation]
\label{p5:prop:faithful-global-sm-group}
On \(\mathsf G_5^{\rm fin}\), the finite represented gauge
algebra is
\[
\mathfrak{su}(3)_c\oplus\mathfrak{su}(2)_L\oplus\mathfrak u(1)_Y.
\]
Its faithful connected global form on this matter carrier is
\[
\boxed{
G_{\rm SM}^{\rm faithful}
=
\frac{SU(3)_c\times SU(2)_L\times U(1)_Y}{\Gamma},
\qquad
\Gamma\cong Z_6 .
}
\]
Here ``finite'' means Lie-algebra integration on the fixed internal carrier.
Spacetime-local gauge fields use the derived
\(\mathsf J_5^{\rm conn,der}\) family on the \(m_x\)-side; they do not enlarge
the \(m_p=1\) carrier.  The \(U(1)_Y\) coordinate is parametrized by the integer-normalized hypercharge \(q_Y:=6Y\), so that \(e^{i\theta}\) acts as \(e^{i\theta q_Y}\).  With this convention, the kernel is generated by
\[
\left(\omega_3\mathbf 1_3,-\mathbf 1_2,e^{i\pi/3}\right),
\qquad
\omega_3=e^{2\pi i/3}.
\]
The integer charges have greatest common divisor one, so this compact
parameterization is primitive up to inversion.
\end{proposition}

\begin{proof}
\ProofPlan{We compute the complete common central kernel on one generation and quotient by it.}

\begin{proofsteps}
\item \textbf{Necessary central form.}
Let \((g_3,g_2,z)\) lie in the kernel.  Triviality on the color-fundamental
right-quark blocks makes \(g_3\) a scalar matrix, and triviality on the
weak-fundamental lepton block makes \(g_2\) a scalar matrix.  The determinant
conditions therefore give
\[
(\omega_3^a,(-1)^b,z),
\qquad
a\in\mathbb Z/3,\quad b\in\mathbb Z/2,\quad z\in U(1).
\]
Triviality on \(e_R\), whose integer hypercharge is \(-6\), gives
\(z^6=1\).  Write \(z=e^{i\pi n/3}\).

\item \textbf{The color and weak congruences.}
Triviality on \(d_R\), whose integer hypercharge is \(-2\), gives
\(\omega_3^a z^{-2}=1\), hence \(a\equiv n\pmod3\).
Triviality on \(L_L\), whose integer hypercharge is \(-3\), gives
\((-1)^b z^{-3}=1\), hence \(b\equiv n\pmod2\).
The conditions on \(Q_L,u_R,\nu_R\) then hold automatically.

\item \textbf{Exhaustion.}
Thus every kernel element is
\[
(\omega_3,-1,e^{i\pi/3})^n.
\]
Conversely, its six powers act trivially on all six matter sectors (and on
the scalar doublet), so the kernel is exactly \(Z_6\).  Quotienting by this
kernel gives the faithful global group.  The unquotiented product has the
same local Lie algebra.
\end{proofsteps}
\end{proof}

\begin{theorem}[Canonical full local gauge connection and joint Dirac family]
\label{p5:thm:canonical-full-joint-connection}
\label{p5:thm:joint5-occurrence-discharge}
Let \(U_0\Subset U\) be a generated local trivializing chart and let
\[
 P_R\longrightarrow U_0,\qquad \omega_R
\]
be the geometric-right (algebraic-left weak) factor extension of the nonflat Born connection constructed in
Theorem~\ref{p3f:thm:generated-fibre-spin-joint-operator}.  Let
\(\mathbb T_5\) integrate the generated candidate plane
\(\operatorname{span}_{\mathbb R}\{T_R^3,B-L\}\), and let
\[
 \kappa_{\rm cand}:
 SU(2)_L\times SU(3)_c\times\mathbb T_5
 \longrightarrow U(F)
\]
be the already generated finite action.  Put
\[
 G_5^{\rm cand}:=\operatorname{Im}\kappa_{\rm cand}.
\]
Then the product and structure-group extension
\begin{align}
\widehat P_5
 &:=
 P_R\times_{U_0}(U_0\times SU(3)_c)
     \times_{U_0}(U_0\times\mathbb T_5),
&
\widehat\omega_5&:=(\omega_R,0,0),\label{p5:eq:candidate-product-connection}\\
P_5^{\rm cand}
 &:=
 \widehat P_5\times_{\kappa_{\rm cand}}G_5^{\rm cand},
&
\omega_5^{\rm cand}
 &:=
 \operatorname{Ext}_{\kappa_{\rm cand}}(\widehat\omega_5)
\label{p5:eq:candidate-full-bundle}
\end{align}
give a generated principal bundle for the full candidate gauge group and a
nonflat connection on it.  Here full refers to the structure group and its
available connection directions: its canonical anchor has a nonzero weak
component and zero color and candidate-Abelian components.

The associated carrier connection \(\nabla_{5,{\rm B}}^F\) and every later
generated gauge and scalar slot lie in the one affine operator family
\begin{equation}
\label{p5:eq:full-joint-dirac-family}
\mathfrak D_5^{\rm cand}(U_0)
:=
\left\{
D_{5,{\rm B}}+c_g(a)+\gamma_g\otimes\Phi:
\begin{array}{l}
a\in\Omega_c^1(U_0,\operatorname{ad}P_5^{\rm cand}),\\[-1mm]
\Phi\in\Gamma_c(U_0,\mathcal V_H^{\rm odd})
\end{array}
\right\},
\end{equation}
where
\[
D_{5,{\rm B}}
:=
c_g\circ
\left(\nabla^{\mathbb S_g}\otimes1_F+
1_{\mathbb S_g}\otimes\nabla_{5,{\rm B}}^F\right)
\]
and \(\mathcal V_H^{\rm odd}\) is the generated self-adjoint odd
intertwiner family carried by the single weak doublet.  In particular,
\(a=a_2\oplus a_3\oplus a_{\mathbb T}\) contains the complete weak, QCD
color, and candidate-Abelian directions.  The base member
\((a,\Phi)=(0,0)\) is already nonflat.

With the structurally generated hypercharge line, the trivial \(\mathbb T_5\)-factor has the generated
reduction \(U_0\times K_Y\).  The background
\(\widehat\omega_5=(\omega_R,0,0)\) preserves that reduction.  Pushing forward
by the matter action and quotienting its discrete kernel \(\Gamma\cong Z_6\)
therefore gives
\begin{equation}
\label{p5:eq:faithful-sm-principal-connection}
 P_{\rm SM}^{\rm can}\longrightarrow U_0,
 \qquad
 \omega_{\rm SM}^{\rm can},
 \qquad
 \operatorname{StrGrp}(P_{\rm SM}^{\rm can})
 =G_{\rm SM}^{\rm faithful},
 \qquad
 \Omega_{\rm SM}^{\rm can}\ne0 .
\end{equation}
Restricting \(a\) in \eqref{p5:eq:full-joint-dirac-family} to
\(\operatorname{ad}P_{\rm SM}^{\rm can}\) gives the full Standard--Model
joint connection/odd family.  Consequently
\[
\boxed{
\mathsf J_5^{\rm conn,der}
=({\rm JOINT}_5^{\rm der})
}
\]
is theorem output, and the historical
\(({\rm JOINT}_5^{\rm occ})=\mathsf R_5^{\rm conn}\) is not an independent
assumption.

The theorem does not select a gauge potential \(a\), Higgs section \(\Phi\),
Yukawa matrix, vacuum, numerical coupling, action, stationary solution, or
global bundle topology beyond the displayed local chart.
\end{theorem}

\begin{proof}
Theorem~\ref{p3f:thm:generated-fibre-spin-joint-operator} gives both
\(SU(2)^a_L\) and \(SU(2)^a_R\) extensions.  The generated visible-chirality
ancestry identifies the geometric-right with the algebraic-left weak factor; it does not delete the
co-realized geometric-left Hopf-side geometry.  The trivial color and candidate-torus
bundles on the local chart are the identity objects for principal-bundle
product.  Product connections give \eqref{p5:eq:candidate-product-connection},
and the same structure-group extension formula as
\eqref{p3f:eq:extension-connection-form} gives
\eqref{p5:eq:candidate-full-bundle}.  Its curvature is
\[
d\kappa_{\rm cand}(\Omega_R,0,0).
\]
The fundamental \(SU(2)_L\)-action on \(W_L\) is faithful, so this expression
is nonzero because \(\Omega_R\ne0\).
The representation acts trivially on \(W_{R+}\) and \(W_{R-}\), so
\([\rho_W(SU(2)_{\rm alg,L}),Q_H]=0\). This is a statement about
the weak-only associated connection and does not make \(Q_H\) parallel
for the original full diagonal Born connection.

By definition, \(G_5^{\rm cand}\) acts faithfully on \(F\).  Equivalently,
GP2 retains every generated factor direction, while GP3 removes only the
common action kernel and unmarked carrier multiplicity; GP3 neither creates
the connection nor chooses a fluctuation.  After parent-relation hypercharge descent,
Proposition~\ref{p5:prop:faithful-global-sm-group} computes the complete
remaining kernel as \(Z_6\).  A discrete quotient changes no Lie algebra, and
the weak curvature is not in its kernel, proving
\eqref{p5:eq:faithful-sm-principal-connection}.

Associated-bundle transport sends \(\omega_5^{\rm cand}\) to
\(\nabla_{5,{\rm B}}^F\).  The generated local cotangent action, Clifford
contraction, identity extension across \(F=W\otimes Q\otimes G\), and graded
addition then give \eqref{p5:eq:full-joint-dirac-family}.  These operations
form the certified word
\[
(P_{\rm B},\omega_{\rm B},F,\kappa_{\rm cand})
\longmapsto(P_R,\omega_R)
\longmapsto(P_5^{\rm cand},\omega_5^{\rm cand})
\longmapsto\nabla_{5,{\rm B}}^F
\longmapsto\mathfrak D_5^{\rm cand}(U_0).
\]
Every arrow is a licensed principal, associated-bundle, tensor-identity, or
local operator constructor, and every target was declared before evaluation.
Paper~II-A's rooted generated-image theorem therefore co-realizes the output
with its generated inputs.  Since its required base member has nonzero
represented curvature, the zero-member alternative cannot realize the same
word.  Paper~III-F's same-datum theorem then derives the common associated
connections, curvatures, and profile.  This proves the certificate and its
scope boundary.
\end{proof}

% BEGIN M8 v1.78 DYN02
% Results: B04-NOETHER-DYNAMICS, B05-NOETHER-DYNAMICS, B06-GLOBAL-ACTION, B07-NOETHER-DYNAMICS, B07-POSITION-TORUS, B08-NOETHER-DYNAMICS, B09-PARALLEL-COMPARISON, B09-NOETHER-DYNAMICS, B10-CONDITIONAL-ACTION, B10-PHASE-AND-NOETHER, B11-CONNECTION-NOETHER, B11-DYNAMICS-HANDOFF, B13-PHYSICAL-AND-DYNAMICAL-HANDOFF

\section{Rooted triality, spin descent and source-incidence checks}
\label{p5d:sec:B}
\noindent\textit{Technical owner.} These exact inherited calculations
support the corresponding statements in Papers~V-A--V-C;
the appropriate action/observable occurrence conditions remain in those papers.
\subsection{Rooted Hurwitz source alphabet and the sharpened G15 compiler target}
\label{p5:sec:v369-rooted-hurwitz-g15-target}
\noindent\textbf{Historical comparison scope.} The alphabet, offset and
polar-word calculations below are conditional source diagnostics. Their
former G15 member/port frontier is superseded for the rooted charged target
by main-text Corollary~\ref{p5:cor:v397-charged-preterminal-assembly-closure}.
The predeclared zero-order target itself remains in the main text as
Definition~\ref{p5:def:v369-direct-mass-compiler-target}.
The active Paper~I-A rooted-Hurwitz theorem
(\texttt{p1:thm:v556-rooted-hurwitz-family-return-algebra}) now supplies on the generated family carrier \(G\) the coefficient-free
operators
\[
 K_7=2I_G+2J_G,\qquad K_9=I_G,
 \qquad R_G=\operatorname{diag}(1,-i,i),
\]
together with the already generated cyclic transport \(X\).  In particular,
\[
 \boxed{C^*(K_7,K_9,R_G)=\End(G).}
\]
This strengthens the interpretation of
Theorem~\ref{p5:thm:v186-generated-yukawa-dirac-owner-family}: the residual
G15 issue is not absence of a sufficiently large family algebra, but
occurrence of the physical sector-specific member in the same action/source
owner.

The local \(Q=\ell\oplus C\) incidence of Paper~I-A also gives
\(K_\ell=J_G\) and \(K_C=I_G+J_G\).  Combining each with the common identity
return \(K_9\) makes the two coefficient-free descendants
\begin{equation}
\label{p5:eq:v369-candidate-sector-returns}
 K_\ell^{(0)}=K_9+K_\ell=I_G+J_G,
 \qquad
 K_C^{(0)}=K_9+K_C=2I_G+J_G.
\end{equation}
Equation~\eqref{p5:eq:v369-candidate-sector-returns} is an exact geometric
identity, not an occurrence theorem: the present action has not yet proved
that these two returns are the physical lepton and quark chiral source words.



There is one useful symmetry restriction on any cyclic-offset realization.
After a root is chosen, an orientation reversal acts on the cyclic torsor by
an affine reflection
\begin{equation}
\label{p5:eq:v369-affine-cyclic-reflection}
 k\longmapsto c-k\pmod 3
\end{equation}
for a constant \(c\) determined by the relative choice of origins.  Thus
orientation reversal constrains the two offset labels, but it does not by
itself justify the stronger convention \(k\mapsto-k\); that form requires a
separate theorem aligning the two roots.

The active Paper~V-B appendix-level rooted Hurwitz--Clifford audit
(\texttt{p6:app:v562-\allowbreak hurwitz-\allowbreak clifford-\allowbreak mass-\allowbreak candidate}) tests
the most economical phase-refined implementation of this idea.  Its exact
characteristic invariants show genuine spectral addressability and nonzero
discrete CP capacity, but the same calculation proves that a nontrivial
cyclic offset of one fixed finite word does not merely permute the three
Clifford hierarchy orders: its two largest singular values coalesce at the
same leading asymptotic order.  Therefore an \(R^\pm\) compiler intended to
transport one copy of each hierarchy order cannot consist solely of choosing
\(k\) in one fixed phase-refined return.  The remaining G15 target must allow
\(R^\pm\) to enter the finite source word, the action-owned functional
calculus, the relative root alignment, or a combination of these.  No
physical mass member is selected here.

\begin{corollary}[Polar-normalized rooted-return coefficient contraction]
\label{p5:cor:v371-polar-normalized-g15-w-branch}
Paper~V-B, Theorem~\texttt{p6:thm:v564-polar-separated-oriented-return},
proves that on the ordered-return \(W\)-branch the sector-dependent positive
incidence modulus and the oriented cyclic transport separate canonically by
positive polar decomposition.  The raw orientation-odd amplitudes
\(a_qC_{2,G}\), with \(a_\ell=1\) and \(a_C=2\), therefore normalize to the
same transport generator \(C_{2,G}\).  Conditional on that branch being the
source consumed by the fermionic zero-order action, the normalized finite
word has the coefficient-free form
\[
 \boxed{W_{\pm,q}^{\rm pol}=F_\pm(a_q)+C_{2,G}},
\]
and no independent pair \((\lambda_\ell,\lambda_C)\) remains to be selected.
The four lepton/color branches are full rank.

This is a normalization theorem, not a source-occurrence theorem.  The
framework still has to prove that the physical G15 source consumes this
polar-normalized ordered return, or else derive the required physical member
through the remaining action-owned \(\Phi\) and/or relative-root \(A\) ports.
It does not identify the \(24/25\) exclusion owner with the mass source.
\end{corollary}


\begin{remark}[Present G15 port-activation ledger]
\label{p5:rem:v370-g15-port-activation-ledger}
At this rooted-Hurwitz frontier the source has isolated three
provenance-distinct ports at which the additional \(R^\pm\)-dependence may
enter: the finite source word \(W\), the action-owned functional calculus
\(\Phi\), and the relative root alignment \(A\).  With these three typed
constructor roles held fixed, their nonempty activation subsets give seven
bookkeeping cases.  This is an exhaustion of the present three-port ledger,
not a finite enumeration of concrete Yukawa/source words and not a theorem
that the seven syntactic placements remain inequivalent under every allowed
reformulation.  Corollary~\ref{p5:cor:v371-polar-normalized-g15-w-branch}
removes the continuous sign-line coefficient ambiguity on the canonical
polar-normalized ordered-return \(W\)-branch; the seven-port ledger is retained
because same-action occurrence of that branch and possible \(\Phi/A\)
participation are still open.  Paper~II-A supplies no construction-depth or
numerical fit selector beyond generated owner normalization.
\end{remark}
% END ROOTED-HURWITZ-G15 v2.62

\begin{remark}[Controlling physical boundary after the structural reductions]
\label{p5:rem:v393-assembly-boundary}
\noindent\textbf{Historical scope.} This release-stage charged residual is superseded on the rooted $u,d,e$ branch by main-text Corollary~\ref{p5:cor:v397-charged-preterminal-assembly-closure}. It is not a current prerequisite.

Primitive source incidence remains \texttt{CLOSED\_GEOMETRIC}.
Paper V-B now constructs the three-dimensional nonzero supports of the
actual terminal maps, and gives a finite conditional spectral assembly.
For its directed candidate branch the missing theorem is a licensed source
word, on the same occurring parent, with the joint output
\[
 \mathbb T_*\xrightarrow{\ w_{\rm src}\ }
 (A_s,\tau_s,q_s(H_s)A_s),\qquad H_s=A_sA_s^\dagger.
\]
It must derive both the coherent directed response and its terminal-role
matching. These objects are not licensed together merely because they can
be defined separately. Thus
\[
\begin{gathered}
 \texttt{FERMION\_SAME\_ACTION\_SVD\_ASSEMBLY}\\
 =\texttt{HISTORICAL\_OPEN\_PHYSICAL}.
\end{gathered}
\]
This target introduces neither a VEV mass law nor a fitted Yukawa parameter.
\texttt{G16\_FINITE\_\allowbreak SCHEME\_\allowbreak VERTEX} and claim-scoped pole/RG transport remain
downstream. The six matchings per sector classify only the fixed candidate
frame branch, not all possible framework source constructors.
\end{remark}

\begin{remark}[Narrowed physical fermion frontier]
\label{p5:rem:v394-parent-contraction-frontier}
\noindent\textbf{Historical scope.} This release-stage charged residual is superseded on the rooted $u,d,e$ branch by main-text Corollary~\ref{p5:cor:v397-charged-preterminal-assembly-closure}. It is not a current prerequisite.

The historical label
\[
 \texttt{FERMION\_SAME\_ACTION\_SVD\_ASSEMBLY}
\]
is superseded as the active gap name.  Once a particular typed parent action
is generated, Paper~III-A claim-complete action co-consumption automatically
retains its source operator, left/right kinetic metrics, mixed current and all
nonzero role blocks.  The remaining structural obligation is therefore
\[
 \boxed{\texttt{FERMION\_TYPED\_\allowbreak PARENT\_\allowbreak CONTRACTION}
        =\texttt{HISTORICAL\_OPEN\_PHYSICAL}.}
\]
It asks for the particular noncentral fermion coefficient and its inherited
kinetic/current data to be extracted from one licensed parent contraction on
the actual spinor/family embeddings.  It does not ask for a new binding rule.

The closed N56 terminal magnitude packet is not part of this remaining debt.
The relative charged/neutral family frame, including PMNS-type readouts, still
descends only after the relevant noncentral family embedding is generated.
Claim-scoped pole/RG transport remains downstream.
\end{remark}

\begin{corollary}[Retyped charged fermion frontier]
\label{p5:cor:v395-preterminal-frontier}
\noindent\textbf{Historical scope.} This release-stage charged residual is superseded on the rooted $u,d,e$ branch by main-text Corollary~\ref{p5:cor:v397-charged-preterminal-assembly-closure}. It is not a current prerequisite.

The historical first missing charged-fermion CORE arrow is
\[
 \boxed{\begin{gathered}
 \texttt{FERMION\_PRETERMINAL\_PARENT\_CONTRACTION}\\
 =\texttt{HISTORICAL\_OPEN\_PHYSICAL}.
 \end{gathered}}
\]
The historical labels
\texttt{FERMION\_SAME\_ACTION\_SVD\_ASSEMBLY} and the standalone
$J_s$-embedding problem are retired as first-principle gaps.  The immediately
preceding label \texttt{FERMION\_TYPED\_\allowbreak PARENT\_\allowbreak CONTRACTION} is superseded by
this sharper chronology: the particular parent must consume the retained
owner-resolved channels before scalar/support terminalization.  Once that
output exists, Paper II-A generated co-consumption supplies its mixed terms,
metrics and currents automatically.
\end{corollary}

\begin{corollary}[Retyped residual of the charged CORE gap]
\label{p5:cor:v396-charged-residual-retyped}
\noindent\textbf{Historical scope.} This release-stage charged residual is superseded on the rooted $u,d,e$ branch by main-text Corollary~\ref{p5:cor:v397-charged-preterminal-assembly-closure}. It is not a current prerequisite.

No new physical \texttt{SrcVal} postulate is required if the contraction
$\mathfrak K_s$ is derived geometrically.  The residual content of
\texttt{FERMION\_PRETERMINAL\_PARENT\_CONTRACTION} is exactly the derivation
of that particular result-independent contraction; the family host, the
$E_{32}/E_{36}/E_{40}$ hierarchy-to-$G$ block bridge, the sector typing, the
charged masses and the downstream action consumer are no longer parts of the
gap.
\end{corollary}

\subsubsection{Historical charged-lepton plane frontier}
\noindent\textbf{Historical scope.} The following v3.107 statement is superseded by main-text Corollary~\ref{p5:cor:v3108-lepton-incidence-status}.
Accordingly the charged-lepton current frontier is no longer a discrete
$\mu/\tau$ placement ambiguity.  The only remaining incidence datum is the
source-owned oriented rank-one decomposition of $Q_\perp G$ between $e$ and
$\mu$.  Equivalently,
\[
 \boxed{
 \texttt{LEPTON\_ROLE\_TO\_FAMILY\_INCIDENCE}
 =\texttt{HISTORICAL\_OPEN\_INPLANE\_SOURCE\_EVALUATION}.}
\]
The structural joint-quadratic relation $R^-\mapsto R_G^\dagger$ is not
promoted to that physical plane selector: its Paper~I-C theorem explicitly
requires a consuming action constructor, and no charged-current word with that
phase as a required input has yet been frozen.


% END INLINE SOURCE: Paper5A-HistoricalChargedAssembly-Appendix-v1_0.tex


% BEGIN C01 V2174 FIELD DICTIONARY APPENDIX

% BEGIN INLINE SOURCE: Paper5A-FieldCoefficientDictionary-Appendix-v1_0.tex
% v2.174: supplementary normalized coordinate and fibre calculation.
\subsection{Countermodels fixing the finite proof boundary}
\label{p5:app:selector-countermodels}

This appendix records the short non-entailment models behind the
fixed-descendant hypercharge reduct, the unbridged grading reduct, and the
conjugate-occurrence boundary.  The main text
contains the theorem statements and dependencies; these countermodels prevent
GP3 or functoriality from being used beyond their proved scope.

\begin{proposition}[Separated descendants alone do not select the parent-derived Abelian ray]
\label{p5:prop:hypercharge-ray-countermodel}
If one discards the generated parent relation used by PCAD and retains only
the already separated finite descendants in the Abelian candidate plane, then
for every nonzero \((a,b)\in\mathbb R^2\) the line
\[
\mathfrak k_{a,b}:=\mathbb R\bigl(aT_R^3+b(B-L)\bigr)
\]
commutes with the generated \(SU(2)_L\) and \(SU(3)_c\) actions, is generation
blind, and lies in the coordinate-conjugate Abelian plane.  Thus the
\emph{fixed-descendant reduct} alone does not select the hypercharge ray.  This
does not contradict Proposition~\ref{p5:prop:v187-pcad-yfix}, which retains
the common generated parent quadratic before descendant separation and uses
PCAD to descend its stabilizer.
\end{proposition}
\begin{proof}
The two generators act on different commuting factors and both act as the
identity on \(G\).  Lemma~\ref{p5:lem:intrinsic-abelian-centralizer-plane}
shows that every displayed line satisfies the scalar/conjugate covariance
condition.  Hence forgetting the parent relation leaves models with every
\(\mathfrak k_{a,b}\).  The PCAD theorem supplies exactly the omitted
parent-to-descendant relation and therefore is outside this reduct.
\end{proof}

\begin{proposition}[The separate gradings do not select relative parity]
\label{p5:prop:relative-parity-countermodel}
Remove \(({\rm CHI}_5)\) while retaining the external and internal graded
carriers.  Then both \(P_+\) and \(P_-\) define equal-dimensional invariant
co-realizations.  Simultaneous orientation reversal preserves both.
Consequently no invariant of the unbridged graded reduct distinguishes the
two relative parities.
\end{proposition}

\begin{proof}
The projector, rank, commutant, and orientation-reversal calculations are
those of Theorem~\ref{p5:thm:generated-orientation-coherent-chirality}.  Since
the positive and negative eigenspaces in each factor have equal dimension,
there exist nondegenerate even and odd graded isomorphisms
\[
J_+\gamma_x=\Gamma_WJ_+,
\qquad
J_-\gamma_x=-\Gamma_WJ_-.
\]
They define the \(P_+\) and \(P_-\) fibre products, respectively.  Forgetting
the sign of the bridge gives the same separate graded reduct in both cases,
so GP3 has no host-faithful deletion that selects one of them.
\end{proof}

\begin{proposition}[Conjugate availability does not imply direct-sum occurrence]
\label{p5:prop:conjugate-occurrence-countermodel}
Let the finite structural language contain \(F\), its generated operations,
and the functorial conjugation assignment \(F\mapsto\bar F\), but no
occurrence predicate or finite real-structure constructor joining them.  This
language has an expansion in which only the \(F\)-sector occurs and an
expansion in which both sectors occur.  Their finite structural reducts are
identical.
\end{proposition}

\begin{proof}
Functorial conjugation defines the object \(\bar F\) and transports every
linear representation map to its conjugate.  It makes no assertion about an
occurrence predicate.  Interpreting that predicate as \(\{F\}\) or as
\(\{F,\bar F\}\) changes no equation in the reduct.  Hence simultaneous
occurrence is not a first-order consequence of the finite functorial data.
An antiunitary PCT theorem on \(\mathsf R_5^{\rm PCT}\), or an explicitly
supplied finite real structure, adds the missing constructor.
\end{proof}



\subsection{Status of the finite response-count mnemonic}
\label{p5:app:finite-response-complex}

The mnemonic \(12-1=11\) remains a diagnostic and is not the derivation of
the vector coefficient.  The actual proof is now the explicit generated
background-gauge response complex of Paper~III-F together with the finite
determinant powers of Paper~IV and the heat-boundary calculation in
Theorem~\ref{p5:thm:part4-standard-one-loop-bridge}.  In particular,
\[
 \frac{11}{12}=\frac56+\frac1{12}
\]
is the one-form plus gauge-orbit quotient heat response, and the universal
conversion \(b=-4\omega h\) gives \(-11/3\).  The quotient scalar is auxiliary
Jacobian data, not an additional physical particle species.

\subsection{common electroweak Hessian}
\label{p5:sec:v164-common-ew-hessian}

\begin{theorem}[Charged and neutral responses have one action owner]
\label{p5:thm:v164-common-ew-hessian}
On the derived joint electroweak connection and a supplied nonzero Higgs orbit
\(h=2^{-1/2}(0,v)^T\), the physical mass operator is
\[
 \mathcal M_{\rm EW}
 :=\left.\delta_A^2S_H[A,H]\right|_{A=0,H=h}.
\]
Its charged and neutral restrictions give
\[
 m_W^2=\frac14g_2^2v^2,
 \qquad
 m_Z^2=\frac14(g_2^2+g_Y^2)v^2,
 \qquad
 m_\gamma^2=0.
\]
The \(W\), \(Z\), and photon directions are spectral subspaces of this one
Hessian; they do not require separate response normalizations.
\end{theorem}

\begin{corollary}[Source-native common-heat ratio]
\label{p5:cor:v164-source-native-ew-ratio}
At the common-heat matching datum,
\(g_Y^2/g_2^2=3/5\), and therefore
\[
 \boxed{\frac{m_W^2}{m_Z^2}=\frac58.}
\]
The finite \(7:9\) Gram relation remains a counting-metric theorem, but its
promotion to equal physical shell amplitudes is incompatible with this
same-action Hessian.  A later scale at which running couplings happen to give
\(7/9\) is an RG crossing, not shell ancestry.
\end{corollary}

\subsection{Supplementary same-moment family countermodel}
\label{p5:app:v181-same-moment-family-countermodel}
The local scalar action does not reconstruct Yukawa source ancestry from the
two moments it consumes.  For example, the positive family operators
\[
 X=\operatorname{diag}(1,2,3),\qquad
 Y=\operatorname{diag}\!\left(\frac{17}{7},\frac67,\frac{19}{7}\right)
\]
have different spectra but satisfy
\[
 \operatorname{tr}X=\operatorname{tr}Y=6,\qquad
 \operatorname{tr}X^2=\operatorname{tr}Y^2=14.
\]
Taking square roots as Yukawa singular-value operators therefore gives the
same scalar $a_Y,b_Y$, canonical quartic and local radial Hessian while
retaining inequivalent family spectra/provenance.  Weak gauge operations act
as the identity on $G$, so the same weak covariance/Noether identities hold
for both.  This countermodel proves only the source-selection boundary: an
already selected exclusion source is retained by the action, but scalar
moments or weak Noether identities cannot manufacture that missing G15
source word.


% BEGIN OWNER-COMPLETION v1.84 AUDIT CORRECTION
\subsection{contraction of the \texorpdfstring{\(D_4\)}{D4} neutral occurrence gate}
\label{p5:sec:v374-d4-neutral-source-contraction}

Paper~I-A v5.60 now co-realizes the positive central determinant lift on the
same rooted trimer that supplies the physical family carrier
\(G\cong\CC[C_3]\).  Consequently the determinant return, the cube root and
the rank-three family lift are no longer constituents of a new physical
occurrence gate.  Only consumption of that already generated operator by the
actual neutral source remains physical.

\begin{definition}[Generated \(D_4\) family response]
\label{p5:def:v374-generated-d4-family-response}
On the generated family carrier \(G\) put
\[
 \boxed{R_{D_4}^{\rm fam}:=xI_G,\qquad x:=192^{1/3}.}
\]
This is imported as generated mathematical owner content from Paper~I-A
v5.60.  It is not a numerical Yukawa member and does not by itself assert a
nonzero neutrino mass edge.
\end{definition}

\begin{definition}[Contracted terminal-neutral source-consumption relation]
\label{p5:def:v374-n56-source-consumption}
Let
\[
 \boxed{\mathsf{N56}_{D_4}^{\rm consume}}
\]
mean that the actual same-action \(K=56\) neutral fermion source/member uses
the already generated response \(R_{D_4}^{\rm fam}\) through the terminal
neutral response word declared in Paper~V-B.  This relation contains no
additional determinant count, cube-root choice, family-carrier occurrence, or
central-lift datum.
\end{definition}

\begin{proposition}[The v3.73 occurrence gate contracts to source consumption]
\label{p5:prop:v374-n56-gate-contraction}
Relative to the generated Paper~I-A v5.60 family response, the mathematical
content of \(\mathsf{N56}_{D_4}^{\rm occ}\) is discharged.  Its remaining
physical content is exactly
\(\mathsf{N56}_{D_4}^{\rm consume}\), together with the independent
normalization and relative-eigenframe obligations owned downstream.
Moreover this consumption relation is a component condition on the same
weighted Yukawa--Dirac source member
\((Y_u,Y_d,Y_e,Y_\nu)\in\mathfrak Y\) whose physical occurrence is already the
G15 source/member frontier; it is not a second independent source-occurrence
cluster.
\end{proposition}

\begin{proof}
Theorem~\ref{p5:thm:generation-selection} generates the physical family
carrier \(G\) from the complementary oriented trimer, while Paper~I-A v5.60
generates the central operator \(R_{D_4}^{\rm fam}\) on that same carrier.
The generated Yukawa--Dirac owner theorem
\ref{p5:thm:v186-generated-yukawa-dirac-owner-family} places the neutral
family operator in the \(Y_\nu\in\End(G)_\nu\) coordinate of one common source
family.  Therefore a further occurrence token cannot be needed merely to
recreate \(G\), \(192\), or its central cube-root lift.  What remains is the
member/source relation asserting that the actually occurring neutral block
uses this generated response.  Since that block is one coordinate of the
same source tuple, the relation refines rather than duplicates the G15
source-member occurrence problem.
\end{proof}

\begin{remark}[No hidden selection claim]
\label{p5:rem:v374-no-hidden-selection}
The proposition does not prove that the physical source member satisfies
\(\mathsf{N56}_{D_4}^{\rm consume}\).  The complete owner family
\(\mathfrak Y\) remains generated without selecting a member.  Thus v3.74
removes a duplicated mathematical occurrence only; it does not infer physical
source consumption from operator availability.
\end{remark}
% END D4-NEUTRAL-SOURCE-CONTRACTION v3.74


% BEGIN D4-TRIALITY-SOURCE-CONSUMPTION v3.75
\subsection{rooted half-step frontier after triality-source closure}
\label{p5:sec:v376-d4-rooted-halfstep-frontier}

The v3.75 theorem discharges the former same-action primitive-owner occurrence
gate.  Paper~I-A v5.62 now supplies two additional facts: the determinant lift
is centrally forced by same-owner equivariance, and the primitive rooted axial
triality contraction is the identity on the complementary trimer.  Therefore
the remaining finite noncentral source problem is not an arbitrary
\(M_3(\CC)\) member selection.

\begin{proposition}[Even rooted transport is generated]
\label{p5:prop:v376-even-rooted-transport-generated}
For the rooted antiprism trimers the complete unoriented cross-incidence is
\[
 J-I=X_G+X_G^\dagger=:C_{1,G}.
\]
Hence the orientation-even rooted transport is generated by the same finite
incidence geometry.  The primitive axial triality contraction contributes the
identity matching \(I_G\).
\end{proposition}

\begin{definition}[Residual half-step source gate]
\label{p5:def:v376-halfstep-source-gate}
Let
\[
 C_{2,G}:=-i(X_G-X_G^\dagger)
\]
be the unique orientation-odd trimer line.  The remaining finite G15 source
gate is
\[
 \boxed{\begin{gathered}
 \mathsf{HS}_{D_4}^{\rm source}:\ \text{the claim-complete rooted}\\
 \text{source consumes an oriented half-step descendant,}\\
 \text{equivalently }C_{2,G}.
 \end{gathered}}
\]
This is a finite marked-transport question.  It is not a renewed occurrence
problem for the triality parent, rank-\(2,4,6\) owners, or a continuous family
matrix.
\end{definition}

\begin{corollary}[Retyping of the G15 finite frontier]
\label{p5:cor:v376-g15-halfstep-retyping}
After v3.75--v3.76, the finite source ancestry, primitive owner ranks, axial
identity matching and orientation-even transport are generated.  The only
remaining primitive noncentral rooted transport is the sign line
\(\RR C_{2,G}\), or equivalently the choice/use of the two oriented
half-step maps \(X_G\) and \(X_G^\dagger\).  Once a claim-complete source word
contains such an ordered half-step return, the existing positive polar theorem
fixes its normalized sign-line coefficient; no continuous half-step amplitude
is reopened.
\end{corollary}

\begin{remark}[Downstream functional calculus is not a competing source port]
\label{p5:rem:v376-feshbach-downstream}
Feshbach/Schur reduction remains an action-owned descendant of an already
fixed finite source word.  It is not an alternative primitive source selector
competing with the rooted triality incidence.  Its remaining task is to
identify the actual paired complement and its normalization on the relevant
claim domain.
\end{remark}
% END D4-ROOTED-HALFSTEP-FRONTIER v3.76



% BEGIN ROOTED-SOURCE-ORIENTATION-COVARIANCE v3.78
\subsection{retirement of the geometric \texorpdfstring{$HS_{D_4}$}{HS-D4} ambiguity}
\label{p5:sec:v380-hsd4-retyping}

Paper~I-A v5.63 and Paper~I-C v1.28 now identify the two rooted half-step
transports with the finite inversion pair of the same-root $q=5$ rotation
plane, while Paper~II-C v2.69 identifies that inversion with the induced
finite Born action.  The former question ``derive or choose the sign line'' is
therefore no longer a live G15 geometric residual.

\begin{theorem}[Geometric closure of the D4 half-step sign line]
\label{p5:thm:v380-hsd4-geometry-closed}
On the real trimer augmentation plane $G_\perp^{\mathbb R}$ there is an exact
same-root isometry
\[
 \Psi_u:G_\perp^{\mathbb R}\xrightarrow{\cong}u^\perp
\]
intertwining the rooted cyclic shift and the $q=5$ rotation,
\[
 \Psi_uX_G=R_u\Psi_u.
\]
Writing
\[
 J_u=\frac{R_u-R_u^{-1}}{\sqrt3},
\]
one has
\[
 \boxed{
 C_{2,G}=-i(X_G-X_G^\dagger)
 =-i\sqrt3\,\Psi_u^{-1}J_u\Psi_u.}
\]
The finite Born involution acts by
\[
 X_G\leftrightarrow X_G^\dagger,
 \qquad C_{2,G}\longmapsto-C_{2,G}.
\]
Hence $\mathbb RC_{2,G}$ is a generated and uniquely normalized geometric
line; no independent family-orientation selector or continuous sign-line
coefficient remains.
\end{theorem}

% END HSD4-BORN-Q5-RETYPE v3.80

% BEGIN G15-C3-SOURCE-CELL v3.81
\subsubsection{Optional finite-bimodule interface}
\label{p5:subsec:ncg-assumption-discharge}

\begin{definition}[Optional conjugate-complete finite bimodule datum]
\label{p5:def:born-hopf-finite-spectral-datum}
Supply a conjugate-complete carrier
\(\mathcal H_F:=F\oplus\bar F\) as part of this optional branch.  An
almost-commutative comparison may
add a real structure \(J_F\), grading \(\Gamma_F\), represented finite algebra,
and a self-adjoint odd \(D_F\) supported on the reconstructed Yukawa blocks
and any optional Majorana block.  Only the availability of the conjugate
representation is derived finitely; simultaneous direct-sum occurrence and
the full spectral-triple data \(J_F,\Gamma_F,D_F\) are supplied optional data
and are not required by the Standard--Model constitution theorem.  Their
optional realization atlas
is
\[
\mathfrak A_{\rm NCG}
:=\{(\mathcal H_F,J_F,\Gamma_F,D_F,\rho):
\text{the stated almost-commutative compatibility axioms hold}\}.
\]
Paper~V-A quantifies over this family; it does not select a physical member.
\end{definition}

\begin{proposition}[Almost-commutative compatibility]
\label{p5:thm:ncg-finite-input-discharge}
\label{p5:thm:born-hopf-inner-fluctuations}
The Paper--V-A carrier does not depend on the order-zero, first-order,
orientability, Poincare-duality, or reconstruction axioms.  For any supplied member of this optional realization atlas satisfying the standard almost-commutative hypotheses, the inner
fluctuation
\[
D_A=D+A+J_FAJ_F^{-1},
\qquad
A=\sum_i\rho(a_i)[D,\rho(b_i)],
\]
has finite gauge and scalar blocks only in the operator modules present in
that supplied realization, provided the order-zero, first-order, grading, and
reality conditions are verified directly.
\end{proposition}

\begin{proof}
The representation, charges, generated scalar module, Yukawa slots, and
neutral right line are fixed before this comparison.  The spectral-triple
identities must be checked directly for the supplied finite matrices.  The
proposition is therefore a compatibility statement, not a derivation of the
finite real structure or a proof that no other scalar module exists.
\end{proof}

\begin{remark}[Comparison boundary]
\label{p5:rem:ncg-theorem-use-firewall}
\label{p5:rem:ccm-novelty-boundary}
\label{p5:rem:ncg-inner-fluctuation-priority-boundary}
Almost-commutative geometry is an optional analytic realization of the
derived representation.  Paper~V-A neither replaces the general
reconstruction and heat-kernel theorems nor claims that GP2 implies their
finite algebraic hypotheses.
\end{remark}

\label{p5:part:qft-bridge}

Assume \(\mathsf R_5^{\rm SSQFT}\).  Paper~III-F supplies the generated
Yang--Mills action, its finite background-gauge quotient response complex, and
the common local heat calculus; Paper~IV fixes the relative bosonic and
paired-Grassmann determinant powers; Paper~I-C supplies the finite heat-shell
flow.  The analytic SSQFT domain still licenses the ordinary four-dimensional
local heat-kernel/background-field calculation.  What is no longer imported
is the numerical triplet of universal determinant weights: it is derived
below as the logarithmic response to displacement of the resolution boundary.
The old finite response-count mnemonic remains only a diagnostic comparison
(Appendix~\ref{p5:app:finite-response-complex}).


\section{Projectors, shell tables and finite spectral certifications}
\label{p5d:sec:C}
\noindent\textit{Technical owner.} These exact inherited calculations
support the corresponding statements in Papers~V-A--V-C;
the appropriate action/observable occurrence conditions remain in those papers.
\subsection{Primitive trimer as a local generic Dirac-flavor chart}
\label{p6:sec:v548-trimer-flavor-chart}

Let $\mathfrak T_{\rm tri}$ be the five-real-dimensional Hermitian operator
space generated in Paper~I-A,
Theorem~\ref{p1:thm:cyclic-trimer-operator-grammar}.  In its oriented cyclic
basis,
\[
 \mathfrak T_{\rm tri}
 =\operatorname{span}_{\RR}\{I,\DtriI,\DtriII,C_1,C_2\},
 \qquad
 C_1=X+X^\dagger,
 \quad C_2=-i(X-X^\dagger).
\]
Write $\mathfrak T_{\rm tri}^{+}$ for its positive-definite locus.

\begin{lemma}[Left-return reduction of the square Dirac quotient]
\label{p6:lem:v548-left-return-reduction}
On the full-rank square-Dirac quark branch, put
\[
 H_u^L=Y_uY_u^\dagger,
 \qquad
 H_d^L=Y_dY_d^\dagger.
\]
Then the classical quark weak-basis quotient is equivalently the simultaneous
unitary-conjugacy quotient
\[
 \boxed{
 (Y_u,Y_d)/\bigl(U(3)_Q\times U(3)_u\times U(3)_d\bigr)
 \cong
 (H_u^L,H_d^L)/U(3)_Q.}
\]
The analogous statement holds for the Dirac-neutrino lepton pair.
\end{lemma}
\begin{proof}
For full-rank $Y_s$, polar decomposition gives
$Y_s=(H_s^L)^{1/2}U_s$ with $U_s$ unitary.  The independent right-handed
basis action removes $U_s$, while the common left-handed basis action sends
$H_s^L\mapsto U_QH_s^LU_Q^\dagger$.  Conversely, two full-rank square
matrices with the same left return differ by a right unitary.  Hence the
fibres of the two quotient descriptions agree.
\end{proof}

\begin{proposition}[Equal-per-sector fixed-slice lower bound]
\label{p6:prop:v548-equal-slice-minimality}
Let $\mathcal S\subset\operatorname{Herm}_3^{+}$ be a fixed smooth
$d$-real-dimensional source slice used independently for each member of a
two-Dirac-sector pair.  If
\[
 \mathcal S\times\mathcal S
 \longrightarrow \mathcal Q_q^{\rm gen}
\]
is locally surjective onto the generic quark weak-basis quotient of
Theorem~\ref{p6:thm:weak-basis-flavor-dimensions}, then
\[
 \boxed{d\ge5.}
\]
Thus a five-dimensional source is dimension-minimal among equal per-sector
fixed smooth source slices.  This does not exclude asymmetric slices with
$d_u+d_d=10$, shared cross-sector variables, or additional dynamical
constraints.
\end{proposition}
\begin{proof}
The generic quark quotient has real dimension ten.  The differential of a
map from $\mathcal S\times\mathcal S$ has rank at most $2d$, so local
surjectivity requires $2d\ge10$.
\end{proof}

\begin{theorem}[Primitive-trimer local submersion onto generic quark flavor]
\label{p6:thm:v548-trimer-local-submersion}
The quotient map
\[
 q:\mathfrak T_{\rm tri}^{+}\times\mathfrak T_{\rm tri}^{+}
 \longrightarrow\mathcal Q_q^{\rm gen}
\]
has rank ten at an explicit positive, simple-spectrum, scalar-stabilizer
point.  In particular the Paper-I-A primitive trimer attains the lower bound of
Proposition~\ref{p6:prop:v548-equal-slice-minimality} and supplies a local
full-rank coordinate chart for generic three-family quark flavor.

One exact witness is
\[
 A_*=10I+\DtriI+2\DtriII+C_1+C_2,
 \qquad
 B_*=12I+2\DtriI-\DtriII+2C_1+3C_2.
\]
For
\[
\begin{aligned}
F(A,B):=(&\operatorname{tr}A,\operatorname{tr}A^2,\operatorname{tr}A^3,
\operatorname{tr}B,\operatorname{tr}B^2,\operatorname{tr}B^3,\\
&\operatorname{tr}AB,\operatorname{tr}AB^2,
\operatorname{tr}A^2B,\operatorname{tr}A^2B^2),
\end{aligned}
\]
exact differentiation in the ordered trimer coefficients gives
\[
 \boxed{
 \det\frac{\partial F}
 {\partial(a_0,\ldots,a_4,b_0,\ldots,b_4)}\bigg|_{(A_*,B_*)}
 =-159690424320\ne0.}
\]
Moreover
\[
 \boxed{
 \chi(A_*,B_*):=\frac1{6i}\operatorname{tr}[A_*,B_*]^3=378\ne0,}
\]
so the full-rank image contains a nonempty open CP-capable neighborhood rather
than lying in the commuting/CP-conserving locus.
\end{theorem}
\begin{proof}
The exact certificate in
Appendix~\ref{p6:app:v548-trimer-local-chart-certificate} shows that $A_*$ and
$B_*$ are positive definite, have simple spectra, and have only the scalar
common commutant, so the witness lies on the generic quotient stratum.  The
ten functions in $F$ are invariant under simultaneous conjugation and hence
factor as $F=\bar F\circ q$.  Their displayed Jacobian has rank ten.  Since
Theorem~\ref{p6:thm:weak-basis-flavor-dimensions} gives
$\dim_{\RR}\mathcal Q_q^{\rm gen}=10$,
\[
 10=\operatorname{rank}d(\bar F\circ q)
 \le\operatorname{rank}dq\le10.
\]
Thus $\operatorname{rank}dq=10$.  Nonvanishing of $\chi$ and positivity are
open conditions, proving the final statement.
\end{proof}

\begin{corollary}[Twenty-dimensional Dirac-neutrino flavor capacity]
\label{p6:cor:v548-dirac-twenty-dimensional-trimer-chart}
Using independent copies of $\mathfrak T_{\rm tri}^{+}$ for
$u,d,e,\nu$, the product source
\[
 (\mathfrak T_{\rm tri}^{+})^4
 \longrightarrow\mathfrak A_{\rm flav}^{D}
\]
has rank twenty at a product of generic witness points.  Hence the primitive
trimer provides exactly enough local geometric coordinates to cover the
generic twenty-dimensional Dirac-neutrino flavor quotient.
\end{corollary}
\begin{proof}
The quark and Dirac-lepton weak-basis actions are independent, and each pair
has the ten-dimensional local-submersion result of
Theorem~\ref{p6:thm:v548-trimer-local-submersion}.  The differential of the
product map is block diagonal and has rank $10+10=20$, matching
Theorem~\ref{p6:thm:weak-basis-flavor-dimensions}.
\end{proof}

\begin{remark}[Capacity, covariance and occurrence firewall]
\label{p6:rem:v548-trimer-capacity-occurrence-firewall}
These results prove local generic \emph{capacity}, not occurrence or
prediction of twenty coefficients.  A covariant family of rooted trimer
returns need not be individually invariant under the residual $C_3$ action;
therefore this local chart does not contradict
Theorem~\ref{p6:thm:cyclic-natural-family-circulant}, which concerns a single
unpointed $C_3$-natural selected operator.  The finite rooted-transport gate is now closed upstream by the Paper-V-A
rooted $C_3$ source-cell co-realization theorem: the source owner retains the
complete presentation $\{I_G,X_G,X_G^{-1}\}$.  The local chart below therefore
tests relative flavor-frame/member dynamics only; it does not select matter
versus antimatter orientation and does not reopen a source-transport selector.
No $2I$ incidence-return
CP-depth claim is used in the proof above.
\end{remark}

\begin{remark}[Coefficient-free shell-return candidate and its scope]
\label{p6:rem:v562-hurwitz-clifford-candidate-scope}
Paper~I-A now supplies a rooted phase and exact shell--trimer return algebra,
so one can test a strictly coefficient-free balanced Hurwitz--Clifford family
without inserting observed mass data.  Appendix~\ref{p6:app:v562-hurwitz-clifford-mass-candidate}
computes its complete characteristic invariants and shows that relative
cyclic alignment is CP-capable on the anchored \(\mathrm{REG}_0\) hierarchy.
The same appendix proves the limiting obstruction relevant to G15: nonzero
cyclic offsets of one fixed phase-refined return create a leading singular
pair of the same hierarchy order rather than transporting one copy of each
of \(E_{32},E_{36},E_{40}\).  It also records a sharper diagnostic: the
unique sign-type direction in the primitive family transport plane can break
the pure-conjugation spectral degeneracy while preserving the three rooted
hierarchy orders for every full-rank coefficient.  The v5.64 appendix then
separates, by positive polar decomposition, the sector-dependent incidence
modulus from the oriented family transport on the ordered-return branch.  On
that branch the normalized sign direction is coefficient one in both the
lepton and color sectors; the earlier \(\lambda_q\)-family is retained only
as a diagnostic deformation family.  This does \emph{not} prove that the
physical Yukawa source consumes the polar-normalized ordered return.
Consequently the action-owned sector-specific source occurrence remains
open.
\end{remark}
% END FLAVOR-QUOTIENT v5.48

\subsection{Packet assignment and conditional raw ledger}
\label{p6:sec:conditional-packet-ledger}

\begin{remark}[Lookup versus physical slots]
\label{p6:rem:selected-slot-map}
Definition~\ref{p6:def:selected-typed-rows} and
Table~\ref{p6:tab:conditional-particle-ledger} together form the selected
lookup component of \(\mathsf{SLOT}\).  Reused supports are read by different
functionals and are not asserted to be disjoint particle projectors.
\end{remark}

\begin{table}[H]
\centering
\caption{Terminal-support and comparison-grade ledger; the nine charged grades are derived in Release v5.91, while other rows retain their stated status.}
\label{p6:tab:conditional-particle-ledger}
\scriptsize
\setlength{\tabcolsep}{3pt}
\renewcommand{\arraystretch}{1.10}
\begin{tabularx}{\textwidth}{@{}L{0.07\textwidth}L{0.13\textwidth}
L{0.31\textwidth}L{0.13\textwidth}Y@{}}
\toprule
Mode & Shell & Raw terminal readout & \(n_{\etaH}\) & Status \\
\midrule
\(g,\gamma\) & --- & gauge carrier / null projection
& \(0\) & declared gauge-null lookup; physical gauge
interpretation downstream \\
\(W\) & \(E_{32}\) & \(\sqrt7\)
& \(1/2\) & weak node \\
\(Z\) & \(E_{32}\) & \(\sqrt9\)
& \(1/2\) & weak node \\
\(H\) & \(E_{32}\) & \(\sqrt{16+1}\)
& \(1/2\) & scalar common column \\
\(t\) & \(E_{32}\) & \(\sqrt7+\sqrt9\)
& \(1/2\) & positive path \\
\(b\) & \(E_{32}\) & \((\sqrt8+\sqrt{16}+1)^{-1}\)
& \(-3/2\) & compound inverse gate \\
\(e\) & \(E_{40}\) & \((1+\sqrt5)^{-1}\)
& \(-3/2\) & inverse boundary gate \\
\(u\) & \(E_{40}\) & \(\sqrt2\)
& \(1/2\) & terminal node \\
\(d\) & \(E_{40}\) & \(\sqrt2+\sqrt3\)
& \(1\) & positive path \\
\(\mu\) & \(E_{36}\) & \(1/\sqrt4\)
& \(1\) & selected inner endpoint \\
\(s\) & \(E_{36}\) & \(1/\sqrt6\)
& \(-1\) & selected exterior endpoint \\
\(c\) & \(E_{36}\) & \(3\sqrt4\)
& \(1\) & positive path \\
\(\tau\) & \(E_{36}\) & \(2\sqrt6+\sqrt{12}\)
& \(1\) & positive path \\
\(\nu_1\) & \(E_{56}\) & \(1/\sqrt8\)
& \(1/2\) & inverse common column \\
\(\nu_2\) & \(E_{56}\) & \(\sqrt5\)
& \(1/2\) & scalar common column \\
\(\nu_3\) & \(E_{56}\) & \(13\)
& \(1/2\) & positive path \\
\bottomrule
\end{tabularx}
\end{table}

\begin{remark}[Coverage boundary]
Table~\ref{p6:tab:conditional-particle-ledger} is total on its declared
particle names but is not a common projector resolution.  The radical
operations and comparison exponent assignments are exactly the data whose
omission produced
the countermodels in Theorem~\ref{p6:thm:positive-dimensional-obstruction}.
\end{remark}

\subsection{Quark contrast and correction order}
\label{p6:sec:quark-contrast-corrections}

Let
\[
R_q=\frac7{16}\DfamI-\frac9{16}\DfamII,
\qquad
D_\pm=\frac{\DfamI\pm\DfamII}{\sqrt2}.
\]

\begin{theorem}[Exact mixed-parity decomposition]
\label{p6:thm:quark-mixed-parity}
Under the formal exchange
\(\DfamI\leftrightarrow\DfamII\),
\[
R_q
=
\frac1{\sqrt2}D_-
-\frac1{8\sqrt2}D_+.
\]
Thus \(R_q\) has a nonzero odd part and a nonzero even part.  It is not odd
under this exchange, under \(P_1\leftrightarrow P_2\), or under either of the
other trimer endpoint transpositions.
\end{theorem}

\begin{proof}
Insert
\[
\DfamI=\frac{D_++D_-}{\sqrt2},
\qquad
\DfamII=\frac{D_+-D_-}{\sqrt2}.
\]
For the literal \(P_1\leftrightarrow P_2\) exchange,
\(\DfamI\) is even and \(\DfamII\) is odd.  Solving the odd-eigenvector
condition for the other two transpositions gives coefficient lines
proportional to \((3,-1)\) and \((3,1)\), neither of which contains
\((7,-9)\).
\end{proof}

\begin{corollary}[Retained quark status]
\label{p6:cor:quark-contrast-status}
The \(7|9\) object may be retained only as a \emph{weighted signed contrast}.
Its two displayed coefficients are coordinates of a selected mixed-parity
operator, not consequences of oddness.  The quark sector still requires
explicit \(H_u,H_d\) before six masses or mixing data are defined.
\end{corollary}

\begin{theorem}[Finite rank does not imply terminality]
\label{p6:thm:no-rank-terminality}
Finite rank and projector idempotence do not force long correction words to
vanish.  For example,
\[
P=
\begin{pmatrix}1&0\\0&0\end{pmatrix},
\qquad
Q=\frac12
\begin{pmatrix}1&1\\1&1\end{pmatrix}
\]
are rank-one orthogonal projectors, but
\[
(PQ)^n=2^{-(n-1)}PQ\ne0
\qquad(n\ge1).
\]
\end{theorem}

\begin{proof}
Direct multiplication gives \((PQ)^2=\frac12PQ\); induction proves the
formula.
\end{proof}

\begin{proposition}[What finite dimension supplies]
\label{p6:prop:finite-recurrence}
The algebra generated by finitely many projectors on \(V\) has dimension at
most \((\dim V)^2\).  Word spans therefore stabilize and yield finite
recurrences.  This permits resummation after convergence is proved, but does
not make higher scalar coefficients vanish.
\end{proposition}

\begin{theorem}[Optional exterior terminality]
\label{p6:thm:exterior-terminality}
If a separately declared correction grammar requires
\[
C_{\chi,k}\in\Lambda^kW_\chi,
\qquad
\dim W_\chi=r_\chi,
\]
then \(C_{\chi,k}=0\) for \(k>r_\chi\).  This exact truncation is a theorem of
the exterior grammar, not of ordinary projector multiplication.
\end{theorem}

\begin{proof}
\(\Lambda^kW_\chi=0\) for \(k>\dim W_\chi\).
\end{proof}

\begin{convention}[Retained-order convention]
\label{p6:conv:retained-order}
This reconstruction does not add the exterior grammar to the physical mass
operator.  Every truncated upstream center, shell, or defect series---in
particular the charged \([4]\)-germ---is therefore tagged as a retained-order
bridge expression with an unresolved exact completion.  Exact unitary
matrix identities and representatives obtained from already declared CKM or
PMNS inputs are exact conditional algebra, not truncated series.  No omitted
term of a genuinely truncated input is asserted to vanish.
\end{convention}

\subsection{Selected finite numerical coordinate ledger}
\label{p6:sec:numerical-ledger}

\begin{definition}[Purpose of the numerical ledger]
\label{p6:def:numerical-ledger-purpose}
The following table evaluates the source \(\mathrm{REG}_0\) member, clock map,
terminal grammar, and cell-label comparison exponents.  It is a reproducible coordinate
ledger for one member of \(\BridgeRaw\), not a benchmark fit and not a proof
that the upstream quotient is a singleton.  The charged-lepton rows are
support diagnostics; the selected retained-order charged absolute map is
Definition~\ref{p6:def:charged-absolute-map}, whose distinct output is
reported in Corollary~\ref{p6:cor:actual-charged-output}.  The quark rows are
native finite coordinates pending \(\mathsf{QRG}\).
The anchors are
\[
\begin{aligned}
E_{32}&=30.9845880638852\,{\rm GeV},&
E_{36}&=0.219578897269190\,{\rm GeV},\\
E_{40}&=0.001556092726698264\,{\rm GeV},&
E_{56}&=3.924776812825316\,{\rm meV}.
\end{aligned}
\]
\end{definition}

\begin{table}[H]
\centering
\caption{Selected overlap-normalized finite coordinates.}
\label{p6:tab:selected-finite-coordinates}
\scriptsize
\hbadness=10000
\renewcommand{\arraystretch}{1.12}
\setlength{\tabcolsep}{3pt}
\begin{tabularx}{\textwidth}{@{}L{0.08\textwidth}L{0.22\textwidth}
L{0.42\textwidth}R{0.20\textwidth}@{}}
\toprule
Mode & Shell/status & Conditional formula & Coordinate \\
\midrule
\(g,\gamma\)& gauge-null & no massive finite readout & \(0\)\\
\(W^\pm\)& \(E_{32}\) &
\(E_{32}\etaH^{1/2}\sqrt7\)&\(80.321232\,{\rm GeV}\)\\
\(Z\)& \(E_{32}\) &
\(E_{32}\etaH^{1/2}\sqrt9\)&\(91.075717\,{\rm GeV}\)\\
\(H\)& \(E_{32}\) &
\(E_{32}\etaH^{1/2}\sqrt{17}\)&\(125.171600\,{\rm GeV}\)\\
\(t\)& \(E_{32}\), native &
\(E_{32}\etaH^{1/2}(\sqrt7+\sqrt9)\)&\(171.396949\,{\rm GeV}\)\\
\(b\)& \(E_{32}\), native &
\(E_{32}\etaH^{-3/2}(\sqrt8+\sqrt{16}+1)^{-1}\)&
\(4.207890\,{\rm GeV}\)\\
\(e\)& \(E_{40}\), diagnostic &
\(E_{40}\etaH^{-3/2}(1+\sqrt5)^{-1}\)&\(0.511224\,{\rm MeV}\)\\
\(u\)& \(E_{40}\), native &
\(E_{40}\etaH^{1/2}\sqrt2\)&\(2.156185\,{\rm MeV}\)\\
\(d\)& \(E_{40}\), native &
\(E_{40}\etaH(\sqrt2+\sqrt3)\)&\(4.700044\,{\rm MeV}\)\\
\(\mu\)& \(E_{36}\), diagnostic &
\(E_{36}\etaH/\sqrt4\)&\(105.397871\,{\rm MeV}\)\\
\(s\)& \(E_{36}\), native &
\(E_{36}\etaH^{-1}/\sqrt6\)&\(93.377822\,{\rm MeV}\)\\
\(c\)& \(E_{36}\), native &
\(3E_{36}\etaH\sqrt4\)&\(1.264774\,{\rm GeV}\)\\
\(\tau\)& \(E_{36}\), diagnostic &
\(E_{36}\etaH(2\sqrt6+\sqrt{12})\)&\(1.762902\,{\rm GeV}\)\\
\bottomrule
\end{tabularx}
\end{table}

\begin{corollary}[Selected finite pole coordinates]
\label{p6:cor:selected-pole-coordinates}
For the additional matrix choice
\(
b=1/120
\)
and ranks \((1,7,1)\),
Definition~\ref{p6:def:conditional-pole-packet} evaluates to
\[
M_W^{\rm mat}\simeq80.3690\,{\rm GeV},
\qquad
M_Z^{\rm mat}\simeq91.1872\,{\rm GeV},
\qquad
M_H^{\rm mat}\simeq125.2023\,{\rm GeV}.
\]
These are eigenvalue coordinates of the declared finite matrices.  The
superscript ``mat'' is retained because physical pole masses require the
downstream interaction and renormalization contract.
\end{corollary}

\subsection{Notation and provenance contract}
\label{p6:sec:notation}

\begin{table}[H]
\centering
\caption{Symbols whose former overloading is removed.}
\label{p6:tab:notation}
\small
\renewcommand{\arraystretch}{1.12}
\begin{tabularx}{\textwidth}{@{}L{0.20\textwidth}L{0.28\textwidth}Y@{}}
\toprule
Symbol & Meaning & Provenance/status \\
\midrule
\(\alpha_{\rm hier}\) & multiplicative hierarchy coordinate &
Derived positive determinant character; its numerical finite part is
conditional on a frozen \(\mathrm{REG}_{\sigma}\) member.  Its neutral dual
is generated; physical energy use requires \((\mathrm{CRP4})\).\\
\(K_{\rm sh}\) & shell-clock coordinate &
All seven scalar levels and the anchors \(32,40,56\) are protected upstream;
historical carrier and particle-role labels, when displayed, use the optional
\(\CLKrem\) and \(\mathsf{SLOT}_{\rm hist}\) atlas.\\
\(\etaH=24/25\) & Hurwitz cell-label exclusion &
Do not confuse with the quark Wolfenstein coordinate.\\
\(\eta_q\) & quark Wolfenstein coordinate &
Flavor invariant supplied by the retained CKM eigenframe class.\\
\(\delta_{\rm geom}\) & family/geometric defect &
Finite carrier coordinate; invariant under \(\mathrm{REG}_{\sigma}\) finite-part changes.\\
\(\sigma_{\rm fin}\) & determinant finite-part selector &
Frozen coordinate of \(\mathrm{REG}_{\sigma}\); source member is \(0\).\\
\(\delta_{\rm std}\) & standard CKM phase &
Not identical to the unitarity-triangle angle.\\
\(\delta_\ell\) & PMNS representative phase &
Coordinate of the supplied PMNS relative-eigenframe data.\\
\(\DfamI,\DfamII\) & centered family-plane operators &
Never used as regularization defects or clock variables.\\
\(M_P\) & series shell rest-energy coordinate &
\(8\hbar c_{\rm op}\sqrt U=8\sqrt{18/\pi}\,\hbar c_{\rm op}/R_{\rm mom}=16\overline M_{\rm Pl}\)
on the active G24 branch; not the reduced Planck mass.\\
\(\overline M_{\rm Pl}\) & reduced Planck mass &
Downstream gravitational and continuum normalization.\\
\(E_N,\Delta_M\) & heavy anchor and splitting &
Symbolic bridge data until assigned by a downstream owner.\\
\(M_B^{\rm mat}\) & finite matrix coordinate &
Not called a physical pole without \(\mathsf{PV}\).\\
\bottomrule
\end{tabularx}
\end{table}

\subsection{assumption supersession for hierarchy and comparison readouts}
\label{p6:sec:post-v125-supersession}

This section supersedes the remaining historical selector language where the
later assumption audit has proved a stronger reduction or no-go.  The finite
support, incidence, shell, radical-classification and weak-basis quotient
theorems are unchanged.  Any counterfamily below that varies a common absolute
reference energy is retained only as a pre-G24/historical minimality witness;
on the active v2.89 branch the native anchor is fixed by
\(E_0=\sqrt{18/\pi}\,\hbar c_{\rm op}/R_{\rm mom}\).

\subsubsection{Energy-owner weight and octahedral placement}
After quotienting the dimensionful reference energy, the general determinant
energy character is
\[
 \frac{E_\lambda(K)}{E_\lambda(0)}
 =\alpha_\sigma^{-\lambda K/32},\qquad \lambda>0.
\]
The former quarter-weight ancestry is the member $\lambda=1$.  Positivity,
monotonicity, $K=8q$, and the logarithmic midpoint
$E_\lambda(36)^2=E_\lambda(32)E_\lambda(40)$ hold for every $\lambda$, so none
selects the quarter weight.

\begin{theorem}[Natural cone automorphisms and the fourth-power density]
\label{p6:thm:physical-cone-fourth-power-density}
At $x\in S^3$, let
\[
 W_x^*=\mathbb R\oplus T_x^*S^3,
 \qquad
 \mathcal C_x^\vee
 =\{(E,p):E\ge c_{\rm op}\|p\|_{h^{-1}}\}
\]
be the generated Paper--III-A scalar-plus-spatial generator owner and its
positive polar cone.  If a real-linear automorphism $A$ is natural under
the full $O(3)$ presentation action and satisfies
$A(\mathcal C_x^\vee)=\mathcal C_x^\vee$, then there are uniquely
$s>0$ and $\varepsilon\in\{+1,-1\}$ such that
\begin{equation}
\label{p6:eq:natural-cone-automorphism-classification}
 A=s\,\operatorname{diag}(1,\varepsilon I_3).
\end{equation}
Consequently the induced action on the positive four-density line is
\begin{equation}
\label{p6:eq:positive-frame-density-fourth-power}
 |\det A|=s^4.
\end{equation}
If $A$ is additionally positive-definite, orientation-preserving, or in the
identity component, then $\varepsilon=+1$ and $A=sI_4$.
\end{theorem}

\begin{proof}
Full $O(3)$-equivariance kills the scalar--vector cross maps and makes the
spatial endomorphism scalar, so $A(E,p)=(aE,bp)$.  Since both $A$ and
$A^{-1}$ preserve the cone, $a>0$, $b\ne0$, and the exact boundary
$E=c_{\rm op}\|p\|_{h^{-1}}$ maps onto itself.  Hence $a=|b|=:s$, which
gives Eq.~\eqref{p6:eq:natural-cone-automorphism-classification}.  Taking the
absolute determinant gives $s^4$.  The branch
$s\operatorname{diag}(1,-I_3)$ is the parity control; each listed additional
condition excludes it.
\end{proof}

Theorem~\ref{p6:thm:physical-cone-fourth-power-density} derives the number
four on the physical first-order owner, but it does not select a determinant
character weight.  Appendix~\ref{p6:app:c13-density-and-counterfamily} gives
the exact positive \(\lambda\)-counterfamily and the independent
contact/filling four-density check.

\begin{theorem}[CRP4 import and C13 discharge]
\label{p6:thm:c13-detframe-closure-condition}
\label{p6:thm:crp4-c13-discharge}
Let $B$ be the scheme--depth base carrying the weight-one dual
determinant-depth associated line
$\bigl(\mathcal L_{\rm det}^{-1},\nabla^{\rm det}\bigr)$, and let $M$ be the
spacetime branch carrying
\(\mathcal D_4:=|\Lambda^4W^*|\).  With projections
\(\pi_B:B\times M\to B\) and \(\pi_M:B\times M\to M\), pull both lines to
$B\times M$.  The one
CRP4 occurrence of
Definition~\ref{p3e:def:crp4-cross-role-positive-parent} induces the
normalization-free associated-line intertwiner
\begin{equation}
\label{p6:eq:c13-detframe-common-base-map}
 \Theta_{\rm C13}:\pi_B^*\mathcal L_{\rm det}^{-1}
 \xrightarrow{\ \cong\ }\pi_M^*\mathcal D_4,
\end{equation}
and it intertwines the two induced connections.  If $H=e^{-\Delta\Phi}$ is
the fundamental determinant transport and $h$ is the scalar first-order
transport, then
\[
 h=H^{1/4},\qquad |\det(hI_4)|=h^4=H,
 \qquad
 \nabla^{\mathcal D_4}=d+d\Phi=\nabla^{\rm det}.
\]
Consequently
\begin{equation}
\label{p6:eq:c13-quarter-weight-consequence}
 \boxed{\lambda_{ED}=1},
 \qquad
 \frac{E(K)}{E(0)}=\alpha_\sigma^{-K/32}.
\end{equation}
Thus C13 is \texttt{DERIVED} on the generated CRP4 ancestry of the
occurring action root.  Paper~III-B's generated-action-response owner and
Paper~III-E's cross-role positive parent supply this antecedent; it is not
an additional independently stipulated source choice.  The historical Paper-V-B notation
\(\mathrm{ANC}_E\equiv\mathrm{CRP4}\) is a deprecated alias, not a second
assumption.  CRP4 by itself fixes no absolute energy value; independently,
the active G24 branch fixes \(E_0=\sqrt{18/\pi}\,\hbar c_{\rm op}/R_{\rm mom}\).
\end{theorem}

\begin{proof}
Paper~III-E,
Theorem~\ref{p3e:thm:crp4-uniqueness}, gives the generator block
$hI_4$ with $h=H^{1/4}$ and connection $d+\tfrac14d\Phi$.
Taking its positive determinant multiplies the connection form by four and
gives transport $h^4=H$, exactly the fundamental line carried by the CRP4
response summand after
Corollary~\ref{p3e:cor:crp4-c12-discharge}.  The parent-induced associated
intertwiner therefore realizes the former DETFRAME criterion without a free
gauge or normalization.  Finally $H=e^{-qL_\sigma}$ and $K=8q$, so the
first-order component is
$h=e^{-qL_\sigma/4}=\alpha_\sigma^{-K/32}$.
The counterfamily in
Appendix~\ref{p6:app:c13-density-and-counterfamily} proves that the conclusion
uses the CRP4 parent and is not a consequence of cone naturality or the number
four alone.
\end{proof}

The former octahedral $K=16$ placement is dropped from the native chain.  As a
$2O$-module,
\[
 \mathbb R[2O/2T]\cong\mathbf1\oplus\mathbf1_{\rm sgn},
\]
whereas the scalar-free protected Clifford carrier has type
$W\cong V_3\oplus V_3\oplus\mathbf1$.  The sign line is absent, so there is no
injective rank-two $2O$-equivariant ancestry into the hierarchy carrier.  The
canonical reduced augmentation quotient is rank one.  The neutral $K=16$ rung
and the quotient-odd heavy carrier remain separately derived, but their
identification is optional rather than native.

\subsubsection{Terminal rows, terminal calculus, and comparison exponents}
The six coarse Standard-Model role blocks are already fixed on the Paper-V-A
role domain.  Structural family projectors are not common mass projectors:
identifying them with all sector mass eigenframes would force commuting sector
operators and trivial mixing.  The historical terminal rows are external
abstract counting copies, not physical subspaces.  Their particle-name table
and optional embeddings are therefore a conditional comparison atlas unless a
later physical family/gauge/vacuum operator reproduces them spectrally.

On the counting supports themselves, the terminal calculus is canonical.  The
positive invariant monoid has unique atom $\Omega_S$, the counting Hilbert
metric is the unmarked metric, common-domain columns are the unique biproduct
tuples, the additive radical is the canonical codiagonal sum, and the
Moore--Penrose inverse is unique.  Relative phase twists require an inserted
diagonal unitary.  Thus no independent TMF law remains once a historical
constructor/list is specified.

The declared orthogonal-label comparison base $\eta_H=24/25$ and the
half-integer grading are theorem level on their typed owners, but physical
use of $\eta_H$ is conditional on the source-word occurrence in
Corollary~\ref{p6:cor:v181-conditional-higgs-exclusion-value}; the charged row exponent table is superseded by the derived v5.91 theorem below; noncharged rows retain their stated status.  Path length, inverse/direct polarity
and additive support cardinality all fail to reproduce the selected ledger.
A common shift of all massive-row exponents is absorbed into the common
reference energy, so only relative exponents matter on the optional numerical
branch.  A theorem-level derivation would require a generated connected
comparison process with an exact additive $\tfrac12\mathbb Z$ cocycle.

\subsubsection{Clock-role semantics}
The generated scalar depth and active anchors remain
$K=32,40,56$.  The value $K=36$ is a derived logarithmic midpoint, not a
generated event: throughout $4<q\le5$ the canonical upper spectral projector
is the same $\Pi_3$.  The historical affine antipodal rule assigning an
interface to $K=36$ is therefore retained, if desired, only as a conditional
role atlas.  A physical interface theorem would require a generated transition
object with an endpoint-exchange involution and a fixed-locus occurrence.

\subsubsection{Scalar coordinate and flavor target}
After canonical field normalization, $\kappa=v/E_{32}$ is a background readout,
not an additional scalar-normalization law.  The factor $\sqrt{24/15}$ is the
derived RMS conversion, while the historical $2\pi$ multiplier is one declared
Hopf-orbit readout among the source-proved alternatives.  No native $v$ is required by the electroweak constitution.  If a conventional single-doublet comparison coordinate is requested, $v_{\rm cmp}$ may be reconstructed from the completed low-energy response, for example $v_{\rm cmp}:=2m_W/g_2$, or evaluated on an explicitly chosen scalar-stationary branch.  In either use, $E_{32}$ is the native dimensionful owner and $\kappa_{\rm cmp}=v_{\rm cmp}/E_{32}$ is not a second scale law.

For generic quark flavor the raw right-handed frames are presentation.  The
low-energy physical quotient is represented by the simultaneous-conjugacy
class
\[
 \bigl[(H_u,H_d)\bigr]_{U(3)_{Q_L}},\qquad
 H_u=Y_uY_u^*,\quad H_d=Y_dY_d^*.
\]
The two positive spectra give six quark singular values and their relative
spectral frame gives CKM.  Theorem~\ref{p6:thm:v548-trimer-local-submersion} now proves that the
five-real-dimensional primitive trimer Hermitian slice is locally sufficient
for generic quark flavor and is dimension-minimal among equal per-sector
fixed smooth slices.  What remains is not additional matrix algebra but the
action-owned occurrence/selection of the sector-specific normalized trimer
coefficients (or equivalent invariant signature).  Appendix~\ref{p6:sec:family-quiver-update}
retains the distinct cyclic-naturality no-go for a single unpointed selected
operator.




% BEGIN HM v1.79 MS01
\subsubsection{Action identities required by the Hurwitz mass comparisons}
\label{p6:sec:v179-action-mass-comparison}

The constructive chain now contains the unital \(P_{17}\), cyclic
projectors, explicit \(E_{36}\) Clifford transition, and incidence polar
map.  To identify one of these sources with a particle row, the generated
comparison must preserve its field grading, gauge action, trace, source
connection, and all derivatives consumed by the same claim-complete action.
For a bosonic row, the compared operator is
\(Z^{-1/2}KZ^{-1/2}\); for a fermionic row it is
\(Z_L^{-1/2}D Z_R^{-1/2}\).  These are supplied on their stated branches by
Theorem~\ref{p3f:thm:v179-normalized-scalar-hessian} and
Corollary~\ref{p5:cor:v179-canonical-fermion-comparison}.

In particular, the historical comparisons in this paper are
\[
 m_H^{\rm comp}=\sqrt{17\etaH}\,E_{32},\qquad
 m_\mu^{\rm comp}=\etaH E_{36}/2.
\]
For the declared orthogonal-label self-exclusion owner,
Theorem~\ref{p6:thm:v181-claim-complete-self-exclusion-owner} gives
$\etaH=24/25$.  Using that value in a physical Higgs/Yukawa word is
conditional on Corollary~\ref{p6:cor:v181-conditional-higgs-exclusion-value}.
On the local canonical one-doublet and fermion-kinetic branch, equality
with those rows would require, at one common matching slice,
\begin{equation}
\label{p6:eq:v179-higgs-muon-action-targets}
 \lambda_Y=\frac{17\etaH E_{32}^2}{2v^2},\qquad
 \sigma_\mu(Y_{e,c})=\frac{\etaH E_{36}}{\sqrt2\,v}.
\end{equation}
These are necessary local matching identities, not selected values of
the flavor family.  A physical pole claim additionally consumes the
complete effective action, its stable background, and the declared
renormalization and comparison maps.  The half-integer comparison grading
and the number \(24/25\) do not determine the particle-row exponents.

GP1--GP3 retain the finite licensed source diagrams without an additional
G0 token; they do not identify different trace owners merely because their
dimensions or numbers agree.  The \(2O\) normalizer contribution to the
hierarchy already has the proved same-owner determinant map of
Theorem~\ref{p3c:thm:occurring-hurwitz-reg-origin}; it fixes
\(\alpha_{\rm hier}=\exp(4\pi^2+1/(6\sqrt2))\) on that branch.
The \(2I\) incidence and \(2T\) split sources now have explicit metrics,
while their physical action maps remain the stated target.
The existing flavor, vacuum and pole interfaces retain those obligations.
% END HM v1.79 MS01


% BEGIN SC v1.80 INCIDENCE-ACTION-BRIDGE
\subsubsection{Harmonic meaning of the rank-ten incidence restriction}
\label{p6:sec:v180-incidence-harmonic-action-bridge}

\begin{theorem}[Rank-ten incidence image is $\mathcal H_0\oplus\mathcal H_2$]
\label{p6:thm:v180-incidence-harmonic-image}
Let $U_B$ be the polar-normalized rank-$17$ incidence map of
Theorem~\ref{p6:thm:v179-incidence-polar-normalization}, let $R_H$ restrict
functions to one marked Hurwitz $2T$ cell, and let $P_0,P_1,P_2$ be the
sampled round harmonic projectors on that cell.  Then
\[
 \boxed{(R_HU_B)(R_HU_B)^\dagger=P_0+\frac4{15}P_2,\qquad
 \operatorname{im}(R_HU_B)=\operatorname{im}(P_0+P_2).}
\]
Thus the restriction has rank $1+9=10$, carries no degree-one harmonic, and
its inherited mode Gram is $\operatorname{diag}(1,(4/15)I_9)$.  Canonical
kinetic normalization multiplies the nine quadratic coordinates by
$\sqrt{15}/2$; the raw factor $4/15$ is not a physical mass coefficient.
\end{theorem}
\begin{proof}
For two Hurwitz-cell vertices $q,r$, with $t=q\cdot r$, the normalized
incidence kernel restricts to $(1+16t^2)/120$.  The round cell projectors are
$(P_0)_{qr}=1/24$ and $(P_2)_{qr}=(4t^2-1)/8$.  Substitution gives the stated
operator identity, whose positive support has ranks one and nine.
\end{proof}

\begin{corollary}[Canonical incidence--Hurwitz return]
\label{p6:cor:v180-incidence-hurwitz-polar-identity}
On the even ten-dimensional band, compress the unital Hurwitz projector
$P_{17}$ of Theorem~\ref{p6:thm:v179-hurwitz-unital-completion} between the
constant-plus-quadratic projectors.  Its response is one on four directions
and $1/2$ on six quadratic directions.  Composing it with the incidence map
gives kinetic Gram spectrum
\[
 1^{(1)}\oplus(4/15)^{(3)}\oplus(1/15)^{(6)}.
\]
After positive polar normalization the resulting source map is identical to
the canonical incidence polar map.  This equality is a same-image action
normalization statement, not an identification of the distinct parent label
constructions.
\end{corollary}

\begin{theorem}[Fourth-power scalar mass-operator readout]
\label{p6:thm:v180-fourth-power-scalar-mass-operator}
On the same local stationary Higgs branch and canonical fermion-pairing
specialization used in the existing mass-moment comparison,
\[
 \Delta V_H=\frac{N_{\rm eff}}{128\pi^2}m_{h,\rm loc}^4,
 \qquad N_{\rm eff}=S_2^2/S_4.
\]
Combining this with the position clock/kinetic identity and the exact scalar
mode action gives the operator form
\[
 \boxed{
 \Delta V_H I=\frac{N_{\rm eff}}{128\pi^2}
 \left(H_{h,\rm free}^2-\sum_iG_i^2\right)^2.}
\]
The spatial multiplicities of the $10$- or $14$-dimensional comparison
spaces do not multiply $N_{\rm eff}$, because $N_{\rm eff}$ is the internal
color-weighted chiral mass trace whereas the spatial modes use the normalized
round volume trace.

For the hierarchy readout,
\[
 E_K=E_0\alpha_{\rm hier}^{-K/32}
 \quad\Longrightarrow\quad
 E_K^4=E_0^4\alpha_{\rm hier}^{-K/8},
\]
which has the same hierarchy exponent as
$I_K=U^2\alpha_{\rm hier}^{-K/8}$.  This establishes the structural fourth-power action bridge.  The generated
line-owner map for $\Delta V_H/E_{32}^4$ is supplied below by
Theorem~\ref{p6:thm:v181-higgs-density-owner-map}; its numerical section
ratio is not selected.  No theorem here sets $\hbar c\Lambda_F=E_{32}$,
selects the physical Yukawa member/effective vacuum, proves the G15
self-exclusion source occurrence, or supplies a physical pole.
\end{theorem}

\begin{theorem}[Generated Higgs four-density on the CRP4 product base]
\label{p6:thm:v181-higgs-density-owner-map}
Let $\mathcal D_4=|\Lambda^4W^*|$ and let
$\Theta_{\rm C13}:\pi_B^*\mathcal L_{\rm det}^{-1}\to\pi_M^*\mathcal D_4$
be the connection-intertwining isomorphism on $B\times M$ of
Theorem~\ref{p6:thm:crp4-c13-discharge}.
On the selected local positive Higgs branch set
$s_H=\Delta V_H\,d\mu_L\in\Gamma(M,\mathcal D_4)$ and define
\[
 \widetilde\sigma_H:=\Theta_{\rm C13}^{-1}(\pi_M^*s_H)
 \in\Gamma(B\times M,\pi_B^*\mathcal L_{\rm det}^{-1}),\qquad
 \widetilde s_{32}:=\Theta_{\rm C13}(\pi_B^*e_{32}^{\otimes4}).
\]
The positive common-line ratio
\[
 \boxed{\widetilde\delta_{32}
 :=\frac{\pi_M^*s_H}{\widetilde s_{32}}>0},\qquad
 \widetilde\sigma_H=\widetilde\delta_{32}\,\pi_B^*e_{32}^{\otimes4}
\]
is exactly typed, with a positive first-order fourth root.
At a specified matching event $x_0\in M$, restriction along
$i_{x_0}:B\to B\times M$ produces a section of
$\mathcal L_{\rm det}^{-1}$ over $B$ and its local scalar ratio
$\delta_{32}^{(x_0)}(b)=\widetilde\delta_{32}(b,x_0)$.
A spacetime-independent $B$-section is \emph{not} implied by local
Higgs stationarity. The physical Yukawa member and numerical source
comparison remain separately selected.
\end{theorem}
\begin{proof}
Both pulled-back densities are sections of the same positive line on
$B\times M$; the CRP4 inverse and their positive ratio are therefore
well-defined. Applying the inverse to the ratio identity proves the
transport formula, and taking the positive fourth root uses the already
constituted first-order character line. Pullback along $i_{x_0}$
produces the claimed section. No descent over all of $M$ is used.
\end{proof}
\begin{proposition}[Exact global descent and spatial-mean alternative]
\label{p6:prop:v382-higgs-crp4-descent-criterion}
If $M$ is connected and $\Theta_{\rm C13}$ intertwines the indicated
pullback connections, then $\widetilde\sigma_H$ descends from $B$
if and only if $\nabla^{\mathcal D_4}s_H=0$ on $M$.
On a chosen round $S^3(R_x)$ slice $\Sigma$ the mean
$\overline\delta_{32}(b)=\int_\Sigma\widetilde\delta_{32}(b,x)\,d\nu_\Sigma$
defines a different, well-typed $B$-section and satisfies
\[
 \int_\Sigma|\widetilde\delta_{32}-\overline\delta_{32}|^2d\nu_\Sigma
 \le\frac{R_x^2}{3}\int_\Sigma|\nabla_\Sigma\widetilde\delta_{32}|^2d\nu_\Sigma.
\]
\end{proposition}
\begin{proof}
Intertwining gives
$\nabla^{\pi_B^*\mathcal L^{-1}_{\rm det}}_{(0,X)}\widetilde\sigma_H
 =\Theta_{\rm C13}^{-1}(\pi_M^*\nabla^{\mathcal D_4}_X s_H)$.
The source pullback connection is constant in $M$ directions exactly
when its section descends on connected $M$. The spatial mean is a scalar
in the common pulled-back fiber, hence defines the claimed section;
the sharp round-three-sphere scalar spectral gap $3/R_x^2$
gives the displayed Poincar\'e inequality.
\end{proof}


\begin{remark}[General physical-unit density ratio]
\label{p6:rem:v382-higgs-frame-comparison}
In any chosen positive density frame, write
$\widetilde s_{32}=E_{32}^4\gamma_{32}\pi_M^*d\mu_L$.
Then $\widetilde\delta_{32}=\Delta V_H/(E_{32}^4\gamma_{32})$;
the original numeric ratio with no $\gamma_{32}$ requires a common
matching frame with $\gamma_{32}=1$ at the evaluated event.
This factor is the coordinate of the already defined line-owner map,
not an independent physical parameter.
\end{remark}

\begin{corollary}[Two-coordinate reduction of the historical Higgs row]
\label{p6:cor:v181-higgs-two-coordinate-reduction}
At a selected matching event with a common physical-unit frame $\gamma_{32}=1$, with $\kappa=v/E_{32}$ and
$N_{\rm eff}$ from Corollary~\ref{p5:cor:v181-yukawa-participation-invariant},
\[
 \boxed{\delta_{32}^{\rm own}
 =\frac{2\pi^2}{N_{\rm eff}}\,\kappa^4.}
\]
If, in addition, the historical owner-relative row
$m_H^2=17(24/25)E_{32}^2$ is imposed on the same matching slice, then
\[
 \boxed{\kappa^2
 =\frac{17(24/25)}{16\pi^2}\,N_{\rm eff}.}
\]
Since $1\le N_{\rm eff}\le24$, this would require
\[
 \frac{\sqrt{102}}{10\pi}\le\kappa\le\frac{6\sqrt{17}}{5\pi}.
\]
These are exact necessary conditions, not a selected vacuum.
\end{corollary}


\begin{remark}[Independent frontier-energy scale]
\label{p6:rem:v191-frontier-e32-firewall}
\leavevmode\par\noindent
The hierarchy owner fixes
$E_{32}=\hbar c_{\rm op}\sqrt{R_p}/\alpha_{\rm hier}$ on its stated branch,
whereas the sharp position frontier owns the independent energy
$E_F=\hbar c_{\rm op}\Lambda_F$.  No current theorem identifies these two
scales.  The equality would require the additional dimensionful condition
$\Lambda_F=\sqrt{R_p}/\alpha_{\rm hier}$, which is not generated by the
comparison-band capture or the fourth-power action bridge.
\end{remark}


% END SC v1.80 INCIDENCE-ACTION-BRIDGE


% BEGIN v2.36 FLAVOR/LIFT/ANALYTIC ATLAS
\subsection{Flavor, neutral-lift and analytic comparison atlas}
\label{p6:sec:v236-flavor-lift-atlas}
\begin{proposition}[Fixed-coupling versus dynamical-flavor branches]
\label{p6:prop:v236-flavor-branches}
Let \(A=Y_uY_u^\dagger\) and \(B=Y_dY_d^\dagger\), and let
\(\Sigma(A,B)\) denote the complete generic weak-basis invariant packet
already constructed in this paper, including the CP orientation
\(\chi=(6i)^{-1}\operatorname{tr}[A,B]^3\).  Two logically distinct physical
branches are retained.
\begin{enumerate}[label=(\roman*),leftmargin=2.2em]
\item On the fixed-coupling branch, the normalized physical Yukawa/flavor
operator is a member of the generated owner family; its numerical orbit is not
selected by minimizing over coefficients that were never declared dynamical.
\item On a genuinely dynamical-flavor branch, a later claim-complete action may
vary the flavor coordinates.  The generic physical flavor orbit is uniquely
discharged exactly when \(\Sigma\) is constant on the retained stable solution
family modulo weak-basis equivalence.
\end{enumerate}
The invariant reconstruction theorem does not select a numerical flavor member.  Framework v2.86 separately closes the finite rooted transport cell; the remaining issue here is only whether downstream observables are constant on, or select within, the retained flavor family.
\end{proposition}

\begin{proposition}[Neutral-lift fibre and observable quotient]
\label{p6:prop:v236-neutral-lift-quotient}
On the high-energy neutral/Majorana branch, the compatible realizations form a
lift fibre \(\mathcal F_{\rm NU}\); a Casas--Ibarra matrix is a coordinate on
that fibre, not a second selector.  For a downstream observable
\(O:\mathcal F_{\rm NU}\to\mathcal O\), only the quotient
\[
 \mathcal F_{\rm NU}/\!\sim_O,
 \qquad R_1\sim_OR_2\Longleftrightarrow O(R_1)=O(R_2),
\]
needs to be fixed.  If \(O\) is constant on the compatible fibre, no physical
lift member is required for that claim.
\end{proposition}

\begin{remark}[Normalizer and global exceptional roles]
\label{p6:rem:v236-o-i-typing}
The octahedral normalizer \(2O=N_{SU(2)}(2T)\), its odd sign line and the
associated heavy carrier are generated descendants of the local tetrahedral
core.  By contrast, the five-core/25-cell \(2I\) geometry is a conditional
global refinement unless a theorem actually consumes a physical \(2I\)
support.  Pure comparison/incidence uses remain comparison geometry and do not
prove a global \(2I\) occurrence.
\end{remark}

\begin{remark}[RG transport versus physical pole]
\label{p6:rem:v236-qrg-pole-split}
The package \(\mathsf{QRG}\) transports an already owned mass/coupling operator
between renormalization coordinates and thresholds.  Once that analytic
package is fixed, the running quantity is generated.  Identification of a
physical pole additionally requires the renormalized two-point/inverse-response
continuation, sheet/domain, isolated-pole and residue conditions.  Finite
spectral rows, running masses and physical poles are therefore distinct typed
readouts of the same upstream physical operator.
\end{remark}
% END v2.36 FLAVOR/LIFT/ANALYTIC ATLAS


% BEGIN OWNER-LIFTED-CHARGED-COMPILER v5.88
\subsection*{canonical exterior comparison and preterminal compiler audit}
\label{p6:sec:v5124-exterior-comparison-compiler-audit}
The exact $24/25$ geometric factor is retained; the $E_{32},E_{36},E_{40}$
hierarchy anchors, existing charged numerical coefficients and
owner-incidence arrows are not refitted.  The material below
separates an independently proved quadratic comparison port from a
\emph{candidate} owner-to-exterior map needed to make its physical
particle-grade assignments action-derived.

\begin{theorem}[Action-band principal-angle transfer and exterior sorts]
\label{p6:thm:v5124-exterior-transfer}
The normalized $14$-mode Hurwitz evaluation frame of Paper~I-C,
Corollary~\ref{p1b:cor:v221-hurwitz-principal-angle}, gives two
isometries with scalar cross amplitude
\[
 t=\sqrt{\eta_H}=\frac{2\sqrt6}{5},\qquad\eta_H=\frac{24}{25}.
\]
For any isometric inclusion of the native family trimer
$G\cong\mathbb C^3$ in that band, its induced $L=tI_G$ has
\[
 \boxed{\Lambda^1L=tI_G,\quad
 \Lambda^2L=t^2I_{\Lambda^2G},\quad
 (\Lambda^2L)^{-1}=t^{-2}I,\quad
 (\det L)^{-1}=t^{-3}.}
\]
The corresponding comparison exponents are $+1/2,+1,-1,-3/2$.
The exterior construction is an exact mathematical realization
of these four classes, independently of any measured particle mass.
\end{theorem}
\begin{proof}
The scalar overlap is an action-band quadratic-frame identity,
not an assigned fermion coefficient.  Exterior cross-Grams are
determinants of $r\times r$ minors of $L$ and hence are $t^rI$;
inverses are taken only on nonzero active support.
\end{proof}

\begin{proposition}[Constructive preterminal comparison on the raw charged words]
\label{p6:prop:v5124-constructive-preterminal-words}
Let $F_p$ denote the \emph{original} direct counting/split-path
primitive or, for an inverse row, the primitive \emph{before}
its Moore--Penrose inverse in Release v5.66.  Give it the indicated
candidate exterior multiplier $\rho_p$, and apply the original
terminal direct or inverse operation only \emph{after} that insertion:
\[
 F_p^{\rm cmp}=\rho_pF_p,\qquad
 \operatorname{Term}_p(F_p^{\rm cmp})
 =\begin{cases}\sigma(F_p^{\rm cmp}),&\text{direct},\\
 \sigma((F_p^{\rm cmp})^+),&\text{inverse}.
 \end{cases}
\]
On each nonzero singular-value support, the following is exact:
\[
\begin{array}{c|c|c|c}
\text{role}&F_p\text{ before inversion}&\rho_p&
\operatorname{Term}_p(F_p^{\rm cmp})/\operatorname{Term}_p(F_p)\\\hline
u&N_2&t&t\\
t&P_{(7,9)}&t&t\\
d&P_{(2,3)}&t^2&t^2\\
c&P_{(4,4,4)}&t^2&t^2\\
\tau&P_{(6,12,6)}&t^2&t^2\\
\mu&N_4\;(\text{inverse})&t^{-2}&t^2\\
s&N_6\;(\text{inverse})&t^2&t^{-2}\\
e&P_{(1,5)}\;(\text{inverse})&t^3&t^{-3}\\
b&P_{(8,16,1)}\;(\text{inverse})&t^3&t^{-3}
\end{array}
\]
The output exponent vector is the existing v5.91 vector in the
order $(u,d,e,t,b,\mu,s,c,\tau)$:
\[
 \boxed{(\tfrac12,1,-\tfrac32,\tfrac12,-\tfrac32,1,-1,1,1).}
\]
No additional continuous parameter is needed to realize this table.
Its \emph{discrete role-to-exterior-sort assignments}, however, have
not been derived by taking variations of the unterminalized charged
source.  The proposition is constructive, not a theorem of physical
owner selection.
\end{proposition}
\begin{proof}
For $\rho>0$, $\sigma(\rho F)=\rho\sigma(F)$ and
$(\rho F)^+=\rho^{-1}F^+$.  Applying these two identities to
all nine rows gives the table directly.  The $\mu$ and $s$ rows
are deliberately different: their underlying positive owners are
multiplied by $t^{-2}$ and $t^2$ \emph{before} inversion,
respectively; the sign of a raw MP inverse alone does not determine
the final comparison grade.
\end{proof}

\begin{proposition}[Post-comparison Hodge-order and fixed-volume test]
\label{p6:prop:v5124-hodge-order}
For the already generated trimer overlap $K_V$ and its native bivector
return $H_{\Lambda^2}=\tfrac14H_V^{-1}$, the ordinary exterior square
of $tK_V$ gives a positive bivector return multiplied by $t^4$.
By contrast \emph{imposing the positive-return Hodge identity after}
comparing the vector return gives
\[
 H_V^{\rm cmp}=t^2H_V\quad\Longrightarrow\quad
 H_{\Lambda^2}^{\rm cmp}
 =\tfrac14(H_V^{\rm cmp})^{-1}
 =t^{-2}H_{\Lambda^2}.
\]
These are different operations.  If the original oriented three-volume
of $K_V$ is held fixed while dualizing $tK_V$, the ratio to its
naturally compared bivector is
\[
 \boxed{(\det(tI_3))^{-1}=t^{-3}=\eta_H^{-3/2}.}
\]
Thus the reciprocal-area and reciprocal-volume classes have exact
Hodge/determinant realizations, but these do not prove that the
$N_4^+,N_6^+$ or outer reciprocal charged roles consume those
returns in the actual action.
\end{proposition}
\begin{proof}
The determinant of $tI_3$ is $t^3$ and
$\Lambda^2(tK_V)=t^2\Lambda^2K_V$.
The nondegenerate $3$-dimensional adjugate identity is
$\Lambda^2 A=(\det A)A^{-T}$ in Hodge coordinates.
The two displayed comparison orders therefore give the asserted
$t^{-2}$ and $t^{-3}$ factors.
\end{proof}

\begin{remark}[Physical-source nonuniqueness and cocycle scope]
\label{p6:rem:v5124-physical-owner-underdetermination}
The family embedding, Clifford hierarchy rung, Hurwitz $24$-cell,
and invariant quadratic comparison scalar do not themselves
choose the left/right chiral-frame assignments or exterior degree
of a particular charged role.  Paper~III-F,
Proposition~\ref{p3f:prop:v221-relative-chiral-comparison-frame},
exhibits two same-rooted-frame choices with canonically normalized
outputs $D_0$ and $tD_0$.  Likewise additional closed scalar returns
$(L^*L)^k$ preserve source/target types while changing the exponent.
Therefore the v5.91 additive cocycle is licensed on its \emph{reduced
scalar-output comparison category}; it must not be extended to
arbitrary matrix compositions of its rank-one $A$.  In fact
$A^2=\eta_HA$, so $\sigma(A^2)=\eta_H^{3/2}$, not $\eta_H$.
The exterior comparison $L=tI_G$ is composable, but its physical
owner assignments remain a separate statement to establish.
\end{remark}

\begin{corollary}[Precise remaining charged comparison certificate]
\label{p6:cor:v5124-owner-selection-open}
A full action-derived particle comparison proof requires an actual
mixed charged-source derivative or equivalent generated parent
contraction showing, before v5.90 terminalization, all of:
(i) the consumed quadratic comparison port, (ii) the chiral relative
frame, (iii) the nine role-to-exterior and Hodge inverse placements,
(iv) the absence of idle comparison loops, and (v) preservation of the
readout through canonical normalization and the physical rest/pole
criterion.  None of (ii)--(v) follows solely from the scalar-frame
identity; their occurrence in the currently printed charged source
is not proved by the table above.  The existing finite spectrum values
remain conditional on the previously chosen charged comparison grades;
no empirical PDG agreement is used as an ownership proof.
\end{corollary}

% BEGIN FRAMEWORK INTEGRATION v3.024: CHARGED PHYSICAL-PORT DIAGNOSTIC
\begin{remark}[Independent native hierarchy and nine-grade source tests]
\label{p6:rem:v3024-charged-source-port-diagnostic}
The mathematical nearest-source theorem of Paper~V-B supplies a
\emph{diagnostic}, not a rule that the physical action minimizes a
Hilbert--Schmidt norm.  On each independently generated
preterminal charged-source member \(p\), first evaluate the
\emph{actual} mixed source and kinetic Grams from the common action,
\[
 \mathcal D_{LR,p}^{\rm pre}
  =\frac{\delta^2S_\star}
  {\delta\bar\psi_{L,p}\,\delta\psi_{R,p}},\qquad
 M_p=Z_{L,p}^{-1/2}\mathcal D_{LR,p}^{\rm pre}Z_{R,p}^{-1/2}.
\]
The hierarchy-incidence test and the comparison-port test are
\emph{independent}.  The former checks the root-sensitive defect
\(\Delta_{\rm hier}(M_p)=0\), its three nonzero oriented blocks,
and, for equality to the source-owned canonical arrows, the three
oriented contractions and the remaining norm equality of
the four-invariant equality test of Paper~V-A.
The latter must extract, \emph{before} terminal mass numbers are
inserted, the source-owned quadratic port, relative chiral placement,
exterior/Hodge operation, direct or inverse readout, and absence of
idle comparison loops.  For a scalar comparison
\((\eta_H^{k_D},\eta_H^{k_L},\eta_H^{k_R})\) in the source and
kinetic blocks, the canonically normalized exponent is
\[
  n_{\rm can}=k_D-\frac12(k_L+k_R).
\]
With \(\vartheta=\sqrt{\eta_H}\), consider the role order
\((u,d,e,t,b,\mu,s,c,\tau)\).
The previously printed \emph{candidate} before-inversion exponents are
\[
  (k_p)=(1,2,3,1,3,-2,2,2,2),\qquad
  \text{candidate input factor }\vartheta^{k_p}.
\]
This list records the operations to be independently verified, not
an evaluated mixed-source selection.  A scalar loop or the vanishing
of \(\Delta_{\rm hier}\) alone cannot certify those grades.
Physical rest/pole occurrence is a further separate condition.
The structural leastness of GP3 is not a numerical nearest-point
postulate.  In particular this diagnostic does not discharge
\(\mathsf{Cert}_{\rm hier}\), \texttt{BH-VA-G-0007}, or
\texttt{BH-G-000113}.
\end{remark}
% END FRAMEWORK INTEGRATION v3.024

% END FRAMEWORK v2.221: EXTERIOR-POWER CHARGED COMPILER AUDIT



% BEGIN MASS-ONTOLOGY-FLAVOR-READOUT v5.92
\subsection{Historical fermion assembly and incidence diagnostics}
\label{p6:app:v2187-historical-fermion-assembly}
These ledgers belong to earlier fixed spectral candidates or intermediate
incidence stages. Their open statuses do not describe the current rooted
charged source. Main-text Release v5.90 and Paper~V-A
Corollary~\ref{p5:cor:v397-charged-preterminal-assembly-closure}
close that source; main-text Release v5.97 evaluates the lepton plane.
Physical source-cyclic roots, particle/rest identification and actual
claimed limits remain separate.

The historical directed-candidate obligation was source-owned left-support incidence
and the coherent response that realizes it. A terminal-role permutation
against a fixed frame is not covered by this covariance.
Appendix~\ref{p6:app:v585-directed-assembly} gives the directed candidate,
its one-parent frame identity, the six-member spectral assembly theorem,
CP preservation, and a general matching test using two additional mixed
moments once the spectrum is fixed.  On the fixed directed native raw
charged packets, Theorems~\ref{p6:thm:v2109-one-moment-matching} and
\ref{p6:thm:v2109-one-moment-full-return} sharpen this to one trace and
reconstruct the entire weighted return.  The thirty-six raw quark pairs are
inequivalent by Theorem~\ref{p6:thm:v2109-thirtysix-inequivalent}.
These remain conditional algebraic results, not a physical source selector.
The release-stage status was
\[
 \texttt{FERMION\_SAME\_ACTION\_SVD\_ASSEMBLY}
 =\texttt{HISTORICAL\_OPEN\_PHYSICAL},
\]
while primitive source incidence remains \texttt{CLOSED\_GEOMETRIC}.


\begin{corollary}[Exact remaining charged-lepton incidence frontier]
\label{p6:cor:v596-lepton-inplane-frontier}
\noindent\textbf{Historical scope.} The open-plane/no-consumer assertion is superseded by main-text Release v5.97 and Paper~V-A Corollary~\ref{p5:cor:v3108-lepton-incidence-status}; the tau/direction results remain active.

The current status contracts to
\[
 \boxed{\texttt{LEPTON\_SINGLET\_PLANE\_PLACEMENT}
 =\texttt{CLOSED}[P_\tau=Q_0],}
\]
\[
 \boxed{\texttt{LEPTON\_E36\_DIRECTIONS}=\texttt{CLOSED}
 [\delta_\tau=-1,\ \delta_\mu=+1],}
\]
and
\[
 \boxed{\texttt{LEPTON\_ROLE\_TO\_FAMILY\_INCIDENCE}
 =\texttt{HISTORICAL\_OPEN\_INPLANE\_SOURCE\_EVALUATION}.}
\]
With $e_\pm$ the normalized $C_3$ character vectors, the surviving family is
\[
 v_e=\cos\theta\,e_+ +e^{i\phi}\sin\theta\,e_-,\qquad
 v_\mu=-e^{-i\phi}\sin\theta\,e_+ +\cos\theta\,e_-,\qquad
 v_\tau=c,
\]
up to role-line phases, and still
\[
 \boxed{\zeta_e=\cos2\theta.}
\]
Thus only one source-owned orientation/helicity in $Q_\perp$ remains.
\end{corollary}

\begin{remark}[Structural $R_G$ phase is not yet the physical plane selector]
\label{p6:rem:v596-rg-consumption-firewall}
\noindent\textbf{Historical scope.} The open-plane/no-consumer assertion is superseded by main-text Release v5.97 and Paper~V-A Corollary~\ref{p5:cor:v3108-lepton-incidence-status}; the tau/direction results remain active.

Paper~I-C proves the structural sector relation
$R^-\mapsto R_G^\dagger$, and the corresponding $Q_\perp$ compressions give
exact coefficient-free candidate lines.  However Paper~I-C also states that
physical use requires an action constructor consuming that joint-parent
phase.  No such charged-current consumption word is currently frozen.
Accordingly $R_G^\dagger$ is not used here to select $(v_e,v_\mu)$, and
\texttt{CHARGED\_LEPTON\_ORIENTATION\_MOMENT} remains
\texttt{OPEN\_SOURCE\_EVALUATION}.
\end{remark}

% END INLINE SOURCE: Paper5B-HistoricalFermionAssembly-Appendix-v1_0.tex



% BEGIN INLINE SOURCE: Paper5B-ElectronArrowLimits-Appendix-v1_0.tex
% Framework v2.191: host-algebra and right-arrow diagnostics; no physical rivals.
\subsection{Exact certificate for the primitive-trimer local chart}
\label{p6:app:v548-trimer-local-chart-certificate}
Use the Paper-I-A cyclic basis
\[
 \DtriI=\operatorname{diag}(2/3,-1/3,-1/3),\qquad
 \DtriII=\operatorname{diag}(0,1,-1),
\]
\[
 X=\begin{pmatrix}0&0&1\\1&0&0\\0&1&0\end{pmatrix},\qquad
 C_1=X+X^\dagger,\qquad C_2=-i(X-X^\dagger).
\]
The five matrices $I,\DtriI,\DtriII,C_1,C_2$ have real vectorized rank five.
For the witness $(A_*,B_*)$ of
Theorem~\ref{p6:thm:v548-trimer-local-submersion}, Sylvester's criterion gives
\[
 A_*:\quad \frac{32}{3},\ \frac{1102}{9},\ \frac{24032}{27},
\qquad
 B_*:\quad \frac{40}{3},\ \frac{1123}{9},\ \frac{30760}{27},
\]
for the leading principal minors.  Both matrices are therefore positive
definite.  Their characteristic discriminants are
\[
 \operatorname{Disc}(A_*)=3240,
 \qquad
 \operatorname{Disc}(B_*)=50240,
\]
so both spectra are simple.  The simultaneous commutant equations
$[T,A_*]=[T,B_*]=0$ have coefficient rank eight in the nine-dimensional
complex matrix space, leaving exactly the scalar line.  Finally
\[
 \det[A_*,B_*]=756i,
 \qquad
 \frac1{6i}\operatorname{tr}[A_*,B_*]^3=378.
\]
Direct exact differentiation of the ten invariant functions displayed in
Theorem~\ref{p6:thm:v548-trimer-local-submersion} gives
\[
 \det J_F=-159690424320.
\]
These identities are an exact arithmetic certificate for the theorem.  They
are not a numerical fit and do not select a physical flavor point.

\subsection{Cyclic-natural family-operator no-go}
\label{p6:sec:family-quiver-update}

Let $G\cong\CC[\mathbb Z_3]$ be the generated unpointed oriented family
trimer and let $X$ be its cyclic shift.

\begin{theorem}[Cyclic-natural family operators are circulant]
\label{p6:thm:cyclic-natural-family-circulant}
If $A\in\End(G)$ is natural under the residual cyclic relabelling, then
\[
 XAX^{-1}=A,
\]
so
\[
 A\in\End_{\mathbb Z_3}(G)=\CC[I,X,X^2].
\]
Every two self-adjoint operators generated from this unpointed cyclic geometry
therefore commute and share the Fourier eigenframe.
\end{theorem}

\begin{corollary}[No CKM/PMNS mixing from one common family functional calculus]
\label{p6:cor:no-single-family-dirac-mixing}
If all sector Hermitian operators satisfy $H_s=f_s(D_G)$ for one normal or
self-adjoint family operator $D_G$, then $[H_s,H_t]=0$ for every pair.  On the
simple-spectrum stratum the relative eigenframes are monomial up to phases,
so the corresponding Jarlskog commutator invariants vanish.  A single common
family operator is therefore not a viable flavor-closure mechanism.
\end{corollary}

The full Weyl algebra $\End(G)=M_3(\CC)$ remains an operator grammar, not a
selection theorem.  The appropriate missing object is a nonabelian transport
on the already generated Yukawa--Majorana slot quiver.  Its return loops
\[
 H_u=Y_uY_u^*,\qquad H_d=Y_dY_d^*,
\]
and their lepton analogues must be generated noncommutingly and the full
quiver representation must be unique only up to the weak-basis gauge group.
This is the structure required to collapse the generic weak-basis quotient;
writing arbitrary matrices in a Weyl basis does not reduce it.

\subsection{C13 contact-density check, counterfamily, and typed boundary}
\label{p6:app:c13-density-and-counterfamily}

This appendix records the exact positive and negative mathematics behind the
main-text C13 discharge.  The contact and cone calculations derive a
four-density degree.  By themselves they do not promote an ambient line
isomorphism to a physical owner map or select a member of the
energy-character family; the CRP4 parent is the additional constitutive input.

\begin{proposition}[Critical contact residue equals the standard four-ball density]
\label{p6:prop:c13-contact-filling-density}
Let $P_H:=\Delta_{b,0}^{1/2}$ for the normalized scalar Hopf contact
sub-Laplacian of Paper~III-C.  Its Heisenberg order is one and the homogeneous
dimension is four, so
\[
 \Delta_{b,0}^{-2}=P_H^{-4}.
\]
For the unit standard ball $B^4\subset\mathbb C^2$, let
$\lambda_{\rm std}$ be the standard Liouville form and
$\omega=d\lambda_{\rm std}$.  In the source normalization,
\begin{equation}
\label{p6:eq:c13-contact-filling-density}
 \operatorname{Res}_H(\Delta_{b,0}^{-2})
 =\frac18\int_{S^3}\eta_0\wedge d\eta_0
 =\int_{B^4}\frac{\omega^2}{2}
 =\frac{\pi^2}{2}.
\end{equation}
Thus the primitive contact residue and the standard symplectic filling agree
as four-densities.  The equality does not include the completed REG character
and is not an energy-owner identification.
\end{proposition}

\begin{proof}
Paper~III-C,
Theorem~\ref{p3:prop:contact-residue-normalization}, gives the first and last
equalities.  Paper~I-A,
Proposition~\ref{p1:prop:standard-contact-fill-tight}, gives
$\lambda_{\rm std}|_{S^3}=\eta_0/2$ at unit radius.  Stokes' theorem then
gives
\[
 \int_{B^4}\frac{\omega^2}{2}
 =\frac12\int_{S^3}\lambda_{\rm std}\wedge d\lambda_{\rm std}
 =\frac18\int_{S^3}\eta_0\wedge d\eta_0.
\]
\end{proof}

\begin{proposition}[Natural positive cone-homothety counterfamily]
\label{p6:prop:c13-lambda-cone-counterfamily}
Fix one REG member, write
$L_\sigma:=\log\alpha_\sigma$,
$\Phi_\sigma(q):=qL_\sigma$, and
$\chi_D(\sigma,q):=e^{-\Phi_\sigma(q)}$.  For every $\lambda>0$, set
\begin{equation}
\label{p6:eq:c13-lambda-cone-counterfamily}
 s_\lambda(\sigma,q):=e^{-\lambda\Phi_\sigma(q)/4},
 \qquad
 \rho_\lambda(\sigma,q):=s_\lambda(\sigma,q)I_{W_x^*}.
\end{equation}
Then $\rho_\lambda$ is positive-definite, $O(3)$-natural, preserves the
exact cone, and is multiplicative in additive depth.  Moreover
\begin{equation}
\label{p6:eq:c13-lambda-frame-determinant}
 |\det\rho_\lambda|
 =s_\lambda^4=e^{-\lambda\Phi_\sigma}=\chi_D^\lambda,
 \qquad
 \frac{E_\lambda(K)}{E_\lambda(0)}
 =\alpha_\sigma^{-\lambda K/32}.
\end{equation}
Hence cone naturality, positivity, the frame dimension, monotonicity, and the
logarithmic midpoint do not select $\lambda=1$.
\end{proposition}

\begin{proof}
Scalar multiplication by $s_\lambda>0$ preserves both sides of the exact
cone inequality and commutes with the full presentation action.  Determinants
multiply the scale four times.  Finally $q=K/8$, so the scalar component has
exponent $-\lambda qL_\sigma/4=-\lambda KL_\sigma/32$.
\end{proof}

The family in
Proposition~\ref{p6:prop:c13-lambda-cone-counterfamily} is an ambient-host
counterfamily, not a family of physically occurring energy laws.  Its role is
to refute selection from the listed mathematical properties alone.

\begin{proposition}[Scheme-neutral density no-go and base-type firewall]
\label{p6:prop:c13-scheme-neutral-density-no-go}
For $q\ne0$, no nonzero linear identification of the completed weight-one
determinant character with a scheme-neutral contact/filling density can be
equivariant under every finite scheme translation.  Moreover, before
pullback, the determinant-depth line is a bundle over the scheme--depth base
$B$, while $\mathcal D_4=|\Lambda^4W^*|$ is a bundle over spacetime $M$;
they are not the same typed functor and cannot be directly identified.
\end{proposition}

\begin{proof}
Under $\sigma\mapsto\sigma+c$, the weight-one character transforms as
$\chi_D\mapsto e^{-qc}\chi_D$, whereas a density built only from the fixed
contact filling has trivial scheme action.  Equivariance of a linear map $T$
would require $T(e^{-qc}v)=T(v)$ for every $c$, forcing $T=0$ when
$q\ne0$.  At $q=0$ the character is the trivial upper anchor and selects
no weight.  The base statement follows from the displayed source and target
bundle projections.  Only the pullbacks
$\pi_B^*\mathcal L_{\rm det}^{-1}$ and
$\pi_M^*\mathcal D_4$ over $B\times M$ are comparable.
\end{proof}

Consequently a bare positive-line isomorphism, a connection obtained by
transport through that freely chosen isomorphism, Born measure balance,
first-symbol scalar conversion, or exact symplecticity cannot silently replace
Theorem~\ref{p6:thm:crp4-c13-discharge}.  This countermodel identifies the
missing physical content as a common occurring first-order parent that
generates both functors and both connection actions before comparison.
The generated occurring action ancestry of Papers~III-B/III-E supplies
that content (the generated-action-response owner and CRP4 positive parent).  The
countermodel therefore proves minimality; it is not evidence that C13 remains
open after the CRP4 import.

\subsection{Projector joins, normalized rays and direct Higgs comparisons}
\label{p6:app:v179-projector-diagnostics}

This appendix records diagnostics of simple physical identifications.
The unital Hurwitz theorem in the main text does not require those
identifications.  Any extra field realization used here is conditional.

\begin{proposition}[Join versus sum and the local reduced septet]
\label{p6:prop:v179-join-sum-diagnostic}
For the projectors of
Theorem~\ref{p6:thm:v179-hurwitz-unital-completion},
\[
 \operatorname{Tr}(P_{16}\Pi_{\rm c})=2/3,\qquad
 \|P_{16}+\Pi_{\rm c}\|_{\rm HS}^2=55/3.
\]
Thus \(P_{16}+\Pi_{\rm c}\) is not the unital-join projector.
The orthogonal complement of \(P_{17}\) is
\(P_8-|\Omega_8\rangle\langle\Omega_8|/8\), a rank-seven reduced cube
space orthogonal to the global seven-point support \(P_7\).
\end{proposition}
\begin{proof}
The global constant has sixteen of its twenty-four equal squared
components in \(\mathcal G_{16}\).  Squaring the sum gives
\(16+1+2(2/3)=55/3\), whereas its trace is \(17\).
The join formula gives its complement, supported on the disjoint cube.
\end{proof}

\begin{proposition}[Normalized scalar rays and the flat-projector test]
\label{p6:prop:v179-flat-projector-higgs-test}
For an admitted linear odd-field ray \(\Phi=xA\), put
\(g_2=\operatorname{Tr}_F A^2>0\), \(g_4=\operatorname{Tr}_F A^4\).
The local scalar action has
\[
 Z_x=2Cg_2,\qquad h=\sqrt{2Cg_2}\,x,\qquad
 V(h)=-\tfrac12\mu_R^2h^2+\tfrac14\lambda_Ah^4,\qquad
 \lambda_A=\frac{g_4}{Cg_2^2}.
\]
If its formal nonzero stationary point exists, the radial local mass
squared is \(2\mu_R^2\).
The purely mathematical doubling
\[
 D(P)=\begin{pmatrix}0&P\\P&0\end{pmatrix}
\]
of a rank-\(r\) projector has \(g_2=g_4=2r\), hence would give
\(\lambda_A=8\pi^2/r\) if admitted with that trace.
It is not a constructed physical Higgs source.
\end{proposition}
\begin{proof}
Substitute the ray in the native action and normalize its kinetic term.
Differentiating the resulting quartic gives the stationary radial mass.
The doubled projector has two rank-\(r\) return blocks in both even
powers, proving its trace identities.
\end{proof}

There is a further exact obstruction to direct unitary identification.
If the three charged matrices \(Y_u,Y_d,Y_e\) all have rank three, the
full chiral block of
Theorem~\ref{p5:thm:v179-canonical-higgs-direction} has rank
\[
 3\operatorname{rank}Y_u+3\operatorname{rank}Y_d
 +\operatorname{rank}Y_e+\operatorname{rank}Y_\nu
 =21+\operatorname{rank}Y_\nu\ge21.
\]
Its self-adjoint completion has rank at least \(42\); \(D(P_{17})\)
has rank \(34\).  They cannot be unitarily identical.  This rules out
that full flat-projector assignment, not every possible partial source
map into the actual flavor family.

The formal quartic minimum is also not a physical vacuum theorem:
Remark~\ref{p3f:rem:v175-local-bosonic-action-boundary} retains the
exact-heat and full-effective-action condition.  The numerical values
of the existing comparison ledger are unchanged by these source
constructions.
% END HM v1.79 HU02



% BEGIN HM v1.79 EX03
% BEGIN EXCEPTIONAL-ROOT-INTEGRATION v2.173

% BEGIN INLINE SOURCE: Paper5B-ExceptionalReadoutComparison-Appendix-v1_0.tex
\subsection{Supplementary exceptional readouts and septet placement}
\label{p6:app:exroot-readout-comparison}

The existing root counts are now accompanied by explicit embedded
configurations and response maps:
Theorem~\ref{p1:thm:exroot-explicit-e876} constructs the specified
golden \(E_8,E_7,E_6\) sections;
Theorem~\ref{p1b:thm:exroot-response-norms} separates vector, scalar
operator and polar normalization. These results are supplementary and
are not new primitive charged mass nodes.

\begin{corollary}[Protected \(E_7\) and the reduced cube septet]
\label{p6:cor:exroot-septet-placement}
Under the fixed real Clifford/Walsh metric comparison of
Proposition~\ref{p3b:prop:exroot-walsh-complement}, the protected
\(E_7\) span maps to
\[
 W_{\rm cube}=P_8-\frac18|\Omega_8\rangle\langle\Omega_8|=I-P_{17}.
\]
This is the reduced cube septet \(7'\). It is distinct from the global
seven-point support \(P_{\mathcal G_7}\), and
\(P_{\mathcal G_7}W_{\rm cube}=0\).
\end{corollary}
\begin{proof}
The metric transport identifies the scalar-free Clifford complement
with the nonconstant cube functions. The global seven points lie
outside that cube, so the counting support projectors are orthogonal.
\end{proof}

The support distinction of
Proposition~\ref{p6:prop:v179-join-sum-diagnostic} and
Proposition~\ref{p6:prop:v566-heavy-hinge-refinement} is retained.
The character comparison in
Proposition~\ref{p3b:prop:exroot-septet-characters} concerns its
specified full-core actions. Its pointed \(C_3\) isometry is not
identified with the actual triality/source incidence, and does not
reopen the existing rooted source-cell closure or impose full-core
symmetry on that source.

\begin{remark}[Status of the 240-label energy comparison]
\label{p6:rem:exroot-240-energy-status}
The exact scalar inverse \(240^{-1/2}\) occurs for both the native
\(H_4\oplus H_4\) pinor presentation and the specified golden \(E_8\)
presentation. Their shared operator Gram does not identify the actions,
the physical source metric or a pole. A common uninserted source scale
cancels by Corollary~\ref{p3f:cor:exroot-common-scale-cancellation};
at fixed action a source amplitude changes residues rather than pole
positions by Proposition~\ref{p3f:prop:exroot-source-scale-poles}.
The generated mode comparison is the full-source condition of
Corollary~\ref{p3f:cor:exroot-residue-criterion}.

Thus the \(E_{32}/\sqrt{240}\) numerical reporting-point candidate
retains exactly the scope of
Corollary~\ref{p6:cor:v582-near-unity-not-invariant}.
No near-unity continuum coefficient, root-count energy divisor,
new \(126\)-based particle, or additional \(7'\) charged node is derived.
The established active action supports \(8,4,3,1\) are not replaced by
the exceptional root ranks.
\end{remark}


% END INLINE SOURCE: Paper5B-ExceptionalReadoutComparison-Appendix-v1_0.tex

% END EXCEPTIONAL-ROOT-INTEGRATION v2.173

\subsection{Incidence restriction and exclusion projectability}
\label{p6:app:v179-incidence-descent}

These calculations explain the limits of the natural maps accompanying
the positive incidence normalization.  They introduce no new field owner.

\begin{proposition}[Restriction to one Hurwitz cell has rank ten]
\label{p6:prop:v179-incidence-restriction-rank}
Let \(R_H:\mathbb C^{120}\to\mathbb C^{24}\) restrict functions to one
selected tetrahedral cell.  Then
\[
 \operatorname{rank}(R_H B)=\operatorname{rank}(R_HU_B)=10 .
\]
\end{proposition}
\begin{proof}
The columns of \(B\) are the matrix-coefficient functions of the
five-point \(A_5\) permutation representation \(\mathbf1\oplus V_4\).
Its coefficient-space dimension is \(1^2+4^2=17\).  On a point
stabilizer \(A_4\), the representation restricts to
\(\mathbf1\oplus\mathbf1\oplus V_3\).  The distinct coefficient types
are only \(\mathbf1,V_3\); repeated trivial summands give the same
constant coefficient function.  Their dimension is \(1^2+3^2=10\).
The central binary lift duplicates values without changing rank.
The positive inverse square root in \(U_B\) is invertible on the
effective source support and has the same image as \(B\).
For a coset cell, translation gives the same rank.
\end{proof}

\begin{proposition}[Single-label exclusion does not descend through incidence]
\label{p6:prop:v179-label-exclusion-nondescent}
Put \(P_{\rm ker}=I-P_{\rm eff}\), let \(e\) be one unit cell-label
vector, and \(Q_{\rm exc}=I-|e\rangle\langle e|\).
There is no linear vertex operator \(Q_{\rm vert}\) satisfying
\(Q_{\rm vert}B=BQ_{\rm exc}\) on the full label carrier.
\end{proposition}
\begin{proof}
A linear operation descends through \(B\) only if it preserves
\(\ker B\).  Set \(k=P_{\rm ker}e\).  Since the diagonal entries of
\(P_{\rm glob}\) and \(P_{\rm con}\) are \(1/25\) and \(16/25\),
\(\langle e,k\rangle=8/25\).  Thus
\[
 Bk=0,\qquad BQ_{\rm exc}k=-\tfrac8{25}Be\ne0 .
\]
This contradicts preservation of the kernel.
\end{proof}

The different inherited comparisons can be displayed exactly.
For normalized \(\Omega_{25}\), the label exclusion has squared norm
\(24/25\); the normalized actual-cell exclusion in the vertex carrier
has value \(4/5\), as in Corollary~\ref{p6:cor:self-exclusion}.
For the whitened cell direction,
\(\|U_Be\|^2=17/25\), and its normalized overlap squared with the global
line is \(1/17\); orthogonal exclusion of that direction retains \(16/17\).
These are different operators, not alternative values of a freely
adjustable physical correction.

The effective compression \(C_{\rm exc}=P_{\rm eff}Q_{\rm exc}P_{\rm eff}\)
is also not an orthogonal projector.  If \(v=P_{\rm eff}e\), then on the
effective support it is \(I-|v\rangle\langle v|\), with
\(\|v\|^2=17/25\).  It has eigenvalue \(8/25\) along \(v\).
Its global expectation is \(24/25\), while its global output squared
norm is \(592/625\).  The expectation alone therefore does not prove
descent of the original exclusion.
% END HM v1.79 EX03


\subsection{Claim-complete exclusion owner and family-compression diagnostics}
\label{p6:app:v181-exclusion-owner-completion}
This appendix records the detailed finite comparison behind
Theorem~\ref{p6:thm:v181-claim-complete-self-exclusion-owner}.  This is a
specified three-state comparison, not the unresolved physical G15 family
source.  Let $G_{\rm cmp}\cong\mathbb C^3$ and choose the three even
permutations
\[
 \tau_1=(1\,3\,4),\qquad
 \tau_2=(1\,4\,2),\qquad
 \tau_3=(1\,2\,3)
\]
in the $A_5$ presentation of the $2I$ vertex carrier.  If $J_G$ inserts the
corresponding three vertex basis vectors, define
\[
 R_I:=J_G^\dagger U_B,\qquad
 G_I:=R_IR_I^\dagger
 =\frac{19}{120}Q_0+\frac2{15}Q_\perp,
\]
where $Q_0=\mathbf1_3\mathbf1_3^T/3$ and $Q_\perp=I_3-Q_0$.  The positive
polar normalization
\[
 \boxed{W:=R_I^\dagger
 \left(\sqrt{\frac{120}{19}}Q_0+\sqrt{\frac{15}{2}}Q_\perp\right)}
 \quad:G_{\rm cmp}\longrightarrow P_{\rm eff}\mathbb C^{25}
\]
satisfies $W^\dagger W=I_3$.  Put
$u_0=(1,1,1)^T/\sqrt3$, $w_0=Wu_0$.  Direct evaluation of this inherited
incidence Gram gives
\[
 W^\dagger\Pi_HW=\frac{243}{475}Q_0,
 \qquad
 \boxed{\langle w_0,Q_{\rm exc}w_0\rangle=\frac{232}{475}},
\]
where $Q_0=\mathbf1_3\mathbf1_3^T/3$.  The original uniform label state and
the family-uniform state are distinct and satisfy
\[
 \boxed{|\langle\Omega_{25},w_0\rangle|^2=\frac3{19}.}
\]

Put $P_W=WW^\dagger$ and
\[
 r_H:=\frac{(I-P_W)e_H}{\sqrt{232/475}}.
\]
Then $WG\oplus\mathbb Cr_H$ is the least reducing completion of $WG$ for the
single projector $\Pi_H$, and in the $(w_0,r_H)$ sector
\[
 Q_{\rm exc}=
 \begin{pmatrix}
 232/475&-\sqrt{(243/475)(232/475)}\\
 -\sqrt{(243/475)(232/475)}&243/475
 \end{pmatrix}.
\]
Retaining the original $\Omega_{25}$ adds one further orthogonal companion
$r_\Omega$, so
\[
 \boxed{\mathcal K_5=\operatorname{span}(W G_{\rm cmp},\Omega_{25},e_H),\qquad
 \dim\mathcal K_5=5.}
\]
The exact decomposition is
\[
 \Omega_{25}
 =\sqrt{\frac3{19}}\,w_0
 -\sqrt{\frac8{551}}\,r_H
 +\sqrt{\frac{24}{29}}\,r_\Omega,
 \qquad 551=19\cdot29,
\]
with $Q_{\rm exc}r_\Omega=r_\Omega$.  Evaluation in this least faithful owner
restores
\[
 \boxed{\langle\Omega_{25},Q_{\rm exc}\Omega_{25}\rangle=\frac{24}{25}.}
\]
The arithmetic identity
\[
 \boxed{\frac{24}{25}
 =2\frac{232}{475}-\frac1{19}\frac8{25}}
\]
is therefore an exact corollary of the completed matrix geometry.  The factor
$2$ is not a primitive Schur/Feshbach weight and must not be promoted to a
universal correction rule.

If instead a claim consumes all $25$ coordinate projectors $\Pi_a$, every
$\Pi_aW G_{\rm cmp}$ is a nonzero coordinate line.  The least invariant owner is then
all of $\mathbb C^{25}$, not $\mathcal K_5$.  These statements concern finite
owner retention and do not select the physical G15 Yukawa/source word.


\subsection{Historical printed scalar readout: same-action no-go}
\label{p6:app:v181-printed-higgs-readout-no-go}
The historical comparison coordinate
\[
 \kappa_{\rm print}=2\pi\sqrt{\frac{24}{15}}=\frac{8\pi}{\sqrt{10}}
\]
is a readout coordinate, not a derived physical vacuum.  If it were promoted
to the same-action canonical Higgs vacuum while also imposing the historical
row and the generated Yukawa quartic, Corollary~\ref{p6:cor:v181-higgs-two-coordinate-reduction}
would require
\[
 \frac{\sqrt{102}}{10\pi}\le\kappa_{\rm print}
 \le\frac{6\sqrt{17}}{5\pi},
\]
which fails.  Equivalently, the historical printed quartic
$\lambda_{\rm print}=51/(40\pi^2)$ would imply
\[
 N_{\rm eff}=\frac{8\pi^2}{\lambda_{\rm print}}
 =\frac{320\pi^4}{51}>24,
\]
contradicting Corollary~\ref{p5:cor:v181-yukawa-participation-invariant}.
Thus that printed coordinate cannot simultaneously be the physical canonical
Higgs VEV, obey the generated same-action Yukawa quartic, and reproduce the
historical $17(24/25)$ row.  The result does not exclude the coordinate from
its declared comparison/readout role.



\subsection{Frontier-scale counterfamily and hierarchy-index no-go diagnostic}
\label{p6:app:v191-frontier-e32-counterfamily}
Fix the hierarchy data $R_p,\alpha_{\rm hier},E_{32}$ and replace the
independently declared sharp-frontier scale by
$\Lambda_F\mapsto s\Lambda_F$, $s>0$.  The shell order, comparison-band
capture, exceptional groups and hierarchy owner are unchanged, while
$a_k\mapsto a_k/s$ and $E_F\mapsto sE_F$.  Hence the current prefix admits a
continuous family with distinct $E_F/E_{32}$ and cannot derive
$E_F=E_{32}$.

Separately, a fourth-power quantity sampled at frontier events scales as
$[k(k+2)]^{-2}$, whereas the protected hierarchy is exponentially spaced in
its owner coordinate.  If one hierarchy ratio $\alpha_{\rm hier}^{-1}$ were
purely a fourth-power redshift step, it would require
$\Delta N=(1/4)\log\alpha_{\rm hier}$.  Therefore seven consecutive small
position-frontier levels cannot be identified with the seven protected
hierarchy rungs.  This is a diagnostic no-go, not a replacement hierarchy
construction.



\clearpage
\subsection{Flavor-selection obstructions and the alternating-word identity}
\label{p6:app:v199-flavor-selection}
\begin{proposition}[Affine mixed terms and one alternating word are insufficient]
\label{p6:prop:v199-affine-mixed-obstruction}
At fixed simple spectra, an affine function of
\(M_{11},M_{12},M_{21},M_{22}\) has a stationary point in the generic
CP-violating interior only if all four coefficients vanish.  Adding a
constant multiple of \(\operatorname{tr}(ABAB)\) still cannot produce a
nondegenerate angular minimum there.
\end{proposition}
\begin{proof}
The four moments are local coordinates by
Proposition~\ref{p6:prop:v199-invariant-selection-equations}.  The gradient
of an affine function in them is its constant coefficient vector.  For the
second assertion the exact \(3\times3\) identity is
\begin{align}
 \operatorname{tr}(ABAB)={}&M_{11}^2+2a_1M_{12}+2b_1M_{21}-2M_{22}
 -2a_1b_1M_{11}\nonumber\\
 &+\tfrac12(a_2-a_1^2)(b_2-b_1^2).
 \label{p6:eq:v199-alternating-word}
\end{align}
It follows by taking the coefficient of \(t^2\) in the Cayley--Hamilton
trace identity for \(A+tB\),
\[
 p_4=\tfrac43p_1p_3-p_1^2p_2+\tfrac12p_2^2+\tfrac16p_1^4.
\]
At fixed spectra its mixed Hessian has rank one.  Thus at a stationary
interior point three independent mixed directions remain locally flat.
Spectral dependence of the coefficients does not change this angular test.
\end{proof}


\begin{proposition}[The full simple mixed determinant is stationary only at commutation]
\label{p6:prop:v203-full-probe-stationarity}
Let $A,B$ be positive Hermitian matrices with fixed spectra, $z,w>0$,
$D_A=zI+A$, $D_B=wI+B$, $R_A=D_A^{-1}$, and $t=|\eta|^2>0$.
On the positive domain $C=D_B-tR_A>0$, define
\[
 \Phi=-\log\frac{\det\left(\begin{smallmatrix}D_A&\eta I\\
                      \bar\eta I&D_B\end{smallmatrix}\right)}
                         {\det D_A\det D_B}.
\]
An interior point is stationary under every relative unitary rotation of
$B$ if and only if $[A,B]=0$. With simple spectra these points have
permutation mixing up to rephasing and zero CP commutator invariant.
\end{proposition}
\begin{proof}
Schur factorization of the specified probe gives
$\Phi=-\log\det C+\log\det D_B$. This is exactly the positive-branch
function in Paper~1B,
Proposition~\ref{p1b:result:determinant-derivatives-and-relative-frame-stationarity},
which makes relative-frame stationarity equivalent to $[A,B]=0$.
For the simple spectra fixed here, simultaneous eigenvectors then give
permutation mixing up to rephasing and vanishing commutator invariant.
\end{proof}
At $t=0$ the probe is independent of mixing. Keeping all powers of this
same logarithm therefore cannot remove the obstruction. The proposition
does not exclude stationary mixing generated by the full physical action
with its additional, independently generated terms.


If complex conjugation is a physical involution on the admitted owner and
the action and state/source data are CP invariant, the two generic conjugate
solutions have equal action.  A nonzero-\(\chi\) pair can occur, but its
sign needs CP-odd data or an occurring symmetry-breaking state.  This
statement must be reapplied after including any physically supplied
orientation-sensitive term.  It is not a prohibition on such a term and
does not import extra flavor fields.

\subsection{Directed frame candidate and finite spectral assembly}
\label{p6:app:v585-directed-assembly}
\noindent\textbf{Comparison scope.} These fixed-candidate matching and
spectral results are retained as conditional diagnostics. The native charged
source compiler is the main-text mass/incidence synthesis of Release v5.90
and Paper~V-A Release v3.97; terminal masses are not obtained from this
candidate's singular values. No charged SVD assembly prerequisite survives.
This appendix treats an explicitly defined candidate matrix branch. It
supplies exact algebra and finite reconstruction certificates, not a
source-origin theorem. The coherent word, its relative coefficient and
its terminal-role matching are not yet derived as physical action outputs.
The parameter $t$ below is a diagnostic variable, not an added physical
coupling. The positive structural support theorem is in the main text.

\subsubsection{Directed positive return and its coherent cross term}
Let $J$ be the all-ones matrix and retain
\[
 K=aI+J,\quad X=\begin{pmatrix}0&0&1\\1&0&0\\0&1&0\end{pmatrix},
 \quad R=\operatorname{diag}(1,-i,i),\quad
 A_+=KR-itX,\quad A_-=KR^\dagger-itX.
\]
The component branches are $a_\ell=1$ and $a_C=2$; the tested coefficient
is $t=1$. Put $R_+=R$, $R_-=R^\dagger$ and $H_\sigma=A_\sigma A_\sigma^\dagger$.
\begin{proposition}[Full rank and exact candidate invariants]
\label{p6:prop:v585-directed-invariants}
For $a>0$ and real $t$, both directed maps are invertible, with
\[
 \det A_\pm=a^2(a+3)-t^2+it(a+t^2),
\]
and
\[
 H_\sigma=K^2+t^2I+it(KR_\sigma X^\dagger-XR_\sigma^\dagger K).
\]
Their common characteristic polynomial is $z^3-e_1z^2+e_2z-e_3$, where
\begin{align*}
 e_1&=3(a^2+2a+t^2+3),\\
 e_2&=3a^4+12a^3+3a^2t^2+18a^2+8at^2+3t^4+10t^2,\\
 e_3&=[a^2(a+3)-t^2]^2+t^2(a+t^2)^2.
\end{align*}
At $(a,t)=(1,1)$ and $(2,1)$ these are, respectively,
\[
 z^3-21z^2+57z-13,\qquad z^3-36z^2+257z-370.
\]
For $\chi(H_+,H_-)=(6i)^{-1}\operatorname{tr}[H_+,H_-]^3$,
\begin{align*}
 \chi={}&16at^4(2a+3)\bigl(4a^4+a^3t^2+18a^3
                    +2a^2t^2-4at^2-27a-4t^2\bigr),\\
 \chi|_{a=1}&=-400t^4(1+t^2),\qquad
 \chi|_{a=2}=448t^4(77+2t^2).
\end{align*}
The corresponding spectral discriminants are
\[
 13500t^2(t^2+5)^2,\qquad
 784t^2(112t^4+1656t^2+6615).
\]
Thus both component branches have simple spectra and nonzero $\chi$ for
$t\ne0$; at $t=1$, $\chi=-800$ and $35392$.
\end{proposition}
\begin{proof}
Expansion gives the determinant. For $t\ne0$ its imaginary part is
$t(a+t^2)\ne0$; for $t=0$ its real part is $a^2(a+3)>0$.
Multiplication gives the Gram formula. Compute $e_1=\operatorname{tr}H$,
$e_2=((\operatorname{tr}H)^2-\operatorname{tr}H^2)/2$ and
$e_3=|\det A|^2$ to obtain the three coefficients.
Expanding the commutator cube gives the stated polynomial. Substitution
in the cubic discriminant gives the two positive expressions. The
companion exact-check script verifies these polynomial identities and
independently reconstructs the two integer Gaussian witnesses.
\end{proof}
The incoherent sum of the two constituent Grams is just $K^2+t^2I$,
independent of sector. It omits the cross term responsible for the
directed pair's relative frame. The displayed nonzero invariant is CP
capacity of this candidate, not a physical CKM or PMNS prediction.

\subsubsection{One map generates both candidate frames}
\begin{theorem}[Opposite Gram sides of one rooted map]
\label{p6:thm:v585-one-parent-frames}
Let $S$ interchange trimer coordinates 1 and 2, fixing coordinate 0, and
$P=XS=\left(\begin{smallmatrix}0&1&0\\1&0&0\\0&0&1\end{smallmatrix}\right)$.
Right-unitary removal gives
\[
 B_+=A_+R^\dagger=K-itXR^\dagger,
 \qquad B_-=A_-R=K-itXR=PB_+^TP.
\]
Consequently
\[
 H_+=B_+B_+^\dagger,\qquad
 H_-=P\overline{B_+^\dagger B_+}P.
\]
If $B_+=U\Lambda^{1/2}V^\dagger$, the two left frames are $U$ and
$P\overline V$, with common spectral labels $\Lambda$.
\end{theorem}
\begin{proof}
The entries are
\[
 B_+=\begin{pmatrix}a+1&1&1-t\\1-it&a+1&1\\1&1+t&a+1\end{pmatrix},\quad
 B_-=\begin{pmatrix}a+1&1&1+t\\1-it&a+1&1\\1&1-t&a+1\end{pmatrix}.
\]
These prove the transpose identity. Multiplication by the adjoint gives
the conjugate right Gram, and substitution of the SVD gives the frames.
\end{proof}
At $a=2,t=1$, $B_+=\left(\begin{smallmatrix}3&1&0\\1-i&3&1\\1&2&3\end{smallmatrix}\right)$.
This is a family-matrix identity. It does not identify a right-chiral
particle space with a down-type left-chiral space.

\subsubsection{Six-member classification and exact directed reconstruction}
\begin{theorem}[Finite spectral assembly on a fixed simple frame branch]
\label{p6:thm:v585-finite-spectral-assembly}
Let $A_s$ be an invertible $3\times3$ matrix, $H_s=A_sA_s^\dagger$ have
simple eigenvalues $0<\lambda_1<\lambda_2<\lambda_3$, and let the terminal
packet have distinct positive magnitudes $\mu_s(r)$, $r\in\mathcal I_s$.
Set
\[
 E_{si}=\prod_{j\ne i}\frac{H_s-\lambda_jI}{\lambda_i-\lambda_j}.
\]
Exactly six positive Hermitian matrices commute with $H_s$ and have the
terminal squared-magnitude spectrum. They are
\[
 \widehat H_s^{(\tau)}=\sum_i\mu_s(\tau(i))^2E_{si},
 \qquad \tau:\{1,2,3\}\mathrel{\mathop{\longrightarrow}\limits^{\sim}}\mathcal I_s.
\]
For each matching there are unique real polynomials of degree at most two,
\begin{align*}
 f_s^{(\tau)}(z)&=\sum_i\mu_s(\tau(i))^2
                   \prod_{j\ne i}\frac{z-\lambda_j}{\lambda_i-\lambda_j},\\
 q_s^{(\tau)}(z)&=\sum_i\frac{\mu_s(\tau(i))}{\sqrt{\lambda_i}}
                   \prod_{j\ne i}\frac{z-\lambda_j}{\lambda_i-\lambda_j},
\end{align*}
and the directed representative
\[
 \widehat R_s^{(\tau)}=q_s^{(\tau)}(H_s)A_s
 \quad\hbox{satisfies}\quad
 \widehat R_s^{(\tau)}(\widehat R_s^{(\tau)})^\dagger
 =f_s^{(\tau)}(H_s)=\widehat H_s^{(\tau)}.
\]
Equivalently, for $p_s(z)=\det(zI-H_s)$,
$zq_s(z)^2\equiv f_s(z)\pmod{p_s(z)}$.
\end{theorem}
\begin{proof}
A Hermitian matrix commuting with a simple Hermitian $H_s$ acts by a real
scalar on each of its one-dimensional eigenspaces. The prescribed
distinct spectrum can be assigned in exactly $3!=6$ ways. Lagrange
interpolation at the three distinct real nodes proves existence and
uniqueness of $f_s,q_s$. On each projector, $\lambda_i q_s(\lambda_i)^2$
equals $\mu_s(\tau(i))^2$, proving the Gram and congruence identities.
\end{proof}
The directed branches above meet these hypotheses at $t=1$. For two
sectors there are 36 assignment pairs.  In general one must still test
common weak-basis equivalence; for the stated raw quark packets on the
$a=2,t=1$ pair, Theorem~\ref{p6:thm:v2109-thirtysix-inequivalent} proves
that all thirty-six are inequivalent. No continuous
interpolation coefficient remains after a matching. Positivity is claimed
on the finite spectrum, not on the entire real axis. Physical use still
requires a constructor deriving $A_s,\tau_s$ and the weighted response
together; $q_s(H_s)$ is not introduced as a kinetic metric or an RG repair.

\subsubsection{CP preservation under distinct reweighting}
\begin{theorem}[Vandermonde law for the commutator invariant]
\label{p6:thm:v585-cp-reweighting}
Write $\Delta(z)=\prod_{i<j}(z_i-z_j)$. For the candidate pair with common
simple eigenvalues $\lambda_i$ and assigned terminal weights
$x_i=\mu_u(\tau_u(i))^2$, $y_i=\mu_d(\tau_d(i))^2$,
\[
 \chi(\widehat H_u,\widehat H_d)=\chi(H_+,H_-)
                  \frac{\Delta(x)\Delta(y)}{\Delta(\lambda)^2}.
\]
Hence every distinct spectral reweighting preserves
\[
 \mathcal J^2(H_1,H_2)=
 \frac{\chi(H_1,H_2)^2}
 {\operatorname{disc}\det(zI-H_1)\,\operatorname{disc}\det(zI-H_2)}.
\]
At $t=1$ the candidate values are
$|\mathcal J|_{a=1}=2/1215$ and $|\mathcal J|_{a=2}=316/58681$.
\end{theorem}
\begin{proof}
In the first eigenbasis the commutator entries are $(x_i-x_j)B_{ij}$.
Its cubic trace is a fixed signed multiple of $\Delta(x)$ times
$\operatorname{Im}(B_{12}B_{23}B_{31})$. For fixed second eigenvectors,
the latter is a homogeneous cubic in $y$. If two $y_i$ coincide, the
off-diagonal part is a real scalar times a rank-one Hermitian projector,
whose triple product is real. The cubic thus vanishes on all three
coincidence planes and is a constant times $\Delta(y)$. Taking the ratio
with the original spectrum proves the formula with its permutation signs.
The discriminants are the squared Vandermondes, so they cancel in
$\mathcal J^2$. The exact values follow from
Proposition~\ref{p6:prop:v585-directed-invariants}.
\end{proof}
The sign of $\chi$ can change with matching parity. Its normalized
magnitude is fixed by the candidate frames and is not an observed mixing
value. At $a=2,t=1$ the complete low mixed-moment matrix is
\[
 \bigl(\operatorname{tr}(H_+^mH_-^n)\bigr)_{m,n=0}^2
 =\begin{pmatrix}3&36&782\\36&710&18338\\782&18114&481162\end{pmatrix}.
\]
Together with the displayed cubic, it fixes all projector overlaps by
inserting the quadratic projector polynomials. As an exact verification
witness, $(H_++I)^2$ and $(H_-+I)(H_-+2I)$ have traces $857,896$,
$\chi=6180481742848$, and unchanged $|\mathcal J|=316/58681$.
These witness polynomials are not a physical mass law.

\subsubsection{Finite acceptance criterion for an independently generated output}
\begin{corollary}[Two additional moments recover the role matching]
\label{p6:cor:v585-action-binding-test}
Suppose $H_s^{\rm act}$ is an independently source-generated positive
return. If $[H_s^{\rm act},H_s]=0$ and its three characteristic coefficients
equal the terminal packet's, it belongs to the six-member family above.
Its matching is uniquely recovered from
$b_k=\operatorname{tr}(H_s^kH_s^{\rm act})$, $k=0,1,2$, by
\[
 \begin{pmatrix}b_0\\b_1\\b_2\end{pmatrix}
 =\begin{pmatrix}1&1&1\\\lambda_1&\lambda_2&\lambda_3\\
                 \lambda_1^2&\lambda_2^2&\lambda_3^2\end{pmatrix}
   \begin{pmatrix}x_1\\x_2\\x_3\end{pmatrix}.
\]
Since $b_0$ is fixed by the packet trace, only two additional scalar mixed
moments are required jointly.
\end{corollary}
\begin{proof}
The classification theorem supplies the simultaneous eigenprojectors.
Taking traces gives the displayed system. Its Vandermonde determinant is
nonzero, so the assigned weights and their distinct role labels are unique.
\end{proof}
The moments must be derived readouts of the source output. Declaring them
does not select or actualize a source. On this route the remaining physical
theorem is one licensed word producing $(A_s,\tau_s,q_s(H_s)A_s)$ on the
occurring root. The finite reduction applies after fixing this candidate
branch; it does not reduce the entire source-constructor problem to 36 choices.

\subsubsection{Matching and undressed symmetrization diagnostics}
An isolated neutral $D_3$ matching orbit does not determine a relative
charged frame. For example, for $H_u=\operatorname{diag}(1,4,9)$ and
$H_d=\operatorname{diag}(2,3,5)$ the mixed trace is $59$; permuting only
the first two weights of $H_d$ changes it to $56$. Simultaneously
permuting weights and their frame map leaves the operator fixed, as in
Corollary~\ref{p6:cor:v585-left-support-quotient}. Moreover, an operator
commuting with both the simple marked $R$ and cyclic $X$ is scalar:
commutation with $R$ makes it diagonal, and commutation with $X$ equates
all diagonal entries. Arbitrary permutations therefore cannot be quotiented
while freezing all the rooted marks.

For the undressed symmetrization $W_\sigma=A_\sigma+A_\sigma^\dagger$,
the same reflection $S$ satisfies $S\overline{W_\sigma}S=W_\sigma$.
It provides a common antiunitary also for $W_\sigma^2$, so
\[
 \operatorname{tr}[W_+^2,W_-^2]^3=0.
\]
Indeed conjugation by that antiunitary makes the trace equal to its complex
conjugate, while a commutator of Hermitian matrices has a purely imaginary
cubic trace. This is a diagnostic for the undressed pair only; an extra
hierarchy mark may break the common real structure. At $t=1$,
$\operatorname{tr}(AA^\dagger)-\operatorname{tr}(A+A^\dagger)^2
=-a^2-2a-14$, another direct check that the two completions differ.

The accompanying \texttt{Paper00-FermionSVDAssemblyChecks-v1\_1.py}
checks the exact identities, support compressions, polynomial interpolation
and independent Gaussian-integer witnesses. Its finite checks do not
certify physical action origin, matching, scheme/vertex completion or poles.


% BEGIN INLINE SOURCE: Paper5B-OneMomentAssembly-v1_0.tex
% Integrated from Paper00-FermionOneMomentMatching-v1_0, Sections 2--5,
% and Paper00-RemainingGapTransportAnalysis-v1_0, CLS-T02--T03.
% Candidate appendix: fixed t=1 branch; no physical source selection.
\subsubsection{One-moment reconstruction on the native raw charged packets}
\label{p6:app:v2109-one-moment-assembly}
Retain the directed candidate of Theorem~\ref{p6:thm:v585-one-parent-frames}
at \(t=1\), with \(a=1\) for the charged-lepton sector and \(a=2\) for
colour.  Write \(H=A A^\dagger=\sum_i\lambda_iE_i\) with ordered, simple
positive eigenvalues.  The two characteristic polynomials are
\[
 p_\ell(z)=z^3-21z^2+57z-13,\qquad
 p_C(z)=z^3-36z^2+257z-370.
\]
The raw squared terminal packets, in the retained role order, are
\[
\begin{aligned}
 x_u&=(2,36,16+6\sqrt7),\\
 x_d&=\left(5+2\sqrt6,\frac16,\frac{33-20\sqrt2}{289}\right),\\
 x_e&=\left(\frac{3-\sqrt5}{8},\frac14,36+24\sqrt2\right).
\end{aligned}
\]
These are structural factors, not measured pole masses.  Every entry lies
in the degree-sixteen field
\[
K=\mathbb Q(\sqrt2,\sqrt3,\sqrt5,\sqrt7),
\]
and each packet is strictly positive with distinct entries.

\begin{theorem}[One mixed moment separates the native assignments]
\label{p6:thm:v2109-one-moment-matching}
Suppose a positive return \(M\) is independently known to commute with the
fixed \(H\) and to have one of the displayed packet spectra.  The single
mixed trace \(b_1=\operatorname{tr}(HM)\) determines its assignment
\(M=M_\tau=\sum_i x_{\tau(i)}E_i\), \(\tau\in S_3\).  In fact the
separation statement holds for any three distinct weights in \(K\).
\end{theorem}
\begin{proof}
The lepton cubic modulo seven is \(z^3+z+1\), with nonzero residues
\(1,3,4,3,6,5,6\).  The colour cubic modulo three is \(z^3+2z+2\),
with residue two at all three elements.  Both are irreducible over
\(\mathbb Q\).  A root in \(K\) would have degree three dividing sixteen;
therefore both cubics remain irreducible over \(K\).  Their splitting-field
Galois groups over \(K\) act transitively on the three roots and contain a
three-cycle.  If two assignments have equal trace, put
\(d_i=x_{\tau(i)}-x_{\sigma(i)}\in K\).  Applying that cycle to
\(\sum_i\lambda_id_i=0\) gives a cyclic three-by-three system whose
determinant is
\[
 -\frac{\lambda_1+\lambda_2+\lambda_3}{2}
 \bigl((\lambda_1-\lambda_2)^2+(\lambda_2-\lambda_3)^2
 +(\lambda_3-\lambda_1)^2\bigr)\ne0.
\]
Positivity and simplicity follow, without numerical approximation, by the
sign-changing isolating intervals \((0,1),(2,3),(17,18)\) for \(p_\ell\)
and \((1,2),(7,8),(26,27)\) for \(p_C\).  Thus \(d_i=0\), and distinct
weights give \(\tau=\sigma\).
\end{proof}

\begin{theorem}[The allowed trace reconstructs the entire weighted return]
\label{p6:thm:v2109-one-moment-full-return}
Set \(s=\operatorname{tr}H\), \(r=\sum_i x_i\),
\(L=H-sI/3\), \(\ell_i=\lambda_i-s/3\), \(y_i=x_i-r/3\), and
\[
\begin{aligned}
 V_H&=\sum_i\ell_i^2,& C_H&=\sum_i\ell_i^3,\\
 V_x&=\sum_i y_i^2,& C_x&=\sum_i y_i^3,\\
 \Delta_H&=V_H^3/2-3C_H^2,&
 \Delta_x&=V_x^3/2-3C_x^2.
\end{aligned}
\]
For \(z=b_1-sr/3\), the six admissible traces obey
\begin{equation}
 \label{p6:eq:v2109-trace-sextic}
 \left(z^3-\frac34V_HV_xz-\frac32C_HC_x\right)^2
 =\frac14\Delta_H\Delta_x.
\end{equation}
On this native domain define
\[
 w=\frac{2C_HV_xz+V_H^2C_x}{4z^2-V_HV_x},\qquad
 D_0=V_H^3/6-C_H^2=\Delta_H/3>0,
\]
\[
 \alpha=\frac{V_H^2z/6-C_Hw}{D_0},\qquad
 \beta=\frac{V_Hw-C_Hz}{D_0}.
\]
Then
\begin{equation}
 \label{p6:eq:v2109-full-return-reconstruction}
 M=\frac r3 I+\alpha L+\beta\left(L^2-\frac{V_H}{3}I\right),
 \qquad b_2=\operatorname{tr}(H^2M)
 =w+\frac{2s}{3}z+\frac r3\operatorname{tr}(H^2).
\end{equation}
\end{theorem}
\begin{proof}
For zero-sum triples, expanding the three even and three odd assignment
values of \(\sum_i\ell_i y_{\tau(i)}\) gives their two monic orbit cubics
\(F(z)\mp d_Hd_x/2\), where
\(F(z)=z^3-3V_HV_xz/4-3C_HC_x/2\), and \(d_H,d_x\) are the
Vandermonde products.  Their product proves
\eqref{p6:eq:v2109-trace-sextic} because \(d_H^2=\Delta_H\) and
\(d_x^2=\Delta_x\).  This identity can alternatively be checked by direct
expansion of the six linear forms; it makes no branch choice.
Under the formal algebraic variation
\(\delta\ell_i=\ell_i^2-V_H/3\), with the assigned \(y_i\) fixed,
\[
 \delta z=w,\quad \delta V_H=2C_H,\quad
 \delta C_H=V_H^2/2,\quad \delta d_H=0.
\]
The last equality follows from
\(\delta\log d_H=\sum_{i<j}(\ell_i+\ell_j)=0\).
Differentiating the applicable orbit cubic yields the stated formula for
\(w=\operatorname{tr}(L^2(M-rI/3))\).
Its denominator cannot vanish at an allowed trace: the sextic has six
simple roots by Theorem~\ref{p6:thm:v2109-one-moment-matching},
\(F\ne0\) at each root, and its derivative is
\(2FF'=3F(4z^2-V_HV_x)/2\).
Finally \(\operatorname{tr}L^4=V_H^2/2\), so the Gram matrix of the
traceless basis \(L,L^2-V_HI/3\) is
\(\left(\begin{smallmatrix}V_H&C_H\\C_H&V_H^2/6\end{smallmatrix}\right)\),
with determinant \(D_0\).  Solving its two moment equations gives
\(\alpha,\beta\).  The simple spectrum makes this traceless commutant
two-dimensional, proving the reconstruction and the formula for \(b_2\).
\end{proof}
The inherited spectral constants are
\[
 (s,V_H,C_H,\Delta_H)_\ell=(21,180,900,486000),\qquad
 (s,V_H,C_H,\Delta_H)_C=(36,350,2226,6572272).
\]
The sextic labels all six alternatives; choosing one of its roots is not a
physical constructor.  The action must still generate \(H\), packet
membership, commutation and the value of \(b_1\).  Common sector rescaling
preserves separation.  Arbitrary role-dependent hierarchy dressing is not
covered without a new coefficient-field or noncollision proof; the general
two-moment criterion remains valid on its original domain.

\begin{theorem}[Thirty-six inequivalent raw quark-pair assemblies]
\label{p6:thm:v2109-thirtysix-inequivalent}
On the \(a=2,t=1\) directed quark pair, the thirty-six raw weighted pairs
\((M_u^\tau,M_d^\upsilon)\) are inequivalent under common unitary
conjugation.  The single joint invariant
\(\operatorname{tr}(M_u^\tau M_d^\upsilon)\) distinguishes them.
\end{theorem}
\begin{proof}
In the fixed candidate coordinates one may use
\[
 B_+=\begin{pmatrix}3&1&0\\1-i&3&1\\1&2&3\end{pmatrix},\qquad
 B_-=\begin{pmatrix}3&1&2\\1-i&3&1\\1&0&3\end{pmatrix},
 \qquad H_\pm=B_\pm B_\pm^\dagger.
\]
Exact Gaussian-integer multiplication gives, for \(0\le m,n\le2\),
\[
 (\operatorname{tr}(H_+^mH_-^n))_{mn}
 =\begin{pmatrix}3&36&782\\36&710&18338\\782&18114&481162\end{pmatrix}.
\]
Use \(E_{\pm i}=\prod_{j\ne i}(H_\pm-\lambda_jI)/(\lambda_i-\lambda_j)\)
and \(P_{ij}=\operatorname{tr}(E_{+i}E_{-j})\).  Every joint invariant is
\(\sum_{ij}x_{u,\tau(i)}P_{ij}x_{d,\upsilon(j)}\).
The finite certificate in
\texttt{Paper00-CrossChatClosureChecks-v1\_0.py} isolates the three cubic
roots and the packet radicals by rational bisection, evaluates these
projector expressions by rational interval arithmetic, and checks all
thirty-six resulting intervals.  Their consecutive gaps exceed
\(47/1000\); thus no two traces agree.  The script and its complete
interval output are release companions.  Trace invariance under common
unitary conjugation proves the conclusion.
\end{proof}
This classifies the stated raw candidate family only; it neither selects the
physical pair nor bounds all possible same-action source words by thirty-six.


% BEGIN FRAMEWORK v2.240: EXACT DIRECTED-COLOUR CROSS-GRAM AUDIT
\begin{proposition}[Exact mixed Gram and two-variable candidate Stieltjes probe]
\label{p6:prop:v240-directed-colour-double-resolvent}
For the \emph{raw candidate} $H_\pm=B_\pm B_\pm^\dagger$ printed in
Theorem~\ref{p6:thm:v2109-thirtysix-inequivalent}, put
$L_\pm=H_\pm-12I$ and $Q_\pm=L_\pm^2-350I/3$.
Direct exact Gaussian-integer multiplication gives the corrected
traceless cross-Gram matrix
\begin{equation}\label{p6:eq:v240-cross-gram}
 \boxed{\begin{pmatrix}
 \operatorname{tr}(L_+L_-) & \operatorname{tr}(L_+Q_-)\\
 \operatorname{tr}(Q_+L_-) & \operatorname{tr}(Q_+Q_-)
 \end{pmatrix}
 =\begin{pmatrix}278&2282\\2058&39098/3\end{pmatrix}.}
\end{equation}
Let $D(z)=z^3+36z^2+257z+370$. Then the entire candidate
mixed-resolvent series is the explicit rational function
\begin{equation}\label{p6:eq:v240-candidate-double-resolvent}
 \boxed{\operatorname{tr}[(zI+H_+)^{-1}(wI+H_-)^{-1}]
       =\frac{N(z,w)}{D(z)D(w)},}
\end{equation}
where
\[
 \boxed{\begin{aligned}
 N(z,w)={}&3z^2w^2+72z^2w+72zw^2
     +257(z^2+w^2)\\
    &+2006zw+7222z+7446w+23001.
 \end{aligned}}
\]
If $\lambda_1<\lambda_2<\lambda_3$ are the simple eigenvalues
of $H_\pm$, then their exact candidate projector overlaps are
\[
 \boxed{P^{(H)}_{ij}=\frac{N(-\lambda_i,-\lambda_j)}
                    {D'(-\lambda_i)D'(-\lambda_j)}.}
\]
Their oriented raw CP invariant obeys
\[
 \boxed{\frac1{6i}\operatorname{tr}[H_+,H_-]^3=35392,
       \qquad |J_H|=\frac{35392}{6572272}=\frac{316}{58681}.}
\]
These algebraic candidate results do not evaluate the physical
$C_{ud}=Z_u^{-1/2}E_u^\dagger E_dZ_d^{-1/2}$.
\end{proposition}
\begin{proof}
Use the exact mixed moments
$\left(\operatorname{tr}(H_+^mH_-^n)\right)_{m,n=0}^2$
printed in Theorem~\ref{p6:thm:v2109-thirtysix-inequivalent}.
Expanding the centered powers gives the displayed cross Gram.
The adjugate identity
$(zI+H_\pm)^{-1}=\operatorname{adj}(zI+H_\pm)/D(z)$
and exact $3\times3$ products give $N$.
Residues at the distinct spectral poles reconstruct the overlap
projectors. Direct commutator multiplication gives $35392$; divide
by the product of the two ordered eigenvalue Vandermondes, here
$\Delta_H=6572272$, to obtain $|J_H|$.
\end{proof}
\begin{remark}[Polar incidence is not the left spectral overlap]
For an illustrative full-rank matrix $B_s=U_sS_sV_s^\dagger$,
its positive left spectral overlap is governed by
$U_u^\dagger U_d$, whereas the normalized \emph{polar incidence}
$W_s=B_s(B_s^\dagger B_s)^{-1/2}=U_sV_s^\dagger$ gives
\[
 \boxed{W_u^\dagger W_d
       =V_u(U_u^\dagger U_d)V_d^\dagger.}
\]
The right singular-frame factors need not be diagonal phases;
therefore the above $P^{(H)}$ and $J_H$ cannot be substituted for
the actual mixed source-current invariants of
Theorem~\ref{p5:thm:v240-five-quark-current-invariants}.
\end{remark}
\begin{remark}[Exact raw terminal-packet arithmetic]
\label{p6:rem:v240-raw-quark-packet-corrections}
For the raw terminal squared weights already printed in this appendix,
\[
 x_u=(2,36,16+6\sqrt7),\qquad
 x_d=\left(5+2\sqrt6,\frac16,\frac{33-20\sqrt2}{289}\right),
\]
put $r_s=\sum_i x_{s,i}$, $y_{s,i}=x_{s,i}-r_s/3$,
$V_{x_s}=\sum_i y_{s,i}^2$ and $C_{x_s}=\sum_i y_{s,i}^3$.
Direct radical arithmetic gives
\[
 \boxed{\begin{aligned}
 r_u&=54+6\sqrt7,&V_{x_u}&=752-24\sqrt7,\\
 C_{x_u}&=1224-3096\sqrt7,&
 \prod_{i<j}(x_{u,i}-x_{u,j})&=952+1224\sqrt7,\\
 r_d&=\frac{9157+3468\sqrt6-120\sqrt2}{1734}.
 \end{aligned}}
\]
The discriminant identity
$\prod_{i<j}(x_{u,i}-x_{u,j})^2
=V_{x_u}^3/2-3C_{x_u}^2$ holds exactly.
These values correct intermediate arithmetic occasionally inferred from
the raw packet; they do not change the source-origin boundary or
select one of the thirty-six weighted candidate pairs.
\end{remark}

% END FRAMEWORK v2.240: EXACT DIRECTED-COLOUR CROSS-GRAM AUDIT

% END INLINE SOURCE: Paper5B-OneMomentAssembly-v1_0.tex

























% BEGIN CHARGED-LEPTON-INCIDENCE-EVALUATION-CORRECTION v5.95
\subsubsection{Orthogonal cell-label comparison}

Declare the orthogonal label carrier
\[
\mathcal L_{25}=\CC[\{\mathcal C_{ij}\}],
\qquad
\Omega_{25}=\frac15\sum_{i,j}|\mathcal C_{ij}\rangle,
\]
use any representative activated label \(\mathcal C_0\) of the transitive
label orbit, and put
\[
\Pi_{\rm exc}=I-|\mathcal C_0\rangle\langle\mathcal C_0|,
\qquad
A=\Pi_{\rm exc}|\Omega_{25}\rangle\langle\Omega_{25}|.
\]
Then
\[
\etaH:=\frac{24}{25},
\qquad
\sigma_1(A)=\sqrt{\etaH},
\qquad
\sigma_1(A^+)=\etaH^{-1/2}.
\]
All activated labels are related by the transitive label action, so the
identity of \(\mathcal C_0\) is an unmarked representative for these invariant
readouts; it is not a physical cell selector.

\begin{theorem}[Factor-count and word no-go]
\label{p6:thm:comparison-factor-no-go}
The cell design selects neither exactly three independent copies of \(A\)
nor a comparison word for any particle.  For every \(N\ge1\) and every
\(e\in\{-1,0,1\}^N\), an \(N\)-fold tensor construction has active singular
value
\[
\etaH^{\frac12\sum_{j=1}^Ne_j}.
\]
All such choices preserve the same cell action on each factor.
\end{theorem}

\begin{proof}
On the active line, use \(A\), \(I\), or \(A^+\) according as
\(e_j=1,0,-1\).  Tensor-product singular values multiply.  The original
incidence geometry contains no datum fixing \(N\) or the values \(e_j\).
\end{proof}

\begin{theorem}[Universal comparison-output factorization]
\label{p6:thm:comparison-output-factorization}
Let
\[
\mathcal W
=
\coprod_{N\ge1}\{-1,0,1\}^{N}
\]
be the host family of finite comparison words and define
\[
q:\mathcal W\longrightarrow\tfrac12\ZZ,
\qquad
q(N,e)=\frac12\sum_{j=1}^{N}e_j.
\]
The active singular-value evaluation factors uniquely as
\[
\mathcal W\xrightarrow{\ q\ }\tfrac12\ZZ
\xrightarrow{\ n\mapsto\etaH^n\ }\RR_{>0}.
\]
Moreover, because \(0<\etaH<1\), two words have the same current Paper--V-B
output if and only if they have the same value of \(q\).  Hence the quotient
of all host factorizations by current-output equivalence is exactly
\(\tfrac12\ZZ\); factor number and factor names are not identifiable.
\end{theorem}

\begin{proof}
Theorem~\ref{p6:thm:comparison-factor-no-go} gives
\(\sigma(W_e)=\etaH^{q(N,e)}\), proving the factorization.  The exponential
map \(n\mapsto\etaH^n\) is injective because \(0<\etaH<1\).  Finally, every
half-integer is reached by a word of repeated \(+1\) or \(-1\), with the zero
word covering \(n=0\), so the stated quotient is all of \(\tfrac12\ZZ\).
\end{proof}

\begin{definition}[Retained comparison-exponent package]
\label{p6:def:selected-comparison-package}
For every current Paper--V-B readout, the minimized package \(\CMPexp\) retains
only the function
\[
n_{\etaH}:\mathcal P_{\rm mass}\cup\{g,\gamma\}
\longrightarrow\tfrac12\ZZ
\]
printed in Table~\ref{p6:tab:conditional-particle-ledger}, and sets
\(w_p=\etaH^{n_{\etaH}(p)}\).  A three-factor
orientation/completion/boundary word is merely one lift through \(q\) and is
not part of the native package unless a later factor-sensitive observable is
constructed.
\end{definition}


\begin{remark}[Charged comparison package superseded below]
\label{p6:rem:v591-charged-package-supersession-forward}
For the nine charged fermions, Release v5.91 below replaces the historical
particle-by-particle exponent input by a derived owner-provenance polar grade.
The present host-factorization and no-go results remain valid for arbitrary
word lifts and for rows not covered by that charged theorem.
\end{remark}

\begin{theorem}[Conditional comparison grading]
\label{p6:thm:conditional-comparison-grading}
Under Definition~\ref{p6:def:selected-comparison-package}, each current
comparison scalar is the single value
\[
w_p=\etaH^{n_{\etaH}(p)}.
\]
There is no factor-count, factor-name, word, or additional continuous
coefficient in the current-output evaluation after the exponent ledger is
fixed.
\end{theorem}

\begin{proof}
This is the evaluation on the quotient identified in
Theorem~\ref{p6:thm:comparison-output-factorization}.  Any two host lifts of
the same exponent have the same singular-value output.
\end{proof}


\section{Mass transfers, flavor operators, clock and countermodel calculations}
\label{p5d:sec:D}
\noindent\textit{Technical owner.} These exact inherited calculations
support the corresponding statements in Papers~V-A--V-C;
the appropriate action/observable occurrence conditions remain in those papers.
\subsection{Hurwitz walk response and scalar normalization}
\label{p6:sec:scalar-normalization}

In quaternion coordinates write
\[
2T
=
\{\pm1,\pm i,\pm j,\pm k\}
\cup
\left\{
\frac{\varepsilon_0+\varepsilon_1i+\varepsilon_2j+\varepsilon_3k}{2}
:\varepsilon_a\in\{\pm1\}
\right\}.
\]
Its nearest nonidentity shell is
\[
S_T
=
\left\{
\frac{1+\varepsilon_1i+\varepsilon_2j+\varepsilon_3k}{2}
:\varepsilon_a\in\{\pm1\}
\right\},
\qquad |S_T|=8.
\]
Let \(L_s\) denote left translation on \(\CC[2T]\) and define
\[
T:=\frac18\sum_{s\in S_T}L_s .
\]
The shell is inverse closed and conjugacy invariant, so \(T\) is
self-adjoint and central.  A direct class-sum character calculation gives
\pvcite{Kugelmann2026a}
\[
\Spec(T)
=
1^{(1)}\oplus(1/2)^{(4)}\oplus0^{(9)}
\oplus(-1/4)^{(8)}\oplus(-1/2)^{(2)}.
\]
Set
\[
\begin{aligned}
P_{15}
&:=I-E_0\\
&=
-16T^4+12T^3+8T^2-3T,
\end{aligned}
\qquad
\tau_{24}=\frac1{24}\Tr .
\]
The polynomial vanishes at \(0\) and equals one at every nonzero displayed
eigenvalue.

\begin{theorem}[Complete walk-projector classification]
\label{p6:thm:walk-projector-classification}
Every orthogonal projector in \(C^*(T)\) is a sum of spectral projectors.
Among projectors below \(P_{15}\), the possible ranks are the unique subset
sums of \(1,2,4,8\), hence every integer from \(0\) through \(15\) occurs
once.  \(P_{15}\) is the unique maximal nonzero-response support projector,
but it is not the only central nonzero-response projector.
\end{theorem}

\begin{proof}
Functional calculus gives
\(
C^*(T)=\bigoplus_{\lambda\in\Spec(T)}\CC E_\lambda
\).
Idempotents select subsets of the five spectral values.  The nonzero
multiplicities \(1,2,4,8\) give binary expansions of \(0,\ldots,15\).
\end{proof}

\begin{theorem}[State dependence of the response fraction]
\label{p6:thm:response-state-dependence}
\[
\tau_{24}(P_{15})=\frac{15}{24}
\]
is the expectation in the maximally mixed state.  For the invariant vector
\[
\Omega_{24}=\frac1{\sqrt{24}}\sum_{g\in2T}|g\rangle,
\]
one has
\[
T\Omega_{24}=\Omega_{24},
\qquad
P_{15}\Omega_{24}=\Omega_{24},
\]
so its response is \(1\).  A unit kernel vector has response \(0\).
\end{theorem}

\begin{proof}
The trace identity is immediate.  The uniform vector spans the eigenvalue-one
line, while kernel states are annihilated by \(P_{15}\).
\end{proof}

\begin{theorem}[Canonical invariant scalar embedding and state]
\label{p6:thm:canonical-invariant-scalar-state}
For the left regular \(2T\)-action,
\[
\operatorname{Hom}_{2T}(\CC,\CC[2T])
=\CC\Omega_{24}.
\]
Consequently every invariant isometry
\(J:\CC\to\CC[2T]\) is
\(J_\phi(z)=e^{i\phi}z\Omega_{24}\).  Its range projector and induced pure
state are independent of \(\phi\):
\[
J_\phi J_\phi^*=|\Omega_{24}\rangle\langle\Omega_{24}|,
\qquad
\omega_H(A)=\langle\Omega_{24},A\Omega_{24}\rangle.
\]
The range lies in \(\operatorname{Ran}P_{15}\).  Thus the invariant embedding
projector and carrier state are theorem outputs, not scalar-normalization
selector coordinates.
\end{theorem}

\begin{proof}
An invariant vector in the regular permutation carrier is constant on the
single transitive orbit, so the invariant subspace is the line
\(\CC\Omega_{24}\).  Isometry forces the coefficient to have modulus one.
The phase cancels in both \(J_\phi J_\phi^*\) and the vector state.  Finally,
\(P_{15}\Omega_{24}=\Omega_{24}\) by
Theorem~\ref{p6:thm:response-state-dependence}.
\end{proof}

\begin{corollary}[Exact conditional trace conversion]
\label{p6:cor:conditional-trace-conversion}
For \(\xi\in\Ran P_{15}\) and \(A_\xi=|\xi\rangle\langle\xi|\), define
\[
\tau_{15}(B)=\frac1{15}\Tr_{\Ran P_{15}}B.
\]
Then
\[
\sqrt{\tau_{15}(A_\xi)}
=
\sqrt{\frac{24}{15}}\sqrt{\tau_{24}(A_\xi)}.
\]
Thus \(\sqrt{24/15}\) is a global-to-conditional RMS conversion, not a
restricted Hilbert norm.
\end{corollary}

\begin{theorem}[Circle-normalization trichotomy]
\label{p6:thm:circle-normalization}
For a constant-norm Hopf-circle orbit, Haar probability \(d\mu=d\theta/(2\pi)\)
and unit-radius arc length \(ds=d\theta\) give
\[
\|h\|_{L^2(d\mu)}=\|h\|,
\qquad
\|h\|_{L^2(ds)}=\sqrt{2\pi}\|h\|,
\qquad
\int_{S^1}\|h\|\,ds=2\pi\|h\|.
\]
The factors \(1,\sqrt{2\pi},2\pi\) are inequivalent functionals.  The Hopf
period alone does not select one as a scalar VEV.
\end{theorem}

\begin{theorem}[Canonical field-rescaling quotient]
\label{p6:thm:canonical-scalar-field-rescaling}
Let a real scalar coordinate have positive kinetic coefficient \(Z>0\),
\[
S_{\rm kin}[h]=\frac Z2\int(\partial h)^2\,d{\rm vol},
\qquad
v=\kappa E_{32}.
\]
The change of coordinate \(h_c=\sqrt Z\,h\) gives canonical kinetic
coefficient one and canonical readout
\[
v_c=\kappa_cE_{32},
\qquad
\kappa_c=\sqrt Z\,\kappa.
\]
Hence the pair \((Z,\kappa)\) factors through the single invariant coordinate
\(\kappa_c\).  Choosing \(Z=1\) is a coordinate representative, not an
additional physical selector.
\end{theorem}

\begin{proof}
Substitution gives \(\partial h=Z^{-1/2}\partial h_c\), so
\(S_{\rm kin}[h]=\frac12\int(\partial h_c)^2d{\rm vol}\).  The same field
rescaling sends its constant coordinate value to
\(v_c=\sqrt Z\,v=\sqrt Z\,\kappa E_{32}\).  Two pairs with equal
\(\sqrt Z\kappa\) therefore describe the same canonically normalized
coordinate.
\end{proof}

\begin{corollary}[Native hierarchy owner of the electroweak comparison scale]
\label{p6:cor:v571-native-ew-scale-owner}
Paper~III-C's protected octet reflection has the unique fixed shell label
\(K=32\).  On the typed energy owner, the derived shell law is
\[
 E(K)=E_{32}\alpha_{\rm hier}^{-(K-32)/32}.
\]
Hence \(E_{32}\) is the native dimensionful owner of the electroweak
comparison scale.  A conventional scalar coordinate may be written
\[
 v_{\rm cmp}=\kappa_{\rm cmp}E_{32},
\]
but this introduces no second dimensionful owner and does not by itself
assert a scalar stationary vacuum.  The dimensionless \(\kappa_{\rm cmp}\)
is a coordinate/readout unless a later claim derives it from physical output.
\end{corollary}

\begin{definition}[Scalar normalization/readout chart and source coordinate]
\label{p6:def:selected-scalar-normalization}
After quotienting by Theorem~\ref{p6:thm:canonical-scalar-field-rescaling},
the scalar normalization/readout chart is
\[
\boxed{
\mathfrak A_H
=
\{\mathsf H_{\rm norm}(\kappa_c):\kappa_c>0\}.
}
\]
Hereafter write \(\kappa\) for the canonical coordinate \(\kappa_c\).  Then
\(v_\kappa=\kappa E_{32}\), and the corresponding couplings are given in
Theorem~\ref{p6:thm:scalar-counterfamily}.  To preserve the former printed
coordinate, the source member declares the mixed-to-conditional RMS
conversion and the unnormalized
\(L^1\) orbit functional:
\[
\kappa_{\rm print}
=2\pi\sqrt{\frac{24}{15}}
=\frac{8\pi}{\sqrt{10}},
\qquad
v_{\rm read}:=\kappa_{\rm print}E_{32}.
\]
This is a normalized readout coordinate.  It is not a theorem that the
scalar potential has a nonzero vacuum at \(v_{\rm read}\).
\end{definition}

\begin{definition}[Explicit scalar normalization package]
\label{p6:def:explicit-hnorm}
Take the phase-zero representative of the invariant isometry class derived in
Theorem~\ref{p6:thm:canonical-invariant-scalar-state},
\[
J_H:\CC\longrightarrow\Ran P_{15},
\qquad
J_H(z)=z\Omega_{24}.
\]
The carrier state and finite line metric are consequently
\[
\omega_H(A)=\langle\Omega_{24},A\Omega_{24}\rangle,
\qquad
K_H(z,w)=\overline zw,
\qquad
Z_H=1.
\]
Here \(Z_H=1\) is the canonical representative of the field-rescaling quotient
in Theorem~\ref{p6:thm:canonical-scalar-field-rescaling}; it is not a second
normalization parameter.
If the real coordinate \(h\) is promoted to a field on a downstream
spacetime chart, the normalization convention is
\[
S_{H,{\rm kin}}
=
\frac12\int(\partial_\mu h)(\partial^\mu h)\,d{\rm vol}.
\]
The derived embedding/state and canonical coordinate, together with the one
declared functional
\(
\kappa_{\rm print}=2\pi\sqrt{24/15}
\),
constitute the retained comparison/readout coordinate
\(\mathsf H_{\rm norm}(\kappa_{\rm print})\).  It is not a separate physical owner family; its finite coordinate is \(\kappa_{\rm print}\), and it contains no potential, sign of a quadratic term, vacuum, or EWSB assertion.
\end{definition}

\begin{theorem}[One-parameter scalar counterfamily]
\label{p6:thm:scalar-counterfamily}
Retain the conditional terminal row-energy coordinates
\[
(\mathcal E_W^{\rm term})^2=7\etaH E_{32}^2,\qquad
(\mathcal E_Z^{\rm term})^2=9\etaH E_{32}^2,\qquad
(\mathcal E_H^{\rm term})^2=17\etaH E_{32}^2.
\]
For every \(\kappa>0\), define
\[
v_\kappa=\kappa E_{32},\qquad
g_2^2(\kappa)=\frac{28\etaH}{\kappa^2},\qquad
g_1^2(\kappa)=\frac{8\etaH}{\kappa^2},
\]
\[
\lambda(\kappa)=\frac{17\etaH}{2\kappa^2}.
\]
If a common downstream broken-branch owner supplies a physical vacuum,
identifies the three rows with physical poles through
\(\mathfrak m_p^{\rm rest}c_{\rm op}^2=\mathcal E_p^{\rm term}\), and uses
the usual canonically normalized dictionary, these coordinates reproduce the
same three mass-energy values and give
\[
\sin^2\theta_W=\frac29
\]
for every \(\kappa\).  Without that owner/ancestry bridge, the theorem is a
counterfamily of row-energy and coupling coordinates, not a physical mass or
weak-angle prediction.
\end{theorem}

\begin{proof}
Substitute into
\[
(\mathcal E_W^{\rm term})^2=\frac14g_2^2v_\kappa^2,\qquad
(\mathcal E_Z^{\rm term})^2=\frac14(g_1^2+g_2^2)v_\kappa^2,\qquad
(\mathcal E_H^{\rm term})^2=2\lambda v_\kappa^2.
\]
The coupling ratio is \(8/(8+28)=2/9\).  The displayed physical
interpretation uses exactly the extra bridge stated in the theorem.
\end{proof}

\begin{corollary}[Conditional printed scalar coordinate]
\label{p6:cor:printed-scalar-coordinate}
At \(\kappa=\kappa_{\rm print}\),
\[
g_2^2=\frac{21}{5\pi^2},
\qquad
g_1^2=\frac{6}{5\pi^2},
\qquad
\lambda_H=\frac{51}{40\pi^2}.
\]
These are exact consequences of the declared normalization and the
conditional broken-branch dictionary.  Paper~V-B does not use them to prove a
potential, vacuum, physical pole ancestry, or equal-amplitude response; those
are downstream obligations.
\end{corollary}

\subsection{External Standard--Model flavor-coordinate dictionaries}
\label{p6:sec:flavor-operators}

\begin{theorem}[Eigenframe counterfamily]
\label{p6:thm:eigenframe-counterfamily}
For positive diagonal matrices with simple spectra and arbitrary
\(V,U\in\UU(3)\), the operators
\[
H_u=D_u^2,\qquad
H_d=VD_d^2V^\dagger,\qquad
H_e=D_e^2,\qquad
M_\nu=U^*D_\nu U^\dagger
\]
have the same four spectra and yield relative frames \(V\) and \(U\), up to
the usual phase and ordering conventions.  They act on the same structural
family carrier.
\end{theorem}

\begin{proof}
Unitary conjugation preserves Hermitian spectra.  The displayed
\(M_\nu\) is symmetric and satisfies \(U^TM_\nu U=D_\nu\).  Family
multiplicity commutes with the gauge action, so no structural gauge condition
restricts \(U\) or \(V\).
\end{proof}

\begin{definition}[External SM family-orbit and relative-eigenframe coordinate]
\label{p6:def:relative-eigenframe-data}
For quarks let
\[
 \mathcal O_q:=[(H_u,H_d)]\in
 (\operatorname{Herm}_3^+\times\operatorname{Herm}_3^+)/U(3)_Q,
 \qquad H_u=Y_uY_u^\dagger,\quad H_d=Y_dY_d^\dagger,
\]
and analogously let
\(\mathcal O_\ell=[(H_e,H_\nu)]\) on the Dirac-neutrino branch.
Within the conventional Standard--Model matrix-mass comparison dictionary,
this simultaneous-conjugacy orbit is a useful coordinate object and the
relative eigenframe classes \([V_{\rm CKM}]\) and \([U_{\rm PMNS}]\) are
standard representatives.  In the framework-native ontology they are not
selector-level data: the charged and neutral masses are terminal readouts and
the physical flavor objects are the generated current maps \(C_{ud}\) and
\(C_{e\nu}\).  Choosing a standard CKM or PMNS section is therefore comparison
presentation only.
\end{definition}

\begin{theorem}[External SM sector-orbit reconstruction]
\label{p6:thm:v577-family-orbit-contraction}
Assume simple positive spectra.  Choose spectral representatives
\[
 H_u=U_uD_u^2U_u^\dagger,\qquad
 H_d=U_dD_d^2U_d^\dagger.
\]
Then the simultaneous-conjugacy orbit \(\mathcal O_q\) determines the ordered
spectra and the CKM rephasing class
\[
 \boxed{[V_{\rm CKM}]=[U_u^\dagger U_d]}.
\]
Every polynomial weak-CP invariant, including
\(\operatorname{Tr}[H_u,H_d]^3\), is constant on that same orbit.  This is a reconstruction theorem for the external matrix dictionary.  It is
not the framework-native flavor closure and supplies no upstream selector.
The active native quark and lepton flavor objects are instead \(C_{ud}\) and
\(C_{e\nu}\), generated by the weak-current incidence of already-massed role
spaces.  The theorem also does not decide a Majorana interaction occurrence.
\end{theorem}

\begin{proof}
Simultaneous weak-basis conjugation sends \(U_u,U_d)\mapsto(WU_u,WU_d)\), so
\(U_u^\dagger U_d\) is unchanged up to the phase/permutation freedoms of the
simple spectral decompositions.  The ordered eigenvalues are orbit invariants,
and traces of words in \(H_u,H_d\), in particular the commutator-cube CP
invariant, are invariant under simultaneous conjugation.  The lepton argument
is the same.
\end{proof}

\begin{corollary}[Reconstruction is realization]
\label{p6:cor:flavor-realization-status}
Formulas such as
\[
Y_d=\frac{\sqrt2}{v}V_{\rm CKM}D_d,
\qquad
M_\nu=U_{\rm PMNS}^*D_\nu U_{\rm PMNS}^\dagger
\]
construct representatives after the spectra and mixing matrices are
supplied.  They prove existence, not upstream selection.
\end{corollary}

\subsubsection{Legacy exact CKM comparison convention}

We use the standard Cabibbo--Kobayashi--Maskawa section, Wolfenstein
coordinates, and Jarlskog invariant
\pvcite{Cabibbo1963,KobayashiMaskawa1973,Wolfenstein1983,Jarlskog1985}.

\begin{definition}[Exact unbarred coordinates]
\label{p6:def:exact-ckm-coordinates}
In the standard unitary section define
\[
s_{12}=\lambda,\qquad
s_{23}=A\lambda^2,\qquad
s_{13}e^{i\delta_{\rm std}}
=A\lambda^3(\rho+i\eta_q).
\]
\end{definition}

\begin{theorem}[Exact CKM invariants]
\label{p6:thm:exact-ckm-identities}
\[
\boxed{
J_q
=
A^2\lambda^6\eta_q\,c_{12}c_{23}c_{13}^2.
}
\]
Moreover,
\[
\bar\rho+i\bar\eta
=
-\frac{V_{ud}V_{ub}^*}{V_{cd}V_{cb}^*}
=
\frac{c_{12}(\rho+i\eta_q)}
{c_{23}+c_{12}A^2\lambda^4(\rho+i\eta_q)}.
\]
Thus the standard phase \(\delta_{\rm std}\) is not identically the
unitarity-triangle angle
\(\gamma_{\rm UT}=\arg(\bar\rho+i\bar\eta)\).
\end{theorem}

\begin{proof}
Substitute Definition~\ref{p6:def:exact-ckm-coordinates} into
\[
J=s_{12}c_{12}s_{23}c_{23}s_{13}c_{13}^2
\sin\delta_{\rm std}
\]
and into the rephasing-invariant apex ratio.
\end{proof}

\begin{theorem}[Former CKM representatives are inequivalent]
\label{p6:thm:former-ckm-inequivalence}
With
\[
\lambda=0.224790118710538,\quad
A=0.826001673049125,\quad
J_q=3.12045659516410\times10^{-5},
\]
the former leading identity
\(
J_q=A^2\lambda^6\eta_q
\)
gave \(|V_{ub}|=0.0036951451\ldots\), whereas the former exact standard
matrix gave \(|V_{ub}|=0.0037955602\ldots\).  The matrices are not related by
rephasing.
\end{theorem}

\begin{proof}
Rephasing preserves every entry modulus.  The exact identity contains the
factor
\(
c_{12}c_{23}c_{13}^2=0.9735447899\ldots
\),
which the leading formula omitted.
\end{proof}

\begin{definition}[Legacy exact CKM comparison branch]
\label{p6:def:selected-ckm-branch}
For historical comparison with a conventional CKM coordinate section, retain the four benchmark inputs
\[
\lambda=0.224790118710538,\qquad
A=0.826001673049125,
\]
\[
J_q=3.12045659516410\times10^{-5},
\qquad
\gamma_{\rm UT}=64.1666666667^\circ .
\]
It interprets the last number as the exact unitarity-triangle angle.  The
simultaneous small solution of the exact Jarlskog and apex equations is
\[
\begin{aligned}
s_{13}&=0.003794424824531,&
\delta_{\rm std}&=64.2021006467^\circ,\\
\rho&=0.176003116174467,&
\eta_q&=0.364114000362982.
\end{aligned}
\]
The corresponding barred apex is
\[
\bar\rho=0.171816450635092,
\qquad
\bar\eta=0.354892340569536.
\]
\end{definition}

\begin{proposition}[Standard comparison representative of the legacy CKM benchmark]
\label{p6:prop:selected-ckm-matrix}
Substitution in the standard unitary section gives
\begingroup
\scriptsize
\[
V_{\rm CKM}^{\rm sel}\simeq
\begin{pmatrix}
0.974400191 & 0.224788500 &
0.001651326-0.003416252i\\
-0.224661391-0.000138940i&
0.973542590-0.000032053i&
0.041738058\\
0.007774708-0.003325920i&
-0.041041035-0.000767271i&
0.999121382
\end{pmatrix}.
\]
\endgroup
It is unitary to the displayed precision and has
\[
|V_{ub}|=0.003794425,\qquad
|V_{td}|=0.008456230,\qquad
J_q=3.12045659516\times10^{-5}.
\]
\end{proposition}

\subsubsection{Legacy SM-coordinate PMNS packet and compatible comparison basis}

\begin{definition}[Conditional SM-coordinate PMNS comparison packet]
\label{p6:def:selected-pmns-packet}
The retained flavor bridge supplies the invariant packet
\[
s_{12}^2=\frac{11}{36},
\qquad
s_{23}^2=\frac47,
\qquad
s_{13}^2=0.0224119926119664,
\]
\[
\delta_\ell
=
\pi+\frac12\arcsin\frac{\sqrt{11}}6
=196.77865^\circ,
\qquad
J_\ell=-0.00963078985599834.
\]
The standard unitary representative is
\[
U_{\rm PMNS}^{\rm sel}\simeq
\begin{pmatrix}
0.823942&0.546541&-0.143333+0.043217i\\
-0.271582+0.027224i&0.605437+0.018058i&0.747410\\
0.496050+0.023577i&-0.578072+0.015639i&0.647276
\end{pmatrix}.
\]
These values are legacy Standard--Model comparison coordinates.  The native current theorem is already structural, but its generated map $C_{e\nu}$ has not been numerically evaluated into this particular conventional section.
\end{definition}

\begin{definition}[Selected diagonal sector spectra]
\label{p6:def:selected-diagonal-spectra}
The flavor package binds its diagonal spectra to the preceding selected
readouts:
\[
\begin{aligned}
D_u^{\rm nat}
&=
\diag\left(
E_{40}\etaH^{1/2}\sqrt2,\
3E_{36}\etaH\sqrt4,\
E_{32}\etaH^{1/2}(\sqrt7+\sqrt9)
\right),\\
D_d^{\rm nat}
&=
\diag\left(
E_{40}\etaH(\sqrt2+\sqrt3),\
E_{36}\etaH^{-1}/\sqrt6,\
\frac{E_{32}\etaH^{-3/2}}{\sqrt8+\sqrt{16}+1}
\right),\\
D_e^{[4]}
&=
\diag\left(
m(\widetilde n_{L,1}),\
m(\widetilde n_{L,2}),\
m(\widetilde n_{L,3})
\right),\\
D_\nu^{(56)}
&=
\sqrt{\etaH}E_{56}
\diag(1/\sqrt8,\sqrt5,13).
\end{aligned}
\]
\(D_u^{\rm nat},D_d^{\rm nat}\) are native finite coordinates pending
\(\mathsf{QRG}\); \(D_e^{[4]}\) is the retained-order charged branch of
Definition~\ref{p6:def:charged-absolute-map}; and \(D_\nu^{(56)}\) is the
positive radical spectrum.  No alternative diagonal ledger enters the
selected flavor representative.
\end{definition}

\begin{proposition}[External Takagi/PMNS coordinate compatibility]
\label{p6:prop:positive-majorana-frame}
Let \(D_\nu=D_\nu^{(56)}\) and \(D_e=D_e^{[4]}\).  In the common family basis
define
\[
M_{\nu,+}=D_\nu,\qquad
H_e=(U_{\rm PMNS}^{\rm sel})^\dagger
D_e^2U_{\rm PMNS}^{\rm sel}.
\]
Then the neutral Takagi frame is \(U_\nu=I\), the charged-lepton frame is
\(U_e=(U_{\rm PMNS}^{\rm sel})^\dagger\), and
\[
U_e^\dagger U_\nu=U_{\rm PMNS}^{\rm sel}.
\]
Thus positivity of \(M_{\nu,+}\), Majorana symmetry, and the complex PMNS
phase are mutually compatible.
\end{proposition}

\begin{proof}
\(M_{\nu,+}\) is real diagonal and positive.  The displayed \(H_e\) is
Hermitian and is diagonalized by \(U_e\).  The relative-frame identity is
immediate.
\end{proof}

\begin{definition}[External matrix-coordinate flavor representative]
\label{p6:def:operator-first-flavor}
Use the four fixed diagonal spectra of
Definition~\ref{p6:def:selected-diagonal-spectra} and set
\[
\begin{aligned}
H_u&=(D_u^{\rm nat})^2,&
H_d&=V_{\rm CKM}^{\rm sel}(D_d^{\rm nat})^2
(V_{\rm CKM}^{\rm sel})^\dagger,\\
H_e&=(U_{\rm PMNS}^{\rm sel})^\dagger
(D_e^{[4]})^2U_{\rm PMNS}^{\rm sel},&
M_{\nu,+}&=D_\nu^{(56)}.
\end{aligned}
\]
These operators precede diagonalization within the conditional model.
They instantiate the selected flavor data but do not prove that the carrier
uniquely selects them.
\end{definition}

\begin{theorem}[Parity does not select the first CP order]
\label{p6:thm:cp-order-no-go}
If
\[
\theta_{12}=a\delta_{\rm geom}^p,\quad
\theta_{23}=b\delta_{\rm geom}^q,\quad
\theta_{13}=c\delta_{\rm geom}^r
\]
and the phase tends to
\(\phi\notin\pi\ZZ\), then
\[
J=abc\sin\phi\,
\delta_{\rm geom}^{p+q+r}
+o(\delta_{\rm geom}^{p+q+r}).
\]
The invariant vanishes on every two-family face for every positive
\((p,q,r)\).  Hence parity and two-family collapse do not select the former
quark order five or lepton order two.
\end{theorem}

\begin{proof}
Insert the angle expansions into the exact Jarlskog formula.  The cosines
tend to one, and every sine factor vanishes on its two-family face.
\end{proof}

\begin{corollary}[Operator-first CP diagnostic]
\label{p6:cor:operator-first-cp}
For nondegenerate quark spectra,
\[
|\det[H_u,H_d]|
=
2|J_q|
\left|\prod_{i<j}(m_{u_i}^2-m_{u_j}^2)\right|
\left|\prod_{k<\ell}(m_{d_k}^2-m_{d_\ell}^2)\right|.
\]
The first nonzero CP order must be computed from explicit sector-operator
series.  This paper makes no independent CP-order selection claim.
\end{corollary}

\begin{corollary}[Conditional external-SM beta readouts]
\label{p6:cor:conditional-beta-readouts}
If one additionally adopts the external Takagi/Majorana comparison dictionary, the selected radical masses, conventional current-coordinate section, and positive parity branch give
\[
m_\beta^2=\sum_i|U_{ei}|^2m_i^2,
\qquad
m_{\beta\beta}=\left|\sum_iU_{ei}^2m_i\right|,
\]
giving approximately
\[
m_\beta=8.907932\ {\rm meV},
\qquad
m_{\beta\beta}=4.468305\ {\rm meV}.
\]
If positivity of \(M_{\nu,+}\) is not retained, the four real Majorana
sign/CP-parity assignments instead give
\[
4.468305,\quad2.631747,\quad0.943554,\quad2.652503
\quad{\rm meV}.
\]
\end{corollary}

\label{p6:part:interfaces}

\subsection{Rooted Hurwitz--Clifford mass candidate: exact invariants and offset-only obstruction}
\label{p6:app:v562-hurwitz-clifford-mass-candidate}
This section records the exact positive algebra and the negative offset-only
test referenced in
Remark~\ref{p6:rem:v562-hurwitz-clifford-candidate-scope}.  It is an audit of
a generated candidate family, not a physical mass assignment.

Let
\[
 K(a,b)=aI+bJ,
 \qquad
 R_G=\operatorname{diag}(1,-i,i),
\]
and define the Hermitian rooted return
\begin{equation}
\label{p6:eq:v562-phase-refined-return}
 F(a,b):=K(a,b)R_G+R_G^\dagger K(a,b)
 =\begin{pmatrix}
 2(a+b)&b(1-i)&b(1+i)\\
 b(1+i)&0&2ib\\
 b(1-i)&-2ib&0
 \end{pmatrix}.
\end{equation}
A direct determinant gives
\begin{equation}
\label{p6:eq:v562-f-charpoly}
 \chi_F(\lambda)
 =(\lambda+2b)
 \left[\lambda^2-2(a+2b)\lambda+4ab\right],
\end{equation}
and hence
\[
 \operatorname{spec}F(a,b)
 =\{-2b,\ a+2b\pm\sqrt{a^2+4b^2}\}.
\]
For \(a,b>0\) the inertia is therefore \((2+,1-)\).

Now write
\[
 x=E_{32},\qquad y=E_{36},\qquad z=E_{40},
 \qquad y^2=xz,
\]
and use the rooted operator-flag pullback
\[
 \mathcal E_G=\operatorname{diag}(x,y,z).
\]
The midpoint \(y=E_{36}\) is used here as the already derived
operator-transition weight, not as a third Clifford carrier eigenspace.  With
\[
 X=\begin{pmatrix}0&0&1\\1&0&0\\0&1&0\end{pmatrix},\qquad
 F_k=X^kFX^{-k},\qquad
 M_k=\mathcal E_G^{1/2}F_k\mathcal E_G^{1/2},
\]
let \(r_k\) denote respectively \(x,y,z\) for \(k=0,1,2\), and let
\(p_k,q_k\) be the other two entries.  The three elementary symmetric
invariants are
\begin{align}
 e_1(M_k)&=2(a+b)r_k,\label{p6:eq:v562-e1}\\
 e_2(M_k)&=-2b^2\left[r_k(p_k+q_k)+2p_kq_k\right],\label{p6:eq:v562-e2}\\
 e_3(M_k)&=-8ab^2xyz=-8ab^2y^3.\label{p6:eq:v562-e3}
\end{align}
Equivalently,
\begin{align*}
 e_2(M_0)&=-2b^2(xy+xz+2yz),\\
 e_2(M_1)&=-2b^2(xy+2xz+yz),\\
 e_2(M_2)&=-2b^2(2xy+xz+yz).
\end{align*}
Thus, with
\(U_k=r_k(p_k+q_k)+2p_kq_k\),
\begin{equation}
\label{p6:eq:v562-mk-charpoly}
 \boxed{\chi_{M_k}(\lambda)
 =\lambda^3-2(a+b)r_k\lambda^2-2b^2U_k\lambda+8ab^2y^3.}
\end{equation}
For the two coefficient-free descendants
\[
 K_\ell^{(0)}=I+J=K(1,1),\qquad
 K_C^{(0)}=2I+J=K(2,1),
\]
the second invariant is identical while the first and third differ.  If
\(x,y,z\) are distinct, the pair \((a,k)\in\{1,2\}\times\mathbb Z_3\) is
recovered from the characteristic invariants by
\[
 a=-\frac{e_3}{8y^3},\qquad
 r_k=\frac{e_1}{2(a+1)}.
\]
Hence the six candidate labels are spectrally distinguishable.  This is a
capacity statement only; Paper~V-A does not yet prove that either displayed
return occurs as a physical sector member.

\begin{proposition}[Discrete relative alignment is CP-capable on the anchored hierarchy]
\label{p6:prop:v562-discrete-cp-capacity}
Put
\[
 y=x\rho,\qquad z=x\rho^2,
\qquad \rho=\alpha_0^{-1/8},
\qquad
 \alpha_0=\exp\!\left(4\pi^2+\frac1{6\sqrt2}\right),
\]
which is the anchored \(\mathrm{REG}_0\) hierarchy.  For
\(H_k=M_kM_k^\dagger=M_k^2\), the relative placements \(k=1,2\) obey
\begin{align}
 \operatorname{tr}[H_1,H_2]^3\big|_{a=b=1}
 ={}&3072i\,x^{12}\rho^9(1+\rho+\rho^2)^2
 \left(14\rho^4+15\rho^3-3\rho^2-1\right),\label{p6:eq:v562-cp-lepton}\\
 \operatorname{tr}[H_1,H_2]^3\big|_{a=2,b=1}
 ={}&768i\,x^{12}\rho^9
 (3\rho^2+2\rho+3)(3\rho^2+3\rho+2)\nonumber\\
 &\qquad\times
 \left(69\rho^4+80\rho^3-11\rho^2+10\rho-1\right).
 \label{p6:eq:v562-cp-color}
\end{align}
Both expressions are nonzero.  Exchanging the two relative orientations
reverses the sign.  Thus no continuously fitted Yukawa phase is
mathematically required for nonzero finite-family CP capacity.  The result
does not select the physical CKM/PMNS branch and does not address a regulated
quantum determinant or strong-CP angle.
\end{proposition}
\begin{proof}
The identities follow by direct multiplication of the three-by-three
matrices in \eqref{p6:eq:v562-phase-refined-return}.  Since
\(0<\rho<e^{-\pi^2/2}<10^{-2}\),
\[
 14\rho^4+15\rho^3-3\rho^2-1<14\cdot10^{-8}+15\cdot10^{-6}-1<0
\]
and
\[
 69\rho^4+80\rho^3-11\rho^2+10\rho-1
 <69\cdot10^{-8}+80\cdot10^{-6}+10^{-1}-1<0.
\]
All remaining real factors are positive, proving nonvanishing.  The
commutator changes sign under interchange of \(H_1,H_2\).
\end{proof}

\begin{proposition}[Offset-only hierarchy-transmission obstruction]
\label{p6:prop:v562-offset-only-hierarchy-obstruction}
Assume \(a,b>0\), set
\[
 y=x\rho,\qquad z=x\rho^2,
\qquad \rho\downarrow0,
\]
and interpret the rooted hierarchy as one occurrence each of the orders
\(x,y,z\).  For the unshifted return,
\begin{align}
 \lambda_1(M_0)&=2(a+b)x+O(y),\nonumber\\
 \lambda_2(M_0)&=-\frac{b^2}{a+b}y+o(y),\label{p6:eq:v562-k0-asymptotics}\\
 \lambda_3(M_0)&=4az+o(z),\nonumber
\end{align}
so its singular values occupy the three distinct leading orders
\(x,y,z\).  In contrast,
\begin{align}
 \lambda_\pm(M_1)&=\pm\sqrt2\,b\sqrt{xy}+O(y),\label{p6:eq:v562-k1-asymptotics}\\
 \lambda_\pm(M_2)&=\pm2b\sqrt{xy}+O(y),\label{p6:eq:v562-k2-asymptotics}
\end{align}
and therefore, for either nonzero offset,
\begin{equation}
\label{p6:eq:v562-leading-singular-doublet}
 \frac{\sigma_1(M_k)}{\sigma_2(M_k)}\longrightarrow1,
 \qquad k=1,2.
\end{equation}
Consequently a compiler in which \(R^\pm\) acts only by choosing a cyclic
offset of one fixed \(F(a,b)\) cannot represent two distinct right-sector
orientations while merely transporting one copy of each of the three rooted
hierarchy orders.  At least one of the finite source word, its action-owned
functional calculus, or the relative root alignment must carry additional
\(R^\pm\)-dependence.
\end{proposition}
\begin{proof}
Divide by the common factor \(x\).  For \(k=0\), the leading matrix has the
simple eigenvalue \(2(a+b)\); the two smaller roots are then fixed to the
stated leading orders by \eqref{p6:eq:v562-e2} and
\eqref{p6:eq:v562-e3}.  For \(k=1\), the leading nonzero two-by-two block at
order \(\rho^{1/2}\) is
\[
 b\sqrt\rho\begin{pmatrix}0&1-i\\1+i&0\end{pmatrix},
\]
with eigenvalues \(\pm\sqrt2\,b\sqrt\rho\).  For \(k=2\) it is
\[
 b\sqrt\rho\begin{pmatrix}0&2i\\-2i&0\end{pmatrix},
\]
with eigenvalues \(\pm2b\sqrt\rho\).  Restoring \(x\) gives
\eqref{p6:eq:v562-k1-asymptotics}--\eqref{p6:eq:v562-k2-asymptotics} and
hence \eqref{p6:eq:v562-leading-singular-doublet}.  This is an internal
hierarchy-transmission obstruction; no observed fermion mass is used.
\end{proof}


\subsubsection{Orientation-odd family-transport repair of the rooted candidate}
\label{p6:app:v563-oriented-signline-repair}
This subsection extends the preceding candidate audit without promoting any
member to a physical Yukawa source.  It uses the family copy of the rooted
cyclic transport
\[
 X_G=\begin{pmatrix}0&0&1\\1&0&0\\0&1&0\end{pmatrix},
 \qquad
 C_{2,G}:=-i(X_G-X_G^\dagger),
 \qquad
 R_G=\operatorname{diag}(1,-i,i).
\]
The source normalization is the coefficient-one transport normalization,
not Hilbert--Schmidt unit normalization:
\begin{equation}
\label{p6:eq:v563-c2-norm-spectrum}
 \operatorname{tr}(C_{2,G}^2)=6,
 \qquad
 \operatorname{spec}C_{2,G}=\{0,\sqrt3,-\sqrt3\}.
\end{equation}

\begin{proposition}[Unique sign line in the primitive family trimer]
\label{p6:prop:v563-unique-trimer-sign-line}
Under the rooted projective-diagonal stabilizer \(S_3\), the primitive
Hermitian trimer chart decomposes as
\[
 \mathfrak T_{\rm tri}
 \cong
 \mathbf1_I\oplus V_{\rm std}^{\rm diag}
 \oplus\mathbf1_{C_1}\oplus\mathbf1_{\rm sgn,C_2}.
\]
Hence \(\mathbb R C_{2,G}\) is the unique sign subrepresentation in this
five-real-dimensional primitive chart.
\end{proposition}
\begin{proof}
The cyclic generator fixes both
\(C_{1,G}=X_G+X_G^\dagger\) and \(C_{2,G}\).  An orientation-reversing
reflection sends \(X_G\mapsto X_G^\dagger\), and therefore fixes
\(C_{1,G}\) while sending \(C_{2,G}\mapsto-C_{2,G}\).  The traceless
diagonal defect plane is the two-dimensional standard representation and
contains no sign subline.  Together with the scalar identity direction this
gives the stated decomposition and uniqueness.
\end{proof}

For \(a>0\), put
\[
 K(a)=aI+J,
\]
\[
 F_+(a)=K(a)R_G+R_G^\dagger K(a),
 \qquad
 F_-(a)=K(a)R_G^\dagger+R_GK(a).
\]
For a component branch \(q\in\{\ell,C\}\), set
\[
 a_\ell=1,
 \qquad
 a_C=2,
\]
and consider the diagnostic repair
\begin{equation}
\label{p6:eq:v563-signline-repair}
 \boxed{
 W_{\pm,q}=F_\pm(a_q)+\lambda_q C_{2,G},
 \qquad \lambda_q\in\mathbb R.}
\end{equation}
The same \(\lambda_q\) is held fixed under the Hopf flip within one component
branch.  At the level of this unrestricted diagnostic deformation family, no
physical source theorem identifies \(\lambda_\ell\) with \(\lambda_C\).  The
polar-normalized ordered-return subbranch derived below fixes the normalized
sign transport to coefficient one in both sectors, conditional on that branch
being consumed by the physical source.

\begin{proposition}[Exact rank and conjugation-splitting conditions]
\label{p6:prop:v563-signline-rank-split}
For real \(\lambda\),
\begin{align}
 \det W_+(a;\lambda)
 &=-2(\lambda+2)\bigl(a\lambda+2a+3\lambda\bigr),
 \label{p6:eq:v563-det-plus}\\
 \det W_-(a;\lambda)
 &=-2(\lambda-2)\bigl(a\lambda-2a+3\lambda\bigr).
 \label{p6:eq:v563-det-minus}
\end{align}
Thus both orientations are full rank exactly when
\begin{equation}
\label{p6:eq:v563-full-rank-set}
 \lambda\notin
 \left\{-2,2,-\frac{2a}{a+3},\frac{2a}{a+3}\right\}.
\end{equation}
With
\[
 D_\rho=\operatorname{diag}(1,\sqrt\rho,\rho),
 \qquad M_\pm=D_\rho W_\pm D_\rho,
 \qquad H_\pm=M_\pm M_\pm^\dagger,
\]
one has the exact trace difference
\begin{equation}
\label{p6:eq:v563-trace-split}
 \boxed{
 \operatorname{tr}H_+-\operatorname{tr}H_-
 =8\lambda\rho(\rho-1)(2\rho+1).}
\end{equation}
Hence for \(0<\rho<1\), every nonzero \(\lambda\) breaks the
pure-conjugation singular-spectrum degeneracy.
\end{proposition}
\begin{proof}
Direct determinant expansion gives
\eqref{p6:eq:v563-det-plus}--\eqref{p6:eq:v563-det-minus}; their zero sets
give \eqref{p6:eq:v563-full-rank-set}.  Direct multiplication of the three by
three Hermitian matrices gives \eqref{p6:eq:v563-trace-split}.
\end{proof}

\begin{proposition}[Full-rank sign-line repair preserves all three hierarchy orders]
\label{p6:prop:v563-signline-hierarchy-orders}
Assume \(a>0\), \(0<\rho\ll1\), and the corresponding determinant in
Proposition~\ref{p6:prop:v563-signline-rank-split} is nonzero.  Then the three
eigenvalue orders of \(M_\pm\) are
\[
 \boxed{1,\rho,\rho^2}.
\]
More precisely,
\begin{align*}
 \lambda_1(M_\pm)&=2(a+1)+O(\rho),\\
 \lambda_2(M_+)&=-\frac{(\lambda-1)^2+1}{2(a+1)}\rho+O(\rho^2),\\
 \lambda_2(M_-)&=-\frac{(\lambda+1)^2+1}{2(a+1)}\rho+O(\rho^2),\\
 \lambda_3(M_+)&=
 \frac{2(\lambda+2)(a\lambda+2a+3\lambda)}{(\lambda-1)^2+1}
 \rho^2+O(\rho^3),\\
 \lambda_3(M_-)&=
 \frac{2(\lambda-2)(a\lambda-2a+3\lambda)}{(\lambda+1)^2+1}
 \rho^2+O(\rho^3).
\end{align*}
Thus the hierarchy-transmission property is not special to a coefficient-one
choice.
\end{proposition}
\begin{proof}
Insert \(\mu=O(1)\), \(\mu=\rho\nu\), and
\(\mu=\rho^2\nu\) successively into the exact characteristic polynomials of
\(M_\pm\).  The order-\(\rho\) coefficients are
\((\lambda\mp1)^2+1>0\); the displayed nonzero determinants fix the final
order-\(\rho^2\) roots.  The stated coefficients follow by matching the
leading terms.
\end{proof}

\begin{proposition}[Coefficient-one noncommuting-frame witness]
\label{p6:prop:v563-signline-cp-witness}
For the diagnostic value \(\lambda=1\), exact multiplication gives
\begin{align}
 \operatorname{tr}[H_+,H_-]^3
 ={}&-192i\,\rho^6(a+1)(\rho-1)(\rho+1)\nonumber\\
 &\times(8a+13\rho^2+13\rho+4)\nonumber\\
 &\times(5a\rho+5a+4\rho^2+5\rho+5).
 \label{p6:eq:v563-cp-witness}
\end{align}
For \(a>0\) and \(0<\rho<1\) this is nonzero.  This is a
noncommuting-frame capacity witness only.  Within the unrestricted diagnostic
\(\lambda\)-family, coefficient one is not selected by numerical simplicity;
it is canonically induced only on the polar-normalized ordered-return branch
proved below, whose physical source occurrence remains open.  The statement
is not a CKM/PMNS member theorem.
\end{proposition}
\begin{proof}
The factorization follows from direct multiplication.  Under the stated
conditions every real factor after the prefactor is positive except
\(\rho-1<0\), so the expression is nonzero.
\end{proof}

\begin{corollary}[Raw-curvature amplitude is not the color mass-word coefficient]
\label{p6:cor:v563-raw-curvature-color-nogo}
The generated orientation-odd return curvature has sector amplitudes
\(C_{2,G}\) and \(2C_{2,G}\) on the two incidence bases.  Inserting the raw
color amplitude literally into \eqref{p6:eq:v563-signline-repair} gives
\(a_C=2\), \(\lambda_C=2\), and therefore
\[
 \boxed{\det W_-(2;2)=0.}
\]
Hence the generated sign-line \emph{direction} is a viable missing oriented
datum, but its raw color curvature magnitude cannot itself be the full-rank
physical color source coefficient.
\end{corollary}

\begin{theorem}[Polar separation of incidence modulus and oriented family transport]
\label{p6:thm:v564-polar-separated-oriented-return}
Let
\[
 K(a,b)=aI+bJ,\qquad a>0,\ b>0,
\]
and let \(X_G\) be the rooted cyclic family transport.  On the ordered
second-return candidate define
\begin{equation}
\label{p6:eq:v564-ordered-return-pair}
 T_{\rightarrow}(a,b):=bK(a,b)X_G^\dagger,
 \qquad
 T_{\leftarrow}(a,b):=bK(a,b)X_G.
\end{equation}
Because \(K(a,b)\) has eigenvalues \(a+3b,a,a>0\) and is circulant, it is
positive and commutes with \(X_G\).  The
right polar decompositions \(T=U|T|\), with
\(|T|=(T^\dagger T)^{1/2}\), are therefore
\begin{equation}
\label{p6:eq:v564-polar-factors}
 |T_{\rightarrow}|=|T_{\leftarrow}|=bK(a,b),
 \qquad
 U_{\rightarrow}=X_G^\dagger,
 \qquad
 U_{\leftarrow}=X_G.
\end{equation}
For
\[
 C_{2,G}=-i(X_G-X_G^\dagger)=i(X_G^\dagger-X_G),
\]
the raw orientation-odd return is
\begin{equation}
\label{p6:eq:v564-raw-oriented-return}
 i(T_{\rightarrow}-T_{\leftarrow})
 =bK(a,b)C_{2,G}
 =ab\,C_{2,G},
\end{equation}
where the last equality uses \(JC_{2,G}=0\).  By contrast the
polar-normalized oriented transport is
\begin{equation}
\label{p6:eq:v564-polar-oriented-return}
 \boxed{
 C_{2,G}^{\rm pol}:=i(U_{\rightarrow}-U_{\leftarrow})=C_{2,G}.}
\end{equation}
Hence the sector coefficient \(ab\) belongs to the positive incidence
modulus, not to the normalized orientation transport.  On the sign line,
positive-polar normalization acts as \((ab)^{-1}\); for the physical
candidate bases \(b=1\), \(a_\ell=1\), \(a_C=2\), this is respectively
\(1\) and \(1/2\).
\end{theorem}
\begin{proof}
Since \(X_G\) permutes the three family coordinates, \(JX_G=J\) and
\(X_GJ=J\), so \([K(a,b),X_G]=0\).  Equation
\eqref{p6:eq:v564-ordered-return-pair} gives
\[
 T_{\rightarrow}^\dagger T_{\rightarrow}
 =b^2X_GK(a,b)^2X_G^\dagger=b^2K(a,b)^2.
\]
Positivity of \(K(a,b)\) gives
\(|T_{\rightarrow}|=bK(a,b)\), and therefore
\[
 U_{\rightarrow}
 =T_{\rightarrow}|T_{\rightarrow}|^{-1}
 =X_G^\dagger.
\]
The reverse orientation is identical with \(X_G\) and \(X_G^\dagger\)
interchanged.  Next,
\[
 i(T_{\rightarrow}-T_{\leftarrow})
 =ibK(a,b)(X_G^\dagger-X_G)=bK(a,b)C_{2,G}.
\]
Because \(JX_G=JX_G^\dagger=J\), one has \(JC_{2,G}=0\), hence
\(K(a,b)C_{2,G}=aC_{2,G}\).  This proves
\eqref{p6:eq:v564-raw-oriented-return} and
\eqref{p6:eq:v564-polar-oriented-return}.  The positive-polar construction
is the standard canonical separation of a full-rank operator into its
unitary and positive factors \cite{Higham1986Polar}; the same
inverse-positive-factor pattern is used in lattice gauge constructions to
extract a group-valued transport variable from a nonunitary matrix
\cite{ShibataKondoShinohara2010ExactDecomposition}.  These references support
the normalization mechanism only; they do not supply the present G15 q=5 paired-complement transport occurrence.
\end{proof}

\begin{corollary}[Coefficient-free polar-normalized diagnostic \(W\)-branch]
\label{p6:cor:v564-polar-normalized-w-branch}
Conditional on consuming the ordered-return branch through its canonical
positive-polar normalization, the diagnostic sign-line family
\eqref{p6:eq:v563-signline-repair} has no independent \(\lambda_q\)-port:
\begin{equation}
\label{p6:eq:v564-polar-w-branch}
 \boxed{
 W_{\pm,q}^{\rm pol}=F_\pm(a_q)+C_{2,G},
 \qquad a_\ell=1,\quad a_C=2.}
\end{equation}
The determinant formulas of
Proposition~\ref{p6:prop:v563-signline-rank-split} give
\[
 \det W_+^{\rm pol}(a)=-18(a+1),
 \qquad
 \det W_-^{\rm pol}(a)=2(3-a),
\]
so the four displayed lepton/color branches are full rank:
\[
 (\det W_{+,\ell},\det W_{-,\ell})=(-36,4),
 \qquad
 (\det W_{+,C},\det W_{-,C})=(-54,2).
\]
The continuous \(\lambda_q\)-family remains useful as a deformation and
capacity diagnostic, but coefficient one is fixed on this polar-normalized
ordered-return branch by the parent normalization rather than by numerical
simplicity or fit.
\end{corollary}
\begin{proof}
Theorem~\ref{p6:thm:v564-polar-separated-oriented-return} gives the common
normalized sign transport \(C_{2,G}\).  Substitution \(\lambda=1\) into
\eqref{p6:eq:v563-det-plus}--\eqref{p6:eq:v563-det-minus} gives the displayed
formulas and values.
\end{proof}



% BEGIN ROOT-COVARIANT-W-AND-CP v5.72
\begin{proposition}[Full rooted reversal of the diagnostic \(W\)-family]
\label{p6:prop:v572-full-root-reversal-of-W}
For real \(a,\lambda\), the complete rooted reversal acts on
\(W_+(a;\lambda)=F_+(a)+\lambda C_{2,G}\) by
\[
 \boxed{\overline{W_+(a;\lambda)}=W_-(a;-\lambda).}
\]
Consequently
\[
 \boxed{\det W_-(a;-\lambda)=\det W_+(a;\lambda)}
\]
and on every full-rank rooted pair
\[
 \boxed{\det W_+(a;\lambda)\det W_-(a;-\lambda)
       =[\det W_+(a;\lambda)]^2>0.}
\]
For \(a_C=2\), \(\lambda=1\), the two complete root pairings give
\((-54,-54)\) or \((2,2)\), not the mixed diagnostic pair \((-54,2)\).
\end{proposition}
\begin{proof}
Because \(K(a)\) is real, \(\overline{F_+(a)}=F_-(a)\), while full rooted
orientation reversal sends \(C_{2,G}\) to \(-C_{2,G}\).  This proves the
matrix identity.  The determinant identity follows either by conjugation or
by substituting \(-\lambda\) in
\eqref{p6:eq:v563-det-minus}.  Squaring then proves positivity on the
full-rank branch.
\end{proof}

\begin{theorem}[Positive color sign-line core]
\label{p6:thm:v572-positive-color-signline-core}
For
\[
 B_\pm:=2I+J\pm C_{2,G}
\]
one has
\[
 \boxed{\operatorname{spec}B_\pm=\{5,2+\sqrt3,2-\sqrt3\},
 \qquad B_\pm>0,
 \qquad \det B_\pm=5.}
\]
More generally,
\[
 \boxed{\det(aI+bJ+\lambda C_{2,G})
 =(a+3b)(a^2-3\lambda^2).}
\]
\end{theorem}
\begin{proof}
On the uniform family line, \(J\) has eigenvalue three and
\(C_{2,G}\) vanishes.  On its orthogonal trimer plane, \(J=0\) and
\(C_{2,G}\) has eigenvalues \(\pm\sqrt3\).  The displayed spectrum and
general determinant follow immediately.
\end{proof}

\begin{proposition}[Positive determinant is compatible with weak-CP capacity]
\label{p6:prop:v572-positive-determinant-weak-cp-capacity}
Let
\[
 D_\rho=\operatorname{diag}(1,\sqrt\rho,\rho),\qquad 0<\rho<1,
\]
\[
 M_\pm=D_\rho B_\pm D_\rho,
 \qquad H_\pm=M_\pm M_\pm^\dagger.
\]
Then
\[
 \det M_\pm=5\rho^3>0
\]
while exact multiplication gives
\[
 \boxed{
 \operatorname{tr}[H_+,H_-]^3
 =2400i\,\rho^6(\rho-1)^3(\rho+1)
 \bigl(21\rho^4+42\rho^3+58\rho^2+42\rho+21\bigr)^2,}
\]
which is nonzero for \(0<\rho<1\).  Positive determinant orientation thus
does not remove weak-family CP capacity.  This is a capacity theorem, not a
physical CKM member selection.
\end{proposition}
\begin{proof}
The determinant follows from Theorem~\ref{p6:thm:v572-positive-color-signline-core}
and \(\det D_\rho=\rho^{3/2}\).  The commutator identity is direct exact
three-by-three multiplication; every real factor displayed is nonzero on the
stated interval.
\end{proof}

\begin{remark}[The old coefficient-one pair is not a strong-CP invariant]
\label{p6:rem:v572-old-w-pair-not-strong-cp}
The coefficient-one values
\((\det W_{+,C}^{\rm pol},\det W_{-,C}^{\rm pol})=(-54,2)\)
were computed while holding the same sign-line coefficient fixed under only
the Hopf-phase flip.  They remain useful deformation/capacity diagnostics,
but they are not a complete rooted orientation pair and must not be read as
\(\arg\det(D_uD_d)=\pi\).  Full root reversal also flips the orientation-odd
sign line, as in Proposition~\ref{p6:prop:v572-full-root-reversal-of-W}.
\end{remark}
% END ROOT-COVARIANT-W-AND-CP v5.72

\paragraph{Typed occurrence firewall.}
The family transport generator \(C_{2,G}\), the binary-octahedral quotient
sign projector, the joint Hopf--axial sign line, and the single-label
exclusion owner are different typed objects.  Common sign-character language
does not identify them or transport coefficients between them.  In
particular the exact \(24/25\) exclusion owner does not descend through the
incidence map.  Theorem~\ref{p6:thm:v564-polar-separated-oriented-return}
removes the sector-dependent sign-line coefficient only \emph{within the
polar-normalized ordered-return branch}; it does not prove
\(\operatorname{Orig}_\Sigma\) of that normalized return in the physical
zero-order Yukawa source.  Nor does it exclude a physical compiler that uses
one of the other presently typed \(\Phi/A\) ports.  Thus
\eqref{p6:eq:v564-polar-w-branch} is the distinguished coefficient-free
normalized \(W\)-candidate.  Framework v2.86 closes the finite rooted $C_3$ source-cell selector upstream; relative CKM/PMNS frame realization and downstream pole/RG matching remain separately typed readouts and are not primitive transport selectors.

\subsection{Scalar-port projection and source-normalization comparisons}
\label{p6:app:v2168-transfer-diagnostics}
\subsubsection{Why projecting the scalar identity is a zero word}
This diagnostic explains the precise constructor distinction; it does not
assert that the actual charged source uses the failed word.
Under any unitary presentation of a rank-four owner with orthogonal
hierarchy supports $Q_4,Q_5$, its identity is still the identity. Therefore
\[
 \boxed{Q_5I_4Q_4=0,\qquad Q_4I_4Q_5=0.}
\]
Moreover the off-diagonal edge $T(v)=Q_5L_vQ_4$ is traceless, so
\[
 \boxed{\epsilon_{E_4}(T(v))=0,}
\]
where this trace is taken after the typed rank-four presentation is supplied.
Thus neither projecting the primitive scalar unit nor tracing the angular
edge produces the nonzero mixed source coefficient. The combined contraction
instead uses $\epsilon_{E_4}(s_4)$ and $T(v_\mu)$ as retained sibling ports.
It neither projects $I_4$ into the mixed corner nor infers the angular edge
from a scalar trace.

The primitive lift's kernel is harmless for its stated scalar/identity
transfer. It is not harmless if one tries to reconstruct an independently
consumed noncentral family relation from that scalar alone. Such a relation
must remain in the complete owner packet before the contraction.

\subsubsection{Port normalization is not an actual source-export theorem}
This appendix tests only the reduced data consisting of primitive transfers,
hierarchy supports and normalized current rays. It is not a counterexample to
a more complete physical source constructor with additional retained relations.
On that reduced data both geometric comparison outputs
\[
 R_1=\varepsilon_eB_e+\varepsilon_\mu B_\mu+\varepsilon_\tau B_\tau,
\]
\[
 R_2=\varepsilon_eB_e+s\varepsilon_\mu B_\mu+s\varepsilon_\tau B_\tau,
 \qquad s=\sqrt{8/9},
\]
have the same normalized left rays and hierarchy supports. Their uninserted
clock returns differ by the muon/tau factors $s^2=8/9$. The scalar $s$ is
already available as the raw Hopf-compression norm, so the distinction is not
merely an absence of possible geometric numbers. It is which actual source
word consumes which ports and inherited normalization.
This does not authorize $R_2$ as a physical source or add a free coefficient.
It shows why the explicit actual-source identity must be established before
GP2 normalization inheritance or the no-new-choice criterion is invoked.

The v2.166 bounded-required-readout theorem bounds a constituted family on
its admitted stage. Neither boundedness nor injectivity turns a source
amplitude into unit transfer fidelity. The endpoint-fidelity theorem likewise
requires an already generated normalized relation and a proved defect
classification; it explicitly does not select a coupling or mass ratio.


% END INLINE SOURCE: Paper5B-ChargedTransferDiagnostics-Appendix-v1_0.tex


% v2.189: supplementary source-evaluation lineage and comparison scope.

% BEGIN INLINE SOURCE: Paper5B-ElectronExecutionLineage-Appendix-v1_0.tex
% Framework v2.189: supplementary current-source execution lineage.
% No new source occurrence premise, comparison choice or pole closure.

\section{Gauge, scalar and common-action comparison calculations}
\label{p5d:sec:E}
\noindent\textit{Technical owner.} These exact inherited calculations
support the corresponding statements in Papers~V-A--V-C;
the appropriate action/observable occurrence conditions remain in those papers.
\subsection{Renormalizable operator skeleton and coefficient typing}
\label{p7v2:sec:operator-skeleton}

\begin{definition}[Finite interaction boundary]
\label{p7v2:def:finite-interaction-boundary}
The Paper--V-C finite interaction boundary is the typed object
\begin{equation}
\label{p7v2:eq:finite-interaction-boundary}
 \mathfrak B_{5C}
 =
 \left(
  \Istr,\,
  \Irole\big|_{\Urole},\,
  \Brawsix[b_6],\,
  \mathcal C_{\rm alg}^{\rm gauge},\,
  \mathcal C_{\rm alg}^{\rm QCD},\,
  \mathcal M_N^{\rm aff},\,
  \mathfrak I_{\rm QFT},\,
  \mathfrak I_{\rm QCD},\,
  \mathfrak I_{\rm th},\,
  \mathfrak O_{5C}
 \right),
\end{equation}
where \(\Istr\) is structural, named physical roles are restricted to
\(\Urole\), and \(\Brawsix[b_6]\) denotes a frozen compatible member
\(b_6\in\Acompsix\) whenever a spectrum/flavor readout is used.  Physical
energy ratios and absolute values additionally use the exact domains
\eqref{p7v2:eq:energy-import-domains};
the \(\mathcal C_{\rm alg}\) entries are exact finite algebraic data, the
\(\mathfrak I\) entries are explicitly conditional interfaces, and
\(\mathfrak O_{5C}\) is the open-obligation ledger of
Table~\ref{p7v2:tab:open-obligations}.  An interface or obligation is not a
hidden coefficient: it is part of the type of the exported boundary.
\end{definition}

\begin{proposition}[Conditional renormalizable Lagrangian dictionary]
\label{p7v2:prop:lagrangian-dictionary}
Assume \(\Hqft+\Urole\).  Fix a compatible
\(b_6\in\Acompsix\) for every displayed charged Yukawa or heavy-spectrum
coefficient.  A coefficient-level \(Y_\nu\) additionally requires one
\(\Hnu\) fiber member; without \(\Hnu\), only its allowed representation slot
is asserted.  A physical scalar potential and physical strong-CP parameter
respectively add \(\Hvac\) and \(\Hcp\).
Use the inherited signature
\(g_{\rm op}=-c_{\rm op}^2d\widetilde t^{\,2}+h\), the inherited
spinor-adjoint convention, and the same metric to raise every spacetime
index.  Then the renormalizable interaction skeleton may be written as
\begin{align}
\label{p7v2:eq:sm-lagrangian-skeleton}
 \mathcal L_{\rm SM+N}
 ={}&
 -\frac14\sum_{i=1}^{3}F_{i,\mu\nu}^{a_i}F_i^{a_i,\mu\nu}
 +\sum_\psi i\bar\psi\slashed D\psi
 -(D_\mu H)^\dagger D^\mu H
 -V(H)
 \nonumber\\
 &-\left[
   \bar Q\,Y_d H d
  +\bar Q\,Y_u\widetilde H u
  +\bar L\,Y_e H e
  +\bar L\,\YN\widetilde H N
  +\frac12\overline{N^c}D_NN
  +{\rm h.c.}
 \right]
 \nonumber\\
 &+\frac{\theta_3g_3^2}{32\pi^2}
   F^a_{\mu\nu}\widetilde F^{a,\mu\nu}.
\end{align}
The role-typed carrier fixes the representation slots and allowed
intertwiners.  It does not fix their coefficient values.  In particular,
\(D_N\) belongs to the compatible Paper--V-B base, whereas \(\YN\) belongs to
the Paper--V-C neutral-lift fiber; they are not one premise.
\end{proposition}

\begin{proof}
Gauge invariance and power counting list the displayed operators once the
physical-role carrier is imported.  Listing an allowed operator does not fix
its coefficient.  The scalar sign follows from the declared metric: in a
local orthonormal frame, write the doublet as four real components
\(H=(\varphi_1+i\varphi_2,\varphi_3+i\varphi_4)^T/\sqrt2\).
The free derivative term is
\(\frac12\sum_a(\dot\varphi_a^2-|\nabla\varphi_a|^2)\),
whose Hamiltonian derivative energy is
\(\frac12\sum_a(\dot\varphi_a^2+|\nabla\varphi_a|^2)\).
This agrees with Paper~V-A's native local scalar action; a mostly-minus
comparison presentation must transform the whole metric/adjoint dictionary
rather than changing this scalar sign alone.
The \(Y_\nu\) coefficient status follows from
Definition~\ref{p7v2:def:nu-selector-family}; the other coefficient statuses
follow respectively from
Propositions~\ref{p7v2:prop:strong-seed-status} and
\ref{p7v2:prop:threshold-obligation},
Theorems~\ref{p7v2:thm:conditional-potential} and
\ref{p7v2:thm:no-vacuum-selection},
Proposition~\ref{p7v2:prop:pmns-cancels}, and
Proposition~\ref{p7v2:prop:physical-strong-cp}.
\end{proof}

\subsection{Open-data ledger}
\label{p7v2:sec:open-ledger}

\begin{longtable}{@{}L{0.22\textwidth}L{0.31\textwidth}L{0.39\textwidth}@{}}
\caption[Premise-minimized open-data ledger]{Minimal objects required to close
the physical Paper--V-C sectors, with optional model extensions explicitly
marked.}\label{p7v2:tab:open-obligations}\\
\toprule
Target & Missing object & Required discharge\\
\midrule
\endfirsthead
\toprule
Target & Missing object & Required discharge\\
\midrule
\endhead
Absolute \(g_1,g_2,g_3\)
& \(\Hgauge\)
& On the common-heat branch, use the derived certificate
\(\Rconn=(\JOINTfiveocc)\); the common heat trace then gives
\(k_i g_i^2=4\pi^2\).  On the broader kinetic branch add the single
\(\CALgauge\) member.  No second strong-seed/E16 coordinate is required.\\
Precision \(\alpha_{\rm em}(0)\)
& dependent \(\Hem\) infrared certificate
& Use the same physical \(\Gamma_{\rm eff}\) and the already-generated occurring-solder
\(U(1)_Q\) restriction; form the same-\(\Gamma\) quadratic response and prove the
unique \(q^2\to0\) limit/matching.  No scalar-stationary certificate is an
upstream prerequisite.  Any finite operator formula is only a representation
candidate until identified with this response.\\
\(\alpha_{\rm em}(M_Z)\)
& Hadronic current correlator
& Compute it downstream without importing a Paper--V-C-dependent QCD result
back into Paper~V-C.\\
Physical threshold/unification
& Heavy fluctuation operators
& Specify masses, spins, multiplicities, Dynkin indices, scheme constants, and
higher-loop accuracy.\\
Scalar stationary/effective-action claims
& \(\Hvac\), claim-scoped
& Only when a downstream claim explicitly needs scalar equilibrium, prove the
required stationarity, Hessian property, energetic comparison, or relaxation
history.  These claims do not constitute the native electroweak stabilizer or
gauge-mass kernel.\\
Physical Higgs couplings
& Renormalized two- and three-point functions
& Compute vertices, pole residues, scheme conversion, and uncertainties.\\
Ordinary spin--\(Z_6\) global anomaly
& Actual regulated anomaly/background descent
& The bordism obstruction vanishes on
Theorem~\ref{p5:thm:v2210-lit-anomaly-obstruction}'s exact domain;
identify the actual native determinant/measure there.  Boundaries and
other spin--gauge categories retain their own check.\\
Regulated strong-CP completion
& \(\Hcp\), claim-scoped
& Preserve the closed native finite character and certify its complete
regulated action/measure transport, compatible with weak CP.  There is no
second finite sector-phase selector.\\
Current harmonic neutral origin
& Result-independent K56 mass/incidence constructor
& Derive the consumed harmonic/local-port word and source insertion.
The conditional gap ratio in
Theorem~\ref{p6:thm:v2210-lit-harmonic-ratio} does not select that
constructor or its physical phase/detector map.\\
Nonperturbative QCD
& Hadronic state and evolution
& Solve or otherwise construct the mass gap, condensates, PDFs, amplitudes, and
hadronization boundary.\\
Heavy high-energy neutral lift and physical source
& \(\Hnu\) plus the stated QFT/width domain
& Freeze \(\Uroleabs+b_6\), supply one high-energy \(Y_\nu\) fiber member,
and derive its invariants and widths.  Its Casas--Ibarra \(R\) is a coordinate,
and the regulator is then a readout rather than another premise.\\
Optional exact saturation
& \(\Hresopt\)
& Use only for a deliberately saturated benchmark.  The repaired equation is
exact and has the unique branch proved in
Theorem~\ref{p7v2:thm:heavy-calibration-reduction}; saturation is not required
by the standard target.\\
Optional Hopf source model
& \(\Hhopfopt\)
& Supply a typed Hopf operator, state, and a new typed map to the CP source.
Without that map, the exact label/Hilbert averages are diagnostics only.\\
Thermal efficiency
& \(\Hkin\)
& Specify and solve production, oscillation, decoherence, inverse-decay,
washout, \(N_1\), flavor, spectator, and expansion dynamics.\\
Baryogenesis
& \(\Hnu+\Hkin\)
& Compute the microscopic source first and propagate it with certified
numerical or analytic errors.\\
\bottomrule
\end{longtable}

\begin{theorem}[Premise-minimized finite interaction boundary]
\label{p7v2:thm:final-boundary}
On the structural domain \(\Ustr\), Paper~V-C fixes:
\begin{enumerate}[label=(\roman*),leftmargin=2.5em]
  \item the derived active support/clock anchors
  \(2T\mapsto32\), \(C_6\mapsto40\), and scalar reclosure
  \(\mapsto56\), together with their finite split grammars and the terminal
  identity \(P_{\rm term}=P_0=\mathbb E_{\rm IR}\);
  \item the defect-count identities
  \eqref{p7v2:eq:defect-counts}--\eqref{p7v2:eq:defect-sum-difference};
  \item the structural \(3\otimes\bar3=1\oplus8\) decomposition; and
  \item the exact finite distinctions between ordinary and normalized
  determinants and between uniform-label and Hilbert-space averages.
\end{enumerate}
On \(\Urole\), the corresponding physical Standard--Model role names,
Yukawa-support dictionary, and physical color interpretation are licensed.
The local-gauge certificate \(\Rconn=(\JOINTfiveocc)\) is derived; factor
connections, curvatures, profiles, and the typed common-trace relations
require no occurrence premise or second relation selector.

After a compatible \(b_6\in\Acompsix\), the exact role/energy domain needed by
the claim, and the interfaces named in its statement are supplied, Paper~V-C
also gives the relative heavy affine splitting
\eqref{p7v2:eq:heavy-splitting}, the type-I identity and complete positive
nonsingular \(O(3,\CC)\) fiber classification, the electroweak dictionary,
perturbative gauge/QCD evolution, potential reconstruction, rare-process
formulas, thermal field count, and sphaleron conversion.

Independently of a physical realization, ordinary mathematics proves the
resonant bound on \(\mathbb R^2\setminus\{0\}\) and the exact block/full
thermal-algebra classification.  If the optional exactly saturated benchmark
is imposed, its angle is closed form and its remaining scalar equation has the
unique certified root of
Proposition~\ref{p7v2:prop:heavy-calibration-numerics}.  These theorems do not
promote the optional source or saturation input into a structural consequence.

For the declared standard high-energy neutral and kinetic comparison targets,
the conditional Paper--V-C interface ledger is
\begin{equation}
\label{p7v2:eq:minimal-physical-package}
\boxed{
\begin{cases}
\{\Hcp,\ \Hnu,\ \Hkin\},
& \text{source-native common-heat branch},\\
\{\CALgauge,\ \Hcp,\ \Hnu,\ \Hkin\},
& \text{broader kinetic-class branch}.
\end{cases}
}
\end{equation}
This is not a count of independent remaining native inputs: native finite
action, source, phase and state closures are reused on their stated domains.
The scalar-stationary certificate \(\Hvac\) is not a member of this comparison
package.  It is appended only to those claims that explicitly require scalar
equilibrium, scalar-Hessian/Goldstone information, global energetic comparison,
or actual scalar relaxation/history.  Analytic theorem domains, the derived
\(\Rconn\) import, upstream compatible members, calibrations, and downstream
nonperturbative data retain their separately typed statuses.  \(\Hresopt\) and
\(\Hhopfopt\) are optional extensions and are not members of
\eqref{p7v2:eq:minimal-physical-package}.

On the source-native common-heat branch the absolute gauge couplings are fixed
at the heat matching datum; on the broader kinetic branch they still require
\(\CALgauge\).  The precision electromagnetic endpoint is a derived-dependent
same-effective-action/IR-limit certificate rather than a member of the owner
package above.  The structural core or a compatible finite/raw member alone
does not fix its analytic limit, exact gauge unification, scalar equilibrium
or its actual history, physical Higgs vertices, an alternative chiral-continuum
strong-CP realization, a unique
high-energy neutrino sector, a thermal efficiency, or \(\eta_B\).
\end{theorem}

\begin{proof}
The active realization is the structural handoff from Papers~III-C and V-B,
and Lemma~\ref{p7v2:lem:defect-counts} is purely finite.  The abstract color
decomposition is structural; its physical role is restricted exactly as in
\eqref{p7v2:eq:physical-role-import}.  Paper~III-F and Paper~V-A supply the
derived \(\Rconn\) datum and its same-datum closure.  The conditional heavy-splitting and
seesaw statements consume the compatible Paper--V-B base, while
Corollary~\ref{p7v2:cor:casas-ibarra-fiber-coordinate} proves that no second
orthogonality datum remains.

The natural resonant domain and thermal algebra branches are
Theorems~\ref{p7v2:thm:resonant-bound} and
\ref{p7v2:thm:thermal-trace-nonunique}.  The exact saturated reduction and
certified bracket are Theorem~\ref{p7v2:thm:heavy-calibration-reduction} and
Proposition~\ref{p7v2:prop:heavy-calibration-numerics}.  Definition
\ref{p7v2:def:interface-hypotheses}, the counterfamilies in
Theorem~\ref{p7v2:thm:type-separation}, and
Proposition~\ref{p7v2:prop:hopf-optional} remove RES and HOPF from the standard
target while retaining them as optional models.

The conditional results name their hypotheses at their statements.  The
comparison non-closure clauses use the reduced-data tests in
Theorem~\ref{p7v2:thm:type-separation} and the stated sector-specific domains;
they do not establish competing complete native owners after the generated
prefix is retained.  Native source, action and state closures keep their
proved domains, while actual downstream expectations and analytic exports
retain their explicit certificates.
\end{proof}

\begin{corollary}[No-hidden-input statement in its valid form]
\label{p7v2:cor:no-hidden-valid}
Within the status-\(\mathsf T\) claims of
Theorem~\ref{p7v2:thm:final-boundary}, no fitted comparison value or downstream
Paper-VI-A/VII-A--VII-C output is used.  Every statement of another status exposes
its additional data in the same section or in
Table~\ref{p7v2:tab:open-obligations}.
\end{corollary}

\begin{proof}
The structural proofs and named physical representation roles use
\(\Urole=\Ustr\) and finite algebra.  Every numerical spectrum/flavor statement
additionally declares a compatible \(b_6\in\Acompsix\), the required
ratio/absolute-energy domain, and every Paper--V-C interface it consumes.  The dependency order is
\eqref{p7v2:eq:dependency-order}; the conditional and open objects are typed by
Definitions~\ref{p7v2:def:claim-status} and
\ref{p7v2:def:interface-hypotheses}.
\end{proof}


% BEGIN v2.36 RETYPED PHYSICAL INTERFACES
\subsection{Common-heat/shell scale diagnostic}
\label{p7:app:common-heat-shell-scale-diagnostic}
The same-scale incompatibility is exact:
\[
 \left.\frac{m_W^2}{m_Z^2}\right|_{\rm heat}=\frac58,
 \qquad
 \left.\frac{m_W^2}{m_Z^2}\right|_{\rm shell,EQAMP}=\frac79.
\]
Under the threshold-free one-loop hypotheses of
Corollary~\ref{p7:cor:conditional-common-heat-shell-transport}, these occur at
scales separated by \(\log(\mu_{79}/\mu_h)=-88/215\).  If, only as a further
comparison chart, one identifies scale ratios with a fixed Paper--III/V-B
shell coordinate
\[
 \frac{\mu(K+\Delta K)}{\mu(K)}
 =\alpha_{\rm hier}^{-\Delta K/32},
 \qquad \alpha_{\rm hier}>1,
\]
then the implied coordinate separation is
\begin{equation}
 \label{p7:eq:conditional-shell-coordinate-separation}
 \boxed{\Delta K=\frac{2816}{215\log\alpha_{\rm hier}}.}
\end{equation}
This last equation is a diagnostic re-expression, not a proof that RG scale
and shell coordinate share an owner.

\begin{theorem}[Carrier symmetry alone does not select a scalar stationary branch]
\label{p7v2:thm:no-vacuum-selection}
The carrier, gauge representation, GP2, and GP3 do not by themselves select
the sign of a conventional scalar quadratic coefficient or prove that the
symmetric scalar coordinate is unstable.
\end{theorem}
\begin{proof}
For \(m^2,\lambda>0\), the two conventional effective potentials
\[
 V_\pm(H)=\pm m^2H^\dagger H+\lambda(H^\dagger H)^2
\]
have the same carrier and gauge representation but opposite Hessian sign at
\(H=0\).  This is a diagnostic of the optional scalar-stationary problem; it
is not a counterexample to the native occurring-solder electroweak theorem.
\end{proof}

\begin{corollary}[Native positive heat branch is excluded as the complete scalar-stationary generator]
\label{p7v2:cor:native-heat-vacuum-exclusion}
On the allowed nontrivial source branch of this appendix, the exact positive
source-native heat contribution satisfies
\[
 S_{\rm heat}'(r)<0\qquad(0\le r<\infty),\quad r=H^\dagger H,
\]
whenever the squared source operator is nonzero.  Therefore that exact heat
term alone has no finite stationary broken scalar orbit.  This constrains only
the claim-scoped scalar-equilibrium branch; it does not reopen native
electroweak reduction.
\end{corollary}

\subsection{Exact-heat vacuum diagnostic and truncation boundary}
\label{p7:app:exact-heat-vacuum-diagnostic}
On the allowed Dirac-neutrino branch, for any generated Yukawa matrices the
constant-field squared Yukawa operator has positive form
\(rB\), \(r=H^\dagger H\), \(B\ge0\).  The exact source-native positive heat
contribution is
\[
 S_{\rm heat}(r)=A\operatorname{Tr}e^{-rB/\Lambda^2},
 \qquad A>0,
\]
and therefore
\[
 S_{\rm heat}'(r)
 =-\frac{A}{\Lambda^2}
   \operatorname{Tr}\!\left(Be^{-rB/\Lambda^2}\right)<0
\]
for every finite \(r\) whenever \(B\ne0\).  Noncommuting flavor geometry
changes the spectrum of \(B\) but not this sign.  Hence this exact positive
heat term alone has no finite broken scalar stationary orbit.

Nor does its quartic truncation supply a controlled counterexample.  Write
\(A_1=\operatorname{Tr}B\) and \(A_2=\operatorname{Tr}B^2\).  Its formal
stationary point is \(r_*=\Lambda^2A_1/A_2\).  If \(b_{\max}\) is the largest
eigenvalue of \(B\), then
\[
 A_2\le b_{\max}A_1,
 \qquad
 \frac{r_*b_{\max}}{\Lambda^2}\ge1.
\]
At least one mode is therefore outside a uniformly small heat-expansion
domain at the alleged minimum.  A viable vacuum must come from an occurring
renormalized or otherwise physical effective action satisfying an independent
criterion such as
Theorem~\ref{p7:thm:sufficient-stable-broken-vacuum}.

\subsubsection{Appendix: exact free-photon spectral corollaries}
\label{p7v2:app:v245-photon-spectral-corollaries}
On the stationary regular free transverse photon branch write
\(n=L+1\ge2\).  The two helicities give degeneracy
\[
 d_n=2(n^2-1).
\]
Hence the coexact spectral zeta function is
\[
 \boxed{\zeta_\gamma(s)
 =2R_x^{2s}\bigl[\zeta_R(2s-2)-\zeta_R(2s)\bigr].}
\]
Under zeta regularization specifically, the formal zero-point spectral value is
\[
 \boxed{E_{0,\gamma}^{\zeta}
 =\frac{11}{120}\frac{\hbar c_{\rm op}}{R_x}
 =\frac{11}{120}\Eper.}
\]
This is a spectral/Casimir regularization value, not a derived cosmological
vacuum energy or a replacement for the gravitational counterterm owner.

The central antipode on either coexact vector representation at level \(L\)
acts by \((-1)^{L+1}=(-1)^n\), while a temporal half-turn acts by
\(e^{-i\omega_nT_t/2}=(-1)^n\).  Therefore every free photon mode obeys
\[
 \boxed{A_n(t+T_t/2,-x)=A_n(t,x).}
\]
Thus the free mode representation descends through the diagonal
\(\mathbb Z_2\) quotient of \(S_t^1\times S_x^3\).  The group comparison
\((U(1)\times SU(2))/\mathbb Z_2^{\rm diag}\cong U(2)\) is representation
theory only; it does not reconstitute physical spacetime as \(U(2)\).

On a regular bosonic thermal branch the free-photon partition function is
\[
 Z_\gamma=\prod_{n=2}^\infty
 (1-e^{-\beta n\Eper})^{-2(n^2-1)},
\]
with one-particle trace
\[
 \boxed{Z_{1,\gamma}(x)=\frac{2x^2(3-x)}{(1-x)^3},
 \qquad x=e^{-\beta\Eper}.}
\]
The exact thermal energy on that regular branch is
\[
 \boxed{U_\gamma
 =2\Eper\sum_{n=2}^\infty
 \frac{n(n^2-1)}{e^{\beta n\Eper}-1}.}
\]
Finite occupation carriers use the finite-trace boundary-corrected calculus of
Paper~IV instead.  The large-radius regular-branch limit recovers the usual
Planck/Stefan--Boltzmann photon density
\(u=\pi^2k_B^4T^4/(15\hbar^3c_{\rm op}^3)\).



% BEGIN S3 LITERATURE INTEGRATION v3.040 [V-D: MODULAR PHOTON ENERGY]
\begin{corollary}[Exact modular reexpression of regular free-photon energy]
\label{p7v2:cor:v340-s3lit-photon-modular}
On the already constituted stationary \emph{regular bosonic free
transverse-photon branch}, set $t=\beta\Eper>0$,
$y=4\pi^2/t$, and $L_s(y)=\sum_{n\ge1}n^s/(e^{ny}-1)$.
The preceding exact Gibbs energy has the all-temperature identity
\[
 \boxed{\begin{aligned}
 \frac{U_\gamma}{\Eper}={}&\frac{2\pi^4}{15t^4}
 -\frac{\pi^2}{3t^2}+\frac1t-\frac{11}{120}\\
 &+\frac{32\pi^4}{t^4}L_3(y)
 +\frac{8\pi^2}{t^2}L_1(y).
 \end{aligned}}
\]
The final two terms are positive and exponentially small as $t\to0$.
The term $-11\Eper/120$ cancels only if the \emph{separately
zeta-regularized} zero-point contribution from this source is included.
This identity does not extend without further proof to a GP1 finite
occupation carrier, a fixed finite photon band, or a physical
clock/thermal duality transformation.
\end{corollary}
\begin{proof}
The existing energy equals
$2\Eper[L_3(t)-L_1(t)]$.  With
$\tau=it/(2\pi)$, the conventional holomorphic Eisenstein series
obey $E_4(\tau)=1+240L_3(t)$ and
$E_2(\tau)=1-24L_1(t)$.  Apply the modular relation for $E_4$
and the quasimodular relation for $E_2$ at
$-1/\tau$ (whose Fourier parameter is $y=4\pi^2/t$), then
simplify.  No extra thermal oscillator or dynamical parameter appears.
\end{proof}
% END S3 LITERATURE INTEGRATION v3.040 [V-D]

\begin{proposition}[Conditional all-level \(7/9\) transverse sum rule]
\label{p7v2:prop:v247-seven-nine-all-level}
On the explicitly conditional tree-shell branch on which
\(m_W^2/m_Z^2=7/9\), every calibrated transverse harmonic level satisfies
\[
 \boxed{
 \frac{E_{W,n}^2-E_{\gamma,n}^2}
 {E_{Z,n}^2-E_{\gamma,n}^2}=\frac79.}
\]
This is only the level-independent rest-Hessian consequence of the conditional
\(7/9\) mass owner.  It neither derives that owner nor identifies a running
or pole-mass ratio at a named experimental scale.
\end{proposition}

% BEGIN HM v1.79 EW02
\subsubsection{Common-scale heavy-row compatibility}
\label{p7:app:v179-heavy-row-compatibility}

This diagnostic tests equality of particular ledger rows with one
canonical action at one matching slice.  It is not a comparison of
unrun heat coefficients with measured pole values.
If the same heavy factor \(s>0\) were to give
\(m_W=s\sqrt7\), \(m_Z=3s\), \(m_H=s\sqrt{17}\), and
\(m_t=s(\sqrt7+3)\) at that slice, the canonical action dictionary would
require
\[
 \frac{g_Y^2}{g_2^2}=\frac27,\qquad
 \frac{\lambda_Y}{g_2^2}=\frac{17}{56},\qquad
 \frac{y_t^2}{g_2^2}=\frac{8+3\sqrt7}{7}.
\]
These follow by dividing \(m_Z^2,m_H^2,m_t^2\) by
\(m_W^2=g_2^2v^2/4\), using
\(m_H^2=2\lambda_Yv^2\) and \(m_t^2=y_t^2v^2/2\).
The native common-heat gauge ratio \(g_Y^2/g_2^2=3/5\) is different.
A physical use therefore needs the declared running, source-amplitude
and comparison maps at the same slice.  No independent coefficient is
introduced to repair that difference.
% END HM v1.79 EW02

\begingroup
\footnotesize
\raggedright
\sloppy
\setlength{\emergencystretch}{8em}
\setlength{\hfuzz}{1pt}
\endgroup

% BEGIN SAFEGUARD-FLAVOR-INTERACTION-SYNC v2.63

\section{Weak-current and CP matrices, density receivers and response tests}
\label{p5d:sec:F}
\noindent\textit{Technical owner.} These exact inherited calculations
support the corresponding statements in Papers~V-A--V-C;
the appropriate action/observable occurrence conditions remain in those papers.
\subsection{Carrier-reduct comparison diagnostics}
\label{p5c:app:v2187-carrier-comparison-scope}
These examples test what follows from a carrier-only reduct.  The current
native action/source/state prefixes remain in the main text; optional
comparison data cannot be inserted into them without their own constructor.

\begin{theorem}[Carrier-reduct comparison does not determine the omitted readouts]
\label{p7v2:thm:type-separation}
If only the carrier dimensions and automorphism groups are retained, without
the corresponding complete action, operator, state or evolution owner, those
comparison data do not uniquely determine an absolute gauge coupling, an electromagnetic endpoint
determinant, the sign of the Higgs quadratic term, a resonant saturation point,
a thermal state on a reducible carrier, or a late-time baryogenesis efficiency.
\end{theorem}

\begin{proof}
The following mathematical comparison families preserve that reduct while
changing omitted target data.  They are not asserted to be distinct licensed
native complete-action or state/history words.
\begin{enumerate}[label=(\alph*),leftmargin=2.5em]
  \item On the broader kinetic-class branch before the common-heat
  specialization, from the same color counts \(d_A=8\) and \(d_P=9\), every
  \(\alpha_s^{(t)}=(d_Ad_P)^{-t}\), \(t>0\), is positive,
  dimensionless, duality-invariant, and uses no external scale.
  \item The family
  \[
    e^{k a^2}
    \left(1-a^4-a^6+c\,a^8\right)^{-1/8},
    \qquad k>0,\ c\in\RR,
  \]
  preserves the same finite rank and monomial support while changing the
  electromagnetic readout.
  \item For \(m^2,\lambda>0\), both
  \[
    V_\pm(H)=\pm m^2H^\dagger H+\lambda(H^\dagger H)^2
  \]
  have the same carrier and gauge symmetry, but only \(V_-\) has a broken
  minimum.
  \item The standard resonant factor takes every value in
  \([-\tfrac12,\tfrac12]\) as its width-to-splitting ratio varies.
  \item On the block observable algebra
  \(\End(\mathcal V_7)\oplus\End(\mathcal V_9)\), every
  \[
    \tau_w(A)
    =
    w\,\frac{\Tr A_7}{7}
    +(1-w)\,\frac{\Tr A_9}{9},
    \qquad 0\leq w\leq1,
  \]
  is a normalized positive trace invariant under
  \(U(7)\times U(9)\).
  \item With the same finite source, the zero collision kernel, a finite
  washout kernel, and an arbitrarily large washout kernel give different
  late-time asymmetries.
\end{enumerate}
Thus these retained comparison data permit distinct mathematical outputs.
Actual native evaluation instead consumes the licensed corresponding owner
of Definition~\ref{p7v2:def:interface-hypotheses}, reusing its generated
coefficients, state and evolution.  The examples prove no native
underdetermination after that complete owner is retained.
\end{proof}


% BEGIN INLINE SOURCE: Paper5C-NativePairComparisons-Appendix-v1_0.tex
% Framework v2.171: integrated conversation proof on its stated domain.
\subsection{Exact field-slot diagnostic: injectivity is weaker than mixed retention}
\label{p5c:v2171:app:fieldslot}
This appendix gives a comparison chart. It is not a licensed alternative
K16 realization. Its purpose is to expose the exact antecedent in the positive
retention theorem, rather than infer a native obstruction from host mathematics.

Let $\phi_i\in\mathbb C^m$ be unit mode columns and define an isometry
$Je_i=\phi_i\otimes e_i$. Put
\[
 C_{ij}=\phi_i^\dagger\phi_j,\qquad
 K_{ij}=\phi_i^\dagger\overline{\phi_j}.
\]
Then the canonical normal and pair pullbacks are
\[
 \widehat A=C\circ A,\qquad \widehat M=K\circ M.
\]
Here $\circ$ is the entrywise product. On the spin-singlet slot
$\epsilon=\left(\begin{smallmatrix}0&1\\-1&0\end{smallmatrix}\right)$,
using the explicitly normalized spin trace $\Tr_{\rm spin}/2$, the marked
contraction equals
\begin{equation}\label{p5c:v2171:eq:rowoverlap}
 Z_J=\gamma Z,\qquad
 \gamma=\overline{K_{22}}K_{11}C_{21}^2,
 \qquad \operatorname{Im}Z_J=
       \frac{24\operatorname{Re}\gamma-7\operatorname{Im}\gamma}{288}.
\end{equation}
\begin{proof}
Insert $J$ on both normal slots and the conjugate slot on the pair.
The unique nonzero marked normal entry is $A_{21}$. Since $M$ is diagonal,
the trace is
$\overline{M_{22}K_{22}}(M_{11}K_{11})(C_{21}A_{21})^2$.
The normalized spin trace contributes
$\Tr\epsilon^\dagger\epsilon/2=1$. This gives
\eqref{p5c:v2171:eq:rowoverlap} and its real form.
\end{proof}

For the exact two-mode comparison choose
\[
 \phi_0=\phi_1=(1,0)^{\mathsf T},\qquad
 \phi_2=\frac{(\sqrt{14},3+4i)^{\mathsf T}}{\sqrt{39}}.
\]
All entries of $C$ and $K$ are nonzero. Thus the normal map is injective on
all $M_3$, and the pair map is injective on all symmetric family matrices.
Nevertheless
\[
 K_{22}=\frac{7-24i}{39},\quad C_{21}^2=\frac{14}{39},\quad
 \gamma=\frac{98+336i}{1521},\quad
 Z_J=-\frac{4375}{219024}\in\mathbb R.
\]
The marked CP-odd contraction is zero. The multiplication guard fails:
\[
 \pi_J(E_{21})\pi_J(E_{12})=\frac{14}{39}E_{22}
 \ne \pi_J(E_{22}).
\]
The comparison refutes only the generic implication from separate
injectivity to retention of this mixed contraction. The actual K16 grammar
has not licensed this comparison chart. It supplies no native no-go theorem.

\subsection{Full cyclic coefficient table and scope comparison}
\label{p5c:v2171:app:matrixscope}
Write $A=a^\dagger$, $B=b^\dagger$, $C=c^\dagger$, $D=d^\dagger$.
In $\mathcal R$ the following cyclic words have coefficient $+1/8$:
\[
 AdDa,\ ADda,\ AabB,\ AadD,\ AaDd,\ ABba,\
 BbCc,\ BcCb,\ BbcC,\ BCcb,\ CcdD,\ CDdc.
\]
The following have coefficient $-1/8$:
\[
 BdCa,\ BCda,\ ABab,\ BadC,\ BaCd,\ AbaB,\
 AbDc,\ AcDb,\ AbcD,\ ADcb,\ CDcd,\ CdcD.
\]
All other cyclic words have zero coefficient. The table gives the entire
calculated response and equals \eqref{p5c:v2171:eq:matrixcoefficient}.

For a formula-scope comparison take a four-dimensional bosonic coefficient
space with $a=|0\rangle\langle2|$, $b=|0\rangle\langle3|$,
$c=|1\rangle\langle2|$, $d=|1\rangle\langle3|$.
All four lower charge for $Q_{\rm bos}=\operatorname{diag}(0,0,1,1)$, while the
exact normalized trace gives $\mathcal R=-1/16$.
This comparison is not a native coordinate or state selection. It verifies
why charge covariance alone does not make the general matrix formula
positive. The coordinate and complementary-slot results retain their stated
positive domains.

% END INLINE SOURCE: Paper5C-NativePairComparisons-Appendix-v1_0.tex



\subsection{The \(7\mid9\) response-amplitude counterfamily}
\label{p7:app:seven-nine-counterfamily}
For every \(t>0\), the block response
\[
 \mathcal J_t:=tP_7+P_9
\]
has exactly the same carrier, block projectors, ranks, ordinary trace
convention, and block symmetries as \(\mathcal J_{1,1}\), while
\[
 \frac{\|\mathcal J_tP_7\|_{\rm HS}^2}
      {\|\mathcal J_tP_9\|_{\rm HS}^2}
 =\frac79t^2.
\]
Thus no construction invariant under independent scalar rescaling of the two
irreducible response blocks can select \(7/9\).  The exact finite Gram weights
are derived; the equality of response amplitudes must descend from a common
pre-split owner such as \((\mathrm{EQAMP}_{79})\).  This diagnostic is the
countermodel behind
Theorem~\ref{p7:thm:physical-seven-nine-owner-reduction}.

\subsection{Explicit family reality and complete ordered coefficient tests}
\begin{theorem}[Native family reality map and odd normal directions]
\label{p5c:thm:v2217-family-reality-ordered-tests}
On the existing rooted coefficient chart let
\[
 \begin{gathered}
 M=\tfrac12\operatorname{diag}(-1+i,1+i,-1-i),\qquad
 g_e=\tfrac12(-1-i,i,1)^T,\\
 g_\mu=\tfrac1{2\sqrt6}(2,-1+3i,-1-3i)^T,\qquad
 g_\tau=\tfrac1{\sqrt3}(1,1,1)^T,
 \end{gathered}
\]
and $P_\alpha=g_\alpha g_\alpha^*$.
Writing $S_{DI}$ for the permutation of the last two coordinates,
\[
 X=e^{-i\pi/4}S_{DI},\qquad
 X\bar X=I,\quad X\bar M X^T=M,\quad X\bar P_\alpha X^*=P_\alpha.
\]
Among unitary matrices satisfying both pair covariance and preservation
of every labelled projector, $X$ is unique up to sign. For Hermitian $h$,
the odd part $h_-=(h-X\bar hX^*)/2$ is precisely
\[
 h_-=\begin{pmatrix}
 0&\alpha+i\beta&-\alpha+i\beta\\
 \alpha-i\beta&\delta&0\\-\alpha-i\beta&0&-\delta
 \end{pmatrix},\qquad \alpha,\beta,\delta\in\mathbb R.
\]
In particular a diagonal neutral normal return has odd part
$(n_D-n_I)(N_D-N_I)/2$.
\end{theorem}
\begin{proof}
Directly $S_{DI}\bar g_e=-ig_e$ and
$S_{DI}\bar g_\mu=g_\mu$, $S_{DI}\bar g_\tau=g_\tau$.
Also $S_{DI}\bar M S_{DI}=iM$, whose factor $i$ is canceled by
the two phases in $X$. For uniqueness let $U=(g_e,g_\mu,g_\tau)$.
Preserving each ray forces any candidate to be $UDU^T$ with diagonal
unitary $D=\operatorname{diag}(d_e,d_\mu,d_\tau)$. The nonzero entries of
$U^*M\bar U$ are
\[
 B=\begin{pmatrix}
 0&-i/\sqrt6&-i/\sqrt3\\
 -i/\sqrt6&(-1+i)/3&\sqrt2(1-i)/6\\
 -i/\sqrt3&\sqrt2(1-i)/6&(-1+i)/6
 \end{pmatrix}.
\]
The equation $D\bar BD=B$ gives $d_ed_\mu=d_ed_\tau=-1$ and
$d_\mu^2=d_\mu d_\tau=d_\tau^2=-i$. Thus
$D=\pm\operatorname{diag}(e^{-3i\pi/4},e^{-i\pi/4},e^{-i\pi/4})$,
equivalent to $\pm X$. Entrywise subtraction in an arbitrary Hermitian
$h$ gives the three-real-parameter matrix above.
\end{proof}

\begin{proposition}[Forward packet matching, Schur reality and state tests]
\label{p5c:prop:v2217-complete-family-packet}
For real $\eta$ put
$\mathbb H[h,\eta M]=\left(\begin{smallmatrix}h&\eta M\\
\eta M^*&-h^T\end{smallmatrix}\right)$ and
$\mathbb X=\operatorname{diag}(X,\bar X)$.
The coefficient-conjugate forward Hamiltonian
$H_c(t)=\overline{\mathbb H[h(t),\eta(t)M]}$ obeys
\[
 \mathbb X H_c\mathbb X^*
       =\mathbb H[h^\sharp,\eta M],\qquad h^\sharp=X\bar hX^*.
\]
Consequently complete even normals $h=h^\sharp$ give equal inclusive
production in the two forward packets at every strength, provided the
inherited states and readout are transported by this same unitary.
For actual retained/rest Hermitian blocks and coupling $B$ satisfying
$h_R=X_R\bar h_RX_R^*$, $h_A=X_A\bar h_AX_A^*$,
$B=X_R\bar B X_A^*$, their Schur block obeys
\[
 K_R(z)=X_R\overline{K_R(\bar z)}X_R^*,\qquad
 K_R^-(z)=\tfrac12[K_R(z)-X_R\overline{K_R(\bar z)}X_R^*]=0.
\]
For any actual compatible family of initial densities $\mathfrak D$, a
complete evolved readout $A_Q=U^*QU$ is independent of the unselected
state precisely when
$\operatorname{Tr}[(\rho-\sigma)A_Q]=0$ for every
$\rho,\sigma\in\mathfrak D$.
\end{proposition}
\begin{proof}
The first identity follows by block multiplication using the theorem.
It is a time-independent unitary intertwiner between two independently
defined forward Hamiltonians; uniqueness of Schr\"odinger evolution gives
the same intertwiner for propagators. Transporting states and readouts
then gives equal traces. It does not conjugate a forward propagator and
pretend its anti-linear time reversal is forward evolution.
For $K_R(z)=zI-h_R-B(zI-h_A)^{-1}B^*$, conjugate at $\bar z$ and insert
the three block identities; all $X_A$ factors cancel. Finally subtracting
the two state expectations proves the annihilator criterion.
\end{proof}
These are conditional packet statements, not a complete physical CP/PCT
map. The spin, gauge, external, constraint and state factors must also be
transported. QROOT on its stated perturbation modes is not a new empty
ordinary-fermion preparation. The odd Schur test must retain the conjugate
spectral argument and cannot turn a frequency-dependent return into an
autonomous Hamiltonian without its actual time-domain reduction.

\begin{proposition}[Three ordered traces reconstruct the entire odd normal]
\label{p5c:prop:v2217-ordered-trace-reconstruction}
Choose the common-real frame of the preceding coefficient packet. In it
the three labelled projectors are the coordinate projectors and the pair is
\[
 M_0=\begin{pmatrix}
 0&-1/\sqrt6&-1/\sqrt3\\
 -1/\sqrt6&\sqrt2/3&-1/3\\
 -1/\sqrt3&-1/3&\sqrt2/6
 \end{pmatrix},\qquad M_0M_0^*=I/2.
\]
For $D=S+iK$, $S$ real symmetric and
$K=\left(\begin{smallmatrix}0&u&v\\-u&0&w\\-v&-w&0\end{smallmatrix}\right)$,
define $L_AN=AN+NA^T$ and
$t(A,B,C)=\operatorname{Im}\operatorname{Tr}(M_0^*L_AL_BL_CM_0)$.
Then
\begin{align*}
 F=t(D,P_\mu,P_\mu)&=(-4\sqrt3u+4\sqrt2w)/9,\\
 H=t(D,P_\tau,P_\tau)&=(-2\sqrt6v-2\sqrt2w)/9,\\
 G=t(P_e,D,P_\mu)&=(-2\sqrt3u+\sqrt6v-3\sqrt2w)/9.
\end{align*}
The determinant of this real map is $-32/81$, and
\[
 \Xi:=2(u^2+v^2+w^2)
 =\frac94(F+H/2)^2+\frac{81}{16}H^2+\frac92G^2.
\]
All 27 ordered entries with one $D$ and two labelled projectors are
\[
 t(D,P_r,P_s)=E_{rs},\quad t(P_r,D,P_s)=GJ_{rs},\quad
 t(P_r,P_s,D)=-E_{rs},
\]
where
\[
 E=\begin{pmatrix}0&(H-F)/2&(F-H)/2\\
 (H-F)/2&F&-(F+H)/2\\(F-H)/2&-(F+H)/2&H\end{pmatrix},\quad
 J=\begin{pmatrix}0&1&-1\\-1&0&1\\1&-1&0\end{pmatrix}.
\]
\end{proposition}
\begin{proof}
To obtain the stated real frame from the charged frame above, take
$D_0=\operatorname{diag}(e^{-3i\pi/4},e^{-i\pi/4},e^{-i\pi/4})$,
choose $\theta_e=\pi/8$, $\theta_\mu=\theta_\tau=3\pi/8$ and put
\[
 V=U\operatorname{diag}(e^{i\theta_e},e^{i\theta_\mu},e^{i\theta_\tau}).
\]
Then $V^*M\bar V=M_0$ and the rays stay coordinate projectors; the
anti-linear reality becomes scalar conjugation up to one common sign.
For each coordinate projector,
$(L_{P_r}N)_{ab}=(\delta_{ar}+\delta_{br})N_{ab}$.
Substitute the displayed $M_0$ and $D=S+iK$ in the three ordered products.
The real $S$ contributes no imaginary part; the three coefficients of
$(u,v,w)$ are exactly the displayed rows. Their determinant is $-32/81$.
Inverting those three linear rows and summing $2(u^2+v^2+w^2)$ gives
the displayed norm identity. Applying the same entry formula with both
coordinate selectors at the three possible positions gives $E,GJ,-E$,
which specifies every ordered entry, including the zero diagonal of $J$.
\end{proof}
In an ordered history a fixed flavor selector vector $a(t)=f(t)a_0$
annihilates the middle coefficient, since $a(t)^TJa(s)=0$.
Independent relative selector histories are required to test this third
direction at that order. These exact coefficient tests select neither a
physical source arrow, envelopes, CP sign nor a baryon-to-photon ratio.

\begin{corollary}[Exact geometric row arrows and their occurrence boundary]
\label{p5c:cor:v2217-geometric-row-arrows}
Put $c=g_\tau$, $Q_0=cc^*$, $Q_\perp=I-Q_0$ and
$R_-=\operatorname{diag}(1,i,-i)$. Then
$c^*R_-c=1/3$ and $s:=\|Q_\perp R_-c\|=\sqrt{8/9}$.
The coefficient-fixed geometric arrows
\[
 \beta_\mu=s^{-1}Q_\perp R_-Q_0=g_\mu c^*,\quad
 \beta_e=P_eR_-Q_\perp=g_eh_e^*,\quad
 \beta_\tau=s^{-1}Q_0R_-Q_\perp=ch_\tau^*,
\]
with $h_e=R_-^*g_e$, $h_\tau=\bar g_\mu$, satisfy
$\beta_r\beta_s^*=\delta_{rs}P_r$. Hence the comparison
$R_0=\sum_r\epsilon_r\beta_r$, for real $\epsilon_r$, has
$R_0R_0^*=\sum_r\epsilon_r^2P_r$ and $\Xi=F=H=G=0$.
\end{corollary}
\begin{proof}
Substitution gives
$s^{-1}Q_\perp R_-c=(2,-1+3i,-1-3i)^T/(2\sqrt6)=g_\mu$.
Taking the adjoint row gives the displayed $h_\tau$; the $P_e$ product
gives $h_e$, with $Q_\perp h_e=h_e$. The three right vectors
$h_e,c,h_\tau$ are orthonormal by direct multiplication of the listed
rays. Their left vectors $g_e,g_\mu,c$ are also orthonormal. The return
identity and the diagonal comparison follow. In the common-real frame
that diagonal has zero $K$, so all four tests vanish.
\end{proof}
The actual source extraction $\iota_r$ has not been identified with these
geometric arrows by this calculation. That equality, its action weights
and inherited source state remain independent occurrence obligations.

\subsection{Mixed family table and chronological pair comparisons}
\label{p5c:app:v2167-cpbl2-comparisons}
This appendix retains coordinate diagnostics and explicitly conditional
finite histories. It does not select the native K16 normal generator,
state, pulse placement, or cosmological yield.

For the generated labelled packet the imaginary coefficients of
Theorem~\ref{p5c:thm:v2167-mixed-family-coefficient} are
\[
 \begin{array}{c|rrr}
 a&\mathrm{inv,dir}&\mathrm{inv,Gram}&\mathrm{dir,Gram}\\\hline
 e&0&0&0\\
 \mu&1/12&1/36&-1/36\\
 \tau&0&-1/18&1/18
 \end{array}.
\]
Each entry is the exact contraction of the fixed projectors and $M$,
before physical spin/mode and action return factors.

\begin{proposition}[Conditional chronological family kernel]
\label{p5c:prop:v2167-chronological-family-kernel}
On a declared common canonical family port, put
\[
 U_E=\sum_{a=e,\mu,\tau}y_aP_a,\qquad
 U_N=\operatorname{diag}(z_{\rm Gram},z_{\rm dir},z_{\rm inv}),\quad
 K_M(V)=\operatorname{Tr}(M^\dagger V M V^{\mathsf T}).
\]
For unit-modulus phases the exact difference is
\[
 \begin{aligned}
 \delta K&=K_M(U_EU_N)-K_M(U_NU_E)\\
 &=\frac i9(z_{\rm inv}^2-z_{\rm dir}^2)
 (y_\tau-y_\mu)(y_\tau+2y_\mu-3y_e).
 \end{aligned}
\]
\end{proposition}
\begin{proof}
Expand both traces with the diagonal generated $M$ and $U_N$, and insert
$U_E=\sum_a y_aP_a$. Diagonal terms cancel. Collecting the off-diagonal
entries with the fixed charged projectors factors the remaining polynomial
as displayed. The companion rational-complex calculation independently
checks this identity before any phase or history is assigned.
\end{proof}
Equal squared inverse/direct phases or the displayed charged factors
cancel this particular return. Its physical use requires those normal
flows to be generated on the pair's same canonical carrier. Terminal
energy values alone do not license such flows.

On the declared spin-singlet finite CAR comparison, balanced conjugate
vacua and two normalized pair pulses with angles $\theta_1,\theta_2$
give
\[
 Q_{\rm comp}=-2\sin(2\theta_1)\sin(2\theta_2)
                 \operatorname{Re}\delta K.
\]
The conjugate sector uses conjugate coefficients at the same chronological
durations. These assigned pulses and product vacua are comparison inputs,
not the inherited native cyclic state. The reproducible companion checks
also cover a continuous pair history on a non-Gaussian state by independent
Fock and covariance calculations. No comparison number is exported as
$Y_{B-L}$, a native efficiency, or a universal perturbative onset order.

% END INLINE SOURCE: Paper5C-CPBL2Comparisons-Appendix-v1_0.tex




% BEGIN INLINE SOURCE: Paper5C-WeakMemoryDiagnostics-Appendix-v1_0.tex
% Framework v2.168: integrated finite charged-source derivations.
\subsection{Additional diagonal weak constituent polynomials}
\label{p5c:app:v2169-occupation-polynomials}
These calculations require the stated diagonal one-body blocks and actual
occupation cap in addition to the native vertex. They evaluate those complete
constituents and do not assert that the physical full action has this form.\par
\subsubsection{The native six-family characteristic and tagged numerator}
For the $N=1$ constituent, additionally suppose its actual $A,B$ are diagonal
in the already occurring current and neutral-role presentations:
$A=\diag(a_e,a_\mu,a_\tau)$ and
$B+\omega I=\diag(b_1,b_2,b_3)$. This is a dynamical diagonal-block
hypothesis, not a new choice of mass frame. Put $u=\pi^2\kappa$ and
\[
 \mathcal P=(|C_{rj}|^2)=
 \begin{pmatrix}1/4&1/4&1/2\\5/12&5/12&1/6\\1/3&1/3&1/3\end{pmatrix}.
\]
Writing products over the indicated complementary indices, the whole
six-dimensional characteristic polynomial is explicitly
\begin{equation}
 \boxed{\begin{aligned}
 D_6(z)={}&\prod_r(z-a_r)\prod_j(z-b_j)\\
 &-u\sum_{rj}\mathcal P_{rj}
             \prod_{s\ne r}(z-a_s)\prod_{k\ne j}(z-b_k)\\
 &+u^2\sum_{rj}\mathcal P_{rj}(z-a_r)(z-b_j)-u^3 .
 \end{aligned}} \label{p5c:eq:v2169-b13-o6}
\end{equation}
For physical tag $r$, its diagonal cofactor is
\begin{equation}
 \boxed{\begin{aligned}
 N_{5,r}(z)={}&\prod_{s\ne r}(z-a_s)\prod_j(z-b_j)\\
 &-u\sum_{\substack{s\ne r\\j}}\mathcal P_{sj}
             \prod_{t\ne r,s}(z-a_t)\prod_{k\ne j}(z-b_k)\\
 &+u^2\sum_j\mathcal P_{rj}(z-b_j).
 \end{aligned}} \label{p5c:eq:v2169-b13-o7}
\end{equation}

\begin{theorem}[Evaluated physical tagged polynomial on this cap]
\label{p5c:v2169:b13:thm:nativepolynomial}
For the stated constituent,
\begin{equation}
 \boxed{G_{rr}(z)=N_{5,r}(z)/D_6(z),\qquad
 p_r(z)=D_6(z)/\gcd(D_6(z),N_{5,r}(z)),} \label{p5c:eq:v2169-b13-o8}
\end{equation}
where the gcd is monic. The quotient is monic, has simple real roots,
and includes physical visibility on the empty support. In particular the
terminal claim uses this reduced polynomial rather than $D_6$ alone.
\end{theorem}
\begin{proof}
Expand the determinant of the two diagonal blocks and their off-diagonal
$\pi\chi C$ word by equal-size selected row and column subsets. A $k$-pair
term is $(-u)^k|\det C_{I,J}|^2$ times the two remaining diagonal products.
For $k=1$ it gives $\mathcal P_{rj}$. For $k=2$, the complementary-minor
identity of a $3$-by-$3$ unitary gives
$|\det C_{I,J}|^2=|C_{I^c,J^c}|^2$, and $|\det C|=1$ gives $k=3$.
This proves \eqref{p5c:eq:v2169-b13-o6}. Delete charged row and column $r$ and repeat the same
expansion; its size-two minor has missing row $r$, giving \eqref{p5c:eq:v2169-b13-o7}.
Cramer's rule proves the ratio. The Hermitian spectral expansion has
positive tagged residues and only first-order poles, so polynomial
cancellation removes exactly the invisible roots and excess degeneracy.
This proves \eqref{p5c:eq:v2169-b13-o8}, including its square-free property.
\end{proof}

This exact determinant calculation consumes the coherent unitary before its
scalar readout. It explains why on \emph{this} diagonal two-block chart the
tagged pole functions can be expressed in the derived $\mathcal P$.
It does not replace the full coherent matrix source response in other
actions, nor does it infer diagonal return from $\mathcal P$'s rank.
\subsubsection{All occupation returns on a scalar receiver branch}
This additional diagonal branch assumes $B=\xi I_3$ on the whole constituent.
Transport the neutral presentation by the existing $C$; each charged role
then has a scalar Jacobi ladder with
\[
 d_{2k}=a_r+2k\omega,\qquad d_{2k+1}=\xi+(2k+1)\omega,\qquad
 w_n^2=u(\lfloor n/2\rfloor+1).
\]
Its continuants $p_0=1$, $p_1=z-a_r$ obey
\begin{equation}
 \boxed{p_{n+1}=(z-d_n)p_n-w_{n-1}^2p_{n-1}.} \label{p5c:eq:v2169-b13-o9}
\end{equation}
All links are nonzero; an eigenvector with zero first component would
successively have every component zero. Hence the complete $p_{N+1}$ is
the actual tagged polynomial on this branch. The first two evaluated caps are
\[
 p_2=(z-a_r)(z-\xi-\omega)-u,
\]
\[
 p_3=(z-a_r-2\omega)
       [(z-a_r)(z-\xi-\omega)-u]-u(z-a_r).
\]
Thus $p_2(a_r)=-u$ and $p_3(a_r)=2\omega u$.
At $N=1$ the two visible energies are
\[
 \frac{a_r+\xi+\omega}{2}
 \ \pm\sqrt{\frac{(a_r-\xi-\omega)^2}{4}+u}.
\]
These first two cap outputs demonstrate genuine shifts on this additional
chart; higher caps can have exact endpoint interference. They do not refute a different complete-action terminal claim,
which can have a different actual free exporter and all other marks.


% BEGIN FRAMEWORK v2.239: FOUR/FIVE-LEVEL JACOBI ENDPOINT
\begin{proposition}[Higher weak Jacobi endpoints and an exact interacting return]
\label{p5c:prop:v239-jacobi-endpoint-interference}
On the \emph{scalar receiver, finite occupation-cutoff constituent}
with the same $d_n,w_n^2$ as in \eqref{p5c:eq:v2169-b13-o9}, write
$e=a_r$, $u>0$, and $\omega>0$. Continuing that exact recurrence yields
\begin{align*}
 p_4(e)&=2u\,[u+\omega(e-\xi-3\omega)],\\
 p_5(e)&=-4\omega u\,[3u+2\omega(e-\xi-3\omega)].
\end{align*}
For the four-level cap an unshifted but source-visible endpoint occurs
exactly when
\begin{equation}\label{p5c:eq:v239-jacobi-four-level-return}
 \boxed{u=\omega(\xi+3\omega-e)>0.}
\end{equation}
For example, $(e,\xi,\omega,u)=(1,3,1,5)$ gives
\[
 J=\begin{pmatrix}
 1&\sqrt5&0&0\\
 \sqrt5&4&\sqrt5&0\\
 0&\sqrt5&3&\sqrt{10}\\
 0&0&\sqrt{10}&6
 \end{pmatrix},\qquad
 J\begin{pmatrix}1\\0\\-1\\\sqrt{2/5}\end{pmatrix}
 =\begin{pmatrix}1\\0\\-1\\\sqrt{2/5}\end{pmatrix}.
\]
Every link remains strictly positive. The first occupation basis
vector is cyclic and sees the eigenvalue $e=1$ with exact residue
$1/(1+1+2/5)=5/12$, not the free two-pole residue $1/2$.
Thus the nonzero weak interaction does not exclude exact bare-energy
return for \emph{all} finite cutoffs; the source-visible pole equality
and its residue are separate observables. Neither this mathematical
parameter example nor the isotropic receiver hypothesis selects the
complete compiled physical action.
\end{proposition}
\begin{proof}
Apply the three-term continuant recurrence at $z=e$ two further
times, beginning with $p_2(e)=-u$ and $p_3(e)=2\omega u$.
At four levels positivity of all off-diagonal links makes the first
site cyclic, so no eigenvalue can be invisible to it. Direct matrix
multiplication verifies the displayed eigenvector; its squared norm
is $12/5$, giving source residue $5/12$.
\end{proof}
% BEGIN FRAMEWORK v2.245: UNIFORM VERSUS ALTERNATING JACOBI AUDIT
\begin{proposition}[Uniform-ladder endpoint roots are not the native alternating returns]
\label{p5c:prop:v245-uniform-jacobi-four-level-audit}
For the \emph{different}, uniformly spaced Jacobi diagnostic
$\widetilde d_k=\xi+k\omega$ and squared links
$\widetilde w_{k-1}^2=ku$ ($k=1,2,3$), define
$\widetilde p_0=1$, $\widetilde p_1=z-\xi$ and
\[
 \widetilde p_{k+1}=(z-\xi-k\omega)\widetilde p_k
 -ku\widetilde p_{k-1}.
\]
At $z=\xi$ its four-level endpoint polynomial is
$\widetilde p_4(\xi)=3u(u-2\omega^2)$.
At the \emph{different} specified energy $z=1$ with
$(\xi,\omega)=(3,1)$, however,
\[
 \boxed{\widetilde p_4(1)=3u^2-58u+120,
 \qquad u_\pm=(29\pm\sqrt{481})/3>0.}
\]
For either root the first site is cyclic and the tagged pole residue is
\[
 \boxed{\widetilde Z_{1}(u_\pm)
 =\frac{19u_\pm-60}{2(19u_\pm-77)}\in(0,1).}
\]
These roots are not the native alternating charged/neutral ladder
of Proposition~\ref{p5c:prop:v239-jacobi-endpoint-interference},
which at $(e,\xi,\omega)=(1,3,1)$ instead gives $u=5$
and residue $5/12$. None of these comparison roots is inferred to
occur in the full compiled physical Hamiltonian.
\end{proposition}
\begin{proof}
Apply the displayed three-term recurrence through $k=3$;
substitution yields both printed polynomials and the quadratic roots.
For $u>0$ all links are nonzero, so the first site is cyclic.
The diagonal resolvent is $\widetilde q_3(z)/\widetilde p_4(z)$,
where $\widetilde q_3$ is the trailing $3\times3$ minor.
Its simple-pole residue at $z=1$ is
$\widetilde q_3(1)/\partial_z\widetilde p_4(1)$,
which simplifies on $\widetilde p_4(1)=0$ to the stated ratio.
The native alternating result follows from its own recurrence
\eqref{p5c:eq:v2169-b13-o9}; the diagonals and links differ.
\end{proof}
% END FRAMEWORK v2.245: UNIFORM VERSUS ALTERNATING JACOBI AUDIT

% END FRAMEWORK v2.239: FOUR/FIVE-LEVEL JACOBI ENDPOINT

% END INLINE SOURCE: Paper5C-WeakOccupationDiagnostics-Appendix-v1_0.tex



% v2.189: spectral cancellation and coefficient/state diagnostics.

% BEGIN INLINE SOURCE: Paper5C-ChargeChronologyDiagnostics-Appendix-v1_0.tex
% Framework v2.189: extended diagnostics supporting the positive inherited-charge chain.
\subsection{Complete spectral cancellation and the external operator guard}
\label{p5c:v2188:app:chargediagnostics}
These finite reductions test the exact domains of the main proof. Their
diagnostic energies and profiles do not choose a physical cosmological
action, cutoff, state or clock.

\subsubsection{A complete paired Gibbs state cannot be reduced to its normal
occupation matrix}
On an actual two-spin restriction write
\[
 H_c=\epsilon_0n_0+\epsilon_1n_1+v n_0n_1,\quad
 P=g a_0^\dagger a_1^\dagger,\quad H=H_c+P+P^\dagger,\quad Q=-n_0-n_1 .
\]
The quartic contribution is retained in \(E=\epsilon_0+\epsilon_1+v\).
For \(\Omega=\sqrt{E^2/4+|g|^2}\), \(A=e^{-\beta E/2}\), the exact
partition function and expectations are
\[
 Z=e^{-\beta\epsilon_0}+e^{-\beta\epsilon_1}
                  +2A\cosh(\beta\Omega),\quad
 \langle P\rangle=-\frac{|g|^2A\sinh(\beta\Omega)}{\Omega Z},\quad
 \langle N\rangle-1=-\frac{EA\sinh(\beta\Omega)}{\Omega Z}.
\]
Thus the pair-only second jet and its complete normal/pair companion are
\begin{equation}
 q''_{\rm pair}=-\frac{4E|g|^2A\sinh(\beta\Omega)}
                          {\hbar^2\Omega Z},\qquad
 q''_{\rm normal,pair}=+\frac{4E|g|^2A\sinh(\beta\Omega)}
                          {\hbar^2\Omega Z}.
 \label{p5c:v2188:eq:gibbscancellation}
\end{equation}
They cancel exactly. The full paired Gibbs state generally fails to
commute with \(Q\); applying the charge-invariant initial-state formula
to this state drops the cancelling term.
By contrast, for the generated Gibbs state of the actual prepair \(H_c\),
\[
 q''(t_i)=\frac{4|g|^2(e^{-\beta E}-1)}
 {\hbar^2(1+e^{-\beta\epsilon_0}
            +e^{-\beta\epsilon_1}+e^{-\beta E})}.
\]
Its sign is determined by the actual pair transition energy \(E\), not
by an assumed empty vacuum.

\subsubsection{Half occupation with a quantum external coefficient}
For \(P=gX a_0^\dagger a_1^\dagger\) with \(X\) an operator on the
consumed external carrier, the exact pair commutator is
\[
 [P,P^\dagger]=|g|^2\{
 XX^\dagger n_0n_1-X^\dagger X(1-n_0)(1-n_1)\}.
\]
On the expressly product restriction with right reduced state \(I_4/4\),
\begin{equation}
 q''(t_i)=\frac{|g|^2}{\hbar^2}
                    \operatorname{Tr}\sigma[X,X^\dagger].
 \label{p5c:v2188:eq:operatorhalfguard}
\end{equation}
Hence \(C_R=I/2\) alone does not force a zero jet for an operator-valued
pair coefficient. Correlated inherited states require the preceding
full joint expectation rather than even this product formula.
The licensed free coordinate unit in
\ref{p4:eq:v2169-b11-l11}--\ref{p4:eq:v2169-b11-l12} has equal
displayed \(\kappa_\pm\); an ordinary oscillator annihilation vacuum
must not replace it. The exact local density receiver
\ref{p5c:v2171:eq:centeredpair} retains \(K_X=[X,X^\dagger]\) and
its actual marked word. Higher ordered or mixed correlations are not
eliminated by equality of these two scalar second moments.

\subsubsection{The actual rooted CP packet}
An actual comparison must reverse the complete rooted source, canonical
embeddings, metric and connection, every source/contact word, inherited
state, signed charge and readout surface. Complex conjugation of the
coefficient matrix at the same time coordinate is not itself a
chronological CP transport: it also reverses the sign of the Schr\"odinger
factor \(i\). The formal charge phase orbit is not substituted for that
physical map. The static real-amplitude theorem gives zero on each
matched chiral spectral packet satisfying its antecedents. It does not
prove a universal cancellation after arbitrary charged-current
insertions or an unverified state transformation.

\subsubsection{The remaining execution belongs to the inherited owner}
Charge-invariant upstream support is transported exactly by a
charge-intertwining parent propagator and equivariant restriction.
Neutral mean alone is weaker than charge invariance, and charge
invariance is weaker than support in a specified charge eigenspace.
The owned nonzero BL2 pair on the K8--K16 boundary does not commute with
the pair charge. Consequently the existing K16 inheritance identity
does not establish charge invariance of the complete transported
ordinary state. This is evaluation of the existing full parent state
and action word, not an independent new state selector.


\subsubsection{Relative source variation and common-family transport}
\label{p5c:v2189:app:relativecharge}
The completed Hermitian same-current value in
Theorem~\ref{p5c:thm:next187-completed-local-pair-current} is real. That
reality does not eliminate a real mixed derivative along an actual relative
source direction. A common transported family moves the normal source,
pair, embeddings and all contacts together; its contraction is covariant.
Holding the pair fixed while rotating a normal source is a different
relative direction. In particular the complete source return
\(\Sigma=\mathsf B G\mathsf B^\dagger\) differentiates both coupling
factors as well as the inverse. Dropping either coupling jet need not equal
the actual common-family derivative. The positive main theorem retains
these terms; this comparison does not install a phase parameter as the
physical discrete CP reversal.

For the formal pair phase orbit, a transported anomalous state and a
fixed anomalous state are distinct packets. The former has the exact
charge equality of
Theorem~\ref{p5c:thm:next187-matched-charge-orbit}; the latter may change
its relative expectation. No formal orbit, particle-hole automorphism or
scalar imaginary-part test replaces the actual complete chronological
CP source/state/charge map.

% END INLINE SOURCE: Paper5C-ChargeChronologyDiagnostics-Appendix-v1_0.tex



% BEGIN CLOSED-DS-HOPF-PHOTON-TRANSFER v2.90 / framework v2.194
\subsubsection{Conditional global de-Sitter--Hopf photon transfer}
\label{p5c:sec:v290-ds-hopf-photon-transfer}
\begin{proposition}[Conformal rigidity of the native free photon modes]
\label{p5c:prop:v290-conformal-photon-rigidity}
On a smooth round closed FLRW metric
$g=a(\eta)^2(-d\eta^2+d\Omega_3^2)$ in geometric units, assume the
unbroken Maxwell action with constant common gauge normalization, no
nonconformal photon correction and vanishing background electromagnetic
field.  On the physical coexact transverse subspace, the modes inherited
from Theorem~\ref{p7v2:thm:v244-transverse-ew-wave-operators} satisfy
\[
 \boxed{A_n''+n^2A_n=0,\qquad n=L+1\ge2.}
\]
A smooth homogeneous scale factor alone does not mix the free positive-
and negative-conformal-frequency branches; its isolated Bogoliubov
coefficient is $\beta_n=0$ on a regular common conformal chart.
\end{proposition}
\begin{proof}
In four spacetime dimensions the factors $a^4$ in the volume element
and $a^{-2}a^{-2}$ in $F_{\mu\nu}F^{\mu\nu}$ cancel exactly.  Hence
Maxwell's quadratic action is that of the Einstein static cylinder.
The coexact one-form Laplacian on the unit round $S^3$ has eigenvalues
$n^2$, which gives the stated oscillator equation and constant
conformal-frequency mode basis.
\end{proof}
\begin{corollary}[Antipodal transfer across the ideal terminal closed-de-Sitter history]
\label{p5c:cor:v290-terminal-de-sitter-antipode}
On the \emph{idealized} terminal Einstein vacuum with
$\Lambda_{\rm eff}=C_{56}/2$, set
$H_\Lambda=c_{\rm op}\sqrt{C_{56}/6}$ and define dimensionless global
conformal time $\tau\in(-\pi/2,\pi/2)$.  Its metric, in length units,
is
\[
 ds^2=\frac{6}{C_{56}\cos^2\tau}
       (-d\tau^2+d\Omega_3^2),
\]
and the total conformal interval is $\Delta\tau=\pi$.  If the free
coexact Maxwell data admit smooth conformal-boundary limits, the antipodal
pullback $\mathcal P^*$, $\mathcal P(x)=-x$, satisfies
\[
 \boxed{
 A_T(+\pi/2,\cdot)=\mathcal P^*A_T(-\pi/2,\cdot),\qquad
 \partial_\tau A_T(+\pi/2,\cdot)
 =\mathcal P^*\partial_\tau A_T(-\pi/2,\cdot).}
\]
Thus on that free transverse phase space the global conformal transfer
is $\boxed{\mathcal U_\gamma=\mathcal P^*}$.
\end{corollary}
\begin{proof}
Each mode acquires $e^{\pm in\pi}=(-1)^n$ across the full conformal
interval.  The already established coexact antipodal parity is
$\mathcal P^*Y_n=(-1)^nY_n$ (the photon spectral corollaries above).
The same factor multiplies the $\tau$ derivative.  Summing smooth mode
data gives both displayed identities.
\end{proof}
\begin{remark}[Transfer is not a slowed local photon cone]
This corollary composes the \emph{existing} $S^3$ photon spectrum with
the \emph{conditional} exact global de-Sitter terminal solution.  It does
not assert that the realized cosmic history equals global de Sitter from
past to future infinity, that a prepared photon traverses both conformal
boundaries, or that physical local $c_\gamma$ differs from $c_{\rm op}$.
Matter, active scalaron dynamics, nonconformal gauge terms, and boundary
source history can modify the transfer.
\end{remark}
% END CLOSED-DS-HOPF-PHOTON-TRANSFER v2.90



% BEGIN GEOMETRY-FIRST CONSEQUENCES v2.201
\setlength{\emergencystretch}{2em}

\section{Thermal comparisons, rare-process and numerical diagnostics}
\label{p5d:sec:G}
\noindent\textit{Technical owner.} These exact inherited calculations
support the corresponding statements in Papers~V-A--V-C;
the appropriate action/observable occurrence conditions remain in those papers.
\subsection{Rare-process consequences of the imported spectrum}
\label{p7v2:sec:rare-processes}

\begin{theorem}[Decay support from the imported mass and charge packet]
\label{p7v2:thm:decay-support}
For a proposed decay \(A\to f_1+\cdots+f_n\) with \(n\geq2\), the imported exact charges and
representation labels determine whether every unbroken additive charge can be
conserved.  The imported masses determine whether the on-shell phase space is
nonempty:
\begin{equation}
\label{p7v2:eq:decay-support}
 m_A>\sum_{k=1}^{n}m_{f_k}
\end{equation}
for nonzero interior phase-space volume.  These data determine kinematic
support, not a partial width or branching fraction.
\end{theorem}

\begin{proof}
An unbroken additive charge is conserved by every invariant interaction, so a
charge mismatch excludes the channel.  In the rest frame of \(A\), energy
conservation and the mass-shell inequalities give
\(m_A=\sum_kE_k\geq\sum_km_{f_k}\); strict inequality permits a phase-space
interior.  Conversely, when the charges can be conserved and the strict mass
inequality holds, momenta can be chosen continuously near threshold to give
nonempty kinematic phase space.  A rate additionally contains an interaction
matrix element, polarization sums, radiative corrections, and the
phase-space integral.
\end{proof}

\begin{proposition}[Claim-scoped decay and branching-ratio boundary]
\label{p7v2:prop:claim-scoped-decay-boundary}
Finite interaction vertices, local current correlators, amputated loop kernels,
the kinematic support in Theorem~\ref{p7v2:thm:decay-support}, and the algebraic
regulator of Theorem~\ref{p7v2:thm:resonant-bound} consume no \(S\)-matrix
hypothesis.  Interpreting such data as a physical decay amplitude, partial
width, or branching ratio additionally uses \textnormal{(SM1)--(SM3)} and the
phase-space/channel-normalization parts of \textnormal{(SM7)} in
Definition~\ref{p12:def:claim-indexed-smat-package}; it uses
\textnormal{(SM4)} when stable external legs are normalized through LSZ and
\textnormal{(SM5)} when the parent is defined as an unstable continued pole.
A massless or long-range channel also uses one compatible route from
Definition~\ref{p12:def:claim-indexed-infrared-package}.  A stable-parent
approximation must be declared as such and satisfies the corresponding
\textnormal{(SM2)} isolation premise only on that approximation domain.
\end{proposition}

\begin{proof}
The first group consists of finite matrices, local distributions, or algebraic
functions and has no large-time comparison map.  The second group names
outgoing particle channels and their probability per unit time; the wave
operator, optional LSZ/pole identification, phase-space normalization, and
infrared qualifications are exactly the clauses allocated by
Theorem~\ref{p12:thm:scattering-hypothesis-firewall}.
\end{proof}

\begin{proposition}[Minimal charged-lepton flavor violation in native current coordinates]
\label{p7v2:prop:minimal-lfv}
Assume the standard electroweak one-loop QFT comparison domain with only the
three generated light neutral roles in the charged-current loop, together
with the existing pole/rest bridge for their physical masses.  Then, on the
physical branching-ratio domain of
Proposition~\ref{p7v2:prop:claim-scoped-decay-boundary},
\begin{equation}
\label{p7v2:eq:minimal-muegamma}
 {\rm Br}(\mu\to e\gamma)
 =
 \frac{3\alpha_{\rm em}}{32\pi}
 \left|
   \sum_r (C_{e\nu})_{\mu r}^{*}(C_{e\nu})_{er}
   \frac{m_{\nu_r}^2}{m_W^2}
 \right|^2
 +O\!\left(\frac{m_\nu^6}{m_W^6}\right).
\end{equation}
In a conventional comparison section
$U_{\rm PMNS}^{\rm cmp}:=[C_{e\nu}]_{\rm std}$ this is the usual SM formula.
The result is a status-\(\mathsf Q\) downstream consequence, not a new finite
selector or a mass-matrix premise.
\end{proposition}

\begin{proof}
Unitarity of the generated current map cancels the role-independent loop term
for unequal charged-lepton flavors.  Expanding the standard one-loop kernel in
$m_{\nu_r}^2/m_W^2$ gives \eqref{p7v2:eq:minimal-muegamma}; the suppression is
controlled by the neutral role-mass splittings on the stated pole/rest bridge.
\end{proof}

\begin{proposition}[Conditional parent-vector proton decay]
\label{p7v2:prop:conditional-proton-decay}
A dimension-six estimate
\begin{equation}
\label{p7v2:eq:conditional-proton-decay}
 \Gamma(p\to e^+\pi^0)
 \sim
 C_{\rm had}\,\frac{g_X^4m_p^5}{M_X^4}
\end{equation}
is meaningful only after a nonminimal parent vector with mass \(M_X\),
coupling \(g_X\), operator normalization, flavor tensor, RG transport, and
hadronic matrix element \(C_{\rm had}\) are supplied, and only on the physical
rate domain of Proposition~\ref{p7v2:prop:claim-scoped-decay-boundary}.  It is
not part of the minimal finite Standard-Model core.
\end{proposition}

\begin{proof}
None of the quantities \(M_X,g_X,C_{\rm had}\) is fixed by the minimal
Standard-Model representation carrier.  Equation
\eqref{p7v2:eq:conditional-proton-decay} is therefore a conditional
effective-field-theory consequence of a separately specified parent
completion.
\end{proof}

\label{p7v2:part:heavy-thermal}

\subsection{Conditional external type-I seesaw comparison dictionary}
\label{p7v2:sec:heavy-seesaw}
\begin{quote}\textbf{v2.133 ontology scope.} The following type-I seesaw/Casas--Ibarra material is retained only as a conditional external Standard--Model comparison dictionary. It does not generate the framework's terminal neutrino masses or native lepton mixing.\end{quote}

\begin{proposition}[What the heavy affine pattern fixes]
\label{p7v2:prop:heavy-affine-status}
The imported pattern \eqref{p7v2:eq:heavy-affine-import} fixes the relative
coefficients and the exact daughter splitting
\begin{equation}
\label{p7v2:eq:heavy-splitting}
 M_{N_2}-M_{N_3}=\frac{\Delta_M}{28}.
\end{equation}
It does not determine the absolute anchor \(E_N\), the splitting scale
\(\Delta_M\), or a thermal initial abundance.  Positivity on the displayed
ordering requires, for \(\Delta_M>0\),
\begin{equation}
\label{p7v2:eq:heavy-positivity}
 E_N>\frac{57}{56}\Delta_M.
\end{equation}
\end{proposition}

\begin{proof}
Subtracting the last two entries of
\eqref{p7v2:eq:heavy-affine-import} gives
\eqref{p7v2:eq:heavy-splitting}.  The carrier and threshold number are
dimensionless and remain unchanged under
\((E_N,\Delta_M)\mapsto(cE_N,c\Delta_M)\), \(c>0\), so they cannot fix the
absolute scale.  The smallest displayed mass is \(M_{N_3}\); its positivity
gives \eqref{p7v2:eq:heavy-positivity}.
\end{proof}

\begin{definition}[High-energy neutrino realization fiber]
\label{p7v2:def:nu-selector-family}
Fix \(\Uroleabs+\Hvac\) and one compatible Paper--V-B member
\(b_6\in\Acompsix\).  Let \(D_\nu(b_6)\), \(D_N(b_6)\), and \(v(b_6)\) denote
its frozen light-neutral, heavy-spectrum, and scalar comparison coordinates.
Only inside this external type-I dictionary define
\[
 U_{\rm PMNS}^{\rm cmp}:=[C_{e\nu}]_{\rm std},\qquad
 m_\nu^{\rm cmp}(b_6)
 :=(U_{\rm PMNS}^{\rm cmp})^*D_\nu(b_6)
 (U_{\rm PMNS}^{\rm cmp})^\dagger .
\]
The Paper--V-C high-energy neutral lift over that base member is the fiber
\[
\boxed{
\mathfrak A_{\rm NU}(b_6)
=
\left\{
[Y_\nu]_{\rm pres}:
-\frac{v(b_6)^2}{2}
Y_\nu D_N(b_6)^{-1}Y_\nu^T
 =m_\nu^{\rm cmp}(b_6)
\right\}.
}
\]
Here \([Y_\nu]_{\rm pres}\) identifies only presentation-equivalent
representatives on the already frozen light/heavy family bases.  The total
conditional V-C family is the disjoint union of these fibers over frozen
compatible Paper--V-B base members.  Thus \(D_N\) belongs to the imported
Paper--V-B base;
low-energy data do not select a member of the high-energy fiber.
\end{definition}

\begin{theorem}[Type-I seesaw realization identity]
\label{p7v2:thm:seesaw-identity}
Let
\[
 D_\nu=\diag(m_1,m_2,m_3),\qquad
 D_N=\diag(M_{N_1},M_{N_2},M_{N_3}),
\]
with positive entries, and let \(R\in O(3,\CC)\), so \(RR^T=I\).  In the
convention
\begin{equation}
\label{p7v2:eq:casas-ibarra}
 \YN
 =
 \frac{i\sqrt2}{v}\,
 (U_{\rm PMNS}^{\rm cmp})^{*}D_\nu^{1/2}R D_N^{1/2},
\end{equation}
one has the exact identity
\begin{equation}
\label{p7v2:eq:seesaw-identity}
 -\frac{v^2}{2}\YN D_N^{-1}\YN^T
 =
 (U_{\rm PMNS}^{\rm cmp})^{*}D_\nu (U_{\rm PMNS}^{\rm cmp})^{\dagger}.
\end{equation}
\end{theorem}

\begin{proof}
Substitution of \eqref{p7v2:eq:casas-ibarra} cancels the heavy factors and
uses \(RR^T=I\):
\[
 \YN D_N^{-1}\YN^T
 =
 -\frac{2}{v^2}
 (U_{\rm PMNS}^{\rm cmp})^{*}D_\nu^{1/2}
 RR^T
 D_\nu^{1/2}(U_{\rm PMNS}^{\rm cmp})^{\dagger}.
\]
Multiplication by \(-v^2/2\) gives
\eqref{p7v2:eq:seesaw-identity}.  This is the standard realization identity
\cite{CasasIbarra2001}; it parameterizes, rather than removes, the high-energy
freedom in \(R\).
\end{proof}


\begin{corollary}[Casas--Ibarra coordinates on the presentation quotient]
\label{p7v2:cor:casas-ibarra-fiber-coordinate}
On the positive nonsingular branch of
Theorem~\ref{p7v2:thm:seesaw-identity}, with labeled pairwise-distinct
heavy Majorana masses and only residual kinetic-unitary heavy-basis
presentation changes, every representative $Y_\nu$ has inverse
\[
 \boxed{R(Y_\nu)=-\frac{iv}{\sqrt2}
 D_\nu^{-1/2}(U_{\rm PMNS}^{\rm cmp})^T Y_\nu D_N^{-1/2}
 \in O(3,\mathbb C).}
\]
The residual presentation group is
$\Gamma_N=\{\operatorname{diag}(\pm1,\pm1,\pm1)\}\cong(\mathbb Z_2)^3$.
It acts by $Y_\nu\mapsto Y_\nu S$ and $R\mapsto RS$.
Consequently the \emph{presentation-class} fiber is
\[
 \boxed{\mathfrak A_{\rm NU}(b_6)\cong O(3,\mathbb C)/\Gamma_N,}
\]
while a chosen representative $Y_\nu$ corresponds bijectively to a
single $R$. Zero or degenerate masses and other weak-basis quotients
require their own strata. No high-energy lift is physically selected by
this coordinate theorem.
\end{corollary}
\begin{proof}
The displayed inverse and $RR^T=I$ follow by inserting the seesaw
identity, as in Theorem~\ref{p7v2:thm:seesaw-identity}.
A unitary heavy-basis change preserving the diagonal Majorana mass
satisfies $S^T D_NS=D_N$. Combining this with unitarity and pairwise
distinct positive masses forces $S$ to be diagonal real signs, so its
residual group is exactly $\Gamma_N$. The representative inverse obeys
$R(Y_\nu S)=R(Y_\nu)S$, proving that only the coset $R\Gamma_N$
descends to the presentation class. Conversely the forward and inverse
formulas distinguish classes precisely up to this sign action.
\end{proof}
\begin{proposition}[Sign-invariant rank-one neutrino fiber coordinates]
\label{p7v2:prop:v382-neutral-lift-rankone-quotient}
Writing $Y_\nu=[y_1,y_2,y_3]$, define
$P_i=-(v^2/2M_i)y_i y_i^T$. Each $P_i$ is invariant under
$\Gamma_N$, and $\sum_iP_i=m_\nu^{\rm cmp}$.
On the nonsingular labeled branch the triple $(P_1,P_2,P_3)$
distinguishes all residual heavy-Majorana sign classes.
Furthermore $h_\nu=Y_\nu^\dagger Y_\nu$ transforms as
$h_\nu\mapsto Sh_\nu S$; in particular
$\operatorname{Im}(h_{\nu,23}^2)$ is presentation-invariant.
\end{proposition}
\begin{proof}
Changing any column sign leaves its rank-one square unchanged.
Conversely $y_i y_i^T=z_i z_i^T$ with nonzero complex columns implies
$z_i=\pm y_i$; the full ordered triple therefore identifies exactly a
$\Gamma_N$ orbit. The seesaw sum is its defining identity.
The Hermitian Gram transformation gives the final squared-entry invariant.
\end{proof}


\begin{proposition}[Low-energy mixing does not fix the unflavored heavy source]
\label{p7v2:prop:pmns-cancels}
For \eqref{p7v2:eq:casas-ibarra},
\begin{equation}
\label{p7v2:eq:heavy-gram}
 h_\nu:=\YN^\dagger\YN
 =
 \frac{2}{v^2}
 D_N^{1/2}R^\dagger D_\nu R D_N^{1/2}.
\end{equation}
The conventional comparison current section cancels.  Therefore an unflavored invariant such as
\(\operatorname{Im}(h_{\nu,23}^2)\) is not determined by the low-energy
Jarlskog invariant unless a definite member of the \(\Hnu\) neutral-lift
fiber is supplied.  Once that \(Y_\nu\) member is fixed, \(R\) is only its
coordinate by Corollary~\ref{p7v2:cor:casas-ibarra-fiber-coordinate}.
\end{proposition}

\begin{proof}
Use \((U_{\rm PMNS}^{\rm cmp})^T(U_{\rm PMNS}^{\rm cmp})^{*}=I\) in
\(\YN^\dagger\YN\).  Equation~\eqref{p7v2:eq:heavy-gram} follows directly.
Since different complex-orthogonal matrices \(R\) give the same low-energy
matrix \eqref{p7v2:eq:seesaw-identity} but different heavy Gram matrices, the
low-energy invariant cannot determine the heavy one.
\end{proof}

\begin{corollary}[Neutral-lift observable-quotient test]
\label{p7v2:cor:neutral-lift-observable-quotient}
Fix a compatible low-energy base member and consider the neutral-lift fibre
$\mathfrak A_{\rm NU}$.  For any claimed readout $O$ on that fibre, no
high-energy member selector is required if $O$ is constant on the complete
compatible fibre (equivalently, if it factors through the frozen base
projection).  On a smooth path-connected stratum it is sufficient to prove
that $DO$ annihilates every vertical fibre tangent, by Paper~II-A
Corollary~\ref{p2a:cor:local-observable-insensitivity}.

Proposition~\ref{p7v2:prop:pmns-cancels} is a negative test for the unflavoured
heavy Gram source: different complex-orthogonal $R$ can change $h_\nu$, so
those observables remain lift-sensitive unless a stronger same-owner theorem
is proved.
\end{corollary}

\begin{theorem}[Neutral-lift determinant readout is fibre invariant]
\label{p7v2:thm:neutral-lift-determinant-closure}
On the positive nonsingular neutral-lift fibre of
Corollary~\ref{p7v2:cor:casas-ibarra-fiber-coordinate}, let
\[
 h_\nu=Y_\nu^\dagger Y_\nu
 =\frac{2}{v^2}D_N^{1/2}R^\dagger D_\nu R D_N^{1/2},
 \qquad R\in O(3,\CC).
\]
Then the determinant readout is independent of the fibre member:
\[
 \boxed{
 \det h_\nu
 =\left(\frac{2}{v^2}\right)^3
   \det D_N\,\det D_\nu .}
\]
Equivalently, the product of the squared singular values of $Y_\nu$ is fixed
by the frozen light/heavy spectra and $v$.  Hence this particular high-energy
readout is closed by Paper~II-A
Theorem~\ref{p2a:thm:compatible-family-observable-closure} and requires no
Casas--Ibarra member selector, even though matrix entries, traces, off-diagonal
Gram data and leptogenesis sources remain lift-sensitive in general.
\end{theorem}

\begin{proof}
Taking determinants in
\eqref{p7v2:eq:heavy-gram} gives
\[
 \det h_\nu
 =\left(\frac{2}{v^2}\right)^3
 \det D_N\,\det(R^\dagger)\det D_\nu\det R.
\]
Since $RR^T=I$, $(\det R)^2=1$, so $\det R=\pm1$ and therefore
$\det(R^\dagger)\det R=|\det R|^2=1$.  The displayed result follows and is
constant on the complete positive nonsingular fibre.
\end{proof}

\subsection{Resonant kernel: bound versus selector}
\label{p7v2:sec:resonant-bound}

\begin{theorem}[Exact resonant-regulator bound]
\label{p7v2:thm:resonant-bound}
For \((X,Y)\in\mathbb R^2\setminus\{(0,0)\}\), define
\begin{equation}
\label{p7v2:eq:resonant-factor}
 \mathcal R(X,Y)
 =
 \frac{XY}{X^2+Y^2}.
\end{equation}
Then
\begin{equation}
\label{p7v2:eq:resonant-bound}
 |\mathcal R(X,Y)|\leq\frac12,
\end{equation}
with equality if and only if \(|X|=|Y|\neq0\).  For the conventional
two-state regulator one sets
\[
 X=M_i^2-M_j^2,\qquad Y=M_i\Gamma_j.
\]
\end{theorem}

\begin{proof}
The denominator is positive on the stated domain.  The inequality
\((|X|-|Y|)^2\geq0\) is equivalent to
\(2|XY|\leq X^2+Y^2\).  Equality holds exactly when
\(|X|=|Y|\), and the domain then forces their common value to be nonzero.
\end{proof}

\begin{corollary}[Saturation is not a dynamics theorem]
\label{p7v2:cor:no-resonance-selector}
The full regulator family is retained, and the saturation locus
\begin{equation}
\label{p7v2:eq:resonance-saturation-locus}
 \mathfrak A_{\rm res}^{\rm sat}
 =\{(X,Y):|X|=|Y|\neq0\}
\end{equation}
is only one special subset.  The bound
\eqref{p7v2:eq:resonant-bound} does not select this subset.  Endpoint
saturation is a status-\(\mathsf A\) member of \(\Hresopt\) unless an independent
variational or dynamical law proves
\(|M_i^2-M_j^2|=M_i\Gamma_j\).
\end{corollary}

\begin{proof}
On \(Y\ne0\), put \(r=|X/Y|\); the continuous function
\(r/(1+r^2)\) takes every value in \([0,1/2]\), while the coordinate axes
supply the value zero.  Compactness supplies a
maximum but does not imply that the physical system realizes the maximizing
point.  No generic Paper--III endpoint rule can be imported to select this
maximum: if Paper~V-C uses saturation, it owns that member locally through
\(\Hresopt\).
\end{proof}

\subsection{Hopf overlap: trace distinction and open matrix element}
\label{p7v2:sec:hopf-overlap}

\begin{theorem}[Dimension obstruction to the former Hopf coisometry]
\label{p7v2:thm:hopf-coisometry-no-go}
Let the protected carrier \(\mathcal K\) have real dimension \(8\), and let
\[
 \mathcal W
 =
 \bigoplus_{q=1}^{8}\Sym^q(\CC^2).
\]
Then
\begin{equation}
\label{p7v2:eq:hopf-dimensions}
 \dim_\CC\mathcal W
 =
 \sum_{q=1}^{8}(q+1)
 =44.
\end{equation}
Thus \(\dim_\RR\mathcal W=88\).  There is no real-linear coisometry from
\(\mathcal K\) onto \(\mathcal W\); if \(\mathcal K\) is equipped with its
compatible complex structure, the same conclusion holds a fortiori for a
complex-linear map.
\end{theorem}

\begin{proof}
A real coisometry onto \(\mathcal W\) would have real rank \(88\), while every
real-linear map from \(\mathcal K\) has rank at most \(8\).  This is
impossible.  With a compatible complex structure the source has complex
dimension \(4\), whereas the target has complex dimension \(44\), giving the
same contradiction.
\end{proof}

\begin{proposition}[Uniform-label and Hilbert-trace averages are inequivalent]
\label{p7v2:prop:hopf-trace-distinction}
For the inverse-charge function \(q\mapsto1/q\), the uniform average on the
eight charge labels is
\begin{equation}
\label{p7v2:eq:hopf-label-average}
 h_{\rm label}
 =
 \frac18\sum_{q=1}^{8}\frac1q
 =
 \frac{H_8}{8}
 =
 \frac{761}{2240}
 =
 0.339732142857\ldots.
\end{equation}
The normalized Hilbert-space trace on
\(\mathcal W\), which weights each sector by \(\dim\Sym^q(\CC^2)=q+1\), is
\begin{equation}
\label{p7v2:eq:hopf-hilbert-average}
 h_{\rm Hilb}
 =
 \frac{\sum_{q=1}^{8}(q+1)/q}
      {\sum_{q=1}^{8}(q+1)}
 =
 \frac{8+H_8}{44}
 =
 \frac{3001}{12320}
 =
 0.243587662338\ldots.
\end{equation}
Neither measure is obtained from the other by a change of notation.
\end{proposition}

\begin{proof}
The first expression is the arithmetic average of eight labels.  The second is
the ordinary normalized trace with each eigenvalue repeated according to its
module dimension.  Evaluation of the two finite sums gives the displayed
rational numbers, which are unequal.
\end{proof}

\begin{definition}[Open Hopf matrix-element realization family]
\label{p7v2:def:open-hopf-element}
The physical heavy-pair role is a typed operator/state expectation
\[
 \boxed{
 \mathfrak A_{\rm Hopf}
 =\{(\mathcal O_{\rm Hopf},\omega):
      h_{23}=\omega(\mathcal O_{\rm Hopf})\}.
 }
\]
The heavy-pair Hopf factor is written
\begin{equation}
\label{p7v2:eq:open-hopf-factor}
 \kappa_{\rm Hopf}=\delta\,h_{23},
\end{equation}
where \(h_{23}\) is the member selected by \(\Hhopfopt\).  The choice
\(h_{23}=H_8/8\) is the status-\(\mathsf A\) uniform-label ansatz, while
\(h_{23}=(8+H_8)/44\) is the distinct Hilbert-trace ansatz.  Their disagreement
is evidence for the unresolved operator/state family, not two notations for one
expectation.
\end{definition}

\begin{proposition}[The Hopf expectation is optional for the standard target]
\label{p7v2:prop:hopf-optional}
The standard type-I resonant source and kinetic boundary consume
\(\Hnu+\Hkin\).  Neither the seesaw identity nor the kinetic system contains
\(\mathcal O_{\rm Hopf}\).  In the absence of a separately constructed typed
map
\[
 (\mathcal O_{\rm Hopf},\omega)
 \longmapsto \varepsilon_{23},
\]
the finite numbers \(h_{\rm label}\) and \(h_{\rm Hilb}\) do not determine the
CP source.  Therefore \(\Hhopfopt\) is an optional Born-specific model branch,
not an irreducible premise of the standard baryogenesis calculation.
\end{proposition}

\begin{proof}
Equations~\eqref{p7v2:eq:seesaw-identity},
\eqref{p7v2:eq:heavy-gram}, and \eqref{p7v2:eq:kinetic-system} are well typed
without a Hopf expectation.  Conversely, Definition~\ref{p7v2:def:open-hopf-element}
only produces the scalar \(h_{23}\); it supplies no morphism from that scalar
to a Yukawa invariant or collision source.  Composing the two sectors would
therefore require new operator/dynamical data.  Omitting an unused optional
branch loses no consequence of the standard target.
\end{proof}

\subsection{Conditional heavy-cell calibration}
\label{p7v2:sec:heavy-cell}

\begin{definition}[Optional exactly saturated benchmark]
\label{p7v2:def:heavy-calibration}
For a complex \(23\)-rotation
\[
 R_{23}(z)
 =
 \begin{pmatrix}
 1&0&0\\
 0&\cos z&\sin z\\
 0&-\sin z&\cos z
 \end{pmatrix},
 \qquad z=x+iy,
\]
let \(h_\nu\) be \eqref{p7v2:eq:heavy-gram}, and use the tree-level width
\(\Gamma_3=M_{N_3}h_{\nu,33}/(8\pi)\).  The repaired optional benchmark imposes
\begin{align}
\label{p7v2:eq:heavy-calibration}
 \Re h_{\nu,23}&=\Im h_{\nu,23}>0,\\
 \frac12\,
 \frac{\Im(h_{\nu,23}^{\,2})}
      {h_{\nu,22}h_{\nu,33}}
 &=\varepsilon_{23}^{\rm cand},\\
 \frac{M_{N_2}M_{N_3}}{8\pi}h_{\nu,33}
 &=M_{N_2}^2-M_{N_3}^2.
\end{align}
The second equation imposes the candidate CP source.  The third is exactly
\(M_{N_2}\Gamma_3=M_{N_2}^2-M_{N_3}^2\), not the near-degenerate replacement
\(\Gamma_3=2(M_{N_2}-M_{N_3})\).  Thus the system defines an
\(\Hresopt\) calibrated realization; it does not derive the candidate source
or physical saturation.  Here \(\Gamma_3\) is an algebraic benchmark width
coordinate; calling it a physical partial width additionally uses
Proposition~\ref{p7v2:prop:claim-scoped-decay-boundary}.
\end{definition}

\begin{theorem}[Analytic reduction and unique saturated branch]
\label{p7v2:thm:heavy-calibration-reduction}
Assume \(0<m_2<m_3\), \(0<\varepsilon<1/2\), and put
\[
 A=m_2+m_3,\qquad B=m_2-m_3,
 \qquad
 q^2=\frac{4\varepsilon m_2m_3}{1-2\varepsilon}.
\]
If \(q^2<B^2\), the principal small-angle branch
\begin{equation}
\label{p7v2:eq:heavy-angle-closed-form}
 x=\frac12\arcsin\frac qB,
 \qquad
 y=\frac12\operatorname{arsinh}\frac qA
\end{equation}
satisfies the first two equations of
\eqref{p7v2:eq:heavy-calibration}, with
\begin{equation}
\label{p7v2:eq:heavy-C-coefficient}
 h_{\nu,33}=\frac{M_{N_3}}{v^2}\,C,
 \qquad
 C=\sqrt{A^2+q^2}+\sqrt{B^2-q^2}>0.
\end{equation}

Now let \(E_N>0\), \(0<\Delta<E_N/2\), and set
\[
 M_2(\Delta)=E_N-\frac{55}{56}\Delta,
 \qquad
 M_3(\Delta)=E_N-\frac{57}{56}\Delta,
 \qquad
 k=\frac{C}{8\pi v^2}.
\]
Then exact saturation is the single scalar equation
\begin{equation}
\label{p7v2:eq:heavy-monotone-equation}
 G(\Delta)
 :=\frac{\Delta(E_N-\Delta)}{14}
   -kM_2(\Delta)M_3(\Delta)^2
 =0.
\end{equation}
Moreover
\begin{equation}
\label{p7v2:eq:heavy-monotonicity}
 G'(\Delta)
 =\frac{E_N-2\Delta}{14}
 +k\left(
   \frac{55}{56}M_3^2
   +2\frac{57}{56}M_2M_3
 \right)>0
\end{equation}
on \((0,E_N/2)\).  Consequently, any sign-changing bracket in this interval
contains exactly one saturated solution.  No three-variable local-isolation
assumption is required.
\end{theorem}

\begin{proof}
For the \(23\) rotation, direct multiplication in
\eqref{p7v2:eq:heavy-gram} gives
\[
\begin{aligned}
 h_{\nu,23}
 &=\frac{\sqrt{M_2M_3}}{v^2}
 \bigl(B\sin 2x+iA\sinh 2y\bigr),\\
 h_{\nu,22}
 &=\frac{M_2}{v^2}\bigl(A\cosh 2y+B\cos 2x\bigr),\\
 h_{\nu,33}
 &=\frac{M_3}{v^2}\bigl(A\cosh 2y-B\cos 2x\bigr).
\end{aligned}
\]
Equation~\eqref{p7v2:eq:heavy-angle-closed-form} yields
\(B\sin2x=A\sinh2y=q>0\), so
\(\Re h_{\nu,23}=\Im h_{\nu,23}>0\).  The principal branches have
\(\cos2x>0\) and \(\cosh2y>0\), whence
\[
 A\cosh2y-B\cos2x
 =\sqrt{A^2+q^2}+\sqrt{B^2-q^2}=C.
\]
Furthermore
\[
\frac{\tfrac12\Im(h_{\nu,23}^2)}{h_{\nu,22}h_{\nu,33}}
=\frac{q^2}{A^2-B^2+2q^2}
=\frac{q^2}{4m_2m_3+2q^2}
=\varepsilon,
\]
where the final equality is precisely the definition of \(q^2\).  This proves
the closed-form reduction.

The affine masses obey
\(M_2-M_3=\Delta/28\) and
\(M_2+M_3=2(E_N-\Delta)\), hence
\(M_2^2-M_3^2=\Delta(E_N-\Delta)/14\).  Equations
\eqref{p7v2:eq:heavy-C-coefficient} and
\(\Gamma_3=M_3h_{\nu,33}/(8\pi)\) give
\(M_2\Gamma_3=kM_2M_3^2\), proving
\eqref{p7v2:eq:heavy-monotone-equation}.  Differentiation gives
\eqref{p7v2:eq:heavy-monotonicity}; every term is positive on the stated
interval.  Strict monotonicity and the intermediate-value theorem prove the
last assertion.
\end{proof}

\begin{proposition}[Certified numerical enclosure of the benchmark]
\label{p7v2:prop:heavy-calibration-numerics}
Using the rounded inputs
\[
 E_N=2.9588\times10^9\,{\rm GeV},\qquad
 (m_1,m_2,m_3)=(1.36,8.60,49.99)\,{\rm meV},
\]
\[
 v=246.2553\,{\rm GeV},\qquad
 \varepsilon_{23}^{\rm cand}=4.1945421\times10^{-7},
\]
the hypotheses of Theorem~\ref{p7v2:thm:heavy-calibration-reduction} hold.
Directed rational interval evaluation, with a Machin-formula rational
enclosure for \(\pi\) and bounding intervals for the two square roots in
\eqref{p7v2:eq:heavy-C-coefficient}, gives
\[
 G(8040.053958\,\mathrm{GeV})<0,
 \qquad
 G(8040.053959\,\mathrm{GeV})>0.
\]
Therefore the unique physical small-branch root is certified by
\begin{equation}
\label{p7v2:eq:heavy-calibration-root}
 \boxed{
 8040.053958<\frac{\Delta_M}{\mathrm{GeV}}<8040.053959,
 \qquad
 z\simeq-3.2444258956\times10^{-4}
   +2.2919742907\times10^{-4}i.}
\end{equation}
It gives
\begin{equation}
\label{p7v2:eq:heavy-calibration-output}
 M_{N_2}-M_{N_3}=287.1447842\ldots\,{\rm GeV},
 \qquad
 \Gamma_3=574.2895406\ldots\,{\rm GeV},
 \qquad
 \mathcal R_{23}=\frac12.
\end{equation}
The last equality is exact on the enclosed root, not a rounding inference.
\end{proposition}

\begin{proof}
The input decimals are treated as the exact benchmark coordinates.  The
source-level certificate evaluates \(q^2\) rationally, encloses each square
root by adjacent decimal rationals, and derives rational bounds for \(\pi\)
from Machin's identity and alternating arctangent remainders.  It obtains
\[
 -66.2\,\mathrm{GeV}^2<G(8040.053958\,\mathrm{GeV})
 <-66.1\,\mathrm{GeV}^2
\]
and
\[
 145.1\,\mathrm{GeV}^2<G(8040.053959\,\mathrm{GeV})
 <145.3\,\mathrm{GeV}^2.
\]
Theorem~\ref{p7v2:thm:heavy-calibration-reduction} then gives existence and
uniqueness in the displayed interval.  The split is \(\Delta_M/28\), and the
width may be evaluated as
\((M_2^2-M_3^2)/M_2\); interval substitution yields the printed digits.
Finally \(X=Y>0\) in Theorem~\ref{p7v2:thm:resonant-bound}, so
\(\mathcal R_{23}=1/2\) exactly.
\end{proof}

\begin{remark}[What was derived and what remains optional]
The angle formula, scalar reduction, monotonicity, root existence, and root
uniqueness are now theorems.  They eliminate the former numerical-solve and
local-Jacobian assumptions.  They do not turn
\(\varepsilon_{23}^{\rm cand}\) or exact saturation into consequences of the
carrier: changing either optional benchmark input changes the root while
preserving the carrier and seesaw identity.
\end{remark}

\begin{remark}[Washout is not negligible on the calibrated branch]
With the standard radiation-era estimate
\[
 H(T)=1.66\sqrt{g_*}\,\frac{T^2}{M_{\rm Pl}},
 \qquad g_*=106.75,
\]
the same rounded branch gives approximately
\[
 K_1\simeq1.27,\qquad K_2\simeq8.03,\qquad K_3\simeq46.70.
\]
Thus neither an inert \(N_1\) channel nor negligible washout follows from the
calibrated spectrum.  These values are themselves conditional on the standard
cosmological interface and on
\eqref{p7v2:eq:heavy-calibration}.
\end{remark}

\subsection{Independent numerical audit}
\label{p7v2:app:numerical-audit}

\begin{table}[htbp]
\centering
\caption[Independent numerical audit]{Reproducible arithmetic retained in
Paper~V-C.  A correct numerical row does not elevate an ansatz to a theorem.}
\label{p7v2:tab:numerical-audit}
\small
\renewcommand{\arraystretch}{1.16}
\begin{tabularx}{\textwidth}{@{}L{0.28\textwidth}L{0.38\textwidth}Y@{}}
\toprule
Quantity & Reproduced value & Status\\
\midrule
\(\delta,\wp_\star\)
& \(1/\sqrt{72}\), \(1/56\)
& exact finite counts\\
Tree-shell couplings
& \(g_2=0.6523411464\), \(g_Y=0.3486910099\),
  \(g_1=0.4501581581\)
& conditional mass dictionary\\
Tree-shell \(\alpha_{\rm em}^{-1}\)
& \(132.88404292\)
& conditional internal boundary, not \(M_Z\)-scheme observable\\
Electromagnetic strict branch
& \(137.035999137806\)
& finite-sum-rule ansatz\\
Electromagnetic half-octic branch
& \(137.035999176726\) expanded;
  \(137.035999175630\) compact
& inequivalent finite-sum-rule ansatz\\
Hopf label trace
& \(761/2240=0.339732142857\)
& uniform-label ansatz\\
Hopf Hilbert trace
& \(3001/12320=0.243587662338\)
& dimension-weighted alternative\\
Certified saturated heavy root
& \(8040.053958<\Delta_M/{\rm GeV}<8040.053959\),
  \(z\simeq-3.2444259\times10^{-4}+2.2919743\times10^{-4}i\)
& unique root of the optional exact-saturation benchmark\\
Calibrated split and width
& \(287.1447842\,{\rm GeV}\), \(574.2895406\,{\rm GeV}\)
& conditional output; regulator exactly \(1/2\)\\
Baryon benchmark
& \(\eta_B=6.12773944\times10^{-10}\)
& arithmetic benchmark, not prediction\\
\bottomrule
\end{tabularx}
\end{table}


% Paper V-D, after p5c:prop:v239-jacobi-endpoint-interference.
\begin{proposition}[Native four-level residue and consecutive-cutoff obstruction]
\label{p5c:prop:v246-native-jacobi-residue-cutoff}
On the native alternating scalar-receiver ladder of
\eqref{p5c:eq:v2169-b13-o9}, with $u,\omega>0$, the four-level
return at $u=\omega(\xi+3\omega-e)>0$ has first-site residue
\[
 \boxed{Z_e=\frac{u}{2(u+\omega^2)}\in(0,1/2).}
\]
For any fixed positive-link Jacobi ladder, no real energy is an
eigenvalue of two consecutive leading-principal cutoffs. Since the
first site is cyclic, no energy is a tagged pole at both cutoffs.
\end{proposition}
\begin{proof}
At the stated return the native four-level matrix has eigenvector
\[
 v=(1,0,-1,\sqrt2\,\omega/\sqrt u)^T
\]
at eigenvalue $e$, as direct substitution shows. Positivity of every
link makes the first site cyclic and the eigenvalue simple. Hence its
residue is $|v_1|^2/\|v\|^2=1/(2+2\omega^2/u)$.
For the general claim, the continuants satisfy
$p_{n+1}(z)=(z-d_n)p_n(z)-w_{n-1}^2p_{n-1}(z)$.
If $p_n(\lambda)=p_{n+1}(\lambda)=0$, the nonzero link implies
$p_{n-1}(\lambda)=0$. Repetition contradicts $p_0=1$.
\end{proof}
An isolated exact return therefore cannot certify an exact pole
independent of every consecutive cutoff. The result permits convergence
of moving poles and returns at nonconsecutive cutoffs; it makes no
assertion that the full physical Hamiltonian is this constituent.


% v2.248 proposed closure on the already stated actual two-spin source restriction.
\begin{theorem}[Complete interacting two-spin odd-source poles and residues]
\label{p5c:thm:v248-paired-quartic-source-poles}
On the actual two-spin restriction of
Section~\ref{p5c:v2188:app:chargediagnostics}, retain
\[
 H=\epsilon_0n_0+\epsilon_1n_1+v n_0n_1
       +g a_0^\dagger a_1^\dagger+\bar g a_1a_0,
 \qquad \rho_\beta=e^{-\beta H}/Z,
\]
with real $\epsilon_0,\epsilon_1,v$, finite $\beta\ge0$ and
$g\ne0$. Define
\[
 E=\epsilon_0+\epsilon_1+v,\quad
 \Omega=\sqrt{E^2/4+|g|^2},\quad
 \lambda_\eta=E/2+\eta\Omega\quad(\eta=\pm1),
\]
\[
 p_\eta=e^{-\beta\lambda_\eta}/Z,\qquad
 p_0=e^{-\beta\epsilon_0}/Z,\qquad
 p_1=e^{-\beta\epsilon_1}/Z.
\]
For the inherited odd source $a_0^\dagger$, the complete stationary
anticommutator resolvent is
\[
 \boxed{\begin{aligned}
 G_0(z)&=\operatorname{Tr}\rho_\beta
  \{a_0,(z-\operatorname{ad}_H)^{-1}a_0^\dagger\}\\
 &=\sum_{\eta=\pm1}\left[
 \frac{(p_\eta+p_0)(1-\eta E/(2\Omega))/2}
      {z-(\epsilon_0-\lambda_\eta)}
 +\frac{(p_\eta+p_1)(1+\eta E/(2\Omega))/2}
      {z-(\lambda_\eta-\epsilon_1)}\right].
 \end{aligned}}
\]
Every displayed residue is strictly positive, their sum is one, and
coincident frequencies are merged by adding their residues. If
$\delta=(\epsilon_0-\epsilon_1)/2$, its complete tagged polynomial
is the monic squarefree part of
\[
 \boxed{P(z)=
 [((z-\delta)+v/2)^2-\Omega^2]
 [((z-\delta)-v/2)^2-\Omega^2].}
\]
For $v\ne0$ and $\Omega\ne|v|/2$, this is a four-pole interacting
response; no terminal-coordinate pole is imposed in its construction.
\end{theorem}
\begin{proof}
Use the ordered even basis $(|0\rangle,|01\rangle)$ and odd
vectors $|0_{\rm f}\rangle=a_0^\dagger|0\rangle$,
$|1_{\rm f}\rangle=a_1^\dagger|0\rangle$. The even block is
\[
 H_{\rm even}=\begin{pmatrix}0&\bar g\\g&E\end{pmatrix},
\]
with eigenvalues $\lambda_\eta$, while the odd eigenvalues are
$\epsilon_0,\epsilon_1$. Its normalized eigenvectors have squared
empty and doubly occupied components
\[
 w_{\eta,0}=(1-\eta E/(2\Omega))/2,\qquad
 w_{\eta,2}=(1+\eta E/(2\Omega))/2.
\]
The source maps the empty component of each even eigenvector to
$|0_{\rm f}\rangle$ and maps $|1_{\rm f}\rangle$ to its doubly
occupied component. These are all its nonzero energy-eigenbasis
matrix elements. A stationary anticommutator spectral term from
$|n\rangle$ to $|m\rangle$ has frequency $E_m-E_n$ and residue
$(p_m+p_n)|\langle m|a_0^\dagger|n\rangle|^2$.
This gives the boxed response. Since $g\ne0$, $|E/(2\Omega)|<1$;
finite Gibbs probabilities are positive. Summing the residues uses
$\sum_\eta w_{\eta,0}=\sum_\eta w_{\eta,2}=1$ and
$w_{\eta,0}+w_{\eta,2}=1$, giving total weight one.
The four frequencies are $\delta\pm(\Omega\pm v/2)$, whose product
polynomial is $P$. Positive residues preclude cancellation of any
distinct pole; only repeated frequencies are merged. This proves the
tagged-polynomial and distinctness assertions.
\end{proof}

\begin{corollary}[Quadratic and number-preserving limits of the same response]
\label{p5c:cor:v248-paired-source-limits}
On the quadratic subbranch $v=0$, merging the coincident source poles
in the theorem gives the state-independent response
\[
 \boxed{G_0(z)=\frac{z+\epsilon_1}
 {(z-\epsilon_0)(z+\epsilon_1)-|g|^2}.}
\]
On the number-preserving subbranch $g=0$, direct evaluation gives
\[
 \boxed{G_0(z)=\frac{1-\langle n_1\rangle_\beta}{z-\epsilon_0}
 +\frac{\langle n_1\rangle_\beta}{z-\epsilon_0-v}.}
\]
The latter formula also covers the $g=E=0$ coordinate, without
introducing $1/\Omega$.
\end{corollary}
\begin{proof}
For $v=0$, $\epsilon_0-\lambda_\eta=
\lambda_{-\eta}-\epsilon_1$ and
$w_{\eta,0}=w_{-\eta,2}$. The two merged probability sums equal
one, yielding the stated two-pole rational function. For $g=0$,
$a_0^\dagger(1-n_1)$ and $a_0^\dagger n_1$ are exact
Liouvillian eigenoperators at $\epsilon_0$ and $\epsilon_0+v$;
their positive residue operators are $1-n_1$ and $n_1$.
\end{proof}
This closes the source-visible observable on its stated interacting
restriction and supplied full Gibbs state. It neither selects the actual
couplings or preparation nor identifies this restriction with the full
physical electron Hamiltonian. The quartic and anomalous terms are
retained in the same action; no normal-occupation replacement is used.

\begin{proposition}[Exact interacting two-spin history compiler]
\label{p5c:prop:v248-paired-history-compiler}
Let the coefficients of the same two-spin Hamiltonian follow any
supplied integrable finite history, with its inherited clock in units
$\hbar=1$. Let $u_{e,\pm}$ solve the complete $2\times2$ even-block
Schr\"odinger equations on the two histories, and let
$\rho_e,p_0,p_1$ be the corresponding diagonal parity blocks of the
fixed initial density, with
$\operatorname{Tr}\rho_e+p_0+p_1=1$. Then its exact interacting
source influence trace is
\[
 \boxed{\begin{aligned}
 Z[q_+,q_-]={}&\operatorname{tr}_2
     (u_{e,+}\rho_eu_{e,-}^\dagger)\\
 &+p_0\exp\!\left[-i\int(\epsilon_{0,+}-\epsilon_{0,-})\,dt\right]\\
 &+p_1\exp\!\left[-i\int(\epsilon_{1,+}-\epsilon_{1,-})\,dt\right].
 \end{aligned}}
\]
The quartic $v(t)$ is retained in the even entry
$E(t)=\epsilon_0(t)+\epsilon_1(t)+v(t)$.
\end{proposition}
\begin{proof}
Parity reduces each complete Hamiltonian into its displayed even
$2\times2$ block and two scalar odd blocks. The unitary propagators
have the same direct-sum decomposition. Trace the exact doubled
unitary against the fixed initial density. Off-diagonal initial blocks
have zero trace against this block-diagonal unitary product.
\end{proof}


% BEGIN FRAMEWORK v2.321: EXACT PETERSEN MATTER CERTIFICATE
\section{Exact rooted Petersen--matter representation certificates}
\label{p5d:sec:v321-petersen-matter-certificate}

This technical section verifies the finite statements added to
Papers~V-A--V-C without inserting a physical member or mass coefficient.
All vector spaces in the first two propositions are oriented real graph
edge spaces; a complexification is made only for the matter comparison.
Its group-theoretic domain is the occurring $2I$ packet; no assertion is
made that a merely local $2T$ core forces this parent.

\begin{proposition}[Signed induced-edge and rooted regular-orbit characters]
\label{p5d:prop:v321-rooted-edge-characters}
Write Petersen vertices as unordered two-subsets of
$\Omega=\{1,2,3,4,5\}$ and edges as partitions of four labels into two
unordered pairs. The signed $A_5$-action on oriented edges gives
\[
 \mathscr W_P\simeq\operatorname{Ind}_{A_4}^{A_5}\mathbf3,
 \qquad
 \mathscr W_P\!\downarrow_{A_4}
   \simeq\mathbf3\oplus\operatorname{Reg}(A_4).
\]
The characters on the $A_4$ element orders $(1,2,3)$ are
\[
 \begin{array}{c|ccc}
 &1&2&3\\ \hline
 V_i&3&-1&0\\
 V_i^\perp&12&0&0\\
 Q\otimes G&12&0&0
 \end{array}.
\]
In particular the twelve external edges form one \emph{signed} regular
$A_4$ module, while the three internal edges transform as the genuine
vector $\mathbf3$, not as the unsigned three-point permutation module.
\end{proposition}
\begin{proof}
An edge omits exactly one of the five labels. Under its stabilizer $A_4$,
the three omitted-label partitions carry the signed permutation character
$3,-1,0$: a nontrivial double transposition fixes one oriented edge and
reverses the other two, while a three-cycle fixes none. The other twelve
edges are indexed by ordered distinct pairs of the remaining four labels;
$A_4$ acts sharply transitively on them, so their character is regular.
The specified comparison action has $Q=\mathbf1\oplus\mathbf3$ and
$G=\mathbf1\oplus\mathbf1'\oplus\mathbf1''$, whose character product is
also regular. The induction statement is the standard induced-module
construction from the $A_5$ orbit of five roots.
\end{proof}

\begin{proposition}[Four canonical conjugation sheets]
\label{p5d:prop:v321-four-complement-sheets}
For the rooted $A_4\subset A_5$ the set $A_5\setminus A_4$ has four free
conjugation orbits, distinguishable by element order and by the two
icosahedral order-five conjugacy classes. Their sizes are
\[
 (12_{2},12_{3},12_{5,+},12_{5,-}).
\]
The outer twelve Petersen edges map to the first orbit by the
$A_4$-equivariant rule
$\{ab,cd\}\mapsto(ab)(cd)$. This is not an assertion that conjugation
and left multiplication on $A_5\setminus A_4$ are the same group action:
they admit isomorphic regular-module decompositions but different natural
set and operator coordinates.
\end{proposition}
\begin{proof}
The fifteen order-two elements of $A_5$ contain three in the root $A_4$;
the twenty order-three elements contain eight; and none of the twenty-four
order-five elements lie in $A_4$. Each of the four indicated subsets is
stable under $A_4$ conjugation and has twelve elements. For any one of its
elements $g$, the intersection of $A_4$ with its $A_5$ centralizer is
trivial, so conjugation is free. The outside involutions are precisely
the double transpositions moving the root label, equivalently Petersen
edges not omitting it.
\end{proof}

\begin{theorem}[Sector-compression proof independent of Hopf coordinates]
\label{p5d:thm:v321-sector-compression-certificate}
Let $P_i$ project onto the three oriented edges omitting core $i$, and let
$\Pi_r$ be an orthogonal $A_5$-central projector on $\mathscr W_P$ with
rank $d_r$ ($d_r=4,5,3,3$). Then
\[
 \boxed{P_i\Pi_rP_i=\frac{d_r}{15}P_i.}
\]
In particular the local comparison $C$ sector has equal quadratic overlap
with $\mathbf3$ and $\mathbf3'$ and does not detect their projector
difference at this order.
\end{theorem}
\begin{proof}
The restriction to $V_i$ of $\Pi_r$ commutes with the irreducible real
$A_{4,i}$ action. The commutant of the tetrahedral vector representation
is scalar, so $P_i\Pi_rP_i=c_rP_i$. As the five root triples partition the
edges and $A_5$ permutes them transitively, the traces
$\operatorname{Tr}(P_i\Pi_r)$ are independent of $i$ and sum to
$\operatorname{Tr}(\Pi_r)=d_r$. Hence $3c_r=d_r/5$.
\end{proof}

\begin{remark}[Machine-verifiable exact certificate and scope]
\leavevmode\par\noindent
The supplementary file
\texttt{Paper5D-\allowbreak PetersenMatter-\allowbreak v1\_1-\allowbreak verify.py} enumerates the twelve
root-stabilizing even permutations, all five local triples, the complete
forty-eight-element projective complement, and the exact rational/
$\mathbb Q(\sqrt5)$ projectors. It checks the seven independent theorem
groups and two additional local-corollary checks. No experimental masses,
particle identities, operator coefficients or phenomenological fits enter.
\end{remark}
% END FRAMEWORK v2.321



% BEGIN FRAMEWORK v2.348: FLAVOR AND SEESAW INFORMATION BOUNDARIES
\clearpage
\section{Exact orientation and neutral-lift non-identifiability checks}
\label{p5d:sec:v348-information-boundaries}
\begin{proposition}[Mixed Hermitian pair moments cannot fix signed CP]
\label{p5d:prop:v348-mixed-moments-cp-sign-nogo}
For a real diagonal $3\times3$ Hermitian $A$ with distinct
eigenvalues and an arbitrary Hermitian $B$, each mixed moment
$\operatorname{tr}(A^mB^n)$, for nonnegative integers $m,n$,
is real and unchanged by $B\mapsto B^*$ (entrywise complex
conjugation).  However the odd invariant
\[
 \chi(A,B):=\frac1{6i}\operatorname{tr}[A,B]^3
\]
changes sign under the same map.  Whenever $\chi(A,B)\ne0$,
CP-even mixed traces --- including all moments with $m,n\le2$ ---
do not determine the sign of the Jarlskog invariant.  A CP-odd
orientation readout must be supplied as separate data.
\end{proposition}
\begin{proof}
The product trace of two Hermitian matrices is real, so
$\operatorname{tr}(A^mB^n)\in\mathbb R$.  Since $A$ is real,
$\operatorname{tr}(A^m(B^*)^n)
=\overline{\operatorname{tr}(A^mB^n)}$.
Similarly $[A,B^*]=[A,B]^*$, so its cubic trace is the complex
conjugate of the original; the latter is purely imaginary for the
anti-Hermitian commutator.  Hence $\chi(A,B^*)=-\chi(A,B)$.
Nonzero examples exist: the primitive-trimer witness of
Theorem~\ref{p6:thm:v548-trimer-local-submersion} has $\chi=378$,
and a common unitary change of basis can diagonalize its first
Hermitian member without changing $\chi$.
\end{proof}

\begin{proposition}[The neutral-lift determinant is fixed while its trace is unbounded]
\label{p5d:prop:v348-casas-trace-unbounded}
On the positive nonsingular external type-I seesaw comparison fibre
of Theorem~\ref{p7v2:thm:seesaw-identity}, retain the correctly
phased Casas--Ibarra coordinate
$Y_\nu=(i\sqrt2/v)U^*D_\nu^{1/2}RD_N^{1/2}$,
$R\in O(3,\mathbb C)$.  The determinant is invariant,
\[
 \det(Y_\nu^\dagger Y_\nu)
 =\left(\frac2{v^2}\right)^3\det D_N\det D_\nu,
\]
but for fixed positive $D_N,D_\nu$ its trace is unbounded above
along a one-parameter family of compatible fibre members.
\end{proposition}
\begin{proof}
The determinant is the established
Theorem~\ref{p7v2:thm:neutral-lift-determinant-closure}; no squared
mass determinants enter.  For real $t$, set
\[
 R(t)=\begin{pmatrix}
 \cosh t&i\sinh t&0\\
 -i\sinh t&\cosh t&0\\
 0&0&1
 \end{pmatrix}.
\]
Direct multiplication gives $R(t)R(t)^T=I$.  With
$D_\nu=\operatorname{diag}(m_1,m_2,m_3)$ and
$D_N=\operatorname{diag}(M_1,M_2,M_3)$,
\[
 \begin{split}
 \operatorname{tr}(Y_\nu^\dagger Y_\nu)
 =\frac2{v^2}\bigl[&M_1(m_1\cosh^2t+m_2\sinh^2t)\\
 &+M_2(m_1\sinh^2t+m_2\cosh^2t)+M_3m_3\bigr].
 \end{split}
\]
Every mass is positive, so the expression diverges as
$|t|\to\infty$ while the light matrix and determinant stay fixed.
This exhibits a strict observable distinction between fibre-invariant
and fibre-sensitive high-energy readouts, without selecting $R$.
\end{proof}

\begin{remark}[Ownership of existing finite proofs]
\label{p5d:rem:v348-existing-certificate-ownership}
The $25$-to-$17$ incidence rank and single-label non-descent are
already proved in Proposition~\ref{p6:prop:v179-label-exclusion-nondescent};
the exact rank-ten trimer local submersion and the nonzero CP witness
are already Theorem~\ref{p6:thm:v548-trimer-local-submersion}.
Framework v2.348 strengthens their downstream scope, not their
underlying finite calculations.  The role-gap quotient already appears
in Paper~V-B; only the explicit physical pole-bridge countermodel is
new.  No claim here identifies terminal geometric readouts with
observed masses or complete interaction poles.
\end{remark}
% END FRAMEWORK v2.348

\printbibliography

% BEGIN FRAMEWORK v2.365: ORDER-SENSITIVE 25-CELL NORMALIZATIONS
\section{Three noncommuting 25-cell exclusion operations}
\label{p5d:sec:v365-operator-ordered-exclusion}
Retain the previously established transporter incidence
$B$ of rank $17$, its effective-label projector $P_{\rm eff}$,
the normalized equal-label vector $\Omega_{25}$, and a chosen
coordinate unit vector $e_H$.  Set $Q_H=I-|e_H\rangle\langle e_H|$
and $v=P_{\rm eff}e_H$.  The inherited Gram decomposition implies
\[
P_{\rm eff}\Omega_{25}=\Omega_{25},\quad
\|v\|^2=\tfrac{17}{25},\qquad
|\langle\Omega_{25},v\rangle|^2=\tfrac1{25}.
\]
\begin{proposition}[Exact ordered norm and inverse comparisons]
\label{p5d:prop:v365-three-polar-factors}
When the inverse is restricted to the effective support,
\[
\boxed{\langle\Omega_{25},Q_H\Omega_{25}\rangle^{-1}
 =\tfrac{25}{24},\qquad
\langle\Omega_{25},(P_{\rm eff}Q_HP_{\rm eff})^{-1}
 \Omega_{25}\rangle=\tfrac98,}
\]
\[
\boxed{\bigl\|P_{\rm eff}Q_H\Omega_{25}\bigr\|^{-2}
 =\tfrac{625}{592}.}
\]
The last expression also equals the squared inverse output norm
after the polar incidence isometry $U_B$ is applied, because
$U_B^\dagger U_B=P_{\rm eff}$.
\end{proposition}
\begin{proof}
The first expression is $(1-|\langle e_H,\Omega_{25}\rangle|^2)^{-1}$.
On $\operatorname{im}P_{\rm eff}$ the compression is
$I-|v\rangle\langle v|$ with positive eigenvalue
$1-17/25=8/25$ along $v$.  Sherman--Morrison gives
$(P_{\rm eff}Q_HP_{\rm eff})^{-1}=I+(25/8)|v\rangle\langle v|$
on that support; hence its stated expectation is $1+1/8=9/8$.
Finally $P_{\rm eff}Q_H\Omega_{25}=\Omega_{25}-v/5$,
whose squared norm is $1-2/25+17/625=592/625$.
\end{proof}
\begin{remark}[No new physical coupling]
\label{p5d:rem:v365-polar-factors-scope}
These operations are ordered differently and therefore their
rational outputs need not agree.  They are finite incidence
comparators, not three estimates of one interaction vertex.
The physical-current normalization remains the existing
same-action response and whitening problem.  In particular,
$625/592$ is not an inferred gauge coupling, and $25/108$
remains an unselected post-hoc neutral-current hypothesis.
\end{remark}
% END FRAMEWORK v2.365


% BEGIN v3.058 ROOTED HOPF TWO-CONTACT TECHNICAL COMPANION
\section{Fixed-chart verification of the rooted two-contact return}
\label{p5d:sec:v158-two-contact-return-checks}
This technical check uses the established rooted family chart, not a
new physical source, field member, or mass calibration.  Let
$c=3^{-1/2}(1,1,1)^\mathsf T$,
$W=\operatorname{diag}(1,i,-i)$,
$P=cc^\dagger$, $Q=I-P$, $u=QWc$, and
$P_\mu=9uu^\dagger/8$; put $P_\tau=P$ and
$P_e=I-P-P_\mu$.  The labels $E,M,T$ denote the \emph{finite}
positive squared role coefficients, not physical propagator-pole data.

\begin{proposition}[Explicit matrix replay and positive return]
\label{p5d:prop:v158-two-contact-matrix-replay}
On this fixed chart $u=(2,3i-1,-3i-1)^\mathsf T/(3\sqrt3)$,
$\|u\|^2=8/9$, $P_\mu P_\tau=0$, and
\[
K_*=EP_e+MP_\mu+TP_\tau
=E I+(T-E)P+(M-E)P_\mu.
\]
For $x=E-M$ and $y=E+M-2T$ its upper offdiagonal entries are
\[
(K_*)_{12}=\frac{-x-2y+3ix}{12},\quad
(K_*)_{13}=\frac{-x-2y-3ix}{12},\quad
(K_*)_{23}=\frac{2x-2y+3ix}{12}.
\]
Therefore $\operatorname{Tr}(K_*^2)=E^2+M^2+T^2$,
$\operatorname{Im}((K_*)_{12}(K_*)_{23}(K_*)_{31})
=(E-M)(E-T)(M-T)/12$ and the squared Hilbert--Schmidt
distance to the $W$ commutant is
\[
 \boxed{\inf_{[N,W]=0}\|NN^\dagger-K_*\|_{\rm HS}^2
 =\frac{11}{24}(E-M)^2+\frac16(E+M-2T)^2.}
\]
The bound is attained by $N=\sqrt{\operatorname{diag}(K_*)}$
in the $W$ eigenframe.  It is a restricted fixed-common-chart result.
\end{proposition}
\begin{proof}
Directly $c^\dagger Wc=1/3$ and
$u=Wc-c/3$, giving the displayed coordinates and norm.
The orthogonal rank-one projector resolution gives the first identity
and the spectral trace moment.  Expanding the offdiagonal entries
gives the cubic.  Since the eigenvalues of $W$ are distinct,
$[N,W]=0$ forces $N$ to be diagonal.  The diagonal projection of
$K_*$ is positive, and its squared offdiagonal norm is precisely the
claimed minimum; its positive diagonal square root attains it.
\end{proof}

\begin{proposition}[Exact correction-rank strata and chart firewall]
\label{p5d:prop:v158-return-rank-strata-check}
Set $d=\#\{E,M,T\}_{\rm distinct}$ and define the
Hermitian commuting-baseline class by
\[
\mathfrak A_W^{\rm sa}=\{D=D^\dagger:[D,W]=0\}.
\]
Then
\[
 \boxed{\min_{D\in\mathfrak A_W^{\rm sa}}\operatorname{rank}(K_*-D)=d-1.}
\]
This counts rank of a \emph{return correction}, not of physical
interaction vertices.  Reversing the rooted orientation
$W\mapsto W^\dagger$ reverses the cubic's sign but does not alter
its absolute rank obstruction.  The conclusion for a source $N$
requires that $[N,W]=0$ be evaluated in the same common left/right
chart as its target; it is not a theorem for independently rotated
left and right kinetic frames.
\end{proposition}
\begin{proof}
For $d=3$ a rank-one Hermitian correction $a vv^\dagger$ to
a diagonal matrix would have real cyclic offdiagonal product,
contrary to the nonzero cubic.  The scalar baseline $EI$ gives
the upper bound two.  If $d=2$, the target differs from the repeated
eigenvalue times $I$ by a nonzero multiple of one of the three
rooted rank-one projectors, giving rank at most one; every such
projector has a nonzero offdiagonal entry, excluding a diagonal
$K_*$.  If $d=1$, the target is scalar.  Orientation reversal is
complex conjugation in the fixed real rooted chart.
\end{proof}
\begin{remark}[Evidence and physical-status boundary]
\label{p5d:rem:v158-two-contact-proof-boundary}
These are exact finite calculations for the comparison target,
not an evaluation of $D_{\rm ch}^{\rm act}$ or proof of
$B_{\rm act}=K_*$.  The two-action-contraction acceptance test is
the already established theorem
\ref{p3f:thm:v267-two-scalar-return-defect}; the sector-owned
$\mathsf{Cert}_{\rm hier}$ and interacting pole/readout obligations
retain their original proof gates.
\end{remark}
% END v3.058 ROOTED HOPF TWO-CONTACT TECHNICAL COMPANION

\end{document}
